\documentclass[fs85,seriesWit]{korfarm-base}
\usepackage{korfarm-witt}

% ============================================================
% passage-note.tex (v11)
% 지문 박스 + 첫 페이지 우측 메모란 (동적 높이)
% 정책: PAGE_BREAK_POLICY.md §3, §4 참조
%
% 변경 (v11):
%  · 지문 박스 폭 확대: enlarge right by=-3.7cm  (메모 3.4cm + gap 3mm)
%  · 메모란 동적 높이: frame.north east ~ frame.south east 사이 자동 채움
%  · 메모 줄선: clip 영역으로 박스 내부만 채움 (박스 높이 모르고도 작동)
%  · 문단 들여쓰기 활성화: tcolorbox 내부 \parindent=1em, \noindent 제거
%
% 사용:
%   \kfPassageNote{제목}{지문 본문}      → 메모란 포함
%   \kfPassageNoteNT{지문 본문}          → 메모란 포함, 제목 없음
%   \kfPassageFull{제목}{본문}           → 메모란 없음 (활동 내 보기)
% ============================================================

\makeatletter

% 메모란 규격 (cm 고정)
\newlength{\kf@memoW}
\setlength{\kf@memoW}{3.4cm}
\newlength{\kf@memoGap}
\setlength{\kf@memoGap}{3mm}

% --- 메인: 제목 있는 버전 ---
\newcommand{\kf@passagenote@with}[2]{%
  \par\ifkfPassageNewPage\ifdim\pagetotal>0.4\textheight\clearpage\fi\fi\smallskip\needspace{6\baselineskip}%
  \begin{tcolorbox}[
    enhanced, breakable,
    width=\dimexpr\linewidth-4cm\relax,
    colback=kfPassageBg, colframe=kfPrimary!70,
    boxrule=0.6pt, arc=3pt,
    left=\kfBoxPad, right=\kfBoxPad, top=\dimexpr\kfBoxPad+8pt\relax, bottom=\kfBoxPad,
    title={\small\bfseries\color{kfPrimary} #1},
    coltitle=kfPrimary,
    colbacktitle=kfPrimaryLight!60,
    attach boxed title to top left={yshift=-2.4mm, xshift=8pt},
    boxed title style={colframe=kfPrimary!50, boxrule=0.5pt, arc=2pt},
    parbox=false,
    overlay unbroken and first={%
      \begin{scope}
        \draw[kfRule, line width=0.4pt, fill=kfNoteBg, rounded corners=2pt]
          ([xshift=\kf@memoGap]frame.north east) rectangle
          ([xshift=\kf@memoGap+\kf@memoW]frame.south east);
        \node[anchor=north east, font=\footnotesize\bfseries\color{kfMute}, inner sep=3pt]
          at ([xshift=\kf@memoGap+\kf@memoW, yshift=-2pt]frame.north east) {메모};
        \node[anchor=north east, fill=kfPrimary!15, text=kfPrimary,
              font=\scriptsize\bfseries, rounded corners=2pt,
              inner xsep=4pt, inner ysep=2pt]
          at ([xshift=\kf@memoGap+\kf@memoW-3pt, yshift=-19pt]frame.north east) {\kf@memocat};
        \begin{scope}
          \path[clip]
            ([xshift=\kf@memoGap]frame.north east) rectangle
            ([xshift=\kf@memoGap+\kf@memoW]frame.south east);
          \foreach \i in {1,...,40}{%
            \draw[kfRule, line width=0.25pt]
              ([xshift=\kf@memoGap+2mm, yshift=-7mm - \i*5mm]frame.north east) --
              ([xshift=\kf@memoGap+\kf@memoW-2mm, yshift=-7mm - \i*5mm]frame.north east);%
          }
        \end{scope}
      \end{scope}
    }
  ]
    \setlength{\parindent}{1em}%
    \hspace*{1em}\ignorespaces#2%
  \end{tcolorbox}\par\smallskip
}

% --- 제목 없는 버전 ---
\newcommand{\kf@passagenote@notitle}[1]{%
  \par\ifkfPassageNewPage\ifdim\pagetotal>0.4\textheight\clearpage\fi\fi\smallskip\needspace{6\baselineskip}%
  \begin{tcolorbox}[
    enhanced, breakable,
    width=\dimexpr\linewidth-4cm\relax,
    colback=kfPassageBg, colframe=kfPrimary!70,
    boxrule=0.6pt, arc=3pt,
    left=\kfBoxPad, right=\kfBoxPad, top=\dimexpr\kfBoxPad+8pt\relax, bottom=\kfBoxPad,
    parbox=false,
    overlay unbroken and first={%
      \begin{scope}
        \draw[kfRule, line width=0.4pt, fill=kfNoteBg, rounded corners=2pt]
          ([xshift=\kf@memoGap]frame.north east) rectangle
          ([xshift=\kf@memoGap+\kf@memoW]frame.south east);
        \node[anchor=north east, font=\footnotesize\bfseries\color{kfMute}, inner sep=3pt]
          at ([xshift=\kf@memoGap+\kf@memoW, yshift=-2pt]frame.north east) {메모};
        \node[anchor=north east, fill=kfPrimary!15, text=kfPrimary,
              font=\scriptsize\bfseries, rounded corners=2pt,
              inner xsep=4pt, inner ysep=2pt]
          at ([xshift=\kf@memoGap+\kf@memoW-3pt, yshift=-19pt]frame.north east) {\kf@memocat};
        \begin{scope}
          \path[clip]
            ([xshift=\kf@memoGap]frame.north east) rectangle
            ([xshift=\kf@memoGap+\kf@memoW]frame.south east);
          \foreach \i in {1,...,40}{%
            \draw[kfRule, line width=0.25pt]
              ([xshift=\kf@memoGap+2mm, yshift=-7mm - \i*5mm]frame.north east) --
              ([xshift=\kf@memoGap+\kf@memoW-2mm, yshift=-7mm - \i*5mm]frame.north east);%
          }
        \end{scope}
      \end{scope}
    }
  ]
    \setlength{\parindent}{1em}%
    \hspace*{1em}\ignorespaces#1%
  \end{tcolorbox}\par\smallskip
}

\newcommand{\kfPassageNote}[2]{%
  \def\kf@tmptitle{#1}%
  \ifx\kf@tmptitle\@empty
    \kf@passagenote@notitle{#2}%
  \else
    \kf@passagenote@with{#1}{#2}%
  \fi
}
\newcommand{\kfPassageNoteNT}[1]{\kf@passagenote@notitle{#1}}
\makeatother

% ============================================================
% 풀폭 지문 박스 (메모란 없음) — 활동 내 보기 박스
% PAGE_BREAK_POLICY.md §3
% ============================================================
\makeatletter
\newcommand{\kf@passagefull@with}[2]{%
  \par\ifkfPassageNewPage\ifdim\pagetotal>0.4\textheight\clearpage\fi\fi\smallskip\needspace{6\baselineskip}%
  \begin{tcolorbox}[
    enhanced, breakable,
    colback=kfPassageBg, colframe=kfPrimary!70,
    boxrule=0.6pt, arc=3pt,
    left=\kfBoxPad, right=\kfBoxPad, top=\dimexpr\kfBoxPad+8pt\relax, bottom=\kfBoxPad,
    title={\small\bfseries\color{kfPrimary} #1},
    coltitle=kfPrimary,
    colbacktitle=kfPrimaryLight!60,
    attach boxed title to top left={yshift=-2.4mm, xshift=8pt},
    boxed title style={colframe=kfPrimary!50, boxrule=0.5pt, arc=2pt},
    parbox=false
  ]
    \setlength{\parindent}{1em}%
    \hspace*{1em}\ignorespaces#2%
  \end{tcolorbox}\par\smallskip
}
\newcommand{\kf@passagefull@notitle}[1]{%
  \par\ifkfPassageNewPage\ifdim\pagetotal>0.4\textheight\clearpage\fi\fi\smallskip\needspace{6\baselineskip}%
  \begin{tcolorbox}[
    enhanced, breakable,
    colback=kfPassageBg, colframe=kfPrimary!70,
    boxrule=0.6pt, arc=3pt,
    left=\kfBoxPad, right=\kfBoxPad, top=\dimexpr\kfBoxPad+8pt\relax, bottom=\kfBoxPad,
    parbox=false
  ]
    \setlength{\parindent}{1em}%
    \hspace*{1em}\ignorespaces#1%
  \end{tcolorbox}\par\smallskip
}
\newcommand{\kfPassageFull}[2]{%
  \def\kf@tmptitle{#1}%
  \ifx\kf@tmptitle\@empty
    \kf@passagefull@notitle{#2}%
  \else
    \kf@passagefull@with{#1}{#2}%
  \fi
}
\newcommand{\kfPassageFullNT}[1]{\kf@passagefull@notitle{#1}}
\makeatother

% ============================================================
% 작은 문장 박스 (문장 독해용) — PAGE_BREAK_POLICY.md §2.5
% ============================================================
\newcommand{\kfSentenceBox}[2]{%
  \par\smallskip\needspace{3\baselineskip}%
  \noindent\begin{tcolorbox}[
    enhanced, breakable=false,
    colback=kfPassageBg, colframe=kfMute,
    boxrule=0.4pt, arc=2pt,
    left=8pt, right=6pt, top=4pt, bottom=4pt,
    title={\footnotesize\bfseries\color{kfPrimary} #1},
    coltitle=kfPrimary, colbacktitle=kfPaper,
    attach boxed title to top left={yshift=-1.6mm, xshift=4pt},
    boxed title style={colframe=kfMute, boxrule=0.3pt, arc=1pt}
  ]%
    \setstretch{1.3}#2%
  \end{tcolorbox}\par\smallskip%
}

% ============================================================
% 지문 본문 아래 안내/정리 박스 (v16 신규)
% 사용: \kfPassageHint{한 줄 정리}{본문 안내 텍스트}
% ============================================================
\newcommand{\kfPassageHint}[2]{%
  \par\smallskip\needspace{4\baselineskip}%
  \begin{tcolorbox}[
    enhanced, breakable=false,
    colback=kfPrimary!5, colframe=kfPrimary!40,
    boxrule=0.4pt, arc=2pt,
    left=8pt, right=6pt, top=4pt, bottom=4pt,
    title={\footnotesize\bfseries\color{kfPrimary} #1},
    coltitle=kfPrimary, colbacktitle=kfPaper,
    attach boxed title to top left={yshift=-1.5mm, xshift=4pt},
    boxed title style={colframe=kfPrimary!40, boxrule=0.3pt, arc=1pt}
  ]%
    \small\setstretch{1.3}#2%
  \end{tcolorbox}\par\smallskip%
}

% ============================================================
% 환경 버전 (호환성) — 메모란 포함
% ============================================================
\makeatletter
\newenvironment{kfPassageNoteEnv}[1]%
  {\par\needspace{6\baselineskip}%
   \def\kf@psrc{#1}%
   \begin{tcolorbox}[
     enhanced, breakable,
     width=\dimexpr\linewidth-4cm\relax,
     colback=kfPassageBg, colframe=kfPrimary!70,
     boxrule=0.6pt, arc=3pt,
     left=\kfBoxPad, right=\kfBoxPad, top=\dimexpr\kfBoxPad+8pt\relax, bottom=\kfBoxPad,
     title={\small\bfseries\color{kfPrimary} \kf@psrc},
     coltitle=kfPrimary, colbacktitle=kfPrimaryLight!60,
     attach boxed title to top left={yshift=-2.4mm, xshift=8pt},
     boxed title style={colframe=kfPrimary!50, boxrule=0.5pt, arc=2pt},
     parbox=false,
     overlay unbroken and first={%
       \begin{scope}
         \draw[kfRule, line width=0.4pt, fill=kfNoteBg, rounded corners=2pt]
           ([xshift=\kf@memoGap]frame.north east) rectangle
           ([xshift=\kf@memoGap+\kf@memoW]frame.south east);
         \node[anchor=north east, font=\footnotesize\bfseries\color{kfMute}, inner sep=3pt]
           at ([xshift=\kf@memoGap+\kf@memoW, yshift=-2pt]frame.north east) {메모};
         \begin{scope}
           \path[clip]
             ([xshift=\kf@memoGap]frame.north east) rectangle
             ([xshift=\kf@memoGap+\kf@memoW]frame.south east);
           \foreach \i in {1,...,40}{%
             \draw[kfRule, line width=0.25pt]
               ([xshift=\kf@memoGap+2mm, yshift=-7mm - \i*5mm]frame.north east) --
               ([xshift=\kf@memoGap+\kf@memoW-2mm, yshift=-7mm - \i*5mm]frame.north east);%
           }
         \end{scope}
       \end{scope}
     }
   ]%
   \setlength{\parindent}{1em}%
  }%
  {\end{tcolorbox}\par\smallskip}
\makeatother

% ============================================================
% condition-box.tex
% <조건> 박스 컴포넌트 (tcolorbox)
% 사용:
%   \begin{kfCondition}
%     \item 첫째 조건
%     \item 둘째 조건
%   \end{kfCondition}
% ============================================================

\newenvironment{kfCondition}%
  {\par\smallskip\needspace{4\baselineskip}%
   \begin{tcolorbox}[
     enhanced, breakable=false,
     colback=kfPrimaryLight!40, colframe=kfPrimary,
     boxrule=0.5pt, arc=2pt,
     left=\kfBoxPad, right=\kfBoxPad, top=\kfBoxPad, bottom=\kfBoxPad,
     title={\small\bfseries\color{kfPrimary} \,조건\,},
     coltitle=kfPrimary, colbacktitle=kfPaper,
     attach boxed title to top left={yshift=-2mm, xshift=6pt},
     boxed title style={colframe=kfPrimary, boxrule=0.4pt, arc=1pt}
   ]%
   \begin{enumerate}[leftmargin=16pt, itemsep=1pt, topsep=1pt,
                    label=\textbf{\arabic*.}]}%
  {\end{enumerate}\end{tcolorbox}\par\smallskip}

% ============================================================
% evidence-box.tex
% <보기> 박스 컴포넌트
% 사용:
%   \begin{kfEvidence}
%     보기 본문 ...
%   \end{kfEvidence}
% ============================================================

\newenvironment{kfEvidence}%
  {\par\smallskip\needspace{4\baselineskip}%
   \begin{tcolorbox}[
     enhanced, breakable=false,
     colback=kfPaper, colframe=kfMute,
     boxrule=0.5pt, arc=2pt,
     left=\kfBoxPad, right=\kfBoxPad, top=\kfBoxPad, bottom=\kfBoxPad,
     title={\small\bfseries\color{kfInk} \,보기\,},
     coltitle=kfInk, colbacktitle=kfPrimaryLight!50,
     attach boxed title to top left={yshift=-2mm, xshift=6pt},
     boxed title style={colframe=kfMute, boxrule=0.4pt, arc=1pt}
   ]%
   \leavevmode\mbox{}\ignorespaces%
  }%
  {\end{tcolorbox}\par\smallskip}

% ============================================================
% choice-list.tex
% 5선지 (① ② ③ ④ ⑤) 자동 들여쓰기 컴포넌트
% 사용:
%   \begin{kfChoices}
%     \kfChoice 첫째 선택지
%     \kfChoice 둘째 선택지
%     \kfChoice 셋째 선택지
%     \kfChoice 넷째 선택지
%     \kfChoice 다섯째 선택지
%   \end{kfChoices}
% ============================================================

\newcounter{kfChoiceCnt}
\newcommand{\kfChoiceMark}[1]{%
  \ifcase#1\relax%
  \or ①\or ②\or ③\or ④\or ⑤\or ⑥\or ⑦\or ⑧\or ⑨%
  \fi
}

% 환경: 자동 번호
\newenvironment{kfChoices}%
  {\setcounter{kfChoiceCnt}{0}%
   \par\smallskip%
   \begin{list}{}{%
     \setlength{\leftmargin}{20pt}%
     \setlength{\itemindent}{0pt}%
     \setlength{\labelwidth}{18pt}%
     \setlength{\labelsep}{6pt}%
     \setlength{\itemsep}{2pt}%
     \setlength{\topsep}{1pt}%
     \setlength{\parsep}{0pt}%
     \renewcommand{\makelabel}[1]{##1\hss}%
   }}%
  {\end{list}\par\smallskip}

\newcommand{\kfChoice}{%
  \stepcounter{kfChoiceCnt}%
  \item[\textbf{\kfChoiceMark{\value{kfChoiceCnt}}}]\ignorespaces%
}

% --- 한 줄 5선지 (짧은 선택지용) ---
\newcommand{\kfChoicesInline}[5]{%
  \par\smallskip%
  \noindent\textbf{①}~#1\quad\textbf{②}~#2\quad\textbf{③}~#3\quad\textbf{④}~#4\quad\textbf{⑤}~#5%
  \par\smallskip%
}

% ============================================================
% answer-note.tex
% 답안 작성란 (노트 스타일, 줄친 영역)
% 사용:
%   \kfAnswerLines{5}        % 5줄 답안란
%   \kfAnswerNote{서술형 답안란}{8} % 라벨+8줄
% ============================================================

% 단일 줄 (필요 시 인라인)
\newcommand{\kfAnswerLine}{%
  \par\noindent\textcolor{kfRule}{\rule{\linewidth}{0.4pt}}\par
}

% N줄 답안란 (간단)
\newcommand{\kfAnswerLines}[1]{%
  \par\smallskip%
  \begin{minipage}{\linewidth}%
  \foreach \i in {1,...,#1}{%
    \textcolor{kfRule}{\rule{\linewidth}{0.4pt}}\\[6pt]%
  }%
  \end{minipage}%
  \par\smallskip%
}

% 라벨이 있는 노트형 답안란
\newcommand{\kfAnswerNote}[2]{%
  \par\smallskip\needspace{#2\baselineskip}%
  \begin{tcolorbox}[
    enhanced, breakable=false,
    colback=kfPaper, colframe=kfRule,
    boxrule=0.4pt, arc=2pt,
    left=\kfBoxPad, right=\kfBoxPad, top=\kfBoxPad, bottom=\kfBoxPad,
    title={\footnotesize\bfseries\color{kfMute} #1},
    coltitle=kfMute, colbacktitle=kfPaper,
    attach boxed title to top left={yshift=-1.5mm, xshift=6pt},
    boxed title style={colframe=kfRule, boxrule=0.3pt, arc=1pt}
  ]%
  \foreach \i in {1,...,#2}{%
    \textcolor{kfRule}{\rule{\linewidth}{0.3pt}}\\[6pt]%
  }%
  \end{tcolorbox}\par\smallskip%
}

% 짧은 빈칸 (단답형)
\newcommand{\kfBlank}[1][3cm]{\underline{\hspace{#1}}}
% 넓은 빈칸 (서술 한 줄)
\newcommand{\kfBlankWide}[1][6cm]{\underline{\hspace{#1}}}
% 괄호 빈칸
\newcommand{\kfBlankParen}{(\,\hspace{1cm}\,)}
% O/X 빈칸
\newcommand{\kfBlankOX}{(\,\hspace{1.5cm}\,)}

% ============================================================
% question-block.tex
% 문제 블록 (번호 배지 + 발문 + 보기/조건 + 선택지 + 답안란)
% 모든 구성을 하나의 minipage로 묶어 단/페이지 단절 금지
%
% 사용:
%   \begin{kfQBlock}{1}{객관식}
%     \kfQStem{다음 글에 대한 설명으로 가장 적절한 것은?}
%     \begin{kfEvidence} ... \end{kfEvidence}    % 선택
%     \begin{kfCondition} ... \end{kfCondition}  % 선택
%     \begin{kfChoices}
%       \kfChoice ...
%     \end{kfChoices}
%     \kfAnswerLines{2}                            % 선택
%   \end{kfQBlock}
% ============================================================

% 무단절 minipage로 묶고, 발문/보기/조건/선택지/답안란을 자유롭게 배치
% 사용자 피드백 #4: [객관식]/[서술형] 등 텍스트 라벨 제거 — 번호 배지만 표시
% #2(유형 라벨)는 호환성 인자로만 받고 출력하지 않음.
\newenvironment{kfQBlock}[2]%
  {\par\smallskip\needspace{6\baselineskip}%
   \noindent\begin{minipage}{\linewidth}%
   \samepage%
   \par\noindent
   \kfQNum{#1}\hspace{4pt}%
   \begin{minipage}[t]{\dimexpr\linewidth-30pt\relax}%
  }%
  {\end{minipage}\end{minipage}\par\smallskip}

% 문제 발문
\newcommand{\kfQStem}[1]{%
  \par\noindent#1\par\vspace{2pt}%
}

% 짧은 소문항 라벨 (1), (2), (3)
\newcounter{kfSubQ}
\newcommand{\kfSubQReset}{\setcounter{kfSubQ}{0}}
\newcommand{\kfSubQ}{%
  \stepcounter{kfSubQ}%
  \par\noindent\textbf{(\arabic{kfSubQ})}~\ignorespaces%
}

% ============================================================
% activity-table.tex
% 활동: 정리표 / 내용 확인표 (마크다운 표 → tabularx 변환)
% 사용:
%   % 일반 tabular (직접 컬럼 명시)
%   \begin{kfActTable}{|l|X|X|}
%     \kfTHead{항목} & \kfTHead{설명} & \kfTHead{학생 작성} \\\hline
%     ① & ... & \kfTBlank \\\hline
%   \end{kfActTable}
%
%   % adjustbox 자동 축소 (정해진 폭으로 강제)
%   \kfActTableWrap{
%     ... (\begin{tabular}...\end{tabular} 코드) ...
%   }
% ============================================================

% 사용 시 본문 마지막 행은 \\\hline 으로 끝나야 함.
% v17.1: 흰색 mask로 Canva 배경 본문 침범 차단
\newenvironment{kfActTable}[1]%
  {\par\smallskip\needspace{4\baselineskip}%
   \renewcommand{\arraystretch}{1.3}%
   \arrayrulecolor{kfRule}%
   \begin{tcolorbox}[blanker, colback=white, left=2pt, right=2pt, top=2pt, bottom=2pt]%
   \noindent\begin{adjustbox}{max width=\linewidth}%
   \begin{tabular}{#1}%
   \hline}%
  {\end{tabular}%
   \end{adjustbox}%
   \end{tcolorbox}%
   \arrayrulecolor{black}%
   \par\smallskip}

% 헤더 행 헬퍼
\newcommand{\kfTHead}[1]{\textbf{\color{kfPrimary} #1}}
\newcommand{\kfTHeadRow}[1]{\rowcolor{kfPrimaryLight!50}\textbf{\color{kfPrimary} #1}}

% 빈 셀 (학생 작성용)
\newcommand{\kfTBlank}{\rule[-1.6em]{0pt}{3.4em}}
\newcommand{\kfTBlankSm}{\rule[-1.0em]{0pt}{2.4em}}

% adjustbox 강제 축소 래퍼 — Python에서 변환된 raw tabular 블록을 감쌀 때 사용
\newcommand{\kfActTableWrap}[1]{%
  \par\smallskip\needspace{4\baselineskip}%
  \begin{tcolorbox}[blanker, colback=white, left=2pt, right=2pt, top=2pt, bottom=2pt]%
  \noindent\begin{adjustbox}{max width=\linewidth, center}%
  #1%
  \end{adjustbox}%
  \end{tcolorbox}\par\smallskip%
}

% ============================================================
% activity-vocab.tex
% 어휘 학습 표 (3열: 번호 - 뜻 - 빈칸)
% 사용:
%   \begin{kfVocab}
%     \kfVocabItem{1}{눈으로 보고 마음으로 깨달음}
%     \kfVocabItem{2}{사물의 이치를 따져 헤아림}
%   \end{kfVocab}
% ============================================================

% adjustbox로 폭 강제 + 일반 tabular + p{} 열 사용
\newenvironment{kfVocab}%
  {\par\smallskip\needspace{4\baselineskip}%
   \renewcommand{\arraystretch}{1.35}%
   \arrayrulecolor{kfRule}%
   \noindent\begin{adjustbox}{max width=\linewidth, center}%
   \begin{tabular}{|>{\centering\arraybackslash}m{1.0cm}|m{8.5cm}|>{\centering\arraybackslash}m{4.0cm}|}%
   \hline
   \rowcolor{kfPrimaryLight!50}%
   \multicolumn{1}{|c|}{\textbf{\color{kfPrimary}번호}} &
   \multicolumn{1}{c|}{\textbf{\color{kfPrimary}뜻풀이}} &
   \multicolumn{1}{c|}{\textbf{\color{kfPrimary}어휘 (정답 작성)}} \\\hline
   \ignorespaces}%
  {\end{tabular}%
   \end{adjustbox}%
   \arrayrulecolor{black}%
   \par\smallskip}

\newcommand{\kfVocabItem}[2]{%
  \textbf{#1} & #2 & \rule[-1.0em]{0pt}{2.4em}\\\hline%
}

% ============================================================
% activity-blank.tex
% 빈칸 채우기 활동
% 사용:
%   \begin{kfBlankList}
%     \kfBlankItem 첫 번째 문장에서 (\kfBlank)은(는) ...
%     \kfBlankItem 둘째 문장에서 ...
%   \end{kfBlankList}
% ============================================================

\newenvironment{kfBlankList}%
  {\par\smallskip%
   \begin{enumerate}[leftmargin=20pt, itemsep=4pt, topsep=2pt,
                    label=\textbf{\arabic*.}, labelsep=6pt]}%
  {\end{enumerate}\par\smallskip}

\newcommand{\kfBlankItem}{\item}

% 텍스트 안에 박스형 빈칸 (단어 1개 채우기)
\newcommand{\kfBlankBox}[1][2.4cm]{%
  \tikz[baseline=-0.5ex]{%
    \draw[kfRule, line width=0.4pt] (0,0) -- ++(#1,0);%
  }%
}

% ============================================================
% activity-ox.tex (v15)
% O/X 표기 활동 — 그래픽 디자인
% 사용:
%   \begin{kfOXList}
%     \kfOXItem 글쓴이는 자연을 예찬하고 있다.\kfOX
%     \kfOXItem 1연과 4연은 수미상관 구조이다.\kfOX
%   \end{kfOXList}
%
% \kfOX = 우측에 [O] [X] 두 박스 (학생이 하나 동그라미)
% ============================================================

\newenvironment{kfOXList}%
  {\par\smallskip%
   \begin{enumerate}[leftmargin=20pt, itemsep=5pt, topsep=2pt,
                    label=\textbf{\arabic*.}, labelsep=6pt]}%
  {\end{enumerate}\par\smallskip}

\newcommand{\kfOXItem}{\item}

% v16: O/X 그래픽 박스 키움 (18×14 → 24×20pt) + 영역색 강조
\newcommand{\kfOX}{%
  \hfill\nolinebreak
  \tikz[baseline=-0.9ex]{%
    \node[draw=kfPrimary, fill=kfPrimary!5, line width=0.6pt, rounded corners=3pt,
          minimum width=24pt, minimum height=20pt, inner sep=0pt,
          font=\normalsize\bfseries\color{kfPrimary}] (ob) {O};%
    \node[draw=kfMute, fill=kfMute!5, line width=0.6pt, rounded corners=3pt,
          minimum width=24pt, minimum height=20pt, inner sep=0pt,
          font=\normalsize\bfseries\color{kfMute},
          right=4pt of ob] (xb) {X};%
  }%
}

% ============================================================
% activity-connect.tex
% 좌우 선 긋기 활동 (2단 양식 + 중앙 여백)
% 사용:
%   \begin{kfConnect}
%     \kfPair{사르트르}{실존은 본질에 앞선다}
%     \kfPair{니체}{초인 사상}
%     \kfPair{하이데거}{현존재}
%   \end{kfConnect}
% ============================================================

\newenvironment{kfConnect}%
  {\par\smallskip%
   \renewcommand{\arraystretch}{1.6}%
   \begin{tabular}{@{}>{\raggedleft\arraybackslash}m{4.5cm}
                     @{\hspace{2.5cm}}
                     >{\raggedright\arraybackslash}m{6.5cm}@{}}}%
  {\end{tabular}\par\smallskip}

\newcommand{\kfPair}[2]{%
  \tikz[baseline]{\node[draw=kfRule, fill=kfPrimaryLight!30, rounded corners=2pt,
        inner sep=4pt, font=\small]{#1};} \tikz[baseline]{\node[circle, draw=kfMute, fill=kfPaper, inner sep=1.5pt]{};}
  &
  \tikz[baseline]{\node[circle, draw=kfMute, fill=kfPaper, inner sep=1.5pt]{};} \tikz[baseline]{\node[draw=kfRule, fill=kfNoteBg, rounded corners=2pt,
        inner sep=4pt, font=\small]{#2};} \\
}

% ============================================================
% activity-writing.tex
% 글쓰기 활동 (큰 노트 영역)
% 사용:
%   \kfWriting{독후감}{12}     % 라벨 + 12줄 큰 노트
%   \kfWritingPrompt{프롬프트 텍스트}
% ============================================================

\newcommand{\kfWritingPrompt}[1]{%
  \par\smallskip%
  \begin{tcolorbox}[
    enhanced, breakable=false,
    colback=kfPrimaryLight!30, colframe=kfPrimary,
    boxrule=0.4pt, arc=3pt,
    left=\kfBoxPad, right=\kfBoxPad, top=\kfBoxPad, bottom=\kfBoxPad,
  ]%
  \small\textbf{\color{kfPrimary}글쓰기 안내}\\[2pt]%
  #1
  \end{tcolorbox}\par\smallskip%
}

\newcommand{\kfWriting}[2]{%
  \par\smallskip\needspace{\numexpr#2+2\relax\baselineskip}%
  \begin{tcolorbox}[
    enhanced, breakable=false,
    colback=kfPaper, colframe=kfRule,
    boxrule=0.4pt, arc=2pt,
    left=\kfBoxPad, right=\kfBoxPad, top=\kfBoxPad, bottom=\kfBoxPad,
    title={\footnotesize\bfseries\color{kfMute} #1},
    coltitle=kfMute, colbacktitle=kfPaper,
    attach boxed title to top left={yshift=-1.5mm, xshift=6pt},
    boxed title style={colframe=kfRule, boxrule=0.3pt, arc=1pt}
  ]%
  \foreach \i in {1,...,#2}{%
    \textcolor{kfRule}{\rule{\linewidth}{0.3pt}}\\[7pt]%
  }%
  \end{tcolorbox}\par\smallskip%
}

% ============================================================
% activity-structure.tex
% 구조도 (트리 + 빈칸)
% 사용:
%   \begin{kfStructure}
%     \kfNode{0}{서론}{문제 제기}
%     \kfNode{1}{본론}{(\,\quad\,)}
%     \kfNode{1}{본론}{근거 2}
%     \kfNode{0}{결론}{주장 강화}
%   \end{kfStructure}
% ============================================================

\newenvironment{kfStructure}%
  {\par\smallskip\needspace{6\baselineskip}%
   \begin{tcolorbox}[
     enhanced, breakable=false,
     colback=kfPaper, colframe=kfRule,
     boxrule=0.4pt, arc=2pt,
     left=\kfBoxPad, right=\kfBoxPad, top=\kfBoxPad, bottom=\kfBoxPad,
     title={\footnotesize\bfseries\color{kfMute} 글의 구조},
     coltitle=kfMute, colbacktitle=kfPrimaryLight!40,
     attach boxed title to top left={yshift=-1.5mm, xshift=6pt},
     boxed title style={colframe=kfRule, boxrule=0.3pt, arc=1pt}
   ]%
   \begin{tabbing}
     \hspace{1.4cm}\=\hspace{1.4cm}\=\hspace{8cm}\kill}%
  {\end{tabbing}\end{tcolorbox}\par\smallskip}

% 노드: #1=들여쓰기 단계 (0~3), #2=라벨, #3=내용
\newcommand{\kfNode}[3]{%
  \ifcase#1\relax%
    \>\>\textbf{\color{kfPrimary} ▶ #2}\quad #3\\
  \or
    \>\>\quad\textbf{\color{kfMute} ◦ #2}\quad #3\\
  \or
    \>\>\quad\quad\textbf{\color{kfMute} - #2}\quad #3\\
  \or
    \>\>\quad\quad\quad\textbf{\color{kfMute} \textperiodcentered\ #2}\quad #3\\
  \fi
}

% 빈칸 노드
\newcommand{\kfNodeBlank}[2]{\kfNode{#1}{#2}{(\hspace{2.5cm})}}

% ============================================================
% activity-sentence-reading.tex
% 문장 독해 활동 — 문장 박스 1개 + 그에 묶인 문제 N개를 한 minipage로
% 사용:
%   \begin{kfSentenceUnit}{1}{"바람은 내 귀에 속삭이며..."}
%     \kfSentQ{(1)}{서술어 '속삭이며'의 주어는?}
%     \begin{kfChoices}
%       \kfChoice 바람은
%       ...
%     \end{kfChoices}
%   \end{kfSentenceUnit}
% ============================================================

% kfSentenceBox 매크로는 passage-note.tex에 정의됨

\newenvironment{kfSentenceUnit}[2]%
  {\par\smallskip\needspace{6\baselineskip}%
   \noindent\begin{minipage}{\linewidth}%
   \samepage%
   \kfSentenceBox{문장 #1}{#2}%
   \par\smallskip%
  }%
  {\end{minipage}\par\smallskip}

% 같은 문장에 묶인 소문항 (문장 안내 + 발문)
\newcommand{\kfSentQ}[2]{%
  \par\noindent\textbf{#1}~#2\par\smallskip%
}

% 헤더 없이 문장만 있는 묶음 (sentence_ref가 없는 일반 문항)
\newenvironment{kfSentencelessUnit}%
  {\par\smallskip\noindent\begin{minipage}{\linewidth}\samepage}%
  {\end{minipage}\par\smallskip}

% ============================================================
% chapter-cover.tex (v6)
% 챕터 진입 페이지 (표지) — 하단에 영역 목차 추가
% 사용:
%   \kfChapterCover{1}{근대성과 자아}{Chapter 1}{이번 챕터에서는 ...}
%   \begin{kfChapterAreaList}
%     \kfChapterAreaItem{어휘}{kfAreaVocab}{핵심 어휘 12개}
%     ...
%   \end{kfChapterAreaList}
%
%   영역 목차가 없을 때는 \kfChapterCoverEnd 호출
% ============================================================

\newcommand{\kfChapterCover}[4]{%
  \clearpage%
  \thispagestyle{kfBlank}%
  \kfMarkChapter{Chapter #1. #2}%
  \kfChapterCoverBg%
  % v18.5: 챕터 표지 우상단에 로고
  \begin{tikzpicture}[remember picture, overlay]
    \node[anchor=north east, inner sep=0pt] at ($(current page.north east)+(-18mm,-18mm)$) {\kfLogoMd};
  \end{tikzpicture}%
  \vspace*{20mm}%
  \noindent
  \begin{tikzpicture}[remember picture, overlay]
    % 좌측 액센트 바
    \fill[kfPrimary] (current page.west) ++(12mm,25mm) rectangle ++(5mm, -50mm);
    % 챕터 번호 배지 (이중 원, 약간 작게)
    \node[anchor=north west, inner sep=0pt] at ($(current page.north west)+(24mm,-45mm)$) {%
      \begin{tikzpicture}
        \draw[kfPrimary, line width=1pt] (0,0) circle (10mm);
        \draw[kfPrimary, line width=0.4pt] (0,0) circle (8mm);
        \node[font=\normalsize\bfseries\color{kfPrimary}] at (0,2mm) {CHAPTER};
        \node[font=\LARGE\bfseries\color{kfPrimary}] at (0,-3mm) {#1};
      \end{tikzpicture}%
    };
  \end{tikzpicture}%
  \vspace{42mm}%
  \par\noindent\hspace{32mm}%
  \begin{minipage}{0.65\linewidth}%
    {\footnotesize\color{kfMute} #3}\\[4pt]
    {\LARGE\bfseries\color{kfPrimary} #2}\\[8pt]
    {\color{kfRule}\rule{50mm}{0.5pt}}\\[6pt]
    {\small\color{kfInk} #4}
  \end{minipage}%
  \par\vspace{14mm}%
}

% v6 / v18.5 / v18.6: 챕터 표지의 영역 목차 — 흰색 반투명 박스
% v18.6: 배경 가로선/도형과 안 겹치도록 TikZ overlay로 우하단 절대 위치
\makeatletter
% 영역 항목들을 모아 둘 박스
\newsavebox{\kf@arealistbox}
\newenvironment{kfChapterAreaList}{%
  \begin{lrbox}{\kf@arealistbox}%
  \begin{minipage}{0.62\linewidth}%
    \begin{tcolorbox}[%
      enhanced, colback=white, colframe=white,
      opacityback=0.92, opacityframe=0,
      boxrule=0pt, arc=4pt,
      width=\linewidth,
      top=8pt, bottom=8pt, left=10pt, right=10pt,
    ]%
      {\footnotesize\bfseries\color{kfMute} 이 챕터의 영역}\par\vspace{4pt}%
      {\color{kfRule}\hrule height 0.3pt}\par\vspace{4pt}%
      \begin{itemize}[leftmargin=14pt, itemsep=2pt, topsep=0pt, label={\color{kfPrimary!70}\textbullet}]%
}{%
      \end{itemize}%
    \end{tcolorbox}%
  \end{minipage}%
  \end{lrbox}%
  % v18.7: 배경 상단 가로선 ~ 하단 가로선 사이의 중간 영역에 배치
  % 가로선이 내용을 가리지 않도록 박스 중심이 두 선 사이 중점에 위치
  \begin{tikzpicture}[remember picture, overlay]%
    \node[anchor=center, inner sep=0pt]
      at ($(current page.south)+(0mm,90mm)$)
      {\usebox{\kf@arealistbox}};%
  \end{tikzpicture}%
  \vfill%
  \noindent\hfill\kfSeriesBadge\hspace{12mm}%
  \clearpage%
}
\makeatother

% 영역 항목: \kfChapterAreaItem{영역명}{kfArea색}{부제}
\makeatletter
\newcommand{\kfChapterAreaItem}[3]{%
  \item {\small\bfseries\color{#2} #1}%
  \def\kf@cci@sub{#3}%
  \ifx\kf@cci@sub\@empty\else
    {\small\color{kfMute}\, \textemdash\, #3}%
  \fi
}
\makeatother

% 영역 목록 없이 cover만 (legacy)
\newcommand{\kfChapterCoverEnd}{%
  \vfill%
  \noindent\hfill\kfSeriesBadge\hspace{12mm}%
  \clearpage%
}

% ============================================================
% area-header.tex (v16)
% 영역 시작 표제 — 시리즈별 차별화 라벨 + 큰 영역명 + 부제
% PAGE_BREAK_POLICY.md §1.4
%
% v16 (디자인 고도화 Phase 1):
%   - 시리즈별 라벨 박스 추가:
%     소쉬르 = AREA START (영역색 채움 박스)
%     프레게 = SECTION OPEN (영역색 테두리)
%     러셀   = RUSSELL (영역색 라이트 박스)
%     비트   = WITTGENSTEIN (회색 라이트 박스)
%   - 라벨 위치: 영역명 위 좌측
% ============================================================

\makeatletter

% 시리즈 라벨 텍스트
\newcommand{\kf@serieslabel}{%
  \ifcase\kf@series\relax
    AREA START%
  \or
    AREA START%
  \or
    SECTION OPEN%
  \or
    RUSSELL%
  \or
    WITTGENSTEIN%
  \fi
}

% 시리즈 라벨 박스 스타일 (#1 = 영역색)
\newcommand{\kf@labelbox}[1]{%
  \ifcase\kf@series\relax
    % 0/1 = 소쉬르: 영역색 채움
    \tikz[baseline=-2pt]{\node[fill=#1, text=white,
      rounded corners=3pt, inner xsep=6pt, inner ysep=2pt,
      font=\scriptsize\bfseries]{\kf@serieslabel};}%
  \or
    % 1 = 소쉬르: 영역색 채움
    \tikz[baseline=-2pt]{\node[fill=#1, text=white,
      rounded corners=3pt, inner xsep=6pt, inner ysep=2pt,
      font=\scriptsize\bfseries]{\kf@serieslabel};}%
  \or
    % 2 = 프레게: 영역색 테두리
    \tikz[baseline=-2pt]{\node[draw=#1, line width=0.6pt, text=#1,
      rounded corners=3pt, inner xsep=6pt, inner ysep=2pt,
      font=\scriptsize\bfseries]{\kf@serieslabel};}%
  \or
    % 3 = 러셀: 영역색 라이트 박스
    \tikz[baseline=-2pt]{\node[fill=#1!12, text=#1,
      rounded corners=2pt, inner xsep=5pt, inner ysep=1.5pt,
      font=\scriptsize\bfseries]{\kf@serieslabel};}%
  \or
    % 4 = 비트: 회색 미니 박스
    \tikz[baseline=-2pt]{\node[fill=kfMute!10, text=kfMute,
      rounded corners=2pt, inner xsep=5pt, inner ysep=1.5pt,
      font=\scriptsize\bfseries]{\kf@serieslabel};}%
  \fi
}

% v17: 시리즈+영역별 Canva 배경 자동 매핑
% \kf@hasbg=1로 set되면 \kfAreaStart에서 라벨 박스 출력 생략
\def\kf@areaVocab{어휘}\def\kf@areaGrammar{문법}\def\kf@areaConcept{개념}
\def\kf@areaLit{문학}\def\kf@areaNonlit{비문학}
\newif\ifkf@bgon
% 파일 있으면 배경 적용 + boolean true. 없으면 nothing.
\newcommand{\kfBgSet}[1]{%
  \IfFileExists{#1.pdf}{\kfPageBg{#1}\kf@bgontrue}{}%
}
\newcommand{\kf@trybg}[1]{%
  % v18.4: 영역 배경 PDF 비활성화 — 큰 도형이 본문 영역을 침범해
  % "영역 헤더 후 빈 페이지" 문제를 일으킴. 라벨 박스로만 영역 표시.
  \kf@bgonfalse%
}

% Note: 러셀(rus_*)/비트(wit_*) 배경 PDF 미존재 시 \kfPageBg가 에러 → \IfFileExists로 가드
% (현재는 파일이 모두 존재한다고 가정. 부분 제공 시 cls 수정 필요)

\newcommand{\kfAreaStart}[3]{%
  \clearpage%
  \thispagestyle{fancy}%
  \kfMarkArea{#1}%
  \kf@trybg{#1}%
  \vspace*{1mm}%
  % 배경 없으면 라텍스 라벨 박스 출력 (배경 있으면 일러스트가 라벨 역할)
  \ifkf@bgon
    \vspace*{4mm}%
  \else
    \noindent\kf@labelbox{#2}\par\vspace{4pt}%
  \fi
  % 영역명 + 부제
  \noindent
  \begin{tikzpicture}
    \node[anchor=west, inner sep=0pt] (bar) at (0,0) {\color{#2}\rule{3pt}{22pt}};
    \node[anchor=base west, inner sep=0pt, xshift=8pt, yshift=2pt] (lbl) at (bar.east |- 0,0) {%
      {\LARGE\bfseries\color{#2} #1}%
    };
    \def\kf@areasub{#3}%
    \ifx\kf@areasub\@empty\else
      \node[anchor=base west, inner sep=0pt, xshift=8pt, font=\normalsize\color{kfMute}]
        at (lbl.base east) {\,\textperiodcentered\, #3};%
    \fi
  \end{tikzpicture}%
  \par\vspace{2pt}
  {\color{#2!70}\hrule height 1pt}%
  \par\vspace{1pt}%
  {\color{kfRule}\hrule height 0.3pt}%
  \par\vspace{4pt}%
  \needspace{6\baselineskip}\nopagebreak[4]%
  \global\kf@areaJustOpenedtrue% v18.5: 첫 컨텐츠 needspace 무시
}
\makeatother

% 시리즈/영역 단축 매크로
\newcommand{\kfStartVocab}[1]{\kfAreaStart{어휘}{kfAreaVocab}{#1}}
\newcommand{\kfStartGrammar}[1]{\kfAreaStart{문법}{kfAreaGrammar}{#1}}
\newcommand{\kfStartConcept}[1]{\kfAreaStart{개념}{kfAreaConcept}{#1}}
\newcommand{\kfStartLit}[1]{\kfAreaStart{문학}{kfAreaLiterature}{#1}}
\newcommand{\kfStartNonlit}[1]{\kfAreaStart{비문학}{kfAreaNonlit}{#1}}
\newcommand{\kfStartWeekly}[1]{\kfAreaStart{실력 확인}{kfAreaWeekly}{#1}}

% ============================================================
% section-header.tex
% 섹션 시작 라벨 헤더 — [지문] / [활동] / [문제] 라벨
% 좌측 액센트 바 + 라벨 + 부제 (영역 액센트 색)
%
% 사용:
%   \kfSecHeader{kfAreaLiterature}{지문}{빼앗긴 들에도 봄은 오는가 / 이상화}
%   \kfSecHeader{kfAreaConcept}{활동}{내용 확인}
%   \kfSecHeader{kfAreaNonlit}{문제}{}
%
% 헤더 직후 페이지 분할 금지: \needspace{4\baselineskip} + \nopagebreak
% ============================================================

\providecommand{\kfSecHeader}[3]{%
  \par\needspace{10\baselineskip}\filbreak%
  \vspace{3pt}%
  \noindent
  \begin{tikzpicture}
    \node[anchor=west, inner sep=0pt] (bar) at (0,0) {\color{#1}\rule{3pt}{14pt}};
    \node[anchor=west, inner sep=0pt, xshift=6pt, font=\normalsize\bfseries\color{#1}]
       (lbl) at (bar.east) {[\,#2\,]};
    \node[anchor=west, inner sep=0pt, xshift=4pt, font=\small\color{kfMute}]
       at (lbl.east) {#3};
  \end{tikzpicture}%
  \par\vspace{1pt}%
  {\color{#1!50}\hrule height 0.4pt}%
  \par\vspace{2pt}\nopagebreak%
}

% 영역별 단축 매크로
\newcommand{\kfSecPassageVocab}[1]{\kfSecHeader{kfAreaVocab}{지문}{#1}}
\newcommand{\kfSecPassageGrammar}[1]{\kfSecHeader{kfAreaGrammar}{지문}{#1}}
\newcommand{\kfSecPassageConcept}[1]{\kfSecHeader{kfAreaConcept}{지문}{#1}}
\newcommand{\kfSecPassageLit}[1]{\kfSecHeader{kfAreaLiterature}{지문}{#1}}
\newcommand{\kfSecPassageNonlit}[1]{\kfSecHeader{kfAreaNonlit}{지문}{#1}}
\newcommand{\kfSecPassageWeekly}[1]{\kfSecHeader{kfAreaWeekly}{지문}{#1}}

\newcommand{\kfSecActVocab}[1]{\kfSecHeader{kfAreaVocab}{활동}{#1}}
\newcommand{\kfSecActGrammar}[1]{\kfSecHeader{kfAreaGrammar}{활동}{#1}}
\newcommand{\kfSecActConcept}[1]{\kfSecHeader{kfAreaConcept}{활동}{#1}}
\newcommand{\kfSecActLit}[1]{\kfSecHeader{kfAreaLiterature}{활동}{#1}}
\newcommand{\kfSecActNonlit}[1]{\kfSecHeader{kfAreaNonlit}{활동}{#1}}
\newcommand{\kfSecActWeekly}[1]{\kfSecHeader{kfAreaWeekly}{활동}{#1}}

\newcommand{\kfSecQVocab}[1]{\kfSecHeader{kfAreaVocab}{문제}{#1}}
\newcommand{\kfSecQGrammar}[1]{\kfSecHeader{kfAreaGrammar}{문제}{#1}}
\newcommand{\kfSecQConcept}[1]{\kfSecHeader{kfAreaConcept}{문제}{#1}}
\newcommand{\kfSecQLit}[1]{\kfSecHeader{kfAreaLiterature}{문제}{#1}}
\newcommand{\kfSecQNonlit}[1]{\kfSecHeader{kfAreaNonlit}{문제}{#1}}
\newcommand{\kfSecQWeekly}[1]{\kfSecHeader{kfAreaWeekly}{문제}{#1}}


\begin{document}

\thispagestyle{empty}
\vspace*{40mm}
\begin{center}
{\Huge\bfseries\color{kfPrimary} 비트겐슈타인3}\\[10pt]
{\Large\color{kfMute} 학생용 교재 — 제 1 권}\\[20pt]
{\small\color{kfMute} (챕터 1 \textendash{} 4)}
\end{center}
\vfill
\hfill\kfSeriesBadge\hspace{15mm}
\clearpage
\kfChapterCover{1}{비트겐슈타인3 ch01}{Chapter 1}{이번 챕터: 문법 → 문학 5작품 → 비문학 5지문 → 패턴 워크북.}

\kfStartGrammar{이번 챕터 문법 영역 — 음운 / 음운 체계 + 음성·음운 구분 + 최소 대립쌍}


\kfSecHeader{kfAreaGrammar}{문제}{}

\begin{multicols}{2}\raggedcolumns\setlength{\columnsep}{12pt}
\kfPassageFull{문항 1 지문 (concept)}{%
「국어의 음운은 말의 뜻을 구별해 주는 소리의 가장 작은 단위이다. 음운은 분절 음운과 비분절 음운으로 나뉜다. 분절 음운에는 자음과 모음이 있으며, 비분절 음운에는 소리의 길이(장단)가 있다. 자음은 조음 위치에 따라 입술소리(ㅂ, ㅃ, ㅍ, ㅁ), 잇몸소리(ㄷ, ㄸ, ㅌ, ㅅ, ㅆ, ㄴ, ㄹ), 센입천장소리(ㅈ, ㅉ, ㅊ), 여린입천장소리(ㄱ, ㄲ, ㅋ, ㅇ), 목청소리(ㅎ)로 분류된다. 조음 방법에 따라서는 파열음(ㅂ, ㅃ, ㅍ, ㄷ, ㄸ, ㅌ, ㄱ, ㄲ, ㅋ), 파찰음(ㅈ, ㅉ, ㅊ), 마찰음(ㅅ, ㅆ, ㅎ), 비음(ㅁ, ㄴ, ㅇ), 유음(ㄹ)으로 나뉜다. 한편 모음은 발음할 때 입술 모양이나 혀의 위치가 변하지 않는 단모음과, 반모음과 단모음이 결합하여 발음 도중 혀의 위치가 달라지는 이중 모음으로 구분된다.」
}
\begin{kfQBlock}{1}{}
\kfQStem{윗글을 바탕으로 <보기>의 탐구 활동을 수행한 결과로 적절하지 않은 것은?

<보기> [탐구 활동] 다음 단어 쌍에서 의미 차이를 만들어 내는 음운의 특성을 분석해 보자. ㉠ '달'과 '딸'  ㉡ '불'과 '풀'  ㉢ '갈'과 '걸' ㉣ '밤'(夜)과 '밤:'(栗)  ㉤ '살'과 '쌀'}
\begin{kfChoices}
\kfChoice ㉠에서 의미를 구별하는 것은 조음 방법의 차이이며, 'ㄷ'은 예사소리, 'ㄸ'은 된소리이다.
\kfChoice ㉡에서 'ㅂ'과 'ㅍ'은 조음 위치가 같고, 소리의 세기 차이로 뜻을 구별한다.
\kfChoice ㉢에서 'ㅏ'와 'ㅓ'는 혀의 높이는 같으나 혀의 전후 위치 차이로 뜻을 구별한다.
\kfChoice ㉣에서 분절 음운은 동일하나 비분절 음운인 소리의 길이가 다르다.
\kfChoice ㉤에서 'ㅅ'과 'ㅆ'은 조음 위치와 조음 방법이 같고, 소리의 세기가 다르다.
\end{kfChoices}
\end{kfQBlock}

\kfPassageFull{문항 2 지문 (concept)}{%
「음성(phone)은 사람의 발음 기관을 통해 만들어지는 물리적 소리 자체이고, 음운(phoneme)은 의미 구별에 관여하는 추상적 소리 단위이다. 하나의 음운이 환경에 따라 다르게 실현되는 구체적 소리를 변이음(allophone)이라 한다. 예를 들어 국어의 'ㄹ'은 모음 사이에서는 탄설음 [ɾ]로, 'ㄹㄹ'이 연속되거나 어말에서는 설측음 [l]로 실현된다. 이 둘은 물리적으로 다른 소리이지만 국어 화자에게는 같은 음운 'ㄹ'로 인식된다. 반면 영어에서는 [ɾ]와 [l]이 서로 다른 음운이다. 이처럼 같은 소리라도 언어에 따라 별개의 음운으로 기능하기도 하고, 하나의 음운의 변이음이 되기도 한다.」
}
\begin{kfQBlock}{2}{}
\kfQStem{윗글을 읽은 학생들이 <보기>와 같이 대화하였다. ㉠\textasciitilde{}㉤ 중 적절하지 않은 것은?

<보기> 학생 1: 'ㄹ'이 '바라'에서는 [ɾ]로, '달력[달ː력]'에서는 [l]로 소리 나니까, ㉠[ɾ]와 [l]은 국어에서 'ㄹ'의 변이음이야. 학생 2: 맞아. ㉡변이음은 같은 음운이 음성 환경에 따라 다르게 실현된 것이지. 학생 3: 그런데 '불'과 '물'에서 'ㅂ'과 'ㅁ'은 ㉢뜻을 구별해 주니까 서로 다른 음운이야. 학생 1: 그렇지. ㉣만약 두 소리가 변이음 관계라면, 하나를 다른 것으로 바꿔도 의미가 달라지겠지. 학생 2: 아, 그리고 ㉤'국물'에서 'ㄱ'이 [ŋ]으로 발음되는 것도 변이음의 예가 될 수 있어.}
\begin{kfChoices}
\kfChoice ㉠
\kfChoice ㉡
\kfChoice ㉢
\kfChoice ㉣
\kfChoice ㉤
\end{kfChoices}
\end{kfQBlock}

\kfPassageFull{문항 3 지문 (concept)}{%
「최소 대립쌍이란 하나의 음운만 다르고 나머지 음운은 동일한데 의미가 구별되는 두 단어를 말한다. 예를 들어 '달'과 '탈'은 초성 'ㄷ'과 'ㅌ'만 다른 최소 대립쌍이다. 최소 대립쌍을 통해 특정 소리가 해당 언어에서 독자적인 음운 자격을 갖는지를 확인할 수 있다. 한편 국어의 자음 체계에서 파열음·파찰음은 예사소리·된소리·거센소리의 3항 대립을 보인다. 'ㅂ-ㅃ-ㅍ', 'ㄷ-ㄸ-ㅌ', 'ㄱ-ㄲ-ㅋ', 'ㅈ-ㅉ-ㅊ'이 그 예이다. 그러나 마찰음은 'ㅅ-ㅆ'의 2항 대립만 보이고, 비음(ㅁ, ㄴ, ㅇ)과 유음(ㄹ)에서는 이러한 대립이 나타나지 않는다. 이처럼 음운 체계 내에서 특정 자질의 대립이 모든 부류에 균일하게 적용되는 것은 아니다.」
}
\begin{kfQBlock}{3}{}
\kfQStem{윗글을 바탕으로 <보기>의 ⓐ\textasciitilde{}ⓔ를 탐구한 내용으로 적절한 것은?

<보기> 다음 단어 쌍이 최소 대립쌍에 해당하는지 판단하고, 해당하면 대립하는 음운의 자질 차이를 분석하시오. ⓐ '불' — '뿔'     ⓑ '달' — '딸' ⓒ '손' — '존'     ⓓ '살' — '솔' ⓔ '곰' — '공'}
\begin{kfChoices}
\kfChoice ⓐ에서 'ㅂ'과 'ㅃ'은 조음 위치와 조음 방법이 다르므로 최소 대립쌍이 된다.
\kfChoice ⓑ에서 'ㄷ'과 'ㄸ'의 대립은 파열음에서 나타나는 3항 대립 중 일부이다.
\kfChoice ⓒ에서 'ㅅ'과 'ㅈ'은 조음 방법이 같으므로 소리의 세기 차이로 의미를 구별한다.
\kfChoice ⓓ에서 'ㅏ'와 'ㅗ'는 혀의 전후 위치 차이로 의미를 구별하는 최소 대립쌍이다.
\kfChoice ⓔ에서 'ㅁ'과 'ㅇ'은 조음 방법(비음)이 다르므로 최소 대립쌍이 된다.
\end{kfChoices}
\end{kfQBlock}

\kfPassageFull{문항 4 지문 (concept)}{%
「국어의 단모음은 혀의 높이, 혀의 전후 위치, 입술 모양의 세 기준으로 분류된다. 혀의 높이에 따라 고모음(ㅣ, ㅟ, ㅡ, ㅜ), 중모음(ㅔ, ㅚ, ㅓ, ㅗ), 저모음(ㅐ, ㅏ)으로 나뉘고, 혀의 전후 위치에 따라 전설 모음(ㅣ, ㅔ, ㅐ, ㅟ, ㅚ)과 후설 모음(ㅡ, ㅓ, ㅏ, ㅜ, ㅗ)으로 나뉜다. 입술 모양에 따라 원순 모음(ㅟ, ㅚ, ㅜ, ㅗ)과 평순 모음(ㅣ, ㅔ, ㅐ, ㅡ, ㅓ, ㅏ)으로 나뉜다. 그런데 현대 국어에서 'ㅔ'와 'ㅐ'의 발음 구별이 사라지는 경향이 두드러지고, 'ㅟ'와 'ㅚ'를 이중 모음으로 발음하는 것이 허용되면서 실질적으로 7모음 체계를 보이기도 한다. 이는 음운 체계가 고정불변이 아니라 시대에 따라 변화할 수 있음을 보여 준다.」
}
\begin{kfQBlock}{4}{}
\kfQStem{윗글을 바탕으로 <보기>를 탐구한 학생의 반응으로 적절하지 않은 것은?

<보기> [탐구 과제] 다음 모음 쌍의 공통점과 차이점을 분석하시오. (가) ㅣ — ㅡ    (나) ㅓ — ㅗ (다) ㅔ — ㅐ    (라) ㅜ — ㅗ}
\begin{kfChoices}
\kfChoice (가)의 'ㅣ'와 'ㅡ'는 혀의 높이가 같은 고모음이지만 혀의 전후 위치가 다르다.
\kfChoice (나)의 'ㅓ'와 'ㅗ'는 모두 후설 중모음이며, 입술 모양에서 차이가 있다.
\kfChoice (다)의 'ㅔ'와 'ㅐ'는 혀의 전후 위치와 입술 모양이 같고, 혀의 높이만 다르다.
\kfChoice (라)의 'ㅜ'와 'ㅗ'는 입술 모양이 같은 원순 모음이며, 혀의 전후 위치가 다르다.
\kfChoice (다)에서 'ㅔ'와 'ㅐ'의 구별이 사라지는 현상은 음운 체계가 변화할 수 있음을 보여 준다.
\end{kfChoices}
\end{kfQBlock}

\kfPassageFull{문항 5 지문 (concept)}{%
「국어의 음절은 자음과 모음의 결합으로 이루어지며, 'CGVC'(초성-반모음-중성-종성) 구조를 가질 수 있다. 초성과 종성에는 자음이 올 수 있고, 중성에는 모음이 반드시 와야 한다. 음운의 개수를 셀 때 주의할 점이 있다. 예를 들어 '닭'의 음운 개수는 초성 'ㄷ', 중성 'ㅏ', 종성 'ㄹ', 'ㄱ'으로 4개이다. 겹받침은 두 개의 자음 음운으로 센다. 반면 '가'의 초성 'ㄱ'은 자음 음운 1개이고, 'ㅏ'는 모음 음운 1개로, 총 2개이다. 이중 모음은 반모음과 단모음의 결합이므로 음운 2개로 센다. '왜'의 경우 반모음 [w], 단모음 'ㅐ'의 2개 모음 음운에 초성 없음(ㅇ은 음가 없음)이므로 음운은 2개이다.」
}
\begin{kfQBlock}{5}{}
\kfQStem{윗글을 바탕으로 <보기>의 단어들을 분석한 내용으로 적절한 것은?

<보기> ㄱ. 삶     ㄴ. 꿰     ㄷ. 읽다 ㄹ. 외     ㅁ. 흙}
\begin{kfChoices}
\kfChoice ㄱ의 음운 개수는 모두 합쳐 3개로 분석된다.
\kfChoice ㄴ의 음운 개수는 반모음을 포함하여 모두 4개이다.
\kfChoice ㄷ의 첫 음절 '읽'에 든 자음 음운의 개수는 2개이다.
\kfChoice ㄹ에 포함된 음운의 개수는 모두 2개이다.
\kfChoice ㅁ에 포함된 자음 음운의 개수는 모두 2개이다.
\end{kfChoices}
\end{kfQBlock}

\end{multicols}


\kfStartLit{이번 챕터 문학 작품 5편}


\kfSecHeader{kfAreaLiterature}{지문}{간 / 별 헤는 밤 — 윤동주 (현대시(복합지문)) / (가) 윤동주, 「간」(1941) / (나) 윤동주, 「별 헤는 밤」(1941). 둘 다 학평·모평 기출 다수.}

\kfPassageNote{간 / 별 헤는 밤 — 윤동주 (현대시(복합지문)) / (가) 윤동주, 「간」(1941) / (나) 윤동주, 「별 헤는 밤」(1941). 둘 다 학평·모평 기출 다수.}{%
(가) \\[4pt]
바닷가 햇빛 바른 바위 위에 \\[4pt]
습한 간을 펴서 말리우자, \\[14pt]
코카서스 산중에서 도망해 온 토끼처럼 \\[4pt]
둘러리를 빙빙 돌며 간을 지키자, \\[14pt]
내가 오래 기르던 여윈 독수리야! \\[4pt]
와서 뜯어 먹어라, 시름없이 \\[14pt]
너는 살찌고 \\[4pt]
나는 여위어야지, 그러나, \\[14pt]
거북이야! \\[4pt]
다시는 용궁의 유혹에 안 떨어진다. \\[14pt]
프로메테우스 불쌍한 프로메테우스 \\[4pt]
불 도적한 죄로 목에 맷돌을 달고 \\[4pt]
끝없이 침전하는 프로메테우스. \\[14pt]
― 윤동주, 「간」 ― \\[14pt]
(나) \\[4pt]
계절이 지나가는 하늘에는 \\[4pt]
가을로 가득 차 있습니다. \\[14pt]
나는 아무 걱정도 없이 \\[4pt]
가을 속의 별들을 다 헤일 듯합니다. \\[14pt]
가슴 속에 하나 둘 새겨지는 별을 \\[4pt]
이제 다 못 헤는 것은 \\[4pt]
쉬이 아침이 오는 까닭이요, \\[4pt]
내일 밤이 남은 까닭이요, \\[4pt]
아직 나의 청춘이 다하지 않은 까닭입니다. \\[14pt]
별 하나에 추억과 \\[4pt]
별 하나에 사랑과 \\[4pt]
별 하나에 쓸쓸함과 \\[4pt]
별 하나에 동경과 \\[4pt]
별 하나에 시와 \\[4pt]
별 하나에 어머니, 어머니, \\[14pt]
어머님, 나는 별 하나에 아름다운 말 한마디씩 불러 봅니다. 소학교 때 책상을 같이 했던 아이들의 이름과, 패, 경, 옥, 이런 이국 소녀들의 이름과, 벌써 아기 어머니 된 계집애들의 이름과, 가난한 이웃 사람들의 이름과, 비둘기, 강아지, 토끼, 노새, 노루, 「프란시스 잠」, 「라이너 마리아 릴케」, 이런 시인의 이름을 불러 봅니다. \\[14pt]
이네들은 너무나 멀리 있습니다. \\[4pt]
별이 아슬히 먼 듯이, \\[14pt]
어머니, \\[4pt]
그리고 당신은 멀리 북간도에 계십니다. \\[14pt]
나는 무엇인지 그리워 \\[4pt]
이 많은 별빛이 내린 언덕 위에 \\[4pt]
내 이름자를 써 보고 \\[4pt]
흙으로 덮어 버리었습니다. \\[14pt]
딴은 밤을 새워 우는 벌레는 \\[4pt]
부끄러운 이름을 슬퍼하는 까닭입니다. \\[14pt]
― 윤동주, 「별 헤는 밤」 ―
}
\kfSecHeader{kfAreaLiterature}{활동}{작품 정리표 — 간 / 별 헤는 밤(윤동주)}

\noindent\textit{\small 빈칸에 알맞은 말을 채워 넣으시오. (초성 힌트 참고)}\par\medskip
\begin{kfActTable}{|>{\centering\arraybackslash}m{2.6cm}|m{6cm}|m{4cm}|}
\rowcolor{kfPrimaryLight!50}\textbf{\color{kfPrimary}항목} & \textbf{\color{kfPrimary}내용 (빈칸 채우기)} & \textbf{\color{kfPrimary}답안 작성}\\\hline
(가) 갈래 & ㅈㅇㅅ, ㅅㅈㅅ ① & \kfTBlank \\\hline
(가) 성격 & ㅈㅎㅈ, ㅇㅈㅈ, ㅇㅎㅈ ② & \kfTBlank \\\hline
(가) 결합 모티프 & ㄱㅌㅈㅅ + ㅍㄹㅁㅌㅇㅅ ㅅㅎ ③ & \kfTBlank \\\hline
(가) ‘간’의 의미 & ㅇㅅ과 ㅅㄴ의 상징 ④ & \kfTBlank \\\hline
(가) 유혹 거부 & ‘다시는 ㅇㄱ의 ㅇㅎ에 안 떨어진다’ ⑤ & \kfTBlank \\\hline
(가) 자기희생 & 끝없이 ㅊㅈ하는 ㅍㄹㅁㅌㅇㅅ ⑥ & \kfTBlank \\\hline
(나) 표현 기법 & ‘별 하나에 \textasciitilde{}과’의 ㅇㄱㅂ ⑦ & \kfTBlank \\\hline
(나) 별의 의미 & 추억, 사랑, ㅆㅆㅎ, 동경, 시, ㅇㅁㄴ ⑧ & \kfTBlank \\\hline
(나) 자아 행위 & ‘내 ㅇㄹㅈ를 써 보고 / ㅎ으로 덮어 버리었습니다’ ⑨ & \kfTBlank \\\hline
(나) 마지막 정서 & ‘ㅂㄲㄹㅇ ㅇㄹㅇ ㅅㅍㅎㄴ 까닭’ ⑩ & \kfTBlank \\\hline
공통 배경 & ㅇㅈ ㄱㅈㄱ(1941년) ⑪ & \kfTBlank \\\hline
공통 주제 & ㅈㅇ ㅂㅇ과 ㅈㄱㅎㅅ·ㅈㅇ ㅅㅊ ⑫ & \kfTBlank \\\hline
\end{kfActTable}
\kfSecHeader{kfAreaLiterature}{문제}{}

\begin{multicols}{2}\raggedcolumns\setlength{\columnsep}{12pt}
\begin{kfQBlock}{1}{}
\kfQStem{(가)와 (나)에 대한 설명으로 가장 적절한 것은?}
\begin{kfChoices}
\kfChoice (가)와 (나) 모두 마지막 연을 명사로 끝맺어 시적 여운을 남긴다.
\kfChoice (가)는 시간적 표지를 활용해 시상을 전환하지만, (나)는 그러한 표지가 나타나지 않는다.
\kfChoice (가)와 (나) 모두 의문의 방식을 활용하여 시적 의미를 강조한다.
\kfChoice (가)와 (나) 모두 특정한 대상을 부르는 방식을 사용하여 그 대상에 주목하게 한다.
\kfChoice (가)와 (나) 모두 공감각적 이미지를 활용하여 다양한 사물의 역동성을 부각한다.
\end{kfChoices}
\end{kfQBlock}

\begin{kfQBlock}{2}{}
\kfQStem{(가)의 시어·시구에 대한 이해로 적절하지 않은 것은?}
\begin{kfChoices}
\kfChoice ‘습한 간을 펴서 말리우자’에서 ‘말리는’ 행위는 부정적 현실을 극복하고 양심을 회복하려는 의지를 드러낸다.
\kfChoice ‘둘러리를 빙빙 돌며 간을 지키자’는 위기 상황에서도 자신의 양심과 신념을 지키려는 의지를 드러낸다.
\kfChoice ‘여윈 독수리야! / 와서 뜯어 먹어라, 시름없이’는 자기 신념을 지키기 위해 겪어야 할 내적 고통을 감내하겠다는 태도를 드러낸다.
\kfChoice ‘다시는 용궁의 유혹에 안 떨어진다’는 부귀영화나 현실적 안락의 유혹을 거부하겠다는 태도를 드러낸다.
\kfChoice ‘끝없이 침전하는 프로메테우스’는 자기희생을 거부하고 안정된 일상으로 복귀하려는 의지를 드러낸다.
\end{kfChoices}
\end{kfQBlock}

\begin{kfQBlock}{3}{}
\kfQStem{(나)의 ‘내 이름자를 써 보고 / 흙으로 덮어 버리었습니다’가 지닌 함축적 의미로 가장 적절한 것은?}
\begin{kfChoices}
\kfChoice 자기 이름에 대한 자부심과 긍지를 표현한 것이다.
\kfChoice 자기 존재에 대한 부끄러움과 윤리적 성찰의 태도를 형상화한 것이다.
\kfChoice 타인의 이름을 잊으려는 의지를 나타낸 것이다.
\kfChoice 농사일에 몰두하여 자아를 잊는 일상을 표현한 것이다.
\kfChoice 죽음에 대한 근원적 공포를 극복한 초월적 경지를 나타낸 것이다.
\end{kfChoices}
\end{kfQBlock}

\begin{kfQBlock}{4}{}
\kfQStem{<보기>를 바탕으로 (가)와 (나)를 감상한 내용으로 적절하지 않은 것은? [3점]

<보기> (가) 「간」은 토끼가 유혹에 빠져 위기에 처했다가 지혜를 발휘해 간을 지킨 ‘구토지설’과, 인간에게 불을 알려 준 죄로 코카서스 산에서 독수리에게 간을 쪼이는 ‘프로메테우스 신화’를 결합한 작품이다. 화자는 두 설화를 재구성하여 간을 지키려는 토끼의 노력과 프로메테우스의 희생을 연결하고, 일제 강점기에서 세속적 욕망을 추구하지 않고 양심을 지키려는 자기희생의 의지를 드러낸다. (나) 「별 헤는 밤」은 가을밤 별을 매개로 내면의 가치를 하나하나 호명하면서 자아를 성찰하는 작품이다.}
\begin{kfChoices}
\kfChoice (가)의 ‘코카서스 산중에서 도망해 온 토끼’는 두 설화를 연결한 표현으로, 일제 강점기에서 양심을 지키려는 존재로 볼 수 있다.
\kfChoice (가)의 ‘다시는 용궁의 유혹에 안 떨어진다’는 구토지설을 재구성한 것으로, 세속적 욕망을 추구하지 않겠다는 의지를 드러낸다.
\kfChoice (가)의 ‘목에 맷돌을 달고’는 프로메테우스가 벌을 받았다는 신화를 재구성한 것으로, 화자가 감수하고자 하는 자기희생을 상징한다.
\kfChoice (나)의 ‘별 하나에 추억과 / 별 하나에 사랑과’ 등의 열거는 별을 매개로 내면의 가치를 호명하며 자아를 성찰하는 양상으로 볼 수 있다.
\kfChoice (나)의 ‘부끄러운 이름을 슬퍼하는 까닭’은 화자가 이미 자아 성찰을 완전히 끝내고 모든 슬픔에서 벗어났음을 보여 준다.
\end{kfChoices}
\end{kfQBlock}

\begin{kfQBlock}{5}{}
\kfQStem{(가)와 (나)를 비교하여 감상한 내용으로 가장 적절한 것은?}
\begin{kfChoices}
\kfChoice (가)는 두 설화를 우의적으로 결합해 자기희생의 의지를, (나)는 별을 매개로 한 호명을 통해 자아 성찰을 드러낸다.
\kfChoice (가)와 (나) 모두 마지막 연에서 일제에 대한 직접적 분노를 사실적으로 토로한다.
\kfChoice (가)는 신화적 모티프를 거부하며 일상적 감각만을 담백하게 노래하고, (나)는 오히려 신화에 의지하여 자아를 형상화한다.
\kfChoice (가)와 (나) 모두 자연 풍경 묘사에 치중하여 화자의 내면을 거의 드러내지 않는다.
\kfChoice (가)와 (나) 모두 화자가 이미 자기 정체성에 대한 갈등을 완전히 해소한 평온의 상태를 그린다.
\end{kfChoices}
\end{kfQBlock}

\end{multicols}

\kfSecHeader{kfAreaLiterature}{지문}{꽃피는 시절 — 이성복 (현대시) / 이성복, 『뒹구는 돌은 언제 잠 깨는가』(1980)}

\kfPassageNote{꽃피는 시절 — 이성복 (현대시) / 이성복, 『뒹구는 돌은 언제 잠 깨는가』(1980)}{%
봄날은 간다 \\[4pt]
가랑잎 날리던 골목에서 \\[4pt]
네가 떠나 보낸 그 봄날은 간다 \\[4pt]
날은 저문다 \\[4pt]
낡은 철도변에 꽃이 피고 \\[4pt]
기적 소리가 울고 \\[4pt]
먼 데서 실바람이 일고 \\[4pt]
날은 저문다 \\[4pt]
가만히 안개가 살갗에 젖으면 \\[4pt]
모든 것은 흐려지고 \\[4pt]
더이상 소중한 것은 없다 \\[4pt]
다만 꽃피는 시절의 어두운 그늘에 \\[4pt]
눈물짓던 가난한 영혼들을 위해 \\[4pt]
우리는 마지막 한 방울의 눈물도 \\[4pt]
나누어 가져야 했던 것이다
}
\kfSecHeader{kfAreaLiterature}{활동}{작품 정리표 — 꽃피는 시절(이성복)}

\noindent\textit{\small 빈칸에 알맞은 말을 채워 넣으시오. (초성 힌트 참고)}\par\medskip
\begin{kfActTable}{|>{\centering\arraybackslash}m{2.6cm}|m{6cm}|m{4cm}|}
\rowcolor{kfPrimaryLight!50}\textbf{\color{kfPrimary}항목} & \textbf{\color{kfPrimary}내용 (빈칸 채우기)} & \textbf{\color{kfPrimary}답안 작성}\\\hline
갈래 & ㅈㅇㅅ, ㅅㅈㅅ ① & \kfTBlank \\\hline
성격 & ㅇㅅㅈ, ㅅㅊㅈ ② & \kfTBlank \\\hline
제목의 역설 & 아름다운 시절 ↔ ㄱㅌ과 ㅅㅅ ③ & \kfTBlank \\\hline
반복 1 & 'ㅂㄴㅇ ㄱㄷ' — 시간의 경과와 ㅅㅅㄱ ④ & \kfTBlank \\\hline
반복 2 & 'ㄴㅇ ㅈㅁㄷ' — ㅅㄹ의 정서 ⑤ & \kfTBlank \\\hline
자연 이미지 & 꽃, 기적 소리, ㅅㅂㄹ — 동시다발적 ㅂㅎ ⑥ & \kfTBlank \\\hline
은폐 이미지 & ㅇㄱ가 살갗에 젖으면 모든 것은 흐려지고 ⑦ & \kfTBlank \\\hline
절망 표현 & '더이상 ㅅㅈㅎ 것은 없다' ⑧ & \kfTBlank \\\hline
핵심 이미지 & '꽃피는 시절의 ㅇㄷㅇ ㄱㄴ' — 역설적 병치 ⑨ & \kfTBlank \\\hline
연대의 윤리 & '마지막 한 방울의 ㄴㅁ도 / ㄴㄴㅇ 가져야 했던' ⑩ & \kfTBlank \\\hline
주제 1 & ㅈㅈㅇ ㅇㅁㅅ — 아름다움과 고통의 공존 ⑪ & \kfTBlank \\\hline
시대 배경 & ㅊㄱㄷ 1980년대 — 시대적 아픔 ⑫ & \kfTBlank \\\hline
\end{kfActTable}
\kfSecHeader{kfAreaLiterature}{문제}{}

\begin{multicols}{2}\raggedcolumns\setlength{\columnsep}{12pt}
\begin{kfQBlock}{1}{}
\kfQStem{윗시에 대한 설명으로 적절하지 않은 것은?}
\begin{kfChoices}
\kfChoice '봄날은 간다'와 '날은 저문다'의 반복을 통해 시간의 경과와 상실감을 강조하고 있다.
\kfChoice '가랑잎 날리던 골목'은 과거의 기억이 깃든 공간을 환기한다.
\kfChoice '꽃이 피고', '실바람이 일고' 등 자연 이미지가 화자의 내면 정서를 반영하고 있다.
\kfChoice '더이상 소중한 것은 없다'에서 화자는 완전한 행복과 충만함을 느끼고 있다.
\kfChoice '눈물도 / 나누어 가져야 했던 것이다'에서 공동체적 연대 의식이 드러나고 있다.
\end{kfChoices}
\end{kfQBlock}

\begin{kfQBlock}{2}{}
\kfQStem{윗시의 '꽃피는 시절의 어두운 그늘'이 지닌 함축적 의미로 가장 적절한 것은?}
\begin{kfChoices}
\kfChoice 아름다운 시절 이면에 자리한 고통과 슬픔
\kfChoice 꽃이 필 때 드리워진 물리적 그림자의 형상
\kfChoice 지나간 봄날에 함께한 즐거운 추억의 잔상
\kfChoice 자연의 순환을 객관적으로 관찰한 결과
\kfChoice 경제적 성공의 이면에 드리운 그늘진 면모
\end{kfChoices}
\end{kfQBlock}

\begin{kfQBlock}{3}{}
\kfQStem{윗시의 마지막 3행에 대한 해석으로 가장 적절한 것은?}
\begin{kfChoices}
\kfChoice 고통받는 이들에 대한 공감과 연대의 윤리를 당위적으로 제시하고 있다.
\kfChoice 화자가 개인적 이익을 추구해야 한다는 자기중심적 태도를 드러내고 있다.
\kfChoice 눈물의 물리적 양을 분배하는 실용적 제안을 하고 있다.
\kfChoice 과거와의 완전한 단절을 선언하고 있다.
\kfChoice 자연의 아름다움에 대한 찬탄으로 시를 마무리하고 있다.
\end{kfChoices}
\end{kfQBlock}

\begin{kfQBlock}{4}{}
\kfQStem{윗시를 바탕으로 <보기>를 감상한 내용으로 적절하지 않은 것은?

<보기> 이성복의 시는 1980년대 한국 사회의 아픔을 배경으로 한다. 그의 시에서 아름다운 자연 이미지는 종종 폭력과 고통의 이면을 감추는 기제로 작동하며, '꽃'과 '봄'은 순수한 아름다움이 아니라 그 아래 숨겨진 상처를 드러내는 역설적 상징이 된다.}
\begin{kfChoices}
\kfChoice '낡은 철도변에 꽃이 피고'에서 '낡은 철도변'은 훼손된 현실을, '꽃이 피고'는 그 위에 덧씌워진 아름다움의 허상을 암시하겠군.
\kfChoice '기적 소리가 울고'의 '울고'는 기차 소리에 울음의 이미지를 겹쳐 놓아, 아름다운 봄 풍경 속의 슬픔을 드러내겠군.
\kfChoice '안개가 살갗에 젖으면 / 모든 것은 흐려지고'에서 안개는 현실을 은폐하는 기제로 작동하겠군.
\kfChoice '가난한 영혼들'은 시대적 고통 속에서 상처받은 이들을 가리키며, <보기>의 '숨겨진 상처'와 연결되겠군.
\kfChoice '봄날은 간다'는 1980년대의 민주화 운동이 거둔 결실을 기뻐하는 환호의 표현이겠군.
\end{kfChoices}
\end{kfQBlock}

\begin{kfQBlock}{5}{}
\kfQStem{윗시에서 '∼고'라는 연결어미가 연속적으로 사용되는 부분의 표현 효과로 가장 적절한 것은?}
\begin{kfChoices}
\kfChoice 과거와 현재를 뚜렷이 구분 지어 시의 시간적 배경을 확정한다.
\kfChoice 사건 간 논리적 인과 관계를 명확히 제시하여 독자의 이해를 돕는다.
\kfChoice 화자의 기쁨이 점차 고조되는 정서적 흐름을 단계적으로 나타낸다.
\kfChoice 대화체 어조를 활용하여 독자에게 직접 말을 거는 효과를 부여한다.
\kfChoice 자연 이미지를 나란히 늘어놓아 세계 쇠락의 동시성을 형상화한다.
\end{kfChoices}
\end{kfQBlock}

\end{multicols}

\kfSecHeader{kfAreaLiterature}{지문}{겨울 뜸부기 — 오정희 (현대소설) / 오정희, 「겨울 뜸부기」(1979), 핵심 발췌}

\kfPassageNote{겨울 뜸부기 — 오정희 (현대소설) / 오정희, 「겨울 뜸부기」(1979), 핵심 발췌}{%
어머니가 밤새 기침을 했다. 아침에 부엌으로 나가 보니 찬물에 손을 담그고 빨래를 하고 있었다. 손등이 벌겋게 터져 있었다. 나는 말없이 앉아서 어머니가 빨래를 비비는 것을 바라보았다. 마당에는 눈이 내리고 있었다. 처마 끝에 고드름이 달려 있었다.

학교에서 돌아오면 어머니는 없었다. 삼촌이 와서 나를 데리고 국수집에 갔다. 삼촌은 국수를 먹으며 아무 말도 하지 않았다. 나는 국수 위에 떠 있는 기름을 젓가락으로 휘저었다. 밖에서 뜸부기가 울었다. 겨울인데 뜸부기가 울다니. 삼촌이 그건 뜸부기가 아니라 비둘기라고 했다. 나는 뜸부기라고 우겼다.

아버지가 다녀간 것은 봄이었다. 아버지는 내 머리를 쓰다듬으며 금방 돌아올 거라고 했다. 그 뒤로 아버지는 오지 않았다. 어머니는 아버지 이야기를 하지 않았다. 나는 아버지의 얼굴이 점점 흐려지는 것을 느꼈다. 마치 안개 속으로 걸어 들어가는 사람처럼.

밤에 이불 속에서 나는 뜸부기 소리를 들었다. 겨울에는 뜸부기가 울지 않는다고 했다. 하지만 나는 분명히 들었다. 어쩌면 그것은 어머니의 기침 소리였는지도 모른다. 아니면 바람 소리. 아니면 내 안에서 나는 소리. 나는 이불을 뒤집어쓰고 눈을 감았다. 뜸부기는 계속 울고 있었다.
}
\kfSecHeader{kfAreaLiterature}{활동}{작품 정리표 — 겨울 뜸부기(오정희)}

\noindent\textit{\small 빈칸에 알맞은 말을 채워 넣으시오. (초성 힌트 참고)}\par\medskip
\begin{kfActTable}{|>{\centering\arraybackslash}m{2.6cm}|m{6cm}|m{4cm}|}
\rowcolor{kfPrimaryLight!50}\textbf{\color{kfPrimary}항목} & \textbf{\color{kfPrimary}내용 (빈칸 채우기)} & \textbf{\color{kfPrimary}답안 작성}\\\hline
갈래 & 현대 ㄷㅍ ㅅㅅ ① & \kfTBlank \\\hline
시점 & ㅇㄹ 화자의 1인칭 ㅈㅎ 시점 ② & \kfTBlank \\\hline
제목 상징 & 겨울에 울지 않는 ㄸㅂㄱ — ㅂㅈ와 ㄱㅈ의 역설 ③ & \kfTBlank \\\hline
가족 상황 & 아버지의 ㅂㅈ, 어머니의 ㄱㄷ ④ & \kfTBlank \\\hline
감각 이미지 & 눈, ㄱㄷㄹ, ㅊㅁ — 겨울의 차가움 ⑤ & \kfTBlank \\\hline
비유 & 아버지 얼굴의 퇴색 = 'ㅇㄱ 속으로 걸어 들어가는 사람' ⑥ & \kfTBlank \\\hline
인식 차이 & 삼촌: ㅂㄷㄱ / 화자: ㄸㅂㄱ ⑦ & \kfTBlank \\\hline
소리의 모호화 & 뜸부기 → 어머니 ㄱㅊ → ㅂㄹ → 내 안의 소리 ⑧ & \kfTBlank \\\hline
서술 기법 & 감각적 ㅇㅅ의 나열 → 독자가 ㅇㅁ를 구성 ⑨ & \kfTBlank \\\hline
주제 1 & ㄱㅈ ㅎㅊ와 ㅅㅈ ㅂㅇ ⑩ & \kfTBlank \\\hline
주제 2 & ㅂㅈ 속에서도 존재하려는 ㅅ의 의지 ⑪ & \kfTBlank \\\hline
작가 특징 & 오정희 특유의 ㄱㄱㅈ ㅁㅊ ⑫ & \kfTBlank \\\hline
\end{kfActTable}
\kfSecHeader{kfAreaLiterature}{문제}{}

\begin{multicols}{2}\raggedcolumns\setlength{\columnsep}{12pt}
\begin{kfQBlock}{1}{}
\kfQStem{윗글에 대한 설명으로 적절하지 않은 것은?}
\begin{kfChoices}
\kfChoice 어린 화자의 제한된 시점을 통해 사건의 의미가 간접적으로 전달되고 있다.
\kfChoice '눈', '고드름', '찬물' 등 겨울의 감각적 이미지가 배경의 분위기를 형성하고 있다.
\kfChoice 삼촌과의 대화에서 '뜸부기'와 '비둘기'의 인식 차이가 나타나고 있다.
\kfChoice 아버지의 부재가 화자에게 미치는 영향이 '안개 속으로 걸어 들어가는 사람'이라는 비유로 드러나고 있다.
\kfChoice 화자는 가족의 상황을 객관적으로 파악하여 논리적인 어조로 분석하고 있다.
\end{kfChoices}
\end{kfQBlock}

\begin{kfQBlock}{2}{}
\kfQStem{윗글의 '겨울에 우는 뜸부기'가 상징하는 바로 가장 적절한 것은?}
\begin{kfChoices}
\kfChoice 여름새인 뜸부기가 겨울에 자연의 질서에 순응하여 조화롭게 공존하는 모습
\kfChoice 있어야 할 때 없고 없어야 할 때 들리는 존재의 불안정성과 내면 부재의 환청
\kfChoice 겨울 풍경 속 새의 생태에 대한 화자의 객관적 관찰과 학문적 호기심
\kfChoice 삼촌과 화자가 새 이름을 두고 교감하며 깊어지는 가족적 유대감의 표상
\kfChoice 겨울이 끝나고 봄이 올 것이라는 화자의 낙관적 기대와 미래 지향의 희망
\end{kfChoices}
\end{kfQBlock}

\begin{kfQBlock}{3}{}
\kfQStem{윗글의 '아버지의 얼굴이 점점 흐려지는 것을 느꼈다. 마치 안개 속으로 걸어 들어가는 사람처럼'에서 화자가 느끼는 것으로 가장 적절한 것은?}
\begin{kfChoices}
\kfChoice 기억이 시간 속에서 점차 희미해져 가는 상실의 감각
\kfChoice 안개가 두텁게 깔린 날씨 변화에 대한 객관적 관찰
\kfChoice 아버지가 머지않아 다시 돌아올 것이라는 굳건한 확신
\kfChoice 아버지와의 관계가 더욱 돈독해진다고 느끼는 기쁨
\kfChoice 아버지를 의식 속에서 떨쳐 내겠다는 의지적 결단
\end{kfChoices}
\end{kfQBlock}

\begin{kfQBlock}{4}{}
\kfQStem{윗글을 바탕으로 <보기>를 감상한 내용으로 적절하지 않은 것은?

<보기> 오정희 소설에서 어린 화자의 시선은 성인 세계의 의미를 직접 파악하지 못한 채, 감각적 인상의 나열로 사건을 전달한다. 독자는 화자의 서술 사이의 '빈틈'에서 가족의 해체, 빈곤, 부재 등의 주제를 스스로 구성해야 한다.}
\begin{kfChoices}
\kfChoice 화자가 '국수 위에 떠 있는 기름을 젓가락으로 휘저'은 것은 삼촌과의 침묵 속 불편함을 감각적으로 전달하는 장면이겠군.
\kfChoice 화자가 '뜸부기라고 우겼다'는 것은 성인의 합리적 판단을 수용하지 못하는 어린 화자의 고집이자, 부재하는 것을 있다고 느끼는 내면의 표현이겠군.
\kfChoice '어머니는 아버지 이야기를 하지 않았다'는 서술은 화자가 직접 말하지 않는 가족 해체의 빈틈을 독자가 구성하게 하겠군.
\kfChoice 화자는 가족의 해체와 빈곤의 원인을 논리적으로 분석하여 독자에게 직접 설명하고 있겠군.
\kfChoice '손등이 벌겋게 터져 있었다'는 감각적 인상의 나열로, 어머니의 고단한 삶을 간접적으로 전달하겠군.
\end{kfChoices}
\end{kfQBlock}

\begin{kfQBlock}{5}{}
\kfQStem{윗글에서 '뜸부기 소리'의 정체가 점점 모호해지는 서술 방식의 효과로 가장 적절한 것은?}
\begin{kfChoices}
\kfChoice 뜸부기의 생태에 관한 과학적 정보를 자세히 제시하여 독자의 지식을 확장한다.
\kfChoice 소리의 출처가 어머니의 기침·바람·내면으로 흩어지며 존재의 불안정성을 심화한다.
\kfChoice 사건의 결말을 단정적으로 제시하여 독자의 궁금증을 한 번에 해소해 준다.
\kfChoice 화자의 청각이 매우 예민함을 부각하여 인물의 감각적 우위를 강조한다.
\kfChoice 삼촌의 판단이 결국 옳았음을 최종적으로 확인하여 인물의 권위를 부여한다.
\end{kfChoices}
\end{kfQBlock}

\end{multicols}

\kfSecHeader{kfAreaLiterature}{지문}{이생규장전 — 김시습 (고전소설) / 김시습, 「이생규장전」, 『금오신화』(15세기), 핵심 발췌, 현대어 역}

\kfPassageNote{이생규장전 — 김시습 (고전소설) / 김시습, 「이생규장전」, 『금오신화』(15세기), 핵심 발췌, 현대어 역}{%
「이제 저는 인연이 다하였으므로 다시 돌아갈 때가 되었습니다. 낭군께서는 부디 다른 어진 짝을 구하시어 대를 잇는 일을 그르치지 마소서.」

이생이 울며 말하였다.

「내 차라리 그대와 함께 황천으로 가련다. 인간 세상에 홀로 남아 무엇을 하리오.」

최낭이 말하였다.

「낭군의 수명은 아직 남아 있으니, 어찌 억지로 죽음을 구하시렵니까. 인연은 끊어졌어도 정은 다하지 아니하였으니, 만약 저를 잊지 아니하신다면 제 유골을 거두어 장사지내 주소서. 그것이 저의 마지막 소원입니다.」

말을 마치니 최낭의 모습이 차차 흐려지더니 마침내 사라져 보이지 아니하였다. 이생이 슬피 울다가 그 말대로 최낭의 유골을 찾아 장사 지내 주었다. 이생도 그 뒤 세상일에 뜻을 끊고, 이태 만에 세상을 떠났다.
}
\kfSecHeader{kfAreaLiterature}{활동}{작품 정리표 — 이생규장전(김시습)}

\noindent\textit{\small 빈칸에 알맞은 말을 채워 넣으시오. (초성 힌트 참고)}\par\medskip
\begin{kfActTable}{|>{\centering\arraybackslash}m{2.6cm}|m{6cm}|m{4cm}|}
\rowcolor{kfPrimaryLight!50}\textbf{\color{kfPrimary}항목} & \textbf{\color{kfPrimary}내용 (빈칸 채우기)} & \textbf{\color{kfPrimary}답안 작성}\\\hline
갈래 & 한문 ㅈㄱ ㅅㅅ ① & \kfTBlank \\\hline
출전 & 『ㄱㅇㅅㅎ(金鰲新話)』 ② & \kfTBlank \\\hline
작자 & ㄱㅅㅅ — 세조 찬탈에 저항한 지식인 ③ & \kfTBlank \\\hline
제목 의미 & 이생이 ㄷㅈ을 엿보다 → 만남의 계기 ④ & \kfTBlank \\\hline
주요 인물 & ㅇㅅ(남주) + ㅊㄴ(여주) ⑤ & \kfTBlank \\\hline
비극 원인 & ㅎㄱ(전란)으로 최낭 사망 ⑥ & \kfTBlank \\\hline
초월적 사건 & 최낭의 ㅁㄹ과 재회 → ㅅㅅ 초월의 사랑 ⑦ & \kfTBlank \\\hline
핵심 대사 & 'ㅇㅇ은 끊어졌어도 ㅈ은 다하지 아니하였으니' ⑧ & \kfTBlank \\\hline
이생의 행위 & 최낭의 ㅇㄱ을 찾아 ㅈㅅ 지냄 ⑨ & \kfTBlank \\\hline
이생의 결말 & 세상일에 ㄸ을 끊고 ㅇㅌ 만에 사망 ⑩ & \kfTBlank \\\hline
주제 & ㅅㅅ ㅊㅇ의 사랑 + 존재의 ㅇㅎㅅ ⑪ & \kfTBlank \\\hline
시대 배경 & 조선 초기 — ㅅㅈ 찬탈기의 비극적 세계 인식 ⑫ & \kfTBlank \\\hline
\end{kfActTable}
\kfSecHeader{kfAreaLiterature}{문제}{}

\begin{multicols}{2}\raggedcolumns\setlength{\columnsep}{12pt}
\begin{kfQBlock}{1}{}
\kfQStem{윗글에 대한 설명으로 적절하지 않은 것은?}
\begin{kfChoices}
\kfChoice 최낭은 자신의 인연이 다했음을 인정하면서도 이생에 대한 정이 남아 있음을 말하고 있다.
\kfChoice 이생은 최낭과 함께 죽고자 하는 강렬한 심정을 드러내고 있다.
\kfChoice 최낭의 모습이 '차차 흐려지더니 마침내 사라져 보이지 아니하였다'는 서술에서 초월적 존재의 소멸이 나타난다.
\kfChoice 이생은 최낭의 유언을 따르지 않고 즉시 자결하여 최낭과 함께 황천으로 갔다.
\kfChoice 이생이 '세상일에 뜻을 끊고' 세상을 떠난 것은 최낭을 잃은 상실감의 귀결로 볼 수 있다.
\end{kfChoices}
\end{kfQBlock}

\begin{kfQBlock}{2}{}
\kfQStem{윗글에서 '인연은 끊어졌어도 정은 다하지 아니하였으니'라는 최낭의 말이 함축하는 바로 가장 적절한 것은?}
\begin{kfChoices}
\kfChoice 이생과의 관계를 완전히 청산하고 새로 출발하겠다는 단호한 의지
\kfChoice 현세적 인연은 끊어져도 정서적 유대는 그 경계를 넘는다는 인식
\kfChoice 이생에게 빨리 재혼하여 자신을 잊으라고 권하는 냉정한 충고
\kfChoice 인연과 정은 한 몸이므로 둘 다 함께 끝난다는 체념적 수용
\kfChoice 최낭이 머지않아 환생하여 이생과 다시 만날 것이라는 예언
\end{kfChoices}
\end{kfQBlock}

\begin{kfQBlock}{3}{}
\kfQStem{윗글의 결말에서 '이태 만에 세상을 떠났다'가 작품 전체의 주제 의식과 관련하여 의미하는 바로 가장 적절한 것은?}
\begin{kfChoices}
\kfChoice 생사를 초월한 사랑의 절대성이 결국 이생의 삶 자체를 소진시켰음을 보여준다.
\kfChoice 이생이 건강한 삶을 유지하며 행복한 노년을 보냈음을 보여준다.
\kfChoice 최낭의 유언이 이생에게 전혀 영향을 미치지 못했음을 드러낸다.
\kfChoice 이생이 세상에 적극적으로 참여하며 새로운 삶을 시작했음을 나타낸다.
\kfChoice 전기 소설의 관습에 따라 주인공이 반드시 행복한 결말을 맞이해야 함을 보여준다.
\end{kfChoices}
\end{kfQBlock}

\begin{kfQBlock}{4}{}
\kfQStem{윗글을 바탕으로 <보기>를 감상한 내용으로 적절하지 않은 것은?

<보기> 『금오신화』의 작자 김시습은 수양대군의 왕위 찬탈에 저항하여 출가·방랑하며 불교·유교·도교를 넘나든 지식인이었다. 그의 소설에는 현실 세계의 질서가 파괴된 상황에서 초월적 존재와의 만남을 통해 인간적 가치를 탐색하려는 의식이 담겨 있다.}
\begin{kfChoices}
\kfChoice 최낭의 죽음과 망령으로의 재회는 현실 질서가 파괴된 상황에서 초월적 만남이 이루어지는 서사로 볼 수 있겠군.
\kfChoice 이생이 최낭의 유골을 거두는 행위는 죽은 자에 대한 인간적 예의이자 사랑의 실천으로 해석할 수 있겠군.
\kfChoice '인연은 끊어졌어도 정은 다하지 아니하였으니'라는 최낭의 말에서 인간적 가치(정)의 탐색이 드러나겠군.
\kfChoice 김시습의 방랑 경험과 무관하게, 이 작품은 순수한 오락 목적으로만 창작된 것이겠군.
\kfChoice 이생이 세상일에 뜻을 끊은 것은 현실 질서에 대한 회의와 상실감의 표현으로 읽을 수 있겠군.
\end{kfChoices}
\end{kfQBlock}

\begin{kfQBlock}{5}{}
\kfQStem{윗글에서 최낭의 모습이 '차차 흐려지더니 마침내 사라져 보이지 아니하였다'라는 서술의 표현 효과로 가장 적절한 것은?}
\begin{kfChoices}
\kfChoice 사실적 묘사를 통해 자연 현상의 과학적 원리를 설명한다.
\kfChoice 갑작스러운 소멸로 독자에게 충격을 주어 극적 반전을 만든다.
\kfChoice 최낭이 다른 장소로 이동했음을 암시하여 재회의 가능성을 열어 둔다.
\kfChoice 점진적 소멸의 묘사를 통해 이별의 불가피성과 슬픔을 서서히 고조시킨다.
\kfChoice 이생의 시력이 나빠지고 있음을 보여주는 복선이다.
\end{kfChoices}
\end{kfQBlock}

\end{multicols}

\kfSecHeader{kfAreaLiterature}{지문}{자화상 / 고운 심장 — 윤동주 / 신석정 (현대시(복합지문)) / (가) 윤동주, 「자화상」(1939) / (나) 신석정, 「고운 심장」(1933)}

\kfPassageNote{자화상 / 고운 심장 — 윤동주 / 신석정 (현대시(복합지문)) / (가) 윤동주, 「자화상」(1939) / (나) 신석정, 「고운 심장」(1933)}{%
(가) \\[4pt]
산모통이를 돌아 논가 외딴 우물을 홀로 찾아가선 \\[4pt]
가만히 들여다봅니다. \\[14pt]
우물 속에는 달이 밝고 구름이 흐르고 하늘이 펼치고 \\[4pt]
파아란 바람이 불고 가을이 있습니다. \\[14pt]
그리고 한 사나이가 있습니다. \\[4pt]
어쩐지 그 사나이가 미워져 돌아갑니다. \\[14pt]
돌아가다 생각하니 그 사나이가 가엾어집니다. \\[4pt]
도로 가 들여다보니 사나이는 그대로 있습니다. \\[14pt]
다시 그 사나이가 미워져 돌아갑니다. \\[4pt]
돌아가다 생각하니 그 사나이가 그리워집니다. \\[14pt]
우물 속에는 달이 밝고 구름이 흐르고 하늘이 펼치고 \\[4pt]
파아란 바람이 불고 가을이 있고 \\[4pt]
추억처럼 사나이가 있습니다. \\[14pt]
(나) \\[4pt]
꽃이 져도 그만이오 \\[4pt]
풀이 져도 그만이오 \\[14pt]
바람에 고운 심장이 \\[4pt]
울어도 그만이오 \\[14pt]
푸른 산이 날더러 \\[4pt]
너 푸르르다 하거든 \\[14pt]
맑은 물이 날더러 \\[4pt]
너 맑으다 하거든 \\[14pt]
혼자서 일어나 \\[4pt]
하늘을 우러르며 \\[14pt]
살포시 웃어 볼까 \\[4pt]
살포시 웃어 볼까
}
\kfSecHeader{kfAreaLiterature}{활동}{작품 정리표 — 자화상(윤동주) / 고운 심장(신석정)}

\noindent\textit{\small 빈칸에 알맞은 말을 채워 넣으시오. (초성 힌트 참고)}\par\medskip
\begin{kfActTable}{|>{\centering\arraybackslash}m{2.6cm}|m{6cm}|m{4cm}|}
\rowcolor{kfPrimaryLight!50}\textbf{\color{kfPrimary}항목} & \textbf{\color{kfPrimary}내용 (빈칸 채우기)} & \textbf{\color{kfPrimary}답안 작성}\\\hline
(가) 공간 & ㄴㄱ 외딴 ㅇㅁ — 고립과 내면 지향 ① & \kfTBlank \\\hline
(가) 매개체 & ㅇㅁ = 내면의 ㄱㅇ ② & \kfTBlank \\\hline
(가) 감정 변화 & ㅁㅇ → ㄱㅇ → ㄱㄹ ③ & \kfTBlank \\\hline
(가) 마무리 & 'ㅊㅇ처럼 사나이가 있습니다' ④ & \kfTBlank \\\hline
(가) 성찰 방식 & ㅂㄲㄹㅇ을 통한 ㅇㄹㅈ ㅅㅊ ⑤ & \kfTBlank \\\hline
(나) 초연 & '꽃이 져도 ㄱㅁㅇㅇ' — 자연의 변화에 초연 ⑥ & \kfTBlank \\\hline
(나) 매개체 & ㅍㄹ ㅅ + ㅁㅇ ㅁ = 자연의 거울 ⑦ & \kfTBlank \\\hline
(나) 자아 확인 & '너 ㅍㄹㄹㄷ' / '너 ㅁㅇㄷ' ⑧ & \kfTBlank \\\hline
(나) 마무리 & 'ㅅㅍㅅ ㅇㅇ ㅂㄱ'의 반복 ⑨ & \kfTBlank \\\hline
(나) 성찰 방식 & 자연 ㄱㄱ을 통한 ㅈㅈ ㄱㅈ ⑩ & \kfTBlank \\\hline
핵심 대비 & (가) ㄱㄷ·ㅅㅊ vs (나) ㄱㅈ·ㅍㅇ ⑪ & \kfTBlank \\\hline
공통 주제 & ㅈㄱ ㅅㅊ — 자아를 응시하고 인식하는 과정 ⑫ & \kfTBlank \\\hline
\end{kfActTable}
\kfSecHeader{kfAreaLiterature}{문제}{}

\begin{multicols}{2}\raggedcolumns\setlength{\columnsep}{12pt}
\begin{kfQBlock}{1}{}
\kfQStem{(가)와 (나)에 대한 설명으로 적절하지 않은 것은?}
\begin{kfChoices}
\kfChoice (가)의 화자는 우물에 비친 자기 모습을 반복적으로 들여다보며 복잡한 감정을 드러내고 있다.
\kfChoice (가)에서 '미워져', '가엾어', '그리워'의 감정 변화는 자아에 대한 갈등을 보여준다.
\kfChoice (나)의 화자는 '꽃이 져도 그만이오'에서 자연의 변화에 초연한 태도를 드러내고 있다.
\kfChoice (나)의 화자는 자연과 대립하여 갈등하며 고뇌하는 모습을 보이고 있다.
\kfChoice (가)와 (나) 모두 자연물을 매개로 자아를 인식하는 과정이 나타나고 있다.
\end{kfChoices}
\end{kfQBlock}

\begin{kfQBlock}{2}{}
\kfQStem{(가)의 '추억처럼 사나이가 있습니다'에서 '추억처럼'이 함축하는 바로 가장 적절한 것은?}
\begin{kfChoices}
\kfChoice 우물 속 자기 모습이 거리감 있는 아련한 존재로 인식됨을 나타낸다.
\kfChoice 과거의 즐거운 추억을 떠올리며 행복감에 젖어 있음을 나타낸다.
\kfChoice 우물이 실제로 화자의 추억을 저장하는 기능을 한다고 설명한다.
\kfChoice 사나이가 화자와 가까웠던 친구임을 우회적으로 암시한다.
\kfChoice 우물 속 자신의 모습을 화자가 완전히 잊었음을 의미한다.
\end{kfChoices}
\end{kfQBlock}

\begin{kfQBlock}{3}{}
\kfQStem{(가)와 (나)를 비교하여 감상한 내용으로 적절하지 않은 것은?

<보기> 자기 성찰의 시에서 '거울'의 역할을 하는 매개체는 화자가 자아를 어떻게 인식하는지를 결정한다. 내면의 갈등을 비추는 거울은 자아의 분열을 드러내고, 자연이라는 거울은 자아의 본래적 속성을 확인시켜 준다.}
\begin{kfChoices}
\kfChoice (가)와 (나) 모두 자아의 분열과 고뇌를 동일한 강도로 드러내고 있으므로, 두 시의 성찰 방식에는 차이가 없겠군.
\kfChoice (나)의 푸른 산과 맑은 물은 자연이라는 거울로, 화자에게 자아의 순수한 속성을 확인시켜 주는 매개체이겠군.
\kfChoice (가)에서 '미워져', '가엾어', '그리워'의 반복은 <보기>에서 말하는 자아의 분열을 드러내는 양상이겠군.
\kfChoice (나)에서 '살포시 웃어 볼까'의 반복은 자아의 긍정적 속성을 확인한 후의 평안한 반응이겠군.
\kfChoice (가)의 우물은 내면의 갈등을 비추는 거울로, 화자가 자아의 복잡한 감정과 대면하게 하는 매개체이겠군.
\end{kfChoices}
\end{kfQBlock}

\begin{kfQBlock}{4}{}
\kfQStem{(나)에서 '살포시 웃어 볼까 / 살포시 웃어 볼까'의 반복이 의미하는 바로 가장 적절한 것은?}
\begin{kfChoices}
\kfChoice 자연이 비춰 준 자기 존재의 순수함에 대한 조심스러운 긍정
\kfChoice 주변 타인의 부질없는 행위에 대한 조소와 비웃음의 발현
\kfChoice 거친 자연을 정복하겠다는 화자의 강한 의지의 표명
\kfChoice 당대의 사회적 불의를 향한 화자의 냉소적 반응 표출
\kfChoice 자기 자신을 향한 깊은 절망과 체념적 정서의 토로
\end{kfChoices}
\end{kfQBlock}

\begin{kfQBlock}{5}{}
\kfQStem{(가)에서 2연과 6연의 유사한 구절이 반복되면서 달라진 부분의 표현 효과로 가장 적절한 것은?}
\begin{kfChoices}
\kfChoice 6연에서 '사나이' 대신 새로운 인물이 등장하여 시적 대상이 전환되었음을 명시한다.
\kfChoice 2연과 6연이 완전히 같은 표현으로 반복되어 화자의 정서가 변화 없이 고정됨을 드러낸다.
\kfChoice 6연에서 달·구름·하늘 등 자연물이 모두 사라져 화자가 자연에 무관심해졌음을 보여 준다.
\kfChoice 2연의 자연물 나열이 6연에서 거꾸로 배열되어 사건의 시간 순서가 뒤집힘을 강조한다.
\kfChoice 2연 자연물 나열 끝에 '추억처럼 사나이' 행이 더해져 자아 인식이 깊어진다.
\end{kfChoices}
\end{kfQBlock}

\end{multicols}

\kfStartNonlit{이번 챕터 비문학 지문 5편}


\kfSecHeader{kfAreaNonlit}{지문}{하이데거의 존재와 시간 심화 — 시간성과 본래적 실존 (인문)}

\kfPassageNote{하이데거의 존재와 시간 심화 — 시간성과 본래적 실존 (인문)}{%
하이데거의 『존재와 시간』(1927)은 존재의 의미를 '시간성(Zeitlichkeit)'을 통해 해명하려는 기획이다. 하이데거에게 시간성이란 물리적 시계 시간이 아니라, 현존재(Dasein)가 자신의 존재를 이해하는 근원적 지평이다. 현존재는 과거·현재·미래라는 세 계기(Ekstasen)를 동시에 살아가는데, 이 세 계기는 순차적으로 나열되는 것이 아니라 하나의 통일된 구조 속에서 상호 침투한다. 미래는 현존재가 자기 앞의 가능성을 향해 자신을 '기투(Entwurf)'하는 차원이고, 과거는 이미 세계에 '던져져 있음(Geworfenheit)'으로서 현존재의 사실성을 구성하며, 현재는 세계 내 존재자와 관계 맺는 '빠져 있음(Verfallen)'의 차원이다.

이 세 계기 가운데 하이데거는 미래에 근원적 우위를 부여한다. 현존재가 자기의 가장 고유한 가능성, 곧 죽음을 향해 '앞질러 달려감(Vorlaufen)'으로써 자신의 유한성을 떠맡을 때, 비로소 과거의 사실성과 현재의 상황이 본래적으로 열린다. 본래적 미래는 현존재로 하여금 이미 던져진 자신의 과거를 자기의 것으로 '되풀이(Wiederholung)'하게 하고, 현재의 상황 속에서 결의하게 한다. 이것이 본래적 시간성의 구조이다.

죽음을 향한 존재(Sein-zum-Tode)는 본래적 시간성의 열쇠이다. 하이데거에 따르면 죽음은 경험적 사건이 아니라 현존재의 가장 고유하고, 무관계적이며, 추월 불가능한 가능성이다. '가장 고유한'이란 그 누구도 대신할 수 없다는 뜻이고, '무관계적'이란 타인과의 모든 관계가 의미를 잃는 지점이며, '추월 불가능한'이란 그 너머에 어떤 가능성도 없다는 뜻이다. 세인(das Man)은 죽음을 '사람은 언젠가 죽는다'라는 익명적 일반 사실로 은폐하여 현존재를 안심시킨다. 그러나 본래적 현존재는 죽음의 가능성을 가능성으로서 견뎌 내면서, 그 앞에서 자기 자신을 불안 속에 드러낸다.

본래적 실존의 양태를 하이데거는 '선구적 결의성(vorlaufende Entschlossenheit)'이라 부른다. 이 개념은 두 요소의 결합이다. '선구'는 죽음을 향한 앞질러 달려감이고, '결의성'은 양심의 부름에 응답하여 자신의 가장 고유한 '부채(Schuld)', 즉 존재론적 부채성을 떠맡는 태도이다. 여기서 '부채'는 경제적·도덕적 빚이 아니라, 현존재가 자신의 근거를 스스로 놓지 않았다는, 즉 근거의 근거가 될 수 없다는 존재론적 구조를 가리킨다. 현존재는 자기가 세계에 던져진 이유를 스스로 정할 수 없으면서도 자기의 존재를 떠맡아야 한다. 이 '무(Nichts)'로서의 근거를 자각하고 떠맡는 것이 결의성이다.

비본래적 시간성에서 현존재는 미래를 '기대(Gewärtigen)'로, 과거를 '망각(Vergessen)'으로, 현재를 '현재화(Gegenwärtigen)'로 살아간다. 이때 현존재는 세인의 세계에 몰입하여 자기 앞의 가능성을 계획 가능한 것으로 축소하고, 던져진 과거를 잊으며, 현재의 사물에 빠져든다. 본래적 시간성과 비본래적 시간성은 두 종류의 시간이 아니라, 하나의 시간성이 드러나는 두 양태이다. 하이데거는 이 양태 전환이 현존재의 자유에 달려 있다고 본다. 현존재는 언제든 세인의 안일함에서 벗어나 선구적 결의성을 통해 본래적 시간성으로 전환할 수 있으며, 동시에 언제든 다시 비본래적 시간성으로 빠져들 수 있다. 이 가능성의 양방향성이야말로 현존재의 유한한 자유의 표현이다.
}
\kfSecHeader{kfAreaNonlit}{문제}{}

\begin{multicols}{2}\raggedcolumns\setlength{\columnsep}{12pt}
\begin{kfQBlock}{1}{}
\kfQStem{윗글의 내용과 일치하지 않는 것은?}
\begin{kfChoices}
\kfChoice 하이데거의 시간성에서 과거·현재·미래는 순차적으로 나열되지 않고 상호 침투하는 통일적 구조를 이룬다.
\kfChoice 본래적 시간성에서 미래는 현존재가 자기의 가장 고유한 가능성을 향해 기투하는 차원이다.
\kfChoice 죽음은 현존재의 가장 고유하고 무관계적이며 추월 불가능한 가능성이다.
\kfChoice '부채'는 현존재가 타인에게 갚아야 할 도덕적 빚을 가리키는 개념이다.
\kfChoice 본래적 시간성과 비본래적 시간성은 하나의 시간성이 드러나는 두 양태이다.
\end{kfChoices}
\end{kfQBlock}

\begin{kfQBlock}{2}{}
\kfQStem{윗글에 따를 때, 하이데거가 시간성의 세 계기 가운데 '미래'에 근원적 우위를 부여하는 이유로 가장 적절한 것은?}
\begin{kfChoices}
\kfChoice 미래는 아직 오지 않은 것이므로 현존재가 자유롭게 계획할 수 있는 유일한 차원이기 때문이다.
\kfChoice 미래를 향한 기투를 통해 과거의 사실성과 현재의 상황이 비로소 본래적으로 열리기 때문이다.
\kfChoice 과거와 현재는 이미 결정되어 변경 불가능하지만 미래는 열려 있기 때문이다.
\kfChoice 물리적 시간에서도 미래가 과거보다 시간적으로 우선하기 때문이다.
\kfChoice 세인의 세계에서는 현재가 지배적이므로, 그에 대항하여 미래를 강조하기 위해서이다.
\end{kfChoices}
\end{kfQBlock}

\begin{kfQBlock}{3}{}
\kfQStem{윗글을 바탕으로 <보기>를 이해한 내용으로 적절하지 않은 것은?

<보기> B 씨는 정년퇴직 후 무의미한 나날을 보내고 있었다. 어느 날 중병 진단을 받고 자신의 죽음을 실감한 뒤, B 씨는 과거에 이루지 못한 꿈을 되돌아보며 남은 시간 속에서 진정으로 자기가 원하는 일을 시작했다. 그러면서도 간혹 '어차피 다 부질없다'는 생각에 빠져 텔레비전만 보며 하루를 보내기도 했다.}
\begin{kfChoices}
\kfChoice B 씨가 중병 진단 후 자신의 죽음을 실감한 것은 '죽음을 향한 존재'로서 자기의 유한성을 직면한 것으로 볼 수 있겠군.
\kfChoice B 씨가 과거의 꿈을 되돌아본 것은 던져진 과거를 자기 것으로 '되풀이'하는 본래적 시간성의 양태에 해당하겠군.
\kfChoice B 씨가 남은 시간 속에서 원하는 일을 시작한 것은 유한한 시간 속에서의 '선구적 결의성'으로 해석할 수 있겠군.
\kfChoice B 씨가 '어차피 다 부질없다'며 텔레비전만 보는 것은 본래적 시간성에서 비본래적 시간성으로의 전환으로 볼 수 있겠군.
\kfChoice B 씨가 본래적 시간성에서 비본래적 시간성으로 되돌아간 것은 도덕적으로 타락한 것이므로 하이데거의 관점에서 비난의 대상이겠군.
\end{kfChoices}
\end{kfQBlock}

\begin{kfQBlock}{4}{}
\kfQStem{윗글의 '선구적 결의성(vorlaufende Entschlossenheit)'에 대한 설명으로 적절하지 않은 것은?}
\begin{kfChoices}
\kfChoice '선구'는 죽음을 향한 앞질러 달려감을 가리킨다.
\kfChoice '결의성'은 양심의 부름에 응답하여 존재론적 부채성을 떠맡는 태도이다.
\kfChoice 선구적 결의성을 통해 현존재는 자신의 근거를 스스로 정립할 수 있게 된다.
\kfChoice 선구적 결의성은 본래적 실존의 양태를 지칭하는 개념이다.
\kfChoice 선구적 결의성은 '선구'와 '결의성'이라는 두 요소의 결합이다.
\end{kfChoices}
\end{kfQBlock}

\end{multicols}

\kfSecHeader{kfAreaNonlit}{지문}{데카르트의 코기토와 하이데거의 현존재 (인문)}

\kfPassageNote{데카르트의 코기토와 하이데거의 현존재 (인문)}{%
(가) 르네 데카르트는 확실한 지식의 토대를 찾기 위해 '방법적 회의'를 수행했다. 감각 경험은 착각일 수 있고, 수학적 진리조차 악한 천재(genius malignus)가 우리를 속이고 있을 가능성을 배제할 수 없다. 그러나 이 모든 것을 의심하는 바로 그 순간에도 '의심하고 있는 나'의 존재만큼은 부정할 수 없다. 이것이 「나는 생각한다, 고로 나는 존재한다(Cogito, ergo sum)」이다. 데카르트에게 코기토는 일체의 의심을 견뎌 내는 제1원리이며, 여기서 '나'는 연장(extension)을 갖지 않는 순수한 사유 실체(res cogitans)이다.

데카르트는 이 사유 실체로부터 출발하여 외부 세계의 존재를 증명하려 했다. 그 논증의 핵심은 신 존재 증명이다. 나의 의식 속에는 '무한한 완전자'의 관념이 있는데, 유한한 존재인 내가 이 관념을 스스로 산출할 수는 없으므로 실제로 무한한 완전자, 곧 신이 존재해야 한다. 그리고 완전한 신은 기만자가 아니므로, 내가 명석판명하게 인식하는 것은 참이라는 보증이 성립한다. 이로써 외부의 물질 세계, 즉 연장 실체(res extensa)의 존재도 확보된다. 결과적으로 데카르트의 세계는 사유 실체와 연장 실체라는 두 종류의 실체로 구성되며, 이 이원론적 구조 속에서 주관(사유)과 객관(연장)은 본질적으로 분리된다.

(나) 마르틴 하이데거는 데카르트의 코기토가 존재의 의미를 묻지 않은 채 '사유하는 주체'를 자명한 출발점으로 삼은 것이야말로 서양 형이상학의 근본 오류라고 비판했다. 데카르트가 「나는 존재한다」라고 말할 때, 그 '존재한다'의 의미는 무엇인가? 데카르트는 이 물음을 건너뛰고 코기토를 확고한 기점으로 확립했지만, 하이데거에 따르면 존재의 의미를 해명하지 않은 출발점은 사상(砂上)의 기초와 다르지 않다.

하이데거가 제시한 대안이 '현존재(Dasein)' 분석이다. 현존재란 자신의 존재를 문제 삼는 존재자, 곧 인간을 가리키되, 데카르트적 주관과는 근본적으로 다르다. 데카르트의 '나'가 세계로부터 고립된 사유 실체인 반면, 현존재는 처음부터 세계 속에서 사물을 사용하고 타인과 관계 맺으며 살아가는 '세계-내-존재(In-der-Welt-sein)'이다. 여기서 세계란 현존재와 무관하게 존재하는 물리적 용기(容器)가 아니라, 현존재의 관심과 기획에 의해 의미 연관으로 열리는 장(場)이다.

두 입장의 차이는 인식과 존재의 관계에서도 드러난다. 데카르트에게 인식은 주관이 객관을 표상(表象)하는 행위이며, 인식론이 철학의 근본 물음이다. 반면 하이데거는 인식을 현존재의 파생적 양태로 본다. 현존재는 먼저 세계 속에서 사물을 '손안의 존재(Zuhandenheit)'로 다루면서 살아가고, 이론적 관조(觀照)는 이러한 실천적 관계가 중단되었을 때 비로소 파생적으로 성립하는 것이다. 이를테면 망치를 사용하는 동안 우리는 망치를 이론적으로 '인식'하지 않는다. 망치가 부러졌을 때 비로소 그것이 대상으로 드러난다. 하이데거에게 존재론은 인식론에 선행하며, '어떻게 아는가'보다 '어떻게 있는가'가 더 근본적인 물음이다.
}
\kfSecHeader{kfAreaNonlit}{문제}{}

\begin{multicols}{2}\raggedcolumns\setlength{\columnsep}{12pt}
\begin{kfQBlock}{1}{}
\kfQStem{(가)와 (나)의 내용과 일치하지 않는 것은?}
\begin{kfChoices}
\kfChoice 데카르트는 모든 것을 의심하는 방법적 회의의 끝에서 코기토를 제1원리로 확립했다.
\kfChoice 데카르트는 무한자 관념에 근거한 신 존재 증명을 거쳐 연장 실체의 존재를 보증하였다.
\kfChoice 하이데거는 데카르트가 '존재한다'의 의미를 해명한 토대 위에서 코기토를 정립했다고 평가했다.
\kfChoice 하이데거의 현존재는 세계 속에서 사물을 다루고 타인과 관계 맺는 세계-내-존재이다.
\kfChoice 하이데거에게 이론적 관조는 실천적 관계가 중단되었을 때 파생적으로 성립하는 양태이다.
\end{kfChoices}
\end{kfQBlock}

\begin{kfQBlock}{2}{}
\kfQStem{(가)의 데카르트와 (나)의 하이데거를 비교한 내용으로 가장 적절한 것은?}
\begin{kfChoices}
\kfChoice 데카르트는 인식론을 근본 물음으로 보았으나, 하이데거는 존재론이 인식론에 앞선다고 보았다.
\kfChoice 데카르트는 세계를 의미 연관의 장으로 보았으나, 하이데거는 세계를 물리적 용기로 보았다.
\kfChoice 데카르트는 인간을 세계-내-존재로 파악했으나, 하이데거는 인간을 사유 실체로 파악했다.
\kfChoice 데카르트와 하이데거 모두 주관과 객관의 이원론적 분리를 전제하였다.
\kfChoice 데카르트와 하이데거 모두 이론적 관조가 가장 근본적인 인식 방식이라고 보았다.
\end{kfChoices}
\end{kfQBlock}

\begin{kfQBlock}{3}{}
\kfQStem{윗글을 바탕으로 <보기>를 이해한 내용으로 적절하지 않은 것은?

<보기> 기술자 C 씨는 매일 공구를 다루며 기계를 수리한다. 어느 날 애용하던 스패너가 부러지자, C 씨는 비로소 스패너의 재질과 구조에 대해 궁금해지기 시작했다.}
\begin{kfChoices}
\kfChoice C 씨가 공구를 다루며 기계를 수리하는 것은 하이데거가 말하는 '손안의 존재'로서 사물과 관계 맺는 방식에 해당하겠군.
\kfChoice 스패너가 부러지기 전까지 C 씨가 스패너를 이론적으로 인식하지 않는 것은 실천적 관계가 유지되고 있기 때문이겠군.
\kfChoice 스패너가 부러진 후 재질과 구조에 궁금증을 느낀 것은 실천적 관계의 중단에서 이론적 관조가 파생적으로 성립하는 예이겠군.
\kfChoice 데카르트의 관점에서 C 씨의 경험은 사유 실체가 연장 실체를 표상하는 행위로 설명될 수 있겠군.
\kfChoice 하이데거의 관점에서 C 씨는 스패너를 이론적으로 관조한 후에야 비로소 그것을 실천적으로 사용할 수 있게 되었겠군.
\end{kfChoices}
\end{kfQBlock}

\begin{kfQBlock}{4}{}
\kfQStem{(가)에서 데카르트가 '신 존재 증명'을 필요로 한 논리적 이유로 가장 적절한 것은?}
\begin{kfChoices}
\kfChoice 코기토만으로는 외부 세계의 존재를 확보할 수 없었기 때문이다.
\kfChoice 사유 실체의 존재를 증명하기 위해 신의 보증이 필요했기 때문이다.
\kfChoice 하이데거의 비판에 미리 대응하기 위한 논증 전략이었기 때문이다.
\kfChoice 방법적 회의를 통해 신의 존재마저 의심해야 했기 때문이다.
\kfChoice 연장 실체가 사유 실체보다 더 확실한 것임을 보이기 위해서이다.
\end{kfChoices}
\end{kfQBlock}

\end{multicols}

\kfSecHeader{kfAreaNonlit}{지문}{하먼의 객체지향 존재론 (인문)}

\kfPassageNote{하먼의 객체지향 존재론 (인문)}{%
그레이엄 하먼(Graham Harman)은 21세기 초 등장한 '사변적 실재론' 진영의 핵심 인물로, '객체지향 존재론(Object-Oriented Ontology, OOO)'을 제창했다. 하먼의 출발점은 하이데거의 도구 분석이다. 하이데거는 망치를 사용하는 동안 망치는 투명하게 물러서 있다가, 망치가 부러질 때 비로소 대상으로 드러난다고 분석했다. 하먼은 이 '물러섬(withdrawal)'을 인간과 도구의 관계에 한정하지 않고, 모든 객체 사이의 관계로 확장한다. 불이 솜을 태울 때, 불은 솜의 모든 성질에 접근하는 것이 아니라 가연성이라는 특정 측면에만 접촉한다. 솜의 색깔, 무게, 부드러움은 불과의 관계에서 물러서 있다. 하먼은 이처럼 모든 객체가 자신과 관계 맺는 다른 객체에게 자신의 전체를 드러내지 않으며, 언제나 관계로 환원되지 않는 잉여를 보유한다고 주장한다.

하먼의 존재론에서 객체는 네 가지 차원을 갖는다. '실재적 객체(real object)'는 모든 관계 너머에 물러서 있는 객체의 핵심이며, '감각적 객체(sensual object)'는 다른 객체와의 관계 속에서 직접 현전하는 객체의 현상적 측면이다. 마찬가지로 '실재적 성질(real quality)'은 객체가 관계와 무관하게 보유하는 본질적 속성이고, '감각적 성질(sensual quality)'은 관계 속에서 드러나는 우발적 속성이다. 이 네 항의 조합에서 하먼은 '대리 인과(vicarious causation)'라는 독특한 인과 이론을 도출한다.

대리 인과의 핵심은 실재적 객체들이 직접 접촉할 수 없다는 데 있다. 실재적 객체는 모든 관계로부터 물러서 있으므로, 두 실재적 객체가 직접 만나는 일은 불가능하다. 인과적 상호 작용은 오직 감각적 객체를 매개로만 일어난다. 두 실재적 객체가 하나의 감각적 객체 내부에서 '접합(fusion)'할 때, 즉 하나의 경험 공간 안에서 만날 때 비로소 인과가 성립한다. 이것을 하먼은 '대리'라 부르는데, 실재적 객체들이 직접이 아니라 감각적 차원의 매개를 통해 간접적으로만 관계한다는 의미이다.

하먼의 객체지향 존재론은 이른바 '상관주의(correlationism)'에 대한 비판이기도 하다. 칸트 이래 근대 철학은 인간의 인식 조건과 무관한 '사물 자체'에 대해서는 말할 수 없다고 보았다. 우리가 알 수 있는 것은 오직 인간 주관과 세계의 상관관계뿐이라는 것이다. 하먼은 이를 인간 중심적 편향이라 비판하며, 물러섬의 구조가 인간-세계 관계에만 성립하는 것이 아니라 객체 사이 어디에서든 발생한다고 주장한다. 돌과 물, 나무와 바람 사이에서도 물러섬은 동일하게 작동한다. 이로써 하먼은 '평평한 존재론(flat ontology)', 즉 인간이 존재론적으로 특권적 지위를 갖지 않는 세계관을 제시한다. 이 관점에서 인간은 수많은 객체 중 하나일 뿐이며, 존재론은 인간의 경험을 넘어서는 실재의 깊이를 탐구할 수 있어야 한다.
}
\kfSecHeader{kfAreaNonlit}{문제}{}

\begin{multicols}{2}\raggedcolumns\setlength{\columnsep}{12pt}
\begin{kfQBlock}{1}{}
\kfQStem{윗글의 내용과 일치하는 것은?}
\begin{kfChoices}
\kfChoice 하먼은 물러섬의 구조를 인간과 도구의 관계에 한정하여 분석하였다.
\kfChoice 감각적 객체는 모든 관계 너머에 물러서 있는 객체의 핵심이다.
\kfChoice 대리 인과에서 두 실재적 객체는 직접 접촉하여 인과 관계를 형성한다.
\kfChoice 하먼은 상관주의를 지지하여 인간의 인식 조건 내에서만 존재를 논해야 한다고 보았다.
\kfChoice 하먼에 따르면 불이 솜을 태울 때, 불은 솜의 가연성이라는 특정 측면에만 접촉한다.
\end{kfChoices}
\end{kfQBlock}

\begin{kfQBlock}{2}{}
\kfQStem{윗글에 따를 때, 하먼이 '대리 인과'라는 개념을 도입한 이유로 가장 적절한 것은?}
\begin{kfChoices}
\kfChoice 실재적 객체들이 모든 관계로부터 물러서 있어 직접 접촉할 수 없기 때문이다.
\kfChoice 감각적 객체들이 서로 물러서 있어 인과 관계가 성립하지 않기 때문이다.
\kfChoice 인간만이 인과 관계를 인식할 수 있는 특권적 존재이기 때문이다.
\kfChoice 칸트의 사물 자체 개념을 계승하여 인식의 한계를 강조하기 위해서이다.
\kfChoice 하이데거의 도구 분석이 인과 관계를 설명하지 못했기 때문이다.
\end{kfChoices}
\end{kfQBlock}

\begin{kfQBlock}{3}{}
\kfQStem{윗글을 바탕으로 <보기>를 이해한 내용으로 적절하지 않은 것은?

<보기> 비가 내려 바위 표면이 젖고, 오랜 세월에 걸쳐 바위가 침식된다.}
\begin{kfChoices}
\kfChoice 비는 바위의 모든 성질에 접근하는 것이 아니라, 바위의 특정 측면에만 접촉하겠군.
\kfChoice 바위의 색깔이나 경도 같은 일부 성질은 비와의 관계에서 물러서 있겠군.
\kfChoice 비와 바위의 인과적 상호 작용은 감각적 객체를 매개로 성립하겠군.
\kfChoice 바위의 침식은 인간이 관찰해야만 성립하는 인과 관계이겠군.
\kfChoice 비-바위 관계에서도 물러섬의 구조가 작동하므로, 이는 평평한 존재론의 예이겠군.
\end{kfChoices}
\end{kfQBlock}

\begin{kfQBlock}{4}{}
\kfQStem{윗글의 문맥에서 '평평한 존재론'이 의미하는 바로 가장 적절한 것은?}
\begin{kfChoices}
\kfChoice 모든 객체가 크기와 중요성에서 동일하다는 양적 평등주의의 주장
\kfChoice 인간이 특권적 지위를 갖지 않고 존재자 사이에 위계가 없다는 관점
\kfChoice 물질적 세계가 정신적 세계보다 존재론적으로 우월하다는 유물론의 주장
\kfChoice 모든 존재자의 성질이 외부 관계를 통해 완전히 드러난다는 관계주의의 입장
\kfChoice 감각적 객체만이 실재하고 실재적 객체는 환상이라는 현상주의의 결론
\end{kfChoices}
\end{kfQBlock}

\end{multicols}

\kfSecHeader{kfAreaNonlit}{지문}{사변적 실재론 — 메야수의 상관주의 비판 (인문)}

\kfPassageNote{사변적 실재론 — 메야수의 상관주의 비판 (인문)}{%
캉탱 메야수(Quentin Meillassoux)는 2006년 출간한 『유한성 이후(Après la finitude)』에서 칸트 이래 근대 철학을 지배해 온 '상관주의(corrélationnisme)'를 비판하고, 사유와 무관한 절대적 실재에 접근할 수 있는 철학의 길을 모색했다. 메야수에 따르면 상관주의란, 우리가 접근할 수 있는 것은 오직 사유와 존재의 '상관관계(corrélation)'뿐이며, 사유 없는 존재나 존재 없는 사유는 사고 불가능하다는 입장이다. 칸트의 초월적 관념론, 후설의 현상학, 하이데거의 존재론 모두 이 상관주의의 변주라는 것이 메야수의 진단이다.

메야수는 상관주의의 한계를 드러내기 위해 '조상성(ancestralité)' 문제를 제기한다. 조상적 사실이란 인류의 출현 이전, 심지어 모든 생명 이전의 사실을 말한다. 이를테면 지구가 약 45억 년 전에 형성되었다는 과학적 진술이 있다. 상관주의자에게 이 진술은 난제이다. 상관주의에 따르면 존재는 사유와의 상관관계 속에서만 의미를 가지는데, 인간도 생명도 존재하지 않던 시점의 사건에 대해 어떻게 상관관계가 성립할 수 있는가? 상관주의자가 '그것은 현재의 과학적 구성에 불과하다'고 답한다면, 이는 사실상 조상적 사건의 실재성을 부정하는 것이 되어, 자연과학의 가장 기초적인 성과와 충돌한다.

상관주의를 넘어서기 위해 메야수가 제시하는 핵심 원리는 '사실성의 원리(principe de factualité)'이다. 메야수는 모든 것이 우연적이라는 사실, 즉 어떤 사물도 그것이 존재해야 할 필연적 이유를 갖지 않는다는 사실만큼은 필연적이라고 주장한다. 이것을 그는 '우연성의 필연성(nécessité de la contingence)'이라 부른다. 전통적으로 철학은 세계의 근거를 필연적 존재(신, 실체, 주체 등)에서 찾았지만, 메야수는 그 어떤 필연적 근거도 존재하지 않으며, 모든 것은 이유 없이 달라질 수 있다고 본다. 자연 법칙조차 내일 변할 수 있으며, 그것이 변하지 않을 이유는 없다. 다만 이 우연성 자체는 변할 수 없는 유일한 필연이다.

메야수의 논증 구조는 이러하다. 만약 상관주의가 옳다면, 사유 밖의 존재에 대해서는 아무것도 말할 수 없어야 한다. 그러나 상관주의 자체도 '사유와 존재의 상관관계는 필연적이다'라는 주장을 하고 있으며, 이는 사유 밖의 존재에 대한 암묵적 주장이다. 메야수는 이 자기 모순을 파고든다. 상관주의가 일관되려면 상관관계의 필연성마저 포기해야 하는데, 그렇게 되면 사유와 무관한 존재가 가능하다는 것을 인정하게 된다. 이로써 메야수는 절대적인 것, 즉 사유에 의존하지 않는 것에 대한 앎의 가능성을 복원한다. 그가 접근하는 절대적인 것은 수학적으로 형식화 가능한 것, 곧 '제1성질'이다. 갈릴레이-데카르트적 전통에서 연장·운동·수 등의 제1성질은 주관에 의존하지 않는 객관적 속성이며, 메야수는 이것이 사물 자체에 대한 수학적 접근의 근거라 주장한다.
}
\kfSecHeader{kfAreaNonlit}{문제}{}

\begin{multicols}{2}\raggedcolumns\setlength{\columnsep}{12pt}
\begin{kfQBlock}{1}{}
\kfQStem{윗글의 내용과 일치하지 않는 것은?}
\begin{kfChoices}
\kfChoice 메야수에 따르면 칸트, 후설, 하이데거의 철학은 모두 상관주의의 변주에 해당한다.
\kfChoice 조상적 사실이란 인류 출현 이전, 심지어 모든 생명 이전의 사실을 말한다.
\kfChoice 메야수는 자연 법칙이 변할 가능성을 인정하면서도, 우연성 자체는 필연적이라고 주장한다.
\kfChoice 메야수에 따르면 상관주의는 자기 모순 없이 일관되게 유지될 수 있다.
\kfChoice 메야수가 접근하는 절대적인 것은 수학적으로 형식화 가능한 제1성질이다.
\end{kfChoices}
\end{kfQBlock}

\begin{kfQBlock}{2}{}
\kfQStem{윗글에 따를 때, '조상성 문제'가 상관주의에 대한 비판으로 기능하는 논리적 근거로 가장 적절한 것은?}
\begin{kfChoices}
\kfChoice 조상적 사실이 수학적 형식화 자체를 거부하므로 과학적 접근의 근본 한계를 폭로하기 때문이다.
\kfChoice 인류 이전 사건에 사유-존재 상관관계가 성립할 수 없는데도 과학은 실재로 다루기 때문이다.
\kfChoice 조상적 사실은 오직 상관주의가 마련한 인식 틀 안에서만 비로소 해석되는 현상이기 때문이다.
\kfChoice 자연과학이 상관주의를 전제로 성립하므로 조상적 사실 또한 상관주의를 강력히 지지하기 때문이다.
\kfChoice 조상적 사실에 대한 연대 측정의 부정확성이 과학적 진술 전반의 신뢰성에 의문을 제기하기 때문이다.
\end{kfChoices}
\end{kfQBlock}

\begin{kfQBlock}{3}{}
\kfQStem{윗글의 '우연성의 필연성'이 의미하는 바로 가장 적절한 것은?}
\begin{kfChoices}
\kfChoice 모든 것이 우연적이라는 사실 자체가 유일하게 필연적이라는 것
\kfChoice 우연적 사건은 반드시 필연적 결과로 귀결된다는 것
\kfChoice 세계의 근거가 되는 필연적 존재(신, 실체)가 반드시 존재한다는 것
\kfChoice 자연 법칙은 변하지 않으며 영원히 고정되어 있다는 것
\kfChoice 우연성은 인간 인식의 한계에서 비롯된 가상이며 실재는 필연적이라는 것
\end{kfChoices}
\end{kfQBlock}

\begin{kfQBlock}{4}{}
\kfQStem{윗글의 메야수와 칸트를 비교한 내용으로 적절하지 않은 것은?}
\begin{kfChoices}
\kfChoice 칸트는 사유의 조건 밖의 '사물 자체'에 대해서는 알 수 없다고 보았으나, 메야수는 사물 자체에 대한 앎의 가능성을 복원하고자 했다.
\kfChoice 메야수는 칸트의 초월적 관념론을 상관주의의 한 변주로 분류했다.
\kfChoice 칸트와 메야수는 수학적 형식화를 통해 사물 자체에 다가가야 한다는 입장을 함께 취했다.
\kfChoice 메야수는 칸트 이래 근대 철학이 사유-존재의 상관관계 안에 갇혀 있다고 비판했다.
\kfChoice 메야수가 상관주의의 자기 모순을 지적한 것은 칸트적 전통에 대한 비판의 일환이다.
\end{kfChoices}
\end{kfQBlock}

\end{multicols}

\kfSecHeader{kfAreaNonlit}{지문}{인격 동일성의 현대적 논쟁 (인문)}

\kfPassageNote{인격 동일성의 현대적 논쟁 (인문)}{%
인격 동일성(personal identity)의 문제는 '나는 어제의 나와 같은 존재인가?'라는 물음에서 출발한다. 존 로크는 이 물음에 대해 기억 이론을 제시했다. 어제의 경험을 기억하는 한, 어제의 나와 오늘의 나는 동일한 인격이다. 그러나 토마스 리드는 용감한 장교의 역설로 이를 반박했다. 소년 시절 매를 맞은 한 사람이 장교가 되어 적의 깃발을 빼앗았고, 노장군이 되었다. 장교는 소년의 기억을 갖고 있고, 노장군은 장교의 기억을 갖고 있지만, 노장군은 소년의 기억을 갖고 있지 않다. 기억 이론에 따르면 장교는 소년과 동일하고, 노장군은 장교와 동일하지만, 노장군은 소년과 동일하지 않다. 이는 동일성의 이행성(A=B이고 B=C이면 A=C)을 위반하는 모순이다.

데릭 파핏(Derek Parfit)은 이 문제에 대해 혁명적 해답을 제시했다. 파핏에 따르면 인격 동일성은 그 자체로 중요한 것이 아니며, 중요한 것은 '관계 R', 즉 심리적 연결성(직접적 기억, 성격, 의도 등의 인과적 연결)과 심리적 연속성(연결성의 겹치는 사슬)이다. 파핏은 이를 보여주기 위해 순간 이동 사고실험을 제안한다. 지구의 나를 스캐너로 정밀 복제하여 화성에 동일한 물리적·심리적 상태를 가진 존재를 만들고, 지구의 원본은 파괴된다고 하자. 화성의 존재는 나의 모든 기억, 성격, 의도를 갖고 있다. 이 경우 화성의 존재는 '나'인가?

파핏은 한 걸음 더 나아간다. 만약 스캐너의 오류로 지구의 원본이 파괴되지 않았다면 어떨까? 이제 지구와 화성에 각각 나의 심리적 연속체가 존재한다. 둘 다 동일한 관계 R을 보유하지만, 동일성은 일대일 관계이므로 둘 다 '나'일 수는 없다. 파핏은 이 사고실험에서 인격 동일성이라는 개념 자체가 중요한 것이 아님을 도출한다. 중요한 것은 생존(survival)이지, 엄밀한 의미의 동일성이 아니다. 관계 R이 유지되는 한, 비록 동일성이 성립하지 않더라도 '나에게 중요한 것'은 보존된다. 이를 파핏은 인격에 대한 '환원주의(reductionism)'라 부른다. 인격이란 별도의 실체가 아니라 물리적·심리적 사건들의 연속에 불과하다.

파핏의 입장에 반대하는 비환원주의자들은 인격의 동일성을 구성하는 '더 깊은 사실(further fact)'이 있다고 주장한다. 이 더 깊은 사실은 물리적·심리적 연속성으로 환원될 수 없는 별도의 형이상학적 실체이다. 한편 4차원주의자들은 인격을 시간에 걸쳐 펼쳐진 4차원적 존재로 파악한다. 이 관점에서 어제의 나와 오늘의 나는 하나의 시공간적 '벌레(spacetime worm)'의 다른 시간 단면이다. 동일성은 부분들 사이의 관계가 아니라, 전체로서의 4차원 존재가 갖는 속성이다. 또한 서사적 자아 이론은 인격의 동일성을 형이상학적 사실이 아니라, 자기 삶의 이야기를 구성하는 서사적 행위에서 찾는다. 내가 '나'인 것은 과거의 경험들을 하나의 일관된 이야기로 엮어 내기 때문이다.
}
\kfSecHeader{kfAreaNonlit}{문제}{}

\begin{multicols}{2}\raggedcolumns\setlength{\columnsep}{12pt}
\begin{kfQBlock}{1}{}
\kfQStem{윗글의 내용과 일치하지 않는 것은?}
\begin{kfChoices}
\kfChoice 로크의 기억 이론에 따르면, 어제의 경험을 기억하는 한 어제의 나와 오늘의 나는 동일한 인격이다.
\kfChoice 리드의 용감한 장교 역설은 기억 이론이 동일성의 이행성을 위반함을 보여준다.
\kfChoice 파핏에 따르면 인격 동일성 자체가 가장 중요하며, 관계 R은 부차적이다.
\kfChoice 4차원주의에서 어제의 나와 오늘의 나는 하나의 시공간적 존재의 다른 시간 단면이다.
\kfChoice 서사적 자아 이론은 인격의 동일성을 자기 삶의 이야기를 구성하는 서사적 행위에서 찾는다.
\end{kfChoices}
\end{kfQBlock}

\begin{kfQBlock}{2}{}
\kfQStem{윗글에 따를 때, 파핏이 순간 이동 사고실험의 '복제본이 파괴되지 않는 변형'을 통해 도출하고자 한 결론으로 가장 적절한 것은?}
\begin{kfChoices}
\kfChoice 심리적 연속성만 유지되면 분기된 두 존재 모두 동일한 인격이라 보아야 한다.
\kfChoice 인격 동일성이라는 개념 자체가 중요한 것이 아니라 관계 R의 유지가 중요하다.
\kfChoice 원본만이 진정한 인격이므로 분기된 복제본은 별개의 존재로 간주되어야 한다.
\kfChoice 순간 이동은 기술적으로 불가능하므로 인격 동일성 논의 자체가 무의미하다.
\kfChoice 인격 동일성은 물리적 연속성에만 의존하므로 심리적 요소는 부차적인 조건이다.
\end{kfChoices}
\end{kfQBlock}

\begin{kfQBlock}{3}{}
\kfQStem{윗글을 바탕으로 <보기>를 이해한 내용으로 적절하지 않은 것은?

<보기> D 씨는 사고로 10년 전 이전의 기억을 모두 잃었지만, 사고 이후의 기억은 온전하다. D 씨의 가족은 '저 사람은 여전히 우리 가족'이라고 말하고, D 씨 자신도 사고 이후의 경험들로 자기 삶의 이야기를 만들어 가고 있다.}
\begin{kfChoices}
\kfChoice 로크의 기억 이론에 따르면, D 씨는 10년 전의 자신과 동일한 인격이라 보기 어렵겠군.
\kfChoice 파핏의 관점에서 D 씨에게 중요한 것은 인격 동일성보다 관계 R의 유지이겠군.
\kfChoice 비환원주의자라면 D 씨의 기억 상실에도 불구하고 인격을 구성하는 '더 깊은 사실'이 유지된다고 볼 수 있겠군.
\kfChoice 서사적 자아 이론에 따르면, D 씨가 사고 이후의 경험으로 자기 이야기를 구성하는 것이 인격의 근거이겠군.
\kfChoice 4차원주의에 따르면, D 씨의 기억 상실은 D 씨라는 시공간적 존재 자체가 단절되었음을 의미하겠군.
\end{kfChoices}
\end{kfQBlock}

\begin{kfQBlock}{4}{}
\kfQStem{윗글의 파핏과 비환원주의자를 비교한 내용으로 적절하지 않은 것은?}
\begin{kfChoices}
\kfChoice 파핏은 인격을 물리적·심리적 사건들의 연속에 불과하다고 보지만, 비환원주의자는 별도의 형이상학적 실체를 상정한다.
\kfChoice 파핏은 인격 동일성보다 관계 R의 유지를 중시하지만, 비환원주의자는 인격 동일성을 구성하는 '더 깊은 사실'을 중시한다.
\kfChoice 파핏과 비환원주의자 모두 인격이란 별도의 실체라는 점에서는 동의한다.
\kfChoice 파핏의 환원주의는 인격을 경험적 요소로 분석하려 하지만, 비환원주의는 그러한 환원을 거부한다.
\kfChoice 비환원주의자는 파핏의 순간 이동 사고실험에서도 '더 깊은 사실'에 의해 동일성이 결정된다고 주장할 수 있다.
\end{kfChoices}
\end{kfQBlock}

\end{multicols}

\kfStartWeekly{패턴 워크북 — 총 ?문항}


\kfSecHeader{kfAreaWeekly}{지문}{문항 1 지문 / 섞어치기}

\kfPassageNote{문항 1 지문 / 섞어치기}{%
개별적이고 특수한 작업을 수행하는 여러 대의 클라이언트와 연결되어 정보나 서비스를 제공하는 컴퓨터 시스템을 서버라고 한다. 서버는 클라이언트의 요청에 따라 작업을 수행하거나 정보를 검색하고, 그 결과를 클라이언트에 돌려준다. 서버를 구성하는 기계에서 작업을 신속하고 정확하게 처리하기 위해서는 데이터를 안전하게 관리해야 할 뿐 아니라 신속하게 데이터에 접근하여 데이터의 읽기나 쓰기를 실행할 필요가 있다. 이런 필요를 충족시키기 위해 개발된 것이 ㉠분산 파일 관리 체계이다. 분산 파일 관리 체계는 데이터의 불법 변조를 방지하고, 기계적 결함에 대응하여 데이터의 유실로부터 안전을 도모하며, 서버의 부하를 줄여 신속한 연산을 수행하도록 기계적 성능과 저장 용량을 확보하는 데 유리한 전략이다. 단일 기계로 구성된 서버에서 데이터를 관리할 경우에는 하드웨어적으로 램이나 CPU의 성능이 우수한 기계를 확보하는 데 큰 비용이 들지만, 분산 파일 관리 체계에서는 다수의 저렴한 저성능 기계를 연결하여

(중략)
}
\begin{kfQBlock}{1}{}
\kfQStem{<보기>를 참고해 ⓐ의 운영 원리를 이해한 내용으로 적절하지 않은 것은?}
\begin{kfEvidence}<보기> 클러스터 X는 마스터-슬레이브 복제를 사용한다. 읽기 요청은 마스터와 슬레이브 모두에서 처리하고, 쓰기 요청은 마스터에서만 처리한 뒤 슬레이브로 반영한다. 클러스터 Y는 피어 투 피어 복제를 사용한다. 읽기와 쓰기 요청을 바쁘지 않은 피어에서 처리하고, 수정된 데이터는 다른 피어들로 복제한다.\end{kfEvidence}
\begin{kfChoices}
\kfChoice Y에서는 데이터 일관성을 위해 쓰기 요청을 반드시 하나의 중앙 피어에서만 처리해야 한다.
\kfChoice X에서는 쓰기 요청이 마스터에 집중되므로 쓰기 경로 병목 가능성을 점검해야 한다.
\kfChoice X에서는 마스터 반영 후 슬레이브 전파 지연 동안 일시적 불일치가 발생할 수 있다.
\end{kfChoices}
\end{kfQBlock}

\kfSecHeader{kfAreaWeekly}{지문}{문항 2 지문 / 부정상대어}

\kfPassageNote{문항 2 지문 / 부정상대어}{%
(가) 가상 공간에서 영생의 삶을 누린다는 흥미로운 이야기가 SF 영화의 소재로만 그치지 않는다고 주장하는 입장이 있다. 커즈와일이나 보스트룀은 기억, 가치관, 태도, 정서적 성향처럼 인간의 자아를 구성하는 것들을 특정 정보 패턴으로 만들어 보존할 수 있다면 신체적인 죽음 이후에도 인간이 여전히 생존할 수 있다고 주장한다. 이 입장은 인격 동일성에 관한 심리적 연속성 이론과 관련이 있다. 이 이론에서는 ‘나’가 누구인지를 규정하는 것은 ‘나’의 기억, 신념, 생각, 감정, 희망, 두려움 등의 묶음이라고 보며, 이러한 심리적인 특성들의 집합이 유지·보존되느냐에 ‘나’의 지속 여부가 달려 있다고 말한다.

(중략)

커즈와일이나 보스트룀은 마음이나 정신 상태를 의미하는 심적 상태의 본성을 계산 기능주의의 관점에서 이해한다. 계산 기능주의에서는 심적 상태의 독특한 인과적 역할이 두뇌의 물리적 상태에 의해서 이루어진다고 보기 때문에 정신 작용을 두뇌에 구현된 계산 체계의 작동으로 이해하며, 인간의 사고나

(중략)
}
\begin{kfQBlock}{2}{}
\kfQStem{[08] (가)와 (나)를 바탕으로 <보기>에 대해 보인 반응으로 적절하지 않은 것은?}
\begin{kfEvidence}<보기> A는 정신을 물질적인 것의 작용에 불과하다고 말하며 정신적 속성이 물리적 속성에 수반한다고 주장한다. 이 입장에서 A는 원본 인간을 스캔해서 물리적 정보를 취득한 다음 스캔한 정보를 원격지로 전송해서 그곳의 재료로 복제 인간을 만듦으로써 인간을 원격 이동시키는 것이 가능하다고 본다. 이 기술에서 인간으로부터 읽어 내는 정보는 물리적 정보일 뿐이며 심적 상태를 읽어 내는 별도의 절차는 없다. A는 물리적 정보를 통하여 복제된 인간은 기억이나 의도를 포함하여 원본 인간이 가진 정신적 특성을 완벽히 보존한다고 말한다. B는 정보 처리에 관한 기능적 동일성이 심적 상태가 가지는 내용의 동일성을 보장하지 않는다고 말한다. 가령 ‘나’와 모든 신체 기관이 동일한 클론이 외계의 어느 별에서 성장했다고 가정하자. 이 클론이 ‘나’와 한 가지 다른 점은 호흡하는 과정에서 대기 중에 존재하는 가스의 종류를 구별해 낸다는 점이다. 만일 이 클론이 ‘나’와 지구에서 동시에 호흡하고 있다면 호흡에 관한 기능은 유사하지만 호흡을 통해 발생하는 심적 상태의 내용은 전혀 다르다.\end{kfEvidence}
\begin{kfChoices}
\kfChoice 계산 기능주의 입장에서는 A가 인간의 원격 이동이 가능하다고 보는 것에 대해 정신적 속성이 특정한 정보 구조로 만들어질 수 없다는 점을 간과한 것이라고 평가하겠군.
\kfChoice 계산 기능주의 입장에서는 A가 정신을 물질적인 것의 작용으로 보는 것에 대해 심적 상태가 두뇌의 물리적 상태와 관련이 있음을 인정한 입장이라고 평가하겠군.
\kfChoice 커즈와일은 A가 원격지의 복제 인간이 원본 인간의 정신적 특성을 완벽히 보존한다고 보는 것에 대해 정신 작용을 계산 체계의 작동으로 이해할 수 있는 입장이라는 점에서 동의하겠군.
\end{kfChoices}
\end{kfQBlock}

\kfSecHeader{kfAreaWeekly}{지문}{문항 3 지문 / 주체객체}

\kfPassageNote{문항 3 지문 / 주체객체}{%
우리는 외부 세계, 곧 물리적 실재에 둘러싸여 있다. 우리는 감각 기관을 활용하여 물리적 실재에 접촉하고 그에 대하여 지각한다. 물리적 실재가 지각 주체와 독립적으로 존재한다는 것은 자연 과학의 기본 가정이다. 인간의 감각 기관은 물리적 실재에 대한 온전한 정보를 제공해 주지 못하기 때문에 인간은 물리적 실재 개념을 추론을 통해서 구성하게 된다. 그런 점에서 물리적 실재 개념은 최종적일 수 없고 새로운 경험적 증거들이 도출됨에 따라 이것들을 설명하기 위해 개정되기도 한다. 뉴턴이 역학을 통해 물리적 실재 개념의 기초를 놓은 이후 전자기 현상에 대한 패러데이와 맥스웰의 연구에 의해 다른 종류의 물리적 실재 개념이 대두되었다.

(중략)

뉴턴의 물리적 실재 개념은 뉴턴이 역학 이론을 만들면서 공간, 시간, 입자라는 개념과 입자 간의 상호 작용인 힘 개념을 통해 제시되었다. 입자는 형태, 연장, 물성 등을 모두 상실하고 오로지 관성과 병진 운동만을 갖는 물체로 간주되었다. 그렇게 상정된 입자는 질점이

(중략)
}
\begin{kfQBlock}{3}{}
\kfQStem{<보기>를 참고할 때 두 관점 비교로 적절하지 않은 것은?}
\begin{kfEvidence}<보기> 연구자 A는 전기·자기 현상을 전하 사이 직접 작용으로만 해석하며, 공간 자체의 분포 상태는 이론적으로 중요하지 않다고 본다. 연구자 B는 공간 각 점의 장 값을 정의하고, 그 분포 변화를 다변수 방정식으로 기술해 현상을 설명한다.\end{kfEvidence}
\begin{kfChoices}
\kfChoice A의 관점은 원격 작용 중심 뉴턴주의 전자기 해석에 가깝다.
\kfChoice B의 관점은 장 개념을 배제하고 점입자 직접 작용만을 채택한다.
\end{kfChoices}
\end{kfQBlock}

\kfSecHeader{kfAreaWeekly}{지문}{문항 4 지문 / 순서방향}

\kfPassageNote{문항 4 지문 / 순서방향}{%
원자에는 중심에 핵이 있고, 핵 주위에 동심원처럼 퍼져 있는 특정 궤도들을 도는 전자들이 있다. 이 궤도마다 전자가 갖는 에너지 값이 정해져 있는데, 이 값을 에너지 준위라고 한다. 원자가 가진 특정한 몇 개의 궤도에만 전자가 존재하므로, 전자가 가질 수 있는 에너지 준위도 특정한 값으로 정해져 있다. 전자는 에너지를 얻으면 바깥쪽 궤도로, 에너지를 잃으면 안쪽 궤도로 이동한다. 원자가 하나일 때는 전자의 에너지 준위가 특정한 값이지만, 여러 원자가 모이면 여러 원자들의 궤도가 ⓐ중첩되면서 전자들이 서로 영향을 미쳐 에너지 준위에 변화가 생긴다. 이로 인해 에너지 준위는 특정한 값이 아닌 일정한 범위로 나타나게 된다. 이렇게 에너지 값이 범위로 표현된 것을 에너지 밴드라고 한다. 보통의 물질들은 수많은 원자로 구성돼 있기 때문에 대부분 에너지 밴드를 갖고 있다. 전자는 낮은 에너지 상태가 되려는 성질이 있으므로, 발광 물질의 전자에 에너지가 공급되면 전자가 바깥쪽으로 이동한 뒤 다시

(중략)
}
\begin{kfQBlock}{4}{}
\kfQStem{윗글의 Ⓐ와 <보기>의 Ⓑ를 비교하여 이해한 내용으로 적절하지 않은 것은?}
\begin{kfEvidence}태양은 파장이 250\textasciitilde{}2,500nm에 이르는 다양한 빛을 지상으로 보내는데, 실리콘 기반 태양 전지는 이 중 500\textasciitilde{}1,000nm의 빛만 활용할 수 있다. 파장이 1,000nm가 넘는 빛은 이 전지를 통과하고, 500nm 이하의 빛은 흡수되지만 열로 전환돼 날아가 버린다. 이 전지에서 낮은 파장의 빛이 대부분 열로 전환되는 이유는 어떤 물질이 매우 높은 에너지를 받았을 때, 물질이 넓은 에너지 밴드를 갖고 있다면 에너지를 받은 전자가 밴드 내에서 이동하면서 열에너지가 방출되기 때문이다. 하지만 물질이 특정한 에너지 준위를 갖고 있다면 상대적으로 적은 열을 발생하며 전자가 이동하여 전류가 발생한다. 이러한 장점을 가진 양자점은 크기를 조절하면 지상에 도달하는 태양 스펙트럼 중 대부분의 영역을 활용할 수 있기 때문에 Ⓑ양자점 기반 태양 전지가 개발되고 있다. 한편, 양자점의 코어에 셸을 씌우면 빛과 달리 전류는 셸을 뚫고 나오기 힘들므로 양자점 기반 태양 전지에서는 셸을 사용하지 않는 경우가 많다. 따라서 코어 자체를 더 안정한 소재로 만들거나 리간드를 더 붙여 표면을 안정시켜야 한다.\end{kfEvidence}
\begin{kfChoices}
\kfChoice Ⓐ와 Ⓑ 모두 양자점의 전자가 넓은 에너지 밴드가 아닌 매우 좁은 범위의 에너지 준위를 갖고 있다는 점을 활용하는군.
\kfChoice 같은 양의 에너지가 공급될 경우, Ⓐ는 LCD 디스플레이보다, Ⓑ는 실리콘 기반 태양 전지보다 더 많은 열에너지를 방출하겠군.
\end{kfChoices}
\end{kfQBlock}

\kfSecHeader{kfAreaWeekly}{지문}{문항 5 지문 / 인과도치}

\kfPassageNote{문항 5 지문 / 인과도치}{%
리더는 목표 달성을 위해 구성원들을 이끌고 그들에게 영향력을 미치는 조직의 핵심 인물이다. 오래전부터 리더로서의 능력을 나타내는 리더십은 무엇이며, 어떠한 리더십이 효과적인지에 관한 연구가 행해졌다. 피들러는 리더를 업무의 목표 달성을 중시하는 업무 지향적 스타일을 가진 리더와 구성원들과의 관계, 소통 등을 중시하는 관계 지향적 스타일을 가진 리더로 구분하였다. 그리고 리더가 자신의 집단에 영향을 미칠 수 있는 정도를 나타내는 상황의 호의성이 ‘리더와 구성원들의 관계, 업무 구조, 리더의 권한’이라는 세 가지 변수로 결정된다고 보았다. 즉 리더와 구성원들이 서로 믿고 존중하는 관계일수록, 업무가 명확하게 구조화되어 있을수록, 채용 및 승진 등에 대한 리더의 공식적 권한이 클수록 상황의 호의성이 높고, 반대의 경우는 호의성이 낮다고 본 것이다. 그는 상황의 호의성에 따라 적합한 리더십의 종류가 다르다고 보았는데, 상황의 호의성이 매우 낮거나 매우 높을 때는 업무 지향적 리더가, 중간 

(중략)
}
\begin{kfQBlock}{5}{}
\kfQStem{윗글을 바탕으로 <보기>를 이해한 반응으로 적절하지 않은 것은?}
\begin{kfEvidence}<보기>\textbackslash{}n갑이 경영하고 있는 기업 K는 얼마 전 회사에 오랫동안 근무하던 구성원들이 대거 퇴직하여 대부분이 입사한 지 오래되지 않은 이들로 구성되어 있다. 이들은 업무에 임하는 태도는 매우 적극적이지만, 아직 업무에 숙달되지 않은 상태이다. 한편 을이 경영하고 있는 기업 P는 업무에 숙달된 구성원들이 매우 많지만 전반적으로 구성원들이 자신감이 부족하고 근로 의욕이 저하된 상태이다.\textbackslash{}n갑은 각 구성원이 수행해야 할 업무와 성과에 따른 보상과 처벌을 명확히 정하여 제시하였다. 을은 구성원에게 기업이 추구하는 목표를 제시하고 구성원들이 회사의 의사 결정에 적극적으로 참여할 수 있게 하였다. 그리고 목표를 달성한 구성원에게 추가 수당이라는 보상을 제공하였다. 기업 K의 구성원들은 대부분 목표를 달성하지 못했고, 기업 P의 구성원들은 대부분 목표를 쉽게 달성하였다.\end{kfEvidence}
\begin{kfChoices}
\kfChoice 허시와 블랜차드는 갑이 구성원들과 인간적 관계를 형성하면서도 명확한 지시를 내려야 한다고 보겠군.
\kfChoice 데시와 라이언은 기업 P의 구성원들이 목표를 달성한 이유를 을이 구성원에게 제공한 보상 때문이라고 여기겠군.
\end{kfChoices}
\end{kfQBlock}

\kfSecHeader{kfAreaWeekly}{지문}{문항 6 지문 / 의도/목적}

\kfPassageNote{문항 6 지문 / 의도/목적}{%
무선 통신에 대한 수요가 폭발적으로 증가함에 따라 미국 정부는 \uline{ⓐ경매}를 통해 통신 사업자들에게 통신 주파수 대역들을 판매하게 되었다. 경매가 시행되기 전에는 \uline{ⓑ심사}나 \uline{ⓒ제비뽑기}와 같은 방식을 활용하였다. 심사 방식은 여러 통신 사업자들이 사업 계획서를 제출하면 평가자들이 사업 계획서에 근거하여 무선 통신 사업을 가장 잘 운영할 만한 사업자를 선정하고, 그렇게 선정된 사업자가 주파수를 할당받는 것이었다. 그러나 사업자가 고용한 법률가들과 로비스트들이 서로 자기들의 사업 계획이 최선이라고 주장하는 상황에서 사업 계획서에 대한 평가 절차와 이의 제기 과정을 전부 거쳐 주파수를 할당하기까지 수년이 걸리기도 하였다.

1982년 무렵에는 더 이상 심사 방식으로 수많은 주파수를 할당하는 것이 어려운 수준에 이르렀고, 결국 의회는 제비뽑기로 주파수를 판매하는 것에 동의하게 되었다. 제비뽑기 방식은 주파수를 할당하는 데 드는 시간과 노력을 대폭 줄여 주었지만 …(중략)
}
\begin{kfQBlock}{6}{}
\kfQStem{윗글에 대한 이해로 가장 적절하지 않은 것은?}
\begin{kfChoices}
\kfChoice 동시 다중 라운드 경매는 라운드별 정보 공개를 통해 참여자의 전략 수정을 가능하게 했다.
\kfChoice 동시 다중 라운드 경매는 라운드별 정보 공개를 통해 참여자의 전략 수정을 가능하게 했다.
\kfChoice 심사 방식은 주파수 대량 할당에 시간이 거의 들지 않았다.
\kfChoice 제비뽑기 방식에서는 투기 목적 참가가 구조적으로 차단되었다.
\end{kfChoices}
\end{kfQBlock}

\kfSecHeader{kfAreaWeekly}{지문}{문항 7 지문 / 상관관계 반전}

\kfPassageNote{문항 7 지문 / 상관관계 반전}{%
1908년 조지 하디와 빌헬름 바인베르크는 멘델 유전 법칙에 기초한 수학적 모델을 각각 독립적으로 제시했다. 하디는 집단 내 대립 유전자 A와 a의 빈도를 p와 q로 두고 무작위 교배에서 다음 세대의 유전자형 빈도를 분석했다. 바인베르크도 무작위 결합 조건에서 유전자 빈도가 세대를 거쳐 안정될 수 있음을 통계 자료와 함께 제시했다. 이 두 접근이 결합되어 하디-바인베르크 평형 법칙이 정립되었다.

하디-바인베르크 평형은 \uline{㉠집단이 충분히 크고}, \uline{㉡무작위 교배가 일어나며}, \uline{㉢새 돌연변이가 없고}, \uline{㉣이입·이출 같은 이주가 없고}, \uline{㉤자연 선택이 작용하지 않을 때} 유지된다. [A] 예를 들어 p=0.7, q=0.3이면 p+q=1이고 유전자형 빈도는 AA=p²=0.49, Aa=2pq=0.42, aa=q²=0.09로 계산되며 p²+2pq+q²=1이 성립한다. 유전자 빈도가 유지되면 다음 세대에서도 유전자형 빈도는 일정하게 유지된다. …(중략)
}
\begin{kfQBlock}{7}{}
\kfQStem{윗글을 이해한 내용으로 적절한 것은?}
\begin{kfChoices}
\kfChoice 작은 집단일수록 우연 변동 영향이 줄어 평형 유지가 쉽다.
\kfChoice 작은 집단일수록 우연 변동 영향이 줄어 평형 유지가 쉽다.
\kfChoice 유전자 빈도는 특정 대립 유전자 수/전체 유전자 수로 계산한다.
\kfChoice Aa 빈도는 2pq로 계산할 수 있다.
\end{kfChoices}
\end{kfQBlock}

\kfSecHeader{kfAreaWeekly}{지문}{문항 8 지문 / 필요조건}

\kfPassageNote{문항 8 지문 / 필요조건}{%
최근까지 일반적인 컴퓨터 환경에서는 운영 체제를 제어하거나 사용자 인터페이스를 처리하고, 문서를 작성하거나 웹을 탐색하는 등 다양한 범용 작업이 주요 연산 대상이었다. 이러한 작업은 명령어의 종류가 많고 그 순서가 수시로 바뀌며, 연산 간 종속성이 크다는 특징이 있다. 컴퓨터의 연산은 주로 여러 개의 코어로 구성된 \uline{㉠중앙 처리 장치(CPU)}에서 수행된다. 코어는 컴퓨터의 두뇌 역할을 하는 것으로 명령어를 해석하고 각종 연산을 실행한다. CPU는 여러 개의 고성능 코어로 구성되어 복잡한 제어 흐름을 정밀하게 조율할 수 있도록 설계되었다. 그러나 인공 지능 학습, 고해상도 영상 처리, 과학 기술 관련 계산 등 대규모 데이터를 대상으로 동일한 연산을 반복 수행해야 하는 작업이 증가하면서 기존의 연산 환경에 변화가 생겼다. CPU는 코어 수에 한계가 있고 개별 코어가 복잡한 제어 흐름에 맞추어 다양한 연산을 처리하도록 설계되어 있기 때문에 동일한 연산을 대량으로 반복 수행해야 …(중략)
}
\begin{kfQBlock}{8}{}
\kfQStem{<보기>를 바탕으로 GPU 작동을 이해한 내용으로 적절한 것은?}
\begin{kfEvidence}<보기> 연산 유닛 A에는 스레드 블록 4개, B와 C에는 각각 1개가 배치되어 있다. 블록 내부 스레드는 동일 명령어를 실행하며, 같은 메모리 주소를 동시에 접근하는 작업이 일부 포함되어 있다.\end{kfEvidence}
\begin{kfChoices}
\kfChoice 전역 메모리는 연산 중 접근 자체가 불가능하므로 결과 저장 외에는 절대 사용할 수 없다.
\kfChoice 전역 메모리는 연산 중 접근 자체가 불가능하므로 결과 저장 외에는 절대 사용할 수 없다.
\kfChoice A에 작업이 과도하게 몰리면 B·C 유닛보다 대기 가능성이 커질 수 있다.
\kfChoice 블록 내부에서 동일 주소 동시 접근이 많으면 메모리 충돌로 속도가 떨어질 수 있다.
\end{kfChoices}
\end{kfQBlock}

\kfSecHeader{kfAreaWeekly}{지문}{문항 9 지문 / 결합/분리}

\kfPassageNote{문항 9 지문 / 결합/분리}{%
(가) 가상 공간에서 영생의 삶을 누린다는 흥미로운 이야기가 SF 영화의 소재로만 그치지 않는다고 주장하는 입장이 있다. 커즈와일이나 보스트룀은 기억, 가치관, 태도, 정서적 성향처럼 인간의 자아를 구성하는 것들을 특정 정보 패턴으로 만들어 보존할 수 있다면 신체적인 죽음 이후에도 인간이 여전히 생존할 수 있다고 주장한다. 이 입장은 인격 동일성에 관한 심리적 연속성 이론과 관련이 있다. 이 이론에서는 ‘나’가 누구인지를 규정하는 것은 ‘나’의 기억, 신념, 생각, 감정, 희망, 두려움 등의 묶음이라고 보며, 이러한 심리적인 특성들의 집합이 유지·보존되느냐에 ‘나’의 지속 여부가 달려 있다고 말한다.

커즈와일이나 보스트룀은 마음이나 정신 상태를 의미하는 심적 상태의 본성을 계산 기능주의의 관점에서 이해한다. 계산 기능주의에서는 심적 상태의 독특한 인과적 역할이 두뇌의 물리적 상태에 의해서 이루어진다고 보기 때문에 정신 작용을 두뇌에 구현된 계산 체계의 작동으로 이해하며, 인간의 사…(중략)
}
\begin{kfQBlock}{9}{}
\kfQStem{윗글의 내용과 일치하는 것으로 가장 적절하지 않은 것은?}
\begin{kfChoices}
\kfChoice 복수 실현 가능성을 인정해도 물리적 동일성의 자동 보장은 성립하지 않는다고 본다.
\kfChoice 복수 실현 가능성을 인정해도 물리적 동일성의 자동 보장은 성립하지 않는다고 본다.
\kfChoice 계산 기능주의는 정신 상태를 물질 구성과 무관한 순수 도덕 범주로 본다.
\kfChoice 체화된 인지는 신체 종류가 달라도 개념 이해 방식은 항상 같다고 본다.
\end{kfChoices}
\end{kfQBlock}

\kfSecHeader{kfAreaWeekly}{지문}{문항 10 지문 / 조건 왜곡}

\kfPassageNote{문항 10 지문 / 조건 왜곡}{%
프랑스의 기호학자이자 문예 비평가인 롤랑 바르트는 대중 매체의 이미지와 텍스트가 단순한 오락이나 소비 대상이 아니라 특정 사회·문화적 의미를 담고 있으며, 그 수용 과정에서 특정 이념이 대중에게 자연스럽게 받아들여진다고 보았다. 그는 이러한 현상을 ‘현대의 신화’라고 명명하고, 일상 속 기호에 숨어 있는 이념과 권력의 작동 방식을 분석하고자 하였다.

바르트는 소쉬르의 이론을 분석 도구로 활용하였다. 소쉬르는 언어를 기호로 보았고, 기호는 음성이나 문자 등 감각적으로 지각 가능한 형태인 ‘기표’와 그것이 가리키는 개념적 내용인 ‘기의’라는 두 가지 요소로 이루어진다고 보았다. 예를 들어 ‘나무’라는 말은, 말소리 혹은 문자 형태로 주어지는 기표와 ‘줄기와 가지를 가진 목질 식물’이라는 기의가 결합한 하나의 기호이다. 소쉬르는 이러한 기호의 구조에서 기표와 기의 사이의 관계가 본질적으로 자의적이라는 점을 강조하였다. 즉 특정 기표가 반드시 특정 기의를 가리켜야 할 필연성은 없으며, 다…(중략)
}
\begin{kfQBlock}{10}{}
\kfQStem{바르트의 관점에서 <보기>를 이해한 내용으로 옳은 것은?}
\begin{kfEvidence}<보기> 한 식품 회사는 ‘제로 슈거’ 음료 광고를 만들었다. 광고는 운동하는 젊은 모델, 하얀 배경, ‘가벼움’ ‘자기관리’라는 문구를 반복적으로 배치하며, 이 음료를 마시는 생활이 곧 건강하고 바람직한 삶의 방식인 것처럼 보이게 한다. 많은 소비자는 이 광고를 보고 ‘이 음료를 선택하는 것이 곧 자기관리의 증거’라고 느꼈다.\end{kfEvidence}
\begin{kfChoices}
\kfChoice 광고는 ‘가벼움/자기관리’의 가치가 자연스러운 것처럼 느껴지게 만들 수 있어, 특정 생활양식을 보편적 기준처럼 자연화할 위험이 있다.
\kfChoice 광고는 ‘가벼움/자기관리’의 가치가 자연스러운 것처럼 느껴지게 만들 수 있어, 특정 생활양식을 보편적 기준처럼 자연화할 위험이 있다.
\kfChoice 광고는 ‘제로 슈거’라는 객관적 사실만 전달하므로 2차 의미는 형성되지 않는다.
\kfChoice 광고는 이념을 강압적으로 주입하므로 바르트가 말한 신화와는 무관하다.
\end{kfChoices}
\end{kfQBlock}

\kfSecHeader{kfAreaWeekly}{지문}{문항 11 지문 / 긍부정 반전}

\kfPassageNote{문항 11 지문 / 긍부정 반전}{%
\#\#\# [32～34] 다음 글을 읽고 물음에 답하시오.

(가)   청강 녹초변에 소 먹이는 아이들이   석양에 흥이 겨워 피리를 빗기 부니   물 아래 잠긴 용이 잠 깨어 일어날 듯   내 기운에 나온 학이 제 깃을 던져 두고 반공에 솟아 뜰 듯   소선(蘇仙)* 적벽은 추칠월이 좋다 하되   팔월 십오야를 모두 어찌 칭찬하는가   구름이 걷히고 물결이 다 잔 적에   하늘에 돋은 달이 솔 위에 걸렸거든   잡다가 빠진 줄이 적선(謫仙)\textit{이 헌사할샤}   \textit{공산에 쌓인 잎을 삭풍이 거둬 불어}   \textit{떼구름 거느리고 눈조차 몰아오니}   \textit{천공이 호사로워 옥으로 꽃을 지어}   \textit{만수천림을 꾸며곰 낼세이고}   \textit{앞 여울 가리 얼어 독목교(獨木橋) 비꼈는데}   \textit{막대 멘 늙은 중이 어느 절로 간단 말고}   \textit{산옹의 이 부귀를 남더러 자랑 마오}   \textit{경요굴(瓊瑤窟)} 숨은 세계 찾을 이 있을세라   산중에 벗이 없어 서책을 쌓아 두고   만고 인물을 거슬러 혜여하니   성현도 많…(중략)
}
\begin{kfQBlock}{11}{}
\kfQStem{[A]에 대한 이해로 적절한 것은?}
\begin{kfChoices}
\kfChoice 하늘의 이치가 제대로 구현되지 못했음을 ‘시운’의 ‘흥망’에서 발견하고도 모를 일이 많다고 한 것에는, 인물의 담담한 태도가, 이상에 미치지 못하는 현실을 수용하는 것을 통해 드러난다.
\kfChoice 하늘의 이치가 제대로 구현되지 못했음을 ‘시운’의 ‘흥망’에서 발견하고도 모를 일이 많다고 한 것에는, 인물의 담담한 태도가, 이상에 미치지 못하는 현실을 수용하는 것을 통해 드러난다.
\kfChoice ‘삭풍’이 가을 잎을 쓸고 간 자리에 구름을 불러와 ‘공산’을 눈 세상으로 만들었다고 한 것에는, 인물이 거처한 공간의 아름다움에 대한 인식이 계절에 따른 자연의 변화를 통해 드러난다.
\kfChoice ‘앞 여울’을 건너가는 노승을 발견하고 ‘경요굴’이 들키지 않기를 바라는 것에는, 빼어난 경치를 소중하게 여기는 태도가, 숨어 있는 세계가 알려질 것에 대한 염려를 통해 드러난다.
\end{kfChoices}
\end{kfQBlock}

\kfSecHeader{kfAreaWeekly}{지문}{문항 12 지문 / 비교대상 뒤바꿈}

\kfPassageNote{문항 12 지문 / 비교대상 뒤바꿈}{%
\#\#\# [39\textasciitilde{}42] 다음 글을 읽고 물음에 답하시오.

(가)   苦忘亂抽書 잊기를 자주 하여 어지러이 뽑아 놓은 책들   散漫還復整 흩어진 걸 다시 또 정리하자니   曜靈忽西頹 해는 문득 서쪽으로 기울고   江光搖林影 강 위에 숲 그림자 흔들린다.   扶筇下中庭 막대 짚고 마당 가운데 내려서서   矯首望雲嶺 고개 들어 구름 낀 고개 바라보니   漠漠炊烟生 아득히 밥 짓는 연기가 피어나고   蕭蕭原野冷 쓸쓸히 들판은 서늘하구나.   田家近秋穫 농삿집 가을걷이 가까워지니   喜色動臼井 절구질 우물가에 기쁜 빛 돌아   鴉還天機熟 갈까마귀 돌아오니 절기가 무르익고   鷺立風標逈 해오라기 서 있는 모습 우뚝하고 훤하다.   我生獨何爲 내 인생은 홀로 무얼 하는 것인지   宿願久相梗 숙원이 오래도록 풀리질 않네.   無人語此懷 이 회포 털어놓을 사람 아무도 없어   搖琴彈夜靜 거문고만 둥둥 탄다, 고요한 밤에.

* 이황, ｢만보(晩步)｣ -

(나)   밤이다. ㉠ 하늘은 푸르다 못해 …(중략)
}
\begin{kfQBlock}{12}{}
\kfQStem{(가)와 (나)의 공통점으로 가장 적절한 것은?}
\begin{kfChoices}
\kfChoice 시간적 배경의 의미를 활용하여 내적 갈등을 드러내고 있다.
\kfChoice 시간적 배경의 의미를 활용하여 내적 갈등을 드러내고 있다.
\kfChoice 공간의 대비를 통해 일상의 공간에 의미를 부여하고 있다.
\kfChoice 대상과의 문답을 통해 삶에 대한 깨달음을 드러내고 있다.
\end{kfChoices}
\end{kfQBlock}

\kfSecHeader{kfAreaWeekly}{지문}{문항 13 지문 / 주체/객체 혼동}

\kfPassageNote{문항 13 지문 / 주체/객체 혼동}{%
\#\#\# [43 \textasciitilde{} 45] 다음 글을 읽고 물음에 답하시오.

[앞부분의 줄거리] 당나라 때 한림학사가 된 장사운은 옥란과 옥계라는 두 딸이 있었는데, 이 중 옥란을 송 시랑의 아들과 정혼시킨다. 권세를 잡고 있던 환관 강환은 이를 알고도 자신의 아들과 옥란을 강제로 혼인시키려 한다.

장 학사가 강환이 권신이라 독한 마음을 먹을 것을 염려하여 외면으로 말하기를,   “하방 천한 출생에게 대황문족이 구혼하니 감히 감당하지 못하여 허락지 못하오며, 또 이전에 송 시랑과 선약하였으니 이 역시 못할 일이옵니다.”   강환이 그 말을 듣고 크게 노하여,   “저는 서촉의 천한 출생이면서 천행으로 등과하였으면서 남을 환관이라고 업신여기는도다. 저의 생사 내게 맡겼거든 내 말을 어찌 멀리하리오. 소위 송 시랑을 먼저 처치하리라.”   하고 즉시 탑전에 들어가 천자에게 참소하여 송 시랑을 의금부의 신문에 부치니, 송 시랑이 너무 당황하여 어찌할 줄을 모르고 옥중에 내려가 장 학사에게 기별하니, 학…(중략)
}
\begin{kfQBlock}{13}{}
\kfQStem{[A]와 [B]에 대한 설명으로 가장 적절한 것은?}
\begin{kfChoices}
\kfChoice [A]를 통해 전달된 상황에 대한 정보는 [B]를 계기로 신뢰성을 의심받게 된다.
\kfChoice [A]를 통해 전달된 상황에 대한 정보는 [B]를 계기로 신뢰성을 의심받게 된다.
\kfChoice [A]를 통해 전달된 상황의 당위성은 [B]를 계기로 드러나게 된 사실로 강화된다.
\kfChoice [A]를 통해 전달된 인물에 대한 평가는 [B]를 계기로 구체성을 확보하게 된다.
\end{kfChoices}
\end{kfQBlock}

\kfSecHeader{kfAreaWeekly}{지문}{문항 14 지문 / 내용과 방법 혼동}

\kfPassageNote{문항 14 지문 / 내용과 방법 혼동}{%
\#\#\# [27 \textasciitilde{} 30] 다음 글을 읽고 물음에 답하시오.

아버지 입에서 그 집을 사자고 처음 이야기가 나왔을 때, 어머니는 달갑잖은 표정이었다. 직접 가서 집을 둘러볼 때도 마찬가지 표정이었고, 형 역시 못마땅한 내색을 감추지 않았다. 그러나 우리는 모두 지쳐 있었다. 여러 해를 여기저기 잠깐씩 남의 집만 전전하며 살아왔기 때문이다. 전에 우리가 살던, 크고 좋은 집은 아버지 친구 되는 사람한테 어처구니없이 빼앗겨 버렸다. 굉장히 똑똑한 사람이라고 소문난, 아버지의 고향 친구 심 씨였다. 아버지가 사업에 실패하여 한참 빚쟁이들한테 몰리고 있을 때 그 심 씨가 나타나서 묘책을 일러 주었다. 친구의 권고에 따라 아버지는 집문서와 인감도장을 내주었다. 집을 심 씨의 소유로 위장하여 그거라도 건져 보자는 속셈이었다. 곤경에 빠졌을 때 찾아와서 위로하고 충고해 주는 고향 친구가 아버지한테는 친형제만큼이나 살가웠을 것이다. 도장을 넘겨주면서 아버지는 심 씨의 손을 붙잡고 고맙다는 인사를 수…(중략)
}
\begin{kfQBlock}{14}{}
\kfQStem{[A]의 서술상 특징으로 가장 적절하지 않은 것은?}
\begin{kfChoices}
\kfChoice 회상 장면을 삽입하여 인물들이 처한 상황의 이유를 드러내고 있다.
\kfChoice 회상 장면을 삽입하여 인물들이 처한 상황의 이유를 드러내고 있다.
\kfChoice 서술자의 추측을 제시하여 앞으로 일어날 사건을 예고하고 있다.
\kfChoice 공간의 이동에 따라 서술자를 달리하여 사건을 입체적으로 바라보고 있다.
\end{kfChoices}
\end{kfQBlock}

\kfSecHeader{kfAreaWeekly}{지문}{문항 15 지문 / 내용 왜곡}

\kfPassageNote{문항 15 지문 / 내용 왜곡}{%
\#\#\# [39 \textasciitilde{} 42] 다음 글을 읽고 물음에 답하시오.

[앞부분 줄거리] 태종 황제에게 원한을 품은 용왕으로 인해 태종 황제는 죽어서 저승에 가게 된다. 신하 위징은 저승의 지인 최 판관에게 태종 황제를 도와 달라 부탁한다.

십전 명왕이 태종 황제를 향하여 말하기를,   “폐하가 경하 용왕으로 빌미되어 인간의 귀하신 몸이 더러운 곳에 욕되이 몸소 오시니 오랫동안 머무시는 것이 불감당(不堪當)하옵니다.”   태종 황제께서 사양하여 말하기를,   “십전 명왕의 덕으로 명백히 처단하여 주시니 감사하옵기 측량없사오며, 또한 명부(冥府)를 구경하오니 은혜에 감사하는 마음이 황공하여이다.”   염라대왕이 최 판관에게 명하기를,   “태종 황제의 수한(壽限)이 얼마 되시는고? 생사부(生死簿)를 가져오라.”   하였다.   최 판관이 명령을 듣고 생사부가 쌓인 곳에 가서 책을 펴 보니 인간 세상의 세민 황제께서 즉위한 후 정관 13년에 들어올 해이매 다시 살아날 도리가 없었다. 마음에 민망…(중략)
}
\begin{kfQBlock}{15}{}
\kfQStem{㉠과 ㉡에 대한 이해로 가장 적절한 것은?}
\begin{kfChoices}
\kfChoice ㉠은 인물이 자신의 과거 행적을 환기하는 장소이고, ㉡은 인물이 타인의 과거 행적의 결과를 확인하는 장소이다.
\kfChoice ㉠은 인물이 자신의 과거 행적을 환기하는 장소이고, ㉡은 인물이 타인의 과거 행적의 결과를 확인하는 장소이다.
\kfChoice ㉠은 인물 간의 오해가 발생하는 장소이고, ㉡은 인물 간의 오해가 해소되는 장소이다.
\kfChoice ㉠은 인물이 타인의 환심을 사려는 장소이고, ㉡은 인물이 상대의 호의를 거절하는 장소이다.
\end{kfChoices}
\end{kfQBlock}

\kfSecHeader{kfAreaWeekly}{지문}{문항 16 지문 / 과잉 해석}

\kfPassageNote{문항 16 지문 / 과잉 해석}{%
\#\#\# [38 \textasciitilde{} 41] 다음을 읽고 물음에 답하시오.

“아니, 작은 것 한 장도 못 되는 돈 갖고 이 바닥에서 독채 전세를 얻겠다고?”   그러더니 다시 한바탕 해소라도 발작한 것같이 급하게 웃었다. 거금 구십만원을 작은 것 한 장도 안 된다니, 이 노인이 귀가 좀 어두운가 해서 나는 다시 목청을 돋우어 구십만원을 강조했다.   그래도 노인은 탁하고 급한 웃음을 멎을 척도 안 했다. 사무실 앞에 승용차가 나란히 두 대가 멎더니 부인들과 신사들이 섞인 한 떼가 안으로 들이닥쳤다. 이곳도 결코 파리 날리는 한가한 곳이 아니었던 것이다.   “사모님, 지금 보신 그 땅 눈 꽉 감고 잡아놓으십시다. 글쎄 문제없다니까요. 중도금 치르기 전에 평당 오천원 띠기는 누워서 떡 먹기라니까요.”   젊은 신사들이 부인들을 꾀고 노인도 합세했다.   “우리하고 손잡고 이 바닥에서 큰돈 잡은 사모님네들 숱합니다, 숱해.”   나는 그들에게 완전히 잊혀졌다. 영아 기저귀를 갈아주고 다시 업고 나올 때까…(중략)
}
\begin{kfQBlock}{16}{}
\kfQStem{<보기>를 참고하여 윗글을 감상한 내용으로 적절하지 않은 것은? [3점]}
\begin{kfChoices}
\kfChoice ‘나’가 ‘친정 나들이’를 갈 때 주인여자에게 남편이 하는 말을 듣고 ‘절망감을 느’끼는 부분에서, 경제적 격차를 인지하지 못하고 가족의 정착만을 중시했던 태도를 후회하는 것을 확인할 수 있군.
\kfChoice ‘나’가 ‘친정 나들이’를 갈 때 주인여자에게 남편이 하는 말을 듣고 ‘절망감을 느’끼는 부분에서, 경제적 격차를 인지하지 못하고 가족의 정착만을 중시했던 태도를 후회하는 것을 확인할 수 있군.
\kfChoice ‘승용차가 나란히 두 대가 멎’은 후 거기서 내린 젊은 신사들이 ‘그 땅’에 대해 말하는 부분에서, 부동산을 부의 축적 수단으로 인식하던 세태를 짐작할 수 있군.
\kfChoice ‘나’가 ‘첫번째 본 집’을 나와서 ‘셋방살이 신세가 무슨 대역죄’냐고 생각하는 부분에서, 주거 공간을 얻는 과정에서 마주한 현실이 부당하다고 느끼는 것을 짐작할 수 있군.
\end{kfChoices}
\end{kfQBlock}

\kfSecHeader{kfAreaWeekly}{지문}{문항 17 지문 / 범위 오류}

\kfPassageNote{문항 17 지문 / 범위 오류}{%
\#\#\# [25\textasciitilde{}27] 다음 글을 읽고 물음에 답하시오.

[앞부분의 줄거리] 영웅적 면모를 지녔으나 먹고 자기만 하며 허송세월을 보내던 이생은 장인이 죽자 장모에 의해 쫓겨나고, 이생의 부인 양 소저는 시부모의 제사를 지내면서 이생을 기다린다. 한편 집을 나온 이생은 곤경에 빠진 사람에게 자신이 가진 삼백 냥의 돈을 다 내어 준 후 어떤 노인의 도움을 받게 된다.

“그대 오늘 큰 적선을 하였으니 깊이 감격하노라.”   이생이 이렇듯 노인의 신기함을 보고 평범한 인물이 아니리라 생각하며 의아해 마지않았다.   “존옹의 물으심이 무슨 일이니이꼬? 저는 적선한 일이 없소이다.”   노인 말하기를,   “대인은 사람 속이기를 아니하나니라.”   이에 또 말하기를,   “그대 저리 먹는 양에 양식도 없이 어찌 살아가려 하느뇨?”   이생이 답하기를,   “이처럼 얻어먹으면 아니 살아가리이까?”   노인이 크게 웃으며 말하였다.   “㉠ 소년의 말이 사리에 어둡도다. 그러나 나는 그대를 알거…(중략)
}
\begin{kfQBlock}{17}{}
\kfQStem{㉠∼㉣에 대한 설명으로 가장 적절하지 않은 것은?}
\begin{kfChoices}
\kfChoice ㉢에서는 과거의 일을 근거로 자책하고 있고, ㉣에서는 미래의 일을 예측하며 상대를 만류하고 있다.
\kfChoice ㉢에서는 과거의 일을 근거로 자책하고 있고, ㉣에서는 미래의 일을 예측하며 상대를 만류하고 있다.
\kfChoice ㉠에서는 당위를 내세워 상대를 설득하고, ㉡에서는 감정에 호소하여 상대의 동의를 이끌어 내고 있다.
\kfChoice ㉠에서는 상대보다 우월한 신분을 통해, ㉣에서는 상대에 대한 배려를 통해 상대의 태도 변화를 촉구하고 있다.
\end{kfChoices}
\end{kfQBlock}

\kfSecHeader{kfAreaWeekly}{지문}{문항 18 지문 / 시점/화자 혼동}

\kfPassageNote{문항 18 지문 / 시점/화자 혼동}{%
\#\#\# [24 \textasciitilde{} 26] 다음 글을 읽고 물음에 답하시오.

(가)   팽이가 돈다   어린아해이고 어른이고 살아가는 것이 신기로워   물끄러미 보고 있기를 좋아하는 나의 너무 큰 눈 앞에서   아이가 팽이를 돌린다 / 살림을 사는 아해들도 아름다웁듯이   노는 아해도 아름다워 보인다고 생각하면서   손님으로 온 나는 ㉠ 이 집 주인과의 이야기도 잊어버리고   또 한번 팽이를 돌려주었으면 하고 원하는 것이다   도회 안에서 쫓겨다니는 듯이 사는   나의 일이며 / 어느 소설보다도 신기로운 나의 생활이며   모두 다 내던지고   점잖이 앉은 나의 나이와 나이가 준 나의 무게를 생각하면서   ㉡ 정말 속임 없는 눈으로 / 지금 팽이가 도는 것을 본다   그러면 팽이가 까맣게 변하여 서서 있는 것이다   누구 집을 가보아도 나 사는 곳보다는 여유가 있고   바쁘지도 않으니 / 마치 별세계(別世界)같이 보인다   팽이가 돈다 / 팽이가 돈다   팽이 밑바닥에 끈을 돌려 매이니 이상하고   …(중략)
}
\begin{kfQBlock}{18}{}
\kfQStem{㉠ \textasciitilde{} ㉤에 대한 이해로 적절하지 않은 것은?}
\begin{kfChoices}
\kfChoice ㉡ : 자신의 삶에 대한 의구심을 버리고 사물의 본질을 객관적으로 파악하려 하는 화자의 태도를 드러내고 있다.
\kfChoice ㉡ : 자신의 삶에 대한 의구심을 버리고 사물의 본질을 객관적으로 파악하려 하는 화자의 태도를 드러내고 있다.
\kfChoice ㉠ : ‘주인’과 이야기하는 것도 잊을 만큼 팽이를 바라보는 일에 열중하는 화자의 모습을 형상화하고 있다.
\kfChoice ㉢ : 팽이 앞에서 자신의 내면을 드러내는 화자의 모습을 형상화하고 있다.
\end{kfChoices}
\end{kfQBlock}



\clearpage

\kfChapterCover{2}{비트겐슈타인3 ch02}{Chapter 2}{이번 챕터: 문법 → 문학 5작품 → 비문학 5지문 → 패턴 워크북.}

\kfStartGrammar{이번 챕터 문법 영역 — 음운 / 음절의 끝소리 규칙 + 음절 구조 제약}


\kfSecHeader{kfAreaGrammar}{문제}{}

\begin{multicols}{2}\raggedcolumns\setlength{\columnsep}{12pt}
\kfPassageFull{문항 1 지문 (concept)}{%
「국어에서 음절 끝(종성)에 올 수 있는 자음은 'ㄱ, ㄴ, ㄷ, ㄹ, ㅁ, ㅂ, ㅇ'의 7개로 제한된다. 이를 음절의 끝소리 규칙이라 한다. 종성에서 이 7개 이외의 자음은 이 7개 중 하나로 바뀌어 발음된다. 예를 들어, 'ㅅ, ㅆ, ㅈ, ㅊ, ㅌ'은 모두 [ㄷ]으로, 'ㅋ, ㄲ'은 [ㄱ]으로, 'ㅍ'은 [ㅂ]으로 발음된다. 겹받침의 경우, 뒤에 모음으로 시작하는 조사·어미가 오지 않으면 둘 중 하나만 발음한다. 'ㄳ, ㄵ, ㄼ, ㄽ, ㄾ, ㅀ'은 앞의 자음을, 'ㄺ, ㄻ, ㄿ'은 뒤의 자음을 발음하는 것이 원칙이다. 단, 'ㄺ'은 '밟다' 등 일부 용언에서 [ㄹ]이 아닌 [ㄱ]으로 발음되기도 한다.」
}
\begin{kfQBlock}{1}{}
\kfQStem{윗글을 바탕으로 <보기>의 발음을 탐구한 내용으로 적절하지 않은 것은?

<보기> ㄱ. 꽃[꼳]     ㄴ. 부엌[부억]     ㄷ. 닭[닥] ㄹ. 넋[넉]     ㅁ. 삶[삼]}
\begin{kfChoices}
\kfChoice ㄱ: 'ㅊ'은 종성에서 발음될 수 없으므로 [ㄷ]으로 교체된다.
\kfChoice ㄴ: 'ㅋ'은 종성에서 발음될 수 없으므로 [ㄱ]으로 교체된다.
\kfChoice ㄷ: 겹받침 'ㄺ'에서 뒤의 자음 'ㄱ'이 발음되고, 'ㄹ'은 탈락한다.
\kfChoice ㄹ: 겹받침 'ㄳ'에서 앞의 자음 'ㄱ'이 발음되고, 'ㅅ'은 탈락한다.
\kfChoice ㅁ: 겹받침 'ㄻ'에서 뒤의 자음 'ㅁ'이 발음되고, 'ㄹ'은 탈락한다.
\end{kfChoices}
\end{kfQBlock}

\kfPassageFull{문항 2 지문 (concept)}{%
「국어의 음절은 '초성+중성+종성'의 구조를 갖되, 초성과 종성은 비어 있을 수 있고, 중성은 반드시 있어야 한다. 초성에는 자음 하나만 올 수 있으며, 영어처럼 자음군(예: str-)이 초성에 나타나지 않는다. 종성에도 자음 하나만 발음된다. 표기상 겹받침이 있더라도 발음될 때는 하나만 남는 것이 이 제약 때문이다. 이러한 음절 구조 제약은 외래어가 국어에 들어올 때 영향을 미친다. 예를 들어 영어 'strike'는 국어에서 '스트라이크'로 음절을 나누어 차용되며, 'film'은 '필름'으로 종성+초성 구조로 재구성된다. 즉 국어의 음절 구조 제약은 음운 변동의 원인이 되기도 하고, 외래어 수용 방식에도 영향을 미친다.」
}
\begin{kfQBlock}{2}{}
\kfQStem{윗글을 바탕으로 <보기>의 ⓐ\textasciitilde{}ⓔ에 대한 설명으로 적절한 것은?

<보기> 다음 단어의 발음에서 음절 구조 제약이 어떻게 작용하는지 분석해 보자. ⓐ 닭도[닥또]     ⓑ 값[갑]     ⓒ 없다[업따] ⓓ 삼[삼]         ⓔ 읊다[읍따]}
\begin{kfChoices}
\kfChoice ⓐ에서 'ㄺ' 중 'ㄹ'이 탈락하고 뒤 음절 'ㄷ'이 경음화되는 변동이 함께 일어난다.
\kfChoice ⓑ에서 'ㅄ' 중 뒤의 'ㅅ'이 발음으로 남고 앞 자음 'ㅂ'이 탈락하여 [갑]이 된다.
\kfChoice ⓒ에서 겹받침 'ㅄ' 중 뒤의 'ㅅ'이 발음되고 앞의 'ㅂ'이 탈락하여 [업따]가 된다.
\kfChoice ⓓ에서는 음절 구조 제약에 따른 음운 변동이 전혀 적용되지 않는다.
\kfChoice ⓔ에서 겹받침 'ㄿ' 중 앞의 'ㄹ'이 발음되어 결과적으로 [을따]가 된다.
\end{kfChoices}
\end{kfQBlock}

\kfPassageFull{문항 3 지문 (concept)}{%
「겹받침의 발음은 뒤에 오는 음운 환경에 따라 달라진다. 겹받침 뒤에 모음으로 시작하는 조사나 어미가 오면 두 자음 모두 발음되어, 뒤의 자음이 다음 음절 초성으로 연음된다. 예를 들어 '닭이'는 [달기], '삶을'은 [살믈]로 발음된다. 그러나 겹받침 뒤에 자음이 오거나 어말에 위치하면 둘 중 하나만 발음한다. 'ㄳ, ㄵ, ㄼ, ㄽ, ㄾ, ㅀ'은 앞의 자음을, 'ㄺ, ㄻ, ㄿ'은 뒤의 자음을 발음하는 것이 원칙이다. 다만 예외가 있는데, '밟다'는 [밥따], '넓다'는 [널따]로 발음된다. '밟다'에서 'ㄼ'은 원칙(앞 자음 ㄹ)과 달리 'ㅂ'을 발음하고, '넓다'는 원칙대로 'ㄹ'을 발음한다. 그러나 '넓적하다, 넓둥글다'에서는 [넙]으로 발음한다.」
}
\begin{kfQBlock}{3}{}
\kfQStem{윗글을 바탕으로 <보기>의 발음을 판단한 내용으로 적절하지 않은 것은?

<보기> ㉠ 흙만[흥만]   ㉡ 읽고[일꼬]   ㉢ 젊다[점따] ㉣ 밟는[밤는]   ㉤ 삶는[삼는]}
\begin{kfChoices}
\kfChoice ㉠에서 'ㄺ'은 'ㄱ'을 발음하고, 뒤의 'ㅁ' 앞에서 비음화하여 [ㅇ]이 된다.
\kfChoice ㉡에서 'ㄺ'은 'ㄱ'을 발음하고, 뒤의 'ㄱ'은 경음화하여 [ㄲ]이 된다.
\kfChoice ㉢에서 'ㄻ'은 'ㅁ'을 발음하고, 뒤의 'ㄷ'은 경음화하여 [ㄸ]이 된다.
\kfChoice ㉣에서 'ㄼ'은 예외적으로 'ㅂ'을 발음하고, 뒤의 'ㄴ' 앞에서 비음화하여 [ㅁ]이 된다.
\kfChoice ㉤에서 'ㄻ'은 'ㄹ'을 발음하고, 뒤의 'ㄴ' 앞에서 'ㄹ'이 비음화한다.
\end{kfChoices}
\end{kfQBlock}

\kfPassageFull{문항 4 지문 (concept)}{%
「국어에서 받침 뒤에 모음으로 시작하는 조사·어미·접미사가 올 때, 받침이 그 뒤 음절의 초성으로 옮겨 발음되는 현상을 연음(連音)이라 한다. '꽃이'는 종성 'ㅊ'이 뒤 음절로 연음되어 [꼬치]로 발음된다. 이때 연음이 적용되므로 음절의 끝소리 규칙은 적용되지 않는다. 반면 '꽃 위'처럼 합성어나 띄어쓰기가 있는 경우, 음절의 끝소리 규칙이 먼저 적용되어 [꼳]이 된 후 연음되어 [꼬뒤]가 된다. 또한 겹받침의 연음에서는 뒤의 자음만 연음된다. '닭을'은 [달글]이 아니라 종성에 'ㄹ'이 남고 'ㄱ'이 연음되어 [달글]이 된다. 이러한 연음 규칙의 적용 여부와 순서는 정확한 발음의 핵심이다.」
}
\begin{kfQBlock}{4}{}
\kfQStem{윗글을 바탕으로 <보기>의 발음 과정을 분석한 내용으로 적절한 것은?

<보기> ㉠ 옷 안[오단]     ㉡ 옷이[오시] ㉢ 밭 아래[바다래]  ㉣ 밭이[바티] ㉤ 넋이[넉씨]}
\begin{kfChoices}
\kfChoice ㉠은 연음이 바로 적용되어 'ㅅ'이 뒤 음절 초성이 된다.
\kfChoice ㉡은 음절의 끝소리 규칙이 먼저 적용된 후 연음된다.
\kfChoice ㉢은 음절의 끝소리 규칙으로 'ㅌ'이 [ㄷ]이 된 후 연음되어 [바다래]가 된다.
\kfChoice ㉣에서 'ㅌ'은 음절의 끝소리 규칙으로 [ㄷ]이 된 후 연음되어 [바디]가 된다.
\kfChoice ㉤에서 겹받침 'ㄳ'의 두 자음이 모두 뒤 음절로 연음된다.
\end{kfChoices}
\end{kfQBlock}

\kfPassageFull{문항 5 지문 (concept)}{%
「국어의 음절 구조 제약은 여러 음운 변동의 원인이 된다. 첫째, 종성에 두 개의 자음이 올 수 없으므로 겹받침은 자음군 단순화를 겪는다. '삶'→[삼], '닭'→[닥]이 그 예이다. 둘째, 종성에 올 수 있는 자음이 7개로 제한되므로, 그 외의 자음은 7개 중 하나로 교체된다. '꽃'→[꼳], '부엌'→[부억]이 이에 해당한다. 셋째, 초성에 자음군이 올 수 없으므로 외래어 차용 시 모음이 삽입된다. 이러한 음절 구조 제약은 표면형을 결정하는 핵심 원리이며, 다른 음운 규칙(비음화, 경음화 등)의 적용 환경을 만들기도 한다.」
}
\begin{kfQBlock}{5}{}
\kfQStem{윗글을 바탕으로 <보기>의 밑줄 친 단어의 발음에 관여하는 음운 변동을 분석한 내용으로 적절하지 않은 것은?

<보기> ㉠ 삶도 어렵고 죽음도 어렵다. ㉡ 흙만 만지면 마음이 편하다. ㉢ 값도 싸고 품질도 좋다. ㉣ 없는 것이 없다.}
\begin{kfChoices}
\kfChoice ㉠의 '삶도'에서는 자음군 단순화(ㄹ 탈락)와 경음화가 일어나 [삼또]로 발음된다.
\kfChoice ㉡의 '흙만'에서는 자음군 단순화(ㄹ 탈락)와 비음화가 일어나 [흥만]으로 발음된다.
\kfChoice ㉢의 '값도'에서는 자음군 단순화(ㅅ 탈락)와 경음화가 일어나 [갑또]로 발음된다.
\kfChoice ㉣의 '없는'에서는 자음군 단순화(ㅅ 탈락) 후 'ㅂ'이 'ㄴ' 앞에서 비음화하여 [엄는]으로 발음된다.
\kfChoice ㉠\textasciitilde{}㉢에서는 모두 자음군 단순화가 공통적으로 적용된다.
\end{kfChoices}
\end{kfQBlock}

\end{multicols}


\kfStartLit{이번 챕터 문학 작품 5편}


\kfSecHeader{kfAreaLiterature}{지문}{장마 — 윤흥길 (현대소설) / 윤흥길, 「장마」(1973), 핵심 발췌}

\kfPassageNote{장마 — 윤흥길 (현대소설) / 윤흥길, 「장마」(1973), 핵심 발췌}{%
장마가 시작되자 두 할머니 사이의 전쟁도 다시 시작되었다. 친할머니는 방 안에서 태극기를 꺼내 벽에 걸었고, 외할머니는 마루에 앉아 인민군 아들의 무사귀환을 비는 무당굿 노래를 불렀다. 나는 그 사이에서 어쩔 줄을 몰랐다.

어느 날 밤, 마당에 큰 구렁이 한 마리가 나타났다. 친할머니가 「저놈 잡아라!」 하고 소리쳤지만, 외할머니는 「손대지 마라. 그건 우리 삼룡이 영혼이다」라며 울었다. 삼룡이는 외할머니의 아들, 곧 인민군이었다.

장마가 그치자 구렁이는 사라졌다. 친할머니는 「그래, 삼룡이가 갔구나」 하고 중얼거렸다. 외할머니는 울음을 그치고 조용히 고개를 끄덕였다. 두 할머니는 그날 밤 마주 앉아 밥을 먹었다. 나는 그것이 어떤 의미인지 그때는 알지 못했다.
}
\kfSecHeader{kfAreaLiterature}{활동}{작품 정리표 — 장마(윤흥길)}

\noindent\textit{\small 빈칸에 알맞은 말을 채워 넣으시오. (초성 힌트 참고)}\par\medskip
\begin{kfActTable}{|>{\centering\arraybackslash}m{2.6cm}|m{6cm}|m{4cm}|}
\rowcolor{kfPrimaryLight!50}\textbf{\color{kfPrimary}항목} & \textbf{\color{kfPrimary}내용 (빈칸 채우기)} & \textbf{\color{kfPrimary}답안 작성}\\\hline
갈래 & 현대 ㄷㅍ ㅅㅅ / ㅂㄷ 문학 ① & \kfTBlank \\\hline
시점 & ㅇㄹ ㅅㄴ 화자의 1인칭 시점 ② & \kfTBlank \\\hline
대립 구도 & ㅊㅎㅁㄴ(국군) vs ㅇㅎㅁㄴ(인민군) ③ & \kfTBlank \\\hline
갈등 상징 & ㅈㅁ = 전쟁·갈등의 ㅈㅅ ④ & \kfTBlank \\\hline
초자연적 매개 & ㄱㄹㅇ = 삼룡이의 ㅇㅎ ⑤ & \kfTBlank \\\hline
화해 계기 & 친할머니: 'ㅅㄹㅇ가 갔구나' → ㄱㄱ의 수용 ⑥ & \kfTBlank \\\hline
화해 상징 & 두 할머니가 마주 앉아 ㅂ을 먹음 ⑦ & \kfTBlank \\\hline
화자 한계 & '어떤 의미인지 ㅇㅈ ㅁㅎㄷ' ⑧ & \kfTBlank \\\hline
주제 & ㅇㄴ ㄱㄷ을 인간적 ㅈㅅ로 초월 ⑨ & \kfTBlank \\\hline
문학사적 의의 & ㅂㄷ의 비극을 가족 내 갈등으로 형상화 ⑩ & \kfTBlank \\\hline
표현 기법 & ㅅㅈ(장마, 구렁이) + ㅂㅅ(구렁이 등장) ⑪ & \kfTBlank \\\hline
시대 배경 & ㅎㄱ ㅈㅈ 시기 ⑫ & \kfTBlank \\\hline
\end{kfActTable}
\kfSecHeader{kfAreaLiterature}{문제}{}

\begin{multicols}{2}\raggedcolumns\setlength{\columnsep}{12pt}
\begin{kfQBlock}{1}{}
\kfQStem{윗글에 대한 설명으로 적절하지 않은 것은?}
\begin{kfChoices}
\kfChoice 어린 소년 화자의 제한된 시점에서 어른들의 갈등이 서술되고 있다.
\kfChoice '장마'는 전쟁과 갈등이 지속되는 상황을 상징하는 소재로 기능하고 있다.
\kfChoice 구렁이는 죽은 사람의 영혼이 깃든 초자연적 존재로 해석되고 있다.
\kfChoice 두 할머니가 마주 앉아 밥을 먹는 장면은 이념 갈등의 화해를 암시한다.
\kfChoice 화자는 어른들의 이념 갈등을 객관적으로 파악하고 논리적인 어조로 분석하고 있다.
\end{kfChoices}
\end{kfQBlock}

\begin{kfQBlock}{2}{}
\kfQStem{윗글에서 '구렁이'가 지닌 상징적 의미로 가장 적절한 것은?}
\begin{kfChoices}
\kfChoice 친할머니의 이념적 승리를 부각하는 표상적 소재
\kfChoice 장맛비가 빚어낸 자연재해의 공포를 환기하는 소재
\kfChoice 죽은 자의 영혼이 깃든 존재이자 화해의 매개
\kfChoice 화자가 길들여 둔 애완동물이자 가족 평화의 상징
\kfChoice 인민군이 지닌 군사적 위협을 환기하는 전쟁 이미지
\end{kfChoices}
\end{kfQBlock}

\begin{kfQBlock}{3}{}
\kfQStem{윗글에서 친할머니가 '그래, 삼룡이가 갔구나'라고 중얼거린 것이 의미하는 바로 가장 적절한 것은?}
\begin{kfChoices}
\kfChoice 삼룡이의 죽음을 받아들이고 외할머니의 슬픔에 공감하게 되었음을 보여준다.
\kfChoice 구렁이를 반드시 잡아야 한다는 자신의 처음 결심을 다시 확인하는 발화이다.
\kfChoice 삼룡이가 살아서 먼 곳으로 떠났다는 사실에 안도하며 내쉬는 한숨이다.
\kfChoice 외할머니에 대한 적대감이 오히려 한층 더 깊어졌음을 가리키는 표현이다.
\kfChoice 전쟁의 완전한 종결을 선언하는 공식적이고 의례적인 성격의 발화이다.
\end{kfChoices}
\end{kfQBlock}

\begin{kfQBlock}{4}{}
\kfQStem{윗글을 바탕으로 <보기>를 감상한 내용으로 적절하지 않은 것은?

<보기> 윤흥길의 분단 문학에서 갈등의 해결은 이념적 논리가 아니라 인간적 정서, 특히 가족의 슬픔과 공감을 통해 이루어진다. 어린 화자의 시선은 이념의 논리를 벗겨 내고, 그 아래 놓인 인간적 고통의 보편성을 드러내는 장치이다.}
\begin{kfChoices}
\kfChoice 두 할머니가 각각 태극기와 무당굿 노래로 갈등하는 것은 이념에 의한 가족 분열을 보여주겠군.
\kfChoice 구렁이를 통한 화해는 이념적 논리가 아닌 인간적 정서(슬픔, 공감)에 의한 것이겠군.
\kfChoice 화자가 '어쩔 줄을 몰랐다'는 것은 이념의 논리를 벗겨 낸 후 남는 인간적 당혹감이겠군.
\kfChoice 두 할머니의 화해는 이념적 승패가 결정된 결과로, 한쪽이 다른 쪽을 굴복시킨 것이겠군.
\kfChoice 화자가 '알지 못했다'고 한 것은 어린 시선의 제한성을 보여주며, 독자가 의미를 구성하게 하겠군.
\end{kfChoices}
\end{kfQBlock}

\begin{kfQBlock}{5}{}
\kfQStem{윗글에서 '장마가 그치자 구렁이는 사라졌다'라는 서술의 구조적 기능으로 가장 적절한 것은?}
\begin{kfChoices}
\kfChoice 친할머니가 시도한 구렁이 잡기가 마침내 성공으로 마무리되었음을 나타낸다.
\kfChoice 장마철 날씨 변화에 대한 사실적 기록으로 서사 전개와는 무관한 배경 묘사이다.
\kfChoice 구렁이가 다른 장소로 옮겨 갔음을 알려 뒤에 이어질 새로운 갈등의 복선이 된다.
\kfChoice 장마의 끝과 구렁이의 사라짐을 연동시켜 갈등 해소와 망자의 안식을 함께 형상화한다.
\kfChoice 장마와 구렁이의 출현은 우연적으로 겹친 사건이라 어떠한 상징적 의미도 없다.
\end{kfChoices}
\end{kfQBlock}

\end{multicols}

\kfSecHeader{kfAreaLiterature}{지문}{완장 / 아홉 켤레의 구두로 남은 사내 — 윤흥길 (현대소설(복합지문)) / (가) 윤흥길, 「완장」(1979) 핵심 발췌 / (나) 윤흥길, 「아홉 켤레의 구두로 남은 사내」(1977) 핵심 발췌}

\kfPassageNote{완장 / 아홉 켤레의 구두로 남은 사내 — 윤흥길 (현대소설(복합지문)) / (가) 윤흥길, 「완장」(1979) 핵심 발췌 / (나) 윤흥길, 「아홉 켤레의 구두로 남은 사내」(1977) 핵심 발췌}{%
(가) 완장을 차니까 세상이 달라 보였다. 시장 안에서 나를 함부로 대하던 상인들이 고개를 숙이기 시작했다. 나는 그것이 좋았다. 완장 하나에 이렇게 달라질 수 있다는 게 신기했다. 어느 날 내가 단속하던 할머니가 무릎을 꿇었을 때, 나는 잠깐 불편했지만 곧 익숙해졌다. 완장은 나를 다른 사람으로 만들고 있었다.

(나) 권범수는 아홉 켤레의 구두를 남기고 죽었다. 구두 한 켤레 한 켤레에 그의 기술과 정성이 깃들어 있었다. 그는 생전에 이렇게 말했다. 「내가 만든 구두를 신는 사람은 열 년을 신어도 닳지 않을 거요. 그게 내 구두요.」 그러나 아무도 그 구두를 사 가지 않았다. 값이 너무 싸면 업신여기고, 너무 비싸면 아깝다고 했다. 결국 그는 병원비도 없이 죽었고, 아홉 켤레의 구두만이 남았다.
}
\kfSecHeader{kfAreaLiterature}{활동}{작품 정리표 — 완장 / 아홉 켤레의 구두로 남은 사내(윤흥길)}

\noindent\textit{\small 빈칸에 알맞은 말을 채워 넣으시오. (초성 힌트 참고)}\par\medskip
\begin{kfActTable}{|>{\centering\arraybackslash}m{2.6cm}|m{6cm}|m{4cm}|}
\rowcolor{kfPrimaryLight!50}\textbf{\color{kfPrimary}항목} & \textbf{\color{kfPrimary}내용 (빈칸 채우기)} & \textbf{\color{kfPrimary}답안 작성}\\\hline
(가) 핵심 소재 & ㅇㅈ = 미미한 ㄱㄹ의 상징 ① & \kfTBlank \\\hline
(가) 변화 & 양심 ㅁㅂ → 인격 ㅂㅍ ② & \kfTBlank \\\hline
(가) 핵심 장면 & 할머니가 ㅁㄹ을 꿇었을 때 '곧 ㅇㅅㅎㅈㄷ' ③ & \kfTBlank \\\hline
(나) 주인공 & 구두공 ㄱㅂㅅ ④ & \kfTBlank \\\hline
(나) 핵심 소재 & 아홉 켤레의 ㄱㄷ = 존재 증명 ⑤ & \kfTBlank \\\hline
(나) 비극 원인 & 사회의 ㅁㄱㅅ + ㄱㄴ ⑥ & \kfTBlank \\\hline
(나) 결말 & 병원비도 없이 ㅈㅇ → 구두만 남음 ⑦ & \kfTBlank \\\hline
대비 — (가) & ㄱㄹ에 의한 ㄴㅁㅈ 변질 ⑧ & \kfTBlank \\\hline
대비 — (나) & ㄱㄴ에 의한 ㅇㅂㅈ 파멸 ⑨ & \kfTBlank \\\hline
공통 주제 & ㅇㄱ ㅅㅇ와 존엄의 위협 ⑩ & \kfTBlank \\\hline
시대 배경 & 1970년대 한국 사회의 ㄱㅈ ㅁㅅ ⑪ & \kfTBlank \\\hline
작가 & ㅇㅎㄱ — 분단·사회 문제 탐구 ⑫ & \kfTBlank \\\hline
\end{kfActTable}
\kfSecHeader{kfAreaLiterature}{문제}{}

\begin{multicols}{2}\raggedcolumns\setlength{\columnsep}{12pt}
\begin{kfQBlock}{1}{}
\kfQStem{(가)와 (나)에 대한 설명으로 적절하지 않은 것은?}
\begin{kfChoices}
\kfChoice (가)의 화자는 완장을 차기 전과 후의 태도 변화를 서술하고 있다.
\kfChoice (가)에서 '할머니가 무릎을 꿇었을 때' 화자가 느낀 '불편함'은 양심의 작용으로 볼 수 있다.
\kfChoice (나)에서 권범수가 남긴 아홉 켤레의 구두는 그의 기술과 정성을 상징한다.
\kfChoice (나)에서 구두가 팔리지 않은 것은 권범수의 기술이 부족했기 때문이다.
\kfChoice (가)와 (나) 모두 1970년대 한국 사회의 구조적 문제를 반영하고 있다.
\end{kfChoices}
\end{kfQBlock}

\begin{kfQBlock}{2}{}
\kfQStem{(가)의 '완장'이 상징하는 바로 가장 적절한 것은?}
\begin{kfChoices}
\kfChoice 상인들의 자발적 존경을 받는 지위
\kfChoice 시장 질서를 완벽하게 유지하는 정의의 도구
\kfChoice 화자의 도덕적 성장을 보여주는 상징
\kfChoice 미미하지만 사람을 변질시키는 권력의 힘
\kfChoice 화자가 시장에서 구입한 장식품
\end{kfChoices}
\end{kfQBlock}

\begin{kfQBlock}{3}{}
\kfQStem{(가)와 (나)를 비교하여 감상한 내용으로 적절하지 않은 것은?

<보기> 윤흥길의 소설에서 권력과 가난은 인간의 존엄을 위협하는 두 가지 힘이다. (가)는 미미한 권력이 인간을 부패시키는 과정을, (나)는 가난이 성실한 인간을 파멸시키는 과정을 그린다.}
\begin{kfChoices}
\kfChoice (가)의 화자가 권력에 도취되어 가는 것은 인간 존엄의 자발적 상실로 볼 수 있겠군.
\kfChoice (나)의 권범수가 병원비 없이 죽은 것은 가난이라는 구조적 힘에 의한 파멸이겠군.
\kfChoice (가)의 비극은 권력에 의한 내면의 부패이고, (나)의 비극은 사회 구조에 의한 외부적 파멸이겠군.
\kfChoice (가)의 화자와 (나)의 권범수 모두 권력에 도취되어 자신의 존엄을 자발적으로 포기하고 있겠군.
\kfChoice (가)와 (나) 모두 인간의 존엄이 사회적 힘에 의해 위협받는 상황을 그리고 있겠군.
\end{kfChoices}
\end{kfQBlock}

\begin{kfQBlock}{4}{}
\kfQStem{(나)에서 '아홉 켤레의 구두만이 남았다'라는 결말이 함축하는 바로 가장 적절한 것은?}
\begin{kfChoices}
\kfChoice 뛰어난 기술과 정성을 지닌 인간이 사회에서 인정받지 못한 채 소멸한 비극
\kfChoice 구두 사업의 성공으로 권범수의 유산이 풍요로워진 결말
\kfChoice 권범수가 게으름 때문에 실패했음을 보여주는 교훈적 결말
\kfChoice 구두를 통해 권범수가 사후 명성을 얻게 되었음을 나타내는 결말
\kfChoice 구두 장인 정신이 시대를 초월하여 높은 평가를 받았음을 보여주는 결말
\end{kfChoices}
\end{kfQBlock}

\begin{kfQBlock}{5}{}
\kfQStem{(가)에서 '잠깐 불편했지만 곧 익숙해졌다'라는 서술의 표현 효과로 가장 적절한 것은?}
\begin{kfChoices}
\kfChoice 화자가 완장을 곧 반납하기로 결심했음을 우회적으로 암시한다.
\kfChoice 화자의 도덕적 각성이 완성되었음을 단정적으로 드러낸다.
\kfChoice 화자가 할머니를 도와주고 기뻐하는 모습을 부각하여 보여준다.
\kfChoice 시장 질서가 점차 안정되어 가고 있음을 보여주는 배경 묘사이다.
\kfChoice 권력에 의한 양심의 마비 과정을 시간 부사로 압축적으로 보여준다.
\end{kfChoices}
\end{kfQBlock}

\end{multicols}

\kfSecHeader{kfAreaLiterature}{지문}{위로 — 윤동주 (현대시(산문시)) / 윤동주, 「위로」(1940), 학평·모평 기출 다수}

\kfPassageNote{위로 — 윤동주 (현대시(산문시)) / 윤동주, 「위로」(1940), 학평·모평 기출 다수}{%
거미란 놈이 흉한 심보로 병원 뒤뜰 난간과 꽃밭 사이 사람 발이 잘 닿지 않는 곳에 그물을 쳐 놓았다. 옥외 요양을 받는 젊은 사나이가 누워서 치어다보기 바르게―

나비가 한 마리 꽃밭에 날아들다 그물에 걸리었다. 노오란 날개를 파득거려도 파득거려도 나비는 자꾸 감기우기만 한다. 거미가 쏜살같이 가더니 끝없는 끝없는 실을 뽑아 나비의 온몸을 감아 버린다. 사나이는 긴 한숨을 쉬었다.

나이보담 무수한 고생 끝에 때를 잃고 병을 얻은 이 사나이를 위로할 말이―거미줄을 헝클어 버리는 것밖에 위로의 말이 없었다.
}
\kfSecHeader{kfAreaLiterature}{활동}{작품 정리표 — 위로(윤동주)}

\noindent\textit{\small 빈칸에 알맞은 말을 채워 넣으시오. (초성 힌트 참고)}\par\medskip
\begin{kfActTable}{|>{\centering\arraybackslash}m{2.6cm}|m{6cm}|m{4cm}|}
\rowcolor{kfPrimaryLight!50}\textbf{\color{kfPrimary}항목} & \textbf{\color{kfPrimary}내용 (빈칸 채우기)} & \textbf{\color{kfPrimary}답안 작성}\\\hline
갈래 & ㅈㅇㅅ, ㅅㅈㅅ(ㅅㅁㅅ) ① & \kfTBlank \\\hline
성격 & ㅅㄱㅈ, ㅅㅅㅈ, ㅅㅁㅈ ② & \kfTBlank \\\hline
공간 & ㅂㅇ ㄷㄸ ③ & \kfTBlank \\\hline
거미의 행위 & ㅎㅎ ㅅㅂ로 ㄱㅁ을 쳐 놓음 ④ & \kfTBlank \\\hline
나비의 처지 & 노오란 ㄴㄱ를 ㅍㄷㄱㄹ도 자꾸 ㄱㄱ우기만 함 ⑤ & \kfTBlank \\\hline
거미의 결정타 & ㅆㅅㄱㅇ 가더니 ㄲㅇㄴ ㄲㅇㄴ 실로 감음 ⑥ & \kfTBlank \\\hline
사나이의 반응 & ㄱ ㅎㅅ을 쉼 ⑦ & \kfTBlank \\\hline
수법 & ㅇㅎㅈ 수법, ㅇㅇㅎ ⑧ & \kfTBlank \\\hline
상징 1 & 거미 = ㅇㅈ ⑨ & \kfTBlank \\\hline
상징 2 & 나비 = ㅁㄱㄹㅎ ㅇㄹ ㅁㅈ ⑩ & \kfTBlank \\\hline
상징 3 & 젊은 사나이 = ㅁㄱㄹㅎ ㄷㅅㄷㅇ ⑪ & \kfTBlank \\\hline
주제 & 암담한 현실 속 ㅇㄹ와 ㄱㅂㅈ ㅎㄱ에 대한 ㅇㅌㄲㅇ ⑫ & \kfTBlank \\\hline
\end{kfActTable}
\kfSecHeader{kfAreaLiterature}{문제}{}

\begin{multicols}{2}\raggedcolumns\setlength{\columnsep}{12pt}
\begin{kfQBlock}{1}{}
\kfQStem{윗시의 표현상 특징에 대한 설명으로 가장 적절한 것은?}
\begin{kfChoices}
\kfChoice 수미상관 구조를 활용하여 시 전체의 통일성을 강화하고 있다.
\kfChoice 산문적 진술 가운데 ‘\textasciitilde{}다’ 종결과 시구의 반복을 활용하여 운율감을 형성하고 있다.
\kfChoice 도치법을 빈번히 사용하여 화자의 격앙된 정서를 직접 토로하고 있다.
\kfChoice 공감각적 심상을 활용하여 화자의 환상적 분위기를 부각하고 있다.
\kfChoice 구체적 청자에게 직접 말을 건네는 친근한 어조로 일관하여 정서적 친밀감을 강조하고 있다.
\end{kfChoices}
\end{kfQBlock}

\begin{kfQBlock}{2}{}
\kfQStem{윗시에 등장하는 인물·자연물에 대한 이해로 적절하지 않은 것은?}
\begin{kfChoices}
\kfChoice ‘거미’는 사람의 발이 잘 닿지 않는 곳에 ‘흉한 심보’로 그물을 쳐 놓은 부정적 존재로 그려진다.
\kfChoice ‘나비’는 ‘노오란 날개를 파득거려도’ 그물에서 벗어나지 못하는 무력한 존재로 그려진다.
\kfChoice ‘젊은 사나이’는 ‘옥외 요양’ 환자로 거미와 나비의 사건을 누워서 바라보는 인물이다.
\kfChoice ‘젊은 사나이’는 거미줄을 직접 끊어 나비를 구함으로써 부정적 현실을 적극적으로 극복한다.
\kfChoice ‘젊은 사나이’는 사건을 보고 ‘긴 한숨’을 쉬는 데 그쳐 무기력한 인물로 형상화된다.
\end{kfChoices}
\end{kfQBlock}

\begin{kfQBlock}{3}{}
\kfQStem{윗시에서 ‘거미줄을 헝클어 버리는 것밖에 위로의 말이 없었다’의 함축적 의미로 가장 적절한 것은?}
\begin{kfChoices}
\kfChoice 거미줄을 헝클어 버린다는 일시적 위로 외에는 사나이의 고통을 덜어 줄 근본적 방법이 없음을 자각한 것이다.
\kfChoice 거미줄을 헝클어 버리는 행위로 사나이의 모든 고통이 완전히 해결되었음을 선언한 것이다.
\kfChoice 사나이의 오랜 병이 가까운 시일 안에 거짓말처럼 완치될 것이라는 화자의 낙관적 전망을 직접적으로 드러낸 것이다.
\kfChoice 거미를 직접 죽이지 못한 자신의 무능을 분노로 표출한 것이다.
\kfChoice 꽃밭과 그물 사이의 자연적 균형을 회복했음을 강조한 것이다.
\end{kfChoices}
\end{kfQBlock}

\begin{kfQBlock}{4}{}
\kfQStem{<보기>를 바탕으로 윗시를 감상한 내용으로 적절하지 않은 것은? [3점]

<보기> 윤동주의 「위로」는 일제 강점기의 부정적 현실을 우의적으로 형상화한 산문시이다. ‘거미’는 우리 민족을 억압하는 일제, ‘나비’는 그 그물에 걸려 무력하게 감기는 우리 민족, ‘젊은 사나이’는 식민지 현실에 무기력하게 머무는 시인 자신과 동시대인을 상징한다. 화자는 거미줄을 헝클어 버리는 행위를 통해 사나이를 위로하지만, 그 위로가 근본적 해결이 아님을 자각하면서 안타까움도 함께 드러낸다.}
\begin{kfChoices}
\kfChoice ‘흉한 심보’로 그물을 쳐 놓는 ‘거미’는 우리 민족을 억압하는 일제를 상징한다고 볼 수 있다.
\kfChoice ‘파득거려도 파득거려도’ 그물에 감기우기만 하는 ‘나비’의 모습은 일제의 억압 속에서 벗어나지 못하는 우리 민족의 처지를 보여 준다.
\kfChoice ‘긴 한숨’만 쉴 뿐 어떤 적극적 행동도 하지 못하는 ‘젊은 사나이’는 식민지 현실에 무기력하게 머무는 시인과 동시대인을 상징한다.
\kfChoice ‘거미줄을 헝클어 버리는’ 행위에는 일시적 위로를 건네지만 근본적 해결에 이르지 못하는 화자의 안타까움이 담겨 있다.
\kfChoice ‘위로의 말이 없었다’는 진술은 사나이가 이미 모든 고통에서 완전히 벗어났기에 더 이상 위로가 필요 없게 되었음을 보여 준다.
\end{kfChoices}
\end{kfQBlock}

\begin{kfQBlock}{5}{}
\kfQStem{윗시에서 줄표(―)의 사용 효과로 가장 적절한 것은?}
\begin{kfChoices}
\kfChoice 사건의 시간적 순서를 역전시켜 과거 회상의 분위기와 거리감을 동시에 형성한다.
\kfChoice 특정 상황(사나이의 자세, ‘위로할 말’의 한계)을 끊어 보여 의미를 강조한다.
\kfChoice 동일한 시구를 반복하여 시 전체에 일정한 음악적 리듬을 부여하고 정서를 강화한다.
\kfChoice 두 가지 모순된 진술을 결합해 역설적 의미를 형성하여 화자의 모순된 정서를 드러낸다.
\kfChoice 독자에게 직접 질문을 던져 능동적 사고를 유도하고 시 전체의 의미 해석을 위임한다.
\end{kfChoices}
\end{kfQBlock}

\end{multicols}

\kfSecHeader{kfAreaLiterature}{지문}{배꼽을 주제로 한 변주곡 — 이청준 (현대소설) / 이청준, 「배꼽을 주제로 한 변주곡」, 수능·학평 기출}

\kfPassageNote{배꼽을 주제로 한 변주곡 — 이청준 (현대소설) / 이청준, 「배꼽을 주제로 한 변주곡」, 수능·학평 기출}{%
[가] 불편스런 일이 한두 가지가 아니었다. 하지만 허원은 그렇게 스스로 주의하고 고통을 감내해 냈기 때문에 자신의 비밀을 남 앞에 감쪽같이 숨겨 나갈 수 있었다. 아무도 그의 비밀을 눈치챈 사람이 없었다. 비밀이 탄로 나지 않는 한 그의 일상 생활은 더 이상 불편을 겪을 필요도 없었다. 인체 생리나 해부학 서적 같은 걸 뒤져 봐도 성인의 배꼽은 거의 아무런 기능도 수행하지 않음을 알 수 있었다. 적어도 그의 외모나 바깥 생활은 정상을 유지할 수 있었다. 그 점만이라도 무척 다행이었다. 그는 일단 안도의 한숨을 내쉬었다.

[나] ― 그깟 놈의 배꼽, 안 가지고 있음 어때.

그쯤 체념을 하고 될 수 있으면 배꼽에 관한 일들을 잊어버리려 했다. 자신으로부터 배꼽이 사라져 버린 사실을, 그리고 그 때문에 생긴 모든 불편을 잊고, 그 배꼽 없는 생활에 스스로 익숙해져 버리기를 바라 마지않았다. 하지만 문제는 그렇게 간단하지 않았다. 아무리 일상생활에선 드러나게 불편한 점이 없다 해도 그는 역시 배꼽이 없는 자신에 대해 좀처럼 익숙해질 수가 없었다. 그는 자꾸만 허전해서 견딜 수가 없어지곤 했다. 있느니라 여기고 지낼 때는 그처럼 무심스럽던 일이 그런 식으로 한번 의식의 끈을 건드려 오자 허원의 상념은 잠시도 그 잃어 버린 배꼽에서 떠나 있을 수가 없었다.

그는 마침내 회사 출근마저 단념하기에 이르렀다. 그러자 신통하게도 늦잠 버릇이 깨끗이 자취를 감춰 버렸다. 그는 눈만 뜨면 사라져 없어진 배꼽 때문에 기분이 허전했고, 그러면 그 허망감을 쫓기 위해 배꼽에 관한 끝없는 상념들을 쌓기 시작했다.

(중략)

그리하여 배꼽에 관한 허원의 지식과 사념은 자꾸 더 심오하고 추상적인 것이 되어 갔다. 그에게는 어느덧 그 나름의 독특한 배꼽론 같은 것이 윤곽을 지어 가고 있었다. 하지만 그러면 그럴수록 허원은 더욱더 허전해지고, 아무 곳에도 발이 닿아 있는 것 같지 않고, 혼자서 외롭게 허공을 둥둥 떠다니고 있는 것처럼 느껴졌다. 그러면 그는 또 거듭 그 허망감을 쫓기 위해 자신의 배꼽론을 완벽하게 발전시켜 나갔다. 그의 배꼽론은 가령 이런 식으로까지 발전되어 있었다.

[다] ― 우리는 누구나 배꼽을 가지고 있다…… 우리는 우리들의 어머니로부터 탯줄이 끊어지는 순간 이 우주의 한 단자(單子)로서 고독하게 존재하게 되었다. 그러나 우리는 영원히 그 탯줄의 기억을 잊지 않는다. 우리 영혼은 언제까지나 그 어머니의 탯줄과 이어지려 하고, 또다시 그 어머니의 어머니의 탯줄과 이어져 나가면서 우리 존재를 설명하고 근원을 밝혀 나가며, 마침내는 마지막 어머니의 탯줄이 이어지는 우리들의 우주와 만나게 된다…… 우리의 배꼽은 우리가 그 마지막 우주와 만나고자 하는 향수의 표상이며 가능성의 상징이며 존재의 비밀로 나아가는 형이상학이다. 그 비밀의 문이다……

그는 어느덧 배꼽에 대해 당당한 일가견을 이룬 배꼽 전문가가 되어 가고 있었다.

[라] 어느 해 여름이었다. 하니까 그것은 허원이 자신의 배꼽을 잃어버리고 나서 불편하기 그지없는 세 번째의 여름을 맞고 있을 때였다.

[마] 그럴 즈음이었다. 허원은 문득 세상 사람들이 수상쩍어지기 시작했다. 어느 때부턴지는 확실히 알 수 없었지만, 세상 사람들 역시 무슨 이유에선지 이 인간 장기의 한 조그만 흔적에 대해 심상찮은 관심을 나타내기 시작한 것이다. 배꼽 이야기가 일반화의 기미를 엿보이기 시작하자 사람들은 이제 그걸 신호로 아무 흉허물 없이 터놓고 지껄이거나 신문, 잡지 같은 데서 진지하게 논의의 대상을 삼기도 하였다. 배꼽에 관한 논의가 그렇듯 갑자기 시중 일반에까지 성행하기 시작한 것이다. 기묘한 현상이었다.
}
\kfSecHeader{kfAreaLiterature}{활동}{작품 정리표 — 배꼽을 주제로 한 변주곡(이청준)}

\noindent\textit{\small 빈칸에 알맞은 말을 채워 넣으시오. (초성 힌트 참고)}\par\medskip
\begin{kfActTable}{|>{\centering\arraybackslash}m{2.6cm}|m{6cm}|m{4cm}|}
\rowcolor{kfPrimaryLight!50}\textbf{\color{kfPrimary}항목} & \textbf{\color{kfPrimary}내용 (빈칸 채우기)} & \textbf{\color{kfPrimary}답안 작성}\\\hline
갈래 & ㅎㄷ ㄷㅍ ㅅㅅ ① & \kfTBlank \\\hline
성격 & ㅊㅎㅈ, ㅈㅈㅈ, ㅂㅎㅅㅈ ② & \kfTBlank \\\hline
주인공 & ㅍㅂㅎ ㅅㅅㅁ ㅎㅇ ③ & \kfTBlank \\\hline
사건의 출발 & 허원이 자신의 ㅂㄲ을 ㅇㅇㅂㄹ ④ & \kfTBlank \\\hline
초기 반응 & 성인의 배꼽은 ㄱㄴ을 거의 안 함 → ㅇㄷㅇ ㅎㅅ ⑤ & \kfTBlank \\\hline
심리 변화 & ㅇㅅㅇ ㄲ이 ㄱㄷㄹㅈ → 상념이 떠나지 못함 ⑥ & \kfTBlank \\\hline
일상 변화 & ㅎㅅ ㅊㄱ을 단념, ㄴㅈ ㅂㄹ 소실 ⑦ & \kfTBlank \\\hline
사념의 확장 & 허원의 사념이 ㅅㅇㅎㄱ ㅊㅅㅈ인 것이 됨 → ㅂㄲㄹ 정립 ⑧ & \kfTBlank \\\hline
배꼽론의 핵심 1 & 배꼽 = ㅇㅁㄴ의 ㅌㅈ과 이어진 흔적 ⑨ & \kfTBlank \\\hline
배꼽론의 핵심 2 & 배꼽 = 마지막 ㅇㅈ와 만나고자 하는 ㅎㅅㅇ ㅍㅅ ⑩ & \kfTBlank \\\hline
사회적 확장 & 신문·잡지에까지 배꼽 논의가 ㅇㅂㅎ → ㄱㅁㅎ ㅎㅅ ⑪ & \kfTBlank \\\hline
주제 & 비현실적 경험을 통한 ㅇㄱ ㅅㅇ와 ㄱㄷ ⑫ & \kfTBlank \\\hline
\end{kfActTable}
\kfSecHeader{kfAreaLiterature}{문제}{}

\begin{multicols}{2}\raggedcolumns\setlength{\columnsep}{12pt}
\begin{kfQBlock}{1}{}
\kfQStem{윗글의 ‘비밀’이 지닌 서사적 기능으로 가장 적절한 것은?}
\begin{kfChoices}
\kfChoice 자신의 신념을 인물이 돌이켜 본 결과로, 새로운 세계관을 바탕으로 하는 주제를 형성한다.
\kfChoice 얽힌 인간관계를 인물이 성찰하는 전환점으로, 갈등으로 인한 위기감을 완화한다.
\kfChoice 일상적이지 않은 경험을 인물이 의식한다는 표지로, 인물의 심리적 동요를 부른다.
\kfChoice 상충된 이해관계를 인물이 조정하는 단서로, 심화된 사회적 갈등을 해소한다.
\kfChoice 기성의 질서에 인물이 저항한다는 신호로, 돌발적 사건의 발생을 알린다.
\end{kfChoices}
\end{kfQBlock}

\begin{kfQBlock}{2}{}
\kfQStem{‘허원’을 중심으로 윗글을 이해한 내용으로 적절하지 않은 것은?}
\begin{kfChoices}
\kfChoice 허원은 배꼽이라는 ‘실물’과 관련하여 시작된 ‘사념’을 통해 ‘존재’의 의미를 발견해 간다.
\kfChoice 허원은 성인의 배꼽이 몸에서 큰 기능을 하지 않는다는 사실을 알고 일단 안도감을 느낀다.
\kfChoice 허원은 ‘사념’을 방편으로 삼아 자신의 현재 상태에 대해 다른 방향에서 접근하고자 한다.
\kfChoice 허원은 배꼽에 대한 사람들의 ‘심상찮은 관심’의 원인을 궁금해하면서 ‘세상 사람들’에게 주의를 기울이게 된다.
\kfChoice 허원은 배꼽에 대한 인식을 ‘세상 사람들’과 공유하게 되면서, 그간 이어 온 ‘사념’을 더 이상 지속하지 않게 된다.
\end{kfChoices}
\end{kfQBlock}

\begin{kfQBlock}{3}{}
\kfQStem{윗글의 [가]\textasciitilde{}[마]에 대한 서술 방식의 설명으로 가장 적절한 것은?}
\begin{kfChoices}
\kfChoice [가]: 누구의 생각을 누가 말하는지 명시한 표현을 나타내어 서술하고 있다.
\kfChoice [나]: 인물의 생각을 서술자가 평가하며 그 심화된 의미를 함축하여 서술하고 있다.
\kfChoice [다]: 인물의 의식을 인물 자신의 생생한 목소리를 통해 서술하고 있다.
\kfChoice [라]: 인물의 상황에 관련된 정보를 부가하여 서술하고 있다.
\kfChoice [마]: 인물 행동의 진행 과정을 순차적으로 서술하고 있다.
\end{kfChoices}
\end{kfQBlock}

\begin{kfQBlock}{4}{}
\kfQStem{<보기>를 참고하여 윗글을 감상한 내용으로 적절하지 않은 것은? [3점]

<보기> 「배꼽을 주제로 한 변주곡」은 주인공이 배꼽을 잃어버렸다는 허구적 설정으로 시작하여, 이후 배꼽을 둘러싼 희화적 에피소드들이 이어진다. 주인공은 으레 있어야 할 것이 없어져 불편한 생활을 이어 가던 중 배꼽에 관심을 갖는 이들이 늘어나고 있음을 알게 된다. 이 과정에서 배꼽에 관련된 개인적 상황은 물론 인간 존재와 사회 상황에 대한 심층적 의미의 탐색이 이루어진다.}
\begin{kfChoices}
\kfChoice ‘의식의 끈’이 ‘건드려’짐으로써 주인공이 비정상적 문제 상황에 지속적으로 주목하게 된 것이라 할 수 있겠군.
\kfChoice ‘회사 출근’을 포기하게 되고 ‘늦잠 버릇’이 사라진 상황은 주인공의 일상이 변화된 모습을 보여 준다고 할 수 있겠군.
\kfChoice ‘배꼽’을 ‘어머니’의 ‘탯줄’에 연관 지어 이해하는 것은 개인의 신체에 관련된 생각을 ‘우주와 만나’는 ‘심오하고 추상적인’ 형이상학적 사유로 확장해 가는 실마리가 된다고 할 수 있겠군.
\kfChoice ‘그의 사념’이 도달한 ‘배꼽론’의 ‘확고한 경지’는 사소한 것의 심층적 의미를 탐색할 때 이를 수 있으므로, 그 사소한 것에 얽매이지 않는 자유로운 상태에서 실현이 가능해지겠군.
\kfChoice ‘기묘한 현상’은 ‘배꼽 이야기’가 ‘일반화’되는 상황이 뜻밖이지만 ‘사실’로 나타나는 현상을 두고 일컬은 말이라 할 수 있겠군.
\end{kfChoices}
\end{kfQBlock}

\begin{kfQBlock}{5}{}
\kfQStem{윗글에서 ‘배꼽론’이 갖는 의미에 대한 설명으로 가장 적절한 것은?}
\begin{kfChoices}
\kfChoice 허망감을 사념으로 보충하며 인간 존재의 근원을 우주적 차원으로 확장한 사유의 결과물이다.
\kfChoice 배꼽이 사라졌다는 사실 자체를 부정하기 위해 허원이 만든 단순한 자기 합리화의 도구이다.
\kfChoice 주변 사람들에게 성인의 배꼽이 지닌 의학적 기능을 설명하기 위한 실용적 강연용 자료이다.
\kfChoice 허원이 자신의 직장 복귀를 위해 마련한 구체적이고 현실적인 행동 계획표에 해당한다.
\kfChoice 사람들의 ‘배꼽 이야기’ 일반화 흐름을 주도하기 위해 허원이 의도적으로 유포한 담론이다.
\end{kfChoices}
\end{kfQBlock}

\end{multicols}

\kfSecHeader{kfAreaLiterature}{지문}{수궁가 — 작자 미상 (판소리(고전문학)) / 「수궁가」(판소리), 핵심 발췌}

\kfPassageNote{수궁가 — 작자 미상 (판소리(고전문학)) / 「수궁가」(판소리), 핵심 발췌}{%
토끼가 수궁에 들어가니 용왕이 반기며 하는 말이,

「오냐, 네가 토끼란 놈이로구나. 과인이 병이 중하여 네 간을 구하노라. 어서 네 배를 가르겠다.」

토끼 속으로 크게 놀라되 얼굴에는 태연한 체하고 여쭈오되,

「소인의 간은 본디 신묘한 물건이라, 항시 배 속에 두면 더러워지옵기로 깨끗이 씻어 바위 위에 말려 두고 왔사옵니다. 대왕께서 일찍이 기별을 하셨더라면 가져올 것을, 급히 오는 바람에 두고 왔사오니, 소인을 다시 육지로 보내 주시면 가져다 바치겠나이다.」

용왕이 기뻐하여,

「과연 충성스러운 놈이로다. 어서 가서 가져오너라.」

별주부 기가 막혀 아뢰되,

「대왕마마, 간을 빼어 바위에 말려 둔다 함은 만고에 없는 일이옵니다. 토끼의 꾀에 속지 마옵소서.」

용왕이 노하여,

「네가 어찌 과인의 은인을 해하려 하느냐. 물러서라.」

토끼 육지에 오르자마자 껑충 뛰며 하는 말이,

「어허, 이 간을 빼어 말릴 놈이 세상에 어디 있단 말이냐. 바보 용왕아, 꿈 깨거라!」
}
\kfSecHeader{kfAreaLiterature}{활동}{작품 정리표 — 수궁가}

\noindent\textit{\small 빈칸에 알맞은 말을 채워 넣으시오. (초성 힌트 참고)}\par\medskip
\begin{kfActTable}{|>{\centering\arraybackslash}m{2.6cm}|m{6cm}|m{4cm}|}
\rowcolor{kfPrimaryLight!50}\textbf{\color{kfPrimary}항목} & \textbf{\color{kfPrimary}내용 (빈칸 채우기)} & \textbf{\color{kfPrimary}답안 작성}\\\hline
갈래 & ㅍㅅㄹ / 고전 ㅅㅅ ① & \kfTBlank \\\hline
원전 & ㅌㅂㅈ(토별전) ② & \kfTBlank \\\hline
용왕의 역할 & ㅇㄹㅅㅇ 권력자 ③ & \kfTBlank \\\hline
별주부의 역할 & 무시당하는 ㅊㅅ ④ & \kfTBlank \\\hline
토끼의 역할 & ㄱㅈ로 살아남는 ㅇㅈ ⑤ & \kfTBlank \\\hline
토끼의 거짓말 & 간을 ㅂㅇ에 ㅁㄹ 두었다 ⑥ & \kfTBlank \\\hline
용왕의 오판 & 토끼를 '충성스러운 놈'이라 하고 ㅂㅈㅂ의 충언 무시 ⑦ & \kfTBlank \\\hline
결말 & 토끼 탈출 + 'ㅂㅂ 용왕아' ⑧ & \kfTBlank \\\hline
풍자 대상 & ㄱㄹ의 ㅎㅍ + 권력자의 어리석음 ⑨ & \kfTBlank \\\hline
해학 & 웃음을 통한 ㅂㅍ — 판소리 특유 ⑩ & \kfTBlank \\\hline
주제 1 & 약자의 ㅈㄹ이 강자를 이김 ⑪ & \kfTBlank \\\hline
주제 2 & ㄱㄹ에 대한 ㅁㅈㅈ 비판 의식 ⑫ & \kfTBlank \\\hline
\end{kfActTable}
\kfSecHeader{kfAreaLiterature}{문제}{}

\begin{multicols}{2}\raggedcolumns\setlength{\columnsep}{12pt}
\begin{kfQBlock}{1}{}
\kfQStem{윗글에 대한 설명으로 적절하지 않은 것은?}
\begin{kfChoices}
\kfChoice 토끼가 위기 상황에서 기지를 발휘하여 꾀로 용왕을 속이고 있다.
\kfChoice 별주부는 토끼의 거짓말을 간파하고 용왕에게 경고하고 있다.
\kfChoice 용왕은 토끼의 말을 그대로 믿고 별주부의 충언을 무시하고 있다.
\kfChoice 토끼는 육지에 돌아온 후 용왕에 대한 감사와 충성을 표현하고 있다.
\kfChoice 판소리 특유의 대화체와 해학적 표현이 사용되고 있다.
\end{kfChoices}
\end{kfQBlock}

\begin{kfQBlock}{2}{}
\kfQStem{윗글에서 용왕이 별주부의 충언을 무시하고 토끼를 풀어 준 것이 풍자하는 바로 가장 적절한 것은?}
\begin{kfChoices}
\kfChoice 충신의 간언을 물리치고 아첨에 현혹되는 어리석은 권력자의 모습
\kfChoice 현명한 군주가 신하의 의견을 종합하여 올바른 판단을 내리는 모습
\kfChoice 토끼의 뛰어난 의학 지식에 감탄하는 용왕의 학구적 태도
\kfChoice 별주부의 불충을 벌하는 용왕의 정의로운 통치
\kfChoice 수궁의 완벽한 사법 체계를 보여주는 장면
\end{kfChoices}
\end{kfQBlock}

\begin{kfQBlock}{3}{}
\kfQStem{윗글에서 토끼가 '간을 바위에 말려 두었다'고 말한 것이 갖는 서사적 기능으로 가장 적절한 것은?}
\begin{kfChoices}
\kfChoice 위기를 기지로 극복하는 약자의 지략을 보여 주는 서사 전환점이다.
\kfChoice 토끼가 지닌 의학적 지식의 깊이를 입증하는 사실적 서술의 한 사례이다.
\kfChoice 용왕을 향한 토끼의 진심 어린 충성과 헌신을 드러내는 결정적 장면이다.
\kfChoice 별주부의 판단이 처음부터 잘못되었음을 독자에게 확인시켜 주는 부분이다.
\kfChoice 수궁의 의료 체계가 매우 발전되어 있음을 보여 주는 묘사적 장면이다.
\end{kfChoices}
\end{kfQBlock}

\begin{kfQBlock}{4}{}
\kfQStem{윗글을 바탕으로 <보기>를 감상한 내용으로 적절하지 않은 것은?

<보기> 판소리에서 풍자와 해학은 지배층의 어리석음을 웃음으로 비판하는 기능을 한다. 약자(토끼)의 지략과 강자(용왕)의 어리석음이 대비되면서, 권력의 횡포에 대한 민중적 비판 의식이 드러난다.}
\begin{kfChoices}
\kfChoice 용왕이 토끼의 말에 속는 것은 지배층의 어리석음을 풍자하는 것이겠군.
\kfChoice 토끼의 꾀는 약자가 강자의 횡포에 맞서는 지략으로 볼 수 있겠군.
\kfChoice 별주부의 충언이 무시되는 것은 권력자의 판단 오류를 부각하겠군.
\kfChoice '바보 용왕아, 꿈 깨거라!'는 해학적 표현으로 권력에 대한 조롱이겠군.
\kfChoice 이 작품은 용왕의 현명한 통치를 찬양하는 것이 핵심 주제이겠군.
\end{kfChoices}
\end{kfQBlock}

\begin{kfQBlock}{5}{}
\kfQStem{윗글에서 토끼가 '속으로 크게 놀라되 얼굴에는 태연한 체하고'라는 서술의 표현 효과로 가장 적절한 것은?}
\begin{kfChoices}
\kfChoice 내면의 공포와 외면의 침착을 대비해 임기응변을 부각한다.
\kfChoice 토끼가 위기 속에서도 두려움을 전혀 느끼지 않았음을 강조한다.
\kfChoice 토끼의 성격이 본래 소심하고 우유부단함을 우회적으로 드러낸다.
\kfChoice 서술자가 토끼의 진짜 내면을 끝내 알지 못함을 시사하고 있다.
\kfChoice 용왕이 토끼의 미세한 표정을 읽어 내지 못했음을 우회적으로 비판한다.
\end{kfChoices}
\end{kfQBlock}

\end{multicols}

\kfStartNonlit{이번 챕터 비문학 지문 5편}


\kfSecHeader{kfAreaNonlit}{지문}{공리주의의 내부 비판과 의무론적 응답 (인문)}

\kfPassageNote{공리주의의 내부 비판과 의무론적 응답 (인문)}{%
(가) 공리주의는 '최대 다수의 최대 행복'을 도덕의 기준으로 삼는 윤리 이론이다. 벤담(J. Bentham)의 양적 공리주의는 쾌락의 총량을 계산하여 행위의 옳고 그름을 판단한다. 그러나 공리주의 내부에서도 이에 대한 비판이 제기되었다. 밀(J. S. Mill)은 쾌락에 질적 차이가 있다고 주장했다. '배부른 돼지보다 배고픈 소크라테스가 낫다'는 그의 유명한 표현은, 감각적 쾌락보다 지적·도덕적 쾌락이 더 높은 가치를 지닌다는 것을 뜻한다.

밀의 질적 공리주의에 대해서도 비판이 이어졌다. 쾌락의 질을 비교할 객관적 기준이 부재하다는 것이다. 밀은 '유능한 판정자(competent judges)', 즉 두 종류의 쾌락을 모두 경험한 사람의 판단에 호소했지만, 이 판정이 정말 객관적인지는 논쟁적이다. 나아가 공리주의 일반에 대해 '소수자 희생' 문제가 제기된다. 만약 한 명의 무고한 사람을 처형하면 사회 전체의 행복이 극대화되는 경우, 공리주의는 그 처형을 정당화해야 하는가? 이것은 공리주의가 개인의 권리를 충분히 보호하지 못한다는 비판이다.

이에 대응하여 규칙 공리주의는 개별 행위가 아니라 '규칙'의 결과를 평가한다. '무고한 사람을 처형하지 말라'는 규칙이 장기적으로 더 큰 행복을 산출하므로, 그 규칙을 따라야 한다는 것이다. 그러나 규칙 공리주의 역시, 규칙을 깨뜨리는 것이 더 큰 행복을 가져오는 예외 상황에서 규칙을 고수해야 하는 이유를 설명하기 어렵다는 문제가 남는다.

(나) 의무론적 윤리학은 행위의 결과가 아니라 행위 자체의 도덕적 성격에 주목한다. 칸트(I. Kant)는 도덕 법칙의 보편화 가능성을 기준으로 행위의 옳고 그름을 판단한다. 그의 정언 명법(定言命法) 제1정식은 「네 의지의 준칙이 항상 동시에 보편적 입법의 원리로서 타당할 수 있도록 행위하라」이다. 이를테면 '약속을 어겨도 좋다'는 준칙은 보편화되면 약속 제도 자체가 무너지므로 자기 모순에 빠져 도덕 법칙이 될 수 없다.

칸트의 제2정식은 「네 자신과 다른 모든 사람의 인격에서 인간성을, 항상 동시에 목적으로 대우하고 결코 단순한 수단으로만 대우하지 말라」이다. 이 정식은 인간의 존엄성을 절대적으로 보호하며, 아무리 큰 사회적 이익을 위해서라도 개인을 단순한 수단으로 사용할 수 없다는 원칙을 세운다. 공리주의에 대한 의무론의 핵심 비판은 바로 이 지점에 있다. 공리주의가 전체 행복 극대화를 위해 개인을 수단화할 위험이 있는 반면, 의무론은 개인의 권리와 존엄을 결과와 무관하게 보호한다.

그러나 의무론에도 비판이 있다. 결과를 전혀 고려하지 않는 것이 현실적으로 가능한가? 이를테면 살인범에게 쫓기는 사람의 위치를 물을 때, 칸트의 원칙에 따르면 거짓말을 해서는 안 된다. 이것이 도덕적으로 옳은가? 결과주의자는 이 경우 거짓말이 도덕적으로 정당화된다고 주장할 것이다. 이처럼 결과주의와 의무론은 각각의 강점과 약점을 안고 현대 윤리학의 핵심 논쟁을 구성한다.
}
\kfSecHeader{kfAreaNonlit}{문제}{}

\begin{multicols}{2}\raggedcolumns\setlength{\columnsep}{12pt}
\begin{kfQBlock}{1}{}
\kfQStem{(가)와 (나)의 내용과 일치하지 않는 것은?}
\begin{kfChoices}
\kfChoice 밀은 쾌락에 양적 차이뿐 아니라 질적 차이가 있다고 주장했다.
\kfChoice 규칙 공리주의는 개별 행위가 아니라 규칙의 결과를 평가한다.
\kfChoice 칸트의 정언 명법 제1정식은 행위 준칙의 보편화 가능성을 기준으로 삼는다.
\kfChoice 의무론은 전체 행복 극대화를 위해 개인을 수단화하는 것을 허용한다.
\kfChoice 의무론에 대해서도 결과를 전혀 고려하지 않는 것이 현실적인가라는 비판이 제기된다.
\end{kfChoices}
\end{kfQBlock}

\begin{kfQBlock}{2}{}
\kfQStem{(나)에 따를 때, 칸트가 '약속을 어겨도 좋다'는 준칙이 도덕 법칙이 될 수 없다고 본 이유로 가장 적절한 것은?}
\begin{kfChoices}
\kfChoice 약속을 어기는 행위가 사회 전체의 행복 총량을 감소시키기 때문이다.
\kfChoice 그 준칙이 보편화되면 약속 제도 자체가 무너져 자기 모순에 빠지기 때문이다.
\kfChoice 약속 위반이 경험적으로 드문 사건이며 통계적 예외에 불과하기 때문이다.
\kfChoice 약속을 어기는 일이 유능한 판정자의 합리적 판단에 명백히 어긋나기 때문이다.
\kfChoice 약속 위반은 규칙 공리주의가 인정하는 보편 규칙에 어긋나기 때문이다.
\end{kfChoices}
\end{kfQBlock}

\begin{kfQBlock}{3}{}
\kfQStem{윗글을 바탕으로 <보기>의 상황에 대해 각 이론이 내릴 수 있는 판단으로 적절하지 않은 것은?

<보기> 한 도시에서 전염병이 발생했다. 감염자 1명을 강제 격리하면 1,000명의 시민이 감염을 면할 수 있다. 그러나 강제 격리는 그 개인의 자유를 침해한다.}
\begin{kfChoices}
\kfChoice 벤담의 양적 공리주의는 1,000명의 행복이 1명의 불행보다 크므로 강제 격리를 정당화할 수 있다.
\kfChoice 칸트의 의무론에 따르면, 그 개인을 단순한 수단으로 사용하는 것이므로 강제 격리는 정당화되기 어렵다.
\kfChoice 규칙 공리주의는 '전염병 시 감염자를 격리한다'는 규칙이 장기적으로 더 큰 행복을 산출한다면 격리를 지지할 수 있다.
\kfChoice 밀의 질적 공리주의는 격리 대상자의 지적 쾌락이 다른 시민보다 높으면 격리를 반대할 것이다.
\kfChoice 의무론적 관점에서는 행위의 결과보다 개인의 권리 침해 여부가 도덕적 판단의 핵심이다.
\end{kfChoices}
\end{kfQBlock}

\begin{kfQBlock}{4}{}
\kfQStem{(가)의 공리주의와 (나)의 의무론을 비교한 내용으로 가장 적절한 것은?}
\begin{kfChoices}
\kfChoice 공리주의는 행위의 결과를, 의무론은 행위 자체의 도덕적 성격을 판단 기준으로 삼는다.
\kfChoice 공리주의와 의무론 모두 행위의 결과만을 도덕적 판단의 기준으로 삼는다.
\kfChoice 의무론은 전체 행복의 극대화를, 공리주의는 개인의 존엄성 보호를 목표로 한다.
\kfChoice 공리주의는 도덕 법칙의 보편화 가능성을 기준으로 삼고, 의무론은 쾌락의 총량을 기준으로 삼는다.
\kfChoice 공리주의와 의무론 모두 개인을 수단으로 사용하는 것을 절대적으로 금지한다.
\end{kfChoices}
\end{kfQBlock}

\end{multicols}

\kfSecHeader{kfAreaNonlit}{지문}{샤프츠베리의 도덕감과 듀이의 도덕 교육론 (인문)}

\kfPassageNote{샤프츠베리의 도덕감과 듀이의 도덕 교육론 (인문)}{%
도덕성의 기원을 둘러싼 논쟁에서 '선천론'과 '경험론'은 오랜 대립축을 이룬다. 제3대 샤프츠베리 백작(Third Earl of Shaftesbury)은 18세기 초 영국 도덕 철학에서 선천적 도덕감(moral sense) 이론을 제시한 선구적 인물이다. 그에 따르면 인간에게는 아름다운 것을 보면 쾌감을 느끼듯, 선한 행위를 보면 승인의 감정을, 악한 행위를 보면 불승인의 감정을 자연스럽게 느끼는 내적 감각이 있다. 이 도덕감은 이성적 추론의 결과가 아니라, 미적 감각과 유사한 직관적·감정적 능력이다.

샤프츠베리는 도덕감이 인간 본성에 내재하며, 사회적 학습이나 문화적 조건에 의해 만들어지는 것이 아니라고 보았다. 이 점에서 그는 홉스의 이기주의적 인간관에 반대한다. 홉스에 따르면 인간은 본성적으로 이기적이며 도덕은 사회 계약의 산물에 불과하지만, 샤프츠베리에게 도덕은 인간 본성 자체에 뿌리를 둔 것이다. 도덕감은 개인적 이익과 무관하게 타인의 행복에 대한 관심과 공감을 유발한다. 그러나 샤프츠베리의 이론에는 도덕감이 선천적이라면 왜 사람마다 도덕적 판단이 다른지 설명하기 어렵다는 비판이 제기된다.

존 듀이(John Dewey)는 이와 대비되는 입장에서 도덕성을 경험과 교육의 산물로 파악했다. 듀이에 따르면 도덕적 능력은 선천적으로 주어지는 고정된 것이 아니라, 구체적인 사회적 상황에서 문제를 해결하는 과정을 통해 형성되고 성장하는 것이다. 그의 프래그머티즘(pragmatism) 윤리학에서 도덕적 판단은 추상적 원리의 적용이 아니라, 상황 속에서 가장 바람직한 결과를 이끌어 내기 위한 지적 탐구이다.

듀이의 도덕 교육론에서 핵심은 '반성적 경험(reflective experience)'이다. 아이는 사회적 상호 작용 속에서 갈등을 겪고, 그 갈등을 반성적으로 사유하며, 시행착오를 통해 더 나은 해결 방안을 찾아 나간다. 이 과정 자체가 도덕적 성장이다. 듀이는 고정된 도덕 규범을 외부에서 주입하는 교육 방식을 비판하고, 아이가 스스로 문제 상황을 탐구하며 도덕적 판단력을 키울 수 있는 환경을 조성해야 한다고 주장했다. 이처럼 샤프츠베리는 도덕성의 선천적 기반을, 듀이는 경험적 형성 과정을 강조하지만, 두 입장 모두 이성적 추론만으로는 도덕성을 설명할 수 없다는 점에서는 공통점을 보인다.
}
\kfSecHeader{kfAreaNonlit}{문제}{}

\begin{multicols}{2}\raggedcolumns\setlength{\columnsep}{12pt}
\begin{kfQBlock}{1}{}
\kfQStem{윗글의 내용과 일치하는 것은?}
\begin{kfChoices}
\kfChoice 샤프츠베리는 도덕감이 사회적 학습에 의해 후천적으로 형성된다고 보았다.
\kfChoice 샤프츠베리는 홉스의 이기주의적 인간관에 동의하였다.
\kfChoice 듀이는 도덕적 능력이 선천적으로 주어지는 고정된 것이라고 주장했다.
\kfChoice 듀이는 고정된 도덕 규범을 외부에서 주입하는 교육 방식을 지지했다.
\kfChoice 샤프츠베리와 듀이 모두 이성적 추론만으로는 도덕성을 설명할 수 없다고 보았다.
\end{kfChoices}
\end{kfQBlock}

\begin{kfQBlock}{2}{}
\kfQStem{윗글의 샤프츠베리와 듀이를 비교한 내용으로 적절하지 않은 것은?}
\begin{kfChoices}
\kfChoice 샤프츠베리는 도덕성의 선천적 기반을, 듀이는 경험적 형성 과정을 강조한다.
\kfChoice 샤프츠베리의 도덕감은 직관적·감정적 능력이며, 듀이의 도덕적 판단은 지적 탐구이다.
\kfChoice 샤프츠베리와 듀이는 도덕적 판단에서 이성적 추론의 역할을 가장 중요하게 본다는 점에서 일치한다.
\kfChoice 샤프츠베리는 도덕감이 본성에 내재한다고 보고, 듀이는 사회적 상호 작용을 통해 형성된다고 본다.
\kfChoice 샤프츠베리는 홉스의 이기주의에 반대하고, 듀이는 고정된 규범 주입 교육에 반대한다.
\end{kfChoices}
\end{kfQBlock}

\begin{kfQBlock}{3}{}
\kfQStem{윗글을 바탕으로 <보기>를 이해한 내용으로 적절하지 않은 것은?

<보기> E 선생님은 학생들에게 도덕 규범을 암기시키는 대신, 학교에서 발생한 갈등 상황을 제시하고 학생들이 스스로 해결 방안을 토론하게 했다.}
\begin{kfChoices}
\kfChoice 듀이의 관점에서 E 선생님의 방법은 학생들의 반성적 경험을 능동적으로 유도하는 교육이겠군.
\kfChoice 듀이는 학생들이 갈등 상황에서 시행착오를 거쳐 도덕적 판단력을 점차 키울 수 있다고 보겠군.
\kfChoice 샤프츠베리의 관점에서 E 선생님의 방법은 선천적 도덕감을 경유해 경험으로 도덕성을 보완하는 것이겠군.
\kfChoice 듀이는 E 선생님의 방법을 적극 지지하면서 고정된 규범을 암기시키는 방법은 비판하겠군.
\kfChoice 샤프츠베리의 관점에서 학생들의 도덕적 판단은 이러한 토론 교육이 없어도 선천적으로 가능하겠군.
\end{kfChoices}
\end{kfQBlock}

\begin{kfQBlock}{4}{}
\kfQStem{윗글의 문맥에서 '반성적 경험(reflective experience)'이 의미하는 바로 가장 적절한 것은?}
\begin{kfChoices}
\kfChoice 갈등을 겪고 반성적 사유로 더 나은 해결 방안을 찾는 과정
\kfChoice 내재한 선천적 도덕감이 외부 자극 없이 자동으로 발현되는 과정
\kfChoice 고정되어 있는 도덕 규범을 권위자로부터 주입받아 내면화하는 과정
\kfChoice 오직 이성적 추론에 기대어 보편적 도덕 원리를 연역하는 과정
\kfChoice 과거의 잘못을 떠올려 깊이 반성하고 참회하는 종교적 행위
\end{kfChoices}
\end{kfQBlock}

\end{multicols}

\kfSecHeader{kfAreaNonlit}{지문}{다산 정약용의 윤리학 심화 (인문)}

\kfPassageNote{다산 정약용의 윤리학 심화 (인문)}{%
정약용(丁若鏞, 1762\textasciitilde{}1836)은 조선 후기 실학의 집대성자로서, 성리학의 이기론적 심성론을 근본적으로 재해석하여 독자적인 윤리학 체계를 구축했다. 그의 윤리학의 출발점은 '성기호설(性嗜好說)'이다. 성리학에서 '성(性)'은 인간에게 선천적으로 부여된 도덕적 본성, 곧 '리(理)'로 이해되었다. 그러나 정약용은 성을 고정된 리가 아니라 '기호(嗜好)', 즉 선을 좋아하고 악을 싫어하는 경향성으로 파악했다. 인간의 본성은 완성된 덕이 아니라 선을 향한 지향, 즉 도덕적 가능성이다.

이 성기호설에서 핵심적인 것은 '자주지권(自主之權)'이다. 인간은 선을 좋아하는 경향을 갖고 있지만, 동시에 형구(形軀)의 기호, 즉 감각적 욕구도 갖고 있다. 선을 향할 것인지 감각적 욕구를 따를 것인지는 인간의 자유로운 선택, 곧 자주지권에 달려 있다. 정약용에게 도덕성은 본성에 이미 완성되어 있는 것이 아니라, 자유로운 선택을 통해 실현해 나가야 하는 것이다. 이 점에서 그의 윤리학은 성리학의 본성 복원론(본래 선한 본성을 회복하라)과 구별된다.

정약용의 윤리학에서 도덕적 실천의 핵심은 '행사지심(行事之心)'이다. 성리학의 주류가 마음의 덕(미발 상태의 선한 본성)을 강조했다면, 정약용은 실제로 행동으로 옮기는 마음, 즉 실천에서의 도덕성을 중시했다. 그는 「맹자(孟子)」의 사단(四端)을 재해석하여, 측은지심·수오지심·사양지심·시비지심이 단순한 심리적 감정이 아니라 행위로 나아가는 실천적 단서라고 보았다. 측은히 여기는 마음은 실제로 고통받는 사람을 도울 때 비로소 '인(仁)'이 된다.

이러한 실천 중심의 윤리학은 정약용의 경세(經世) 사상과 연결된다. 그는 개인의 내면 수양에 머무르지 않고, 도덕적 실천이 사회·정치적 제도 개혁으로 이어져야 한다고 주장했다. 목민관(牧民官)이 백성을 위해 헌신하는 것, 전제(田制)를 개혁하여 토지를 공정하게 분배하는 것, 이 모든 것이 도덕의 실현이다. 정약용에게 윤리학은 내면의 닦음에서 끝나는 것이 아니라, 구체적인 사회적 실천으로 완성되는 것이다.
}
\kfSecHeader{kfAreaNonlit}{문제}{}

\begin{multicols}{2}\raggedcolumns\setlength{\columnsep}{12pt}
\begin{kfQBlock}{1}{}
\kfQStem{윗글의 내용과 일치하지 않는 것은?}
\begin{kfChoices}
\kfChoice 정약용은 성(性)을 고정된 리가 아니라 선을 좋아하고 악을 싫어하는 경향성으로 파악했다.
\kfChoice 자주지권은 선을 향할 것인지 감각적 욕구를 따를 것인지 자유롭게 선택하는 능력이다.
\kfChoice 정약용은 도덕성이 본성에 이미 완성되어 있으므로 실천은 불필요하다고 보았다.
\kfChoice 정약용은 사단을 행위로 나아가는 실천적 단서로 재해석했다.
\kfChoice 정약용에게 윤리학은 내면 수양에 머무르지 않고 사회적 실천으로 완성되는 것이다.
\end{kfChoices}
\end{kfQBlock}

\begin{kfQBlock}{2}{}
\kfQStem{윗글에 따를 때, 정약용이 성리학의 본성 복원론과 구별되는 이유로 가장 적절한 것은?}
\begin{kfChoices}
\kfChoice 성리학과 정약용 모두 성을 고정된 리로 동일하게 파악하기 때문이다.
\kfChoice 성리학은 선한 본성의 회복을, 정약용은 자유로운 선택의 실현을 강조하기 때문이다.
\kfChoice 정약용이 감각적 욕구를 완전히 제거해야 한다고 강력히 주장했기 때문이다.
\kfChoice 정약용이 오직 내면 수양만을 중시하며 사회적 실천을 경시했기 때문이다.
\kfChoice 성리학이 실천을, 정약용이 내면 성찰을 강조하여 입장이 뒤바뀌기 때문이다.
\end{kfChoices}
\end{kfQBlock}

\begin{kfQBlock}{3}{}
\kfQStem{윗글을 바탕으로 <보기>를 이해한 내용으로 적절하지 않은 것은?

<보기> F 씨는 길에서 쓰러진 노인을 보고 측은한 마음이 들었다. 그러나 약속 시간에 늦을까 봐 망설이다가, 결국 노인을 부축하여 병원에 데려다 주었다.}
\begin{kfChoices}
\kfChoice F 씨가 측은한 마음을 느낀 것은 정약용이 말하는 성기호, 즉 선을 좋아하는 경향성이 발현된 것이겠군.
\kfChoice F 씨가 약속 시간과 노인 사이에서 망설인 것은 도덕적 기호와 형구의 기호 사이의 갈등이겠군.
\kfChoice F 씨가 결국 노인을 도운 것은 자주지권을 통해 선을 선택한 것이겠군.
\kfChoice 정약용의 관점에서 F 씨의 측은한 마음은 실제 행동으로 옮겨졌을 때 비로소 인(仁)이 되겠군.
\kfChoice 정약용의 관점에서 F 씨의 측은한 마음만으로도 이미 인(仁)이 완성되었겠군.
\end{kfChoices}
\end{kfQBlock}

\begin{kfQBlock}{4}{}
\kfQStem{윗글의 '행사지심(行事之心)'이 의미하는 바로 가장 적절한 것은?}
\begin{kfChoices}
\kfChoice 도덕 감정에 머무르지 않고 실제 행동으로 옮기는 실천적 마음
\kfChoice 외부 세계에 대한 이론적 관조를 수행하는 지적 능력
\kfChoice 미발 상태에서 선한 본성이 자동으로 발현되는 과정
\kfChoice 감각적 욕구를 억제하여 마음의 고요함을 유지하는 수양법
\kfChoice 사회 제도의 개혁과 무관한 순수한 내면 수양
\end{kfChoices}
\end{kfQBlock}

\end{multicols}

\kfSecHeader{kfAreaNonlit}{지문}{롤스의 분배 정의와 노직의 소유권 정의 (인문)}

\kfPassageNote{롤스의 분배 정의와 노직의 소유권 정의 (인문)}{%
존 롤스(John Rawls)는 『정의론』(1971)에서 사회 정의의 원칙을 도출하기 위해 '원초적 입장(original position)'과 '무지의 베일(veil of ignorance)'이라는 사고실험을 제시했다. 원초적 입장이란 사회 계약을 맺기 위한 가상적 상황이며, 무지의 베일이란 계약 당사자들이 자신의 사회적 지위, 재능, 성별, 인종 등을 모르는 상태를 말한다. 이 상태에서 합리적 개인들이 합의할 정의의 원칙은 무엇인가?

롤스에 따르면 합의될 원칙은 두 가지이다. 제1원칙은 '평등한 자유의 원칙'으로, 모든 사람은 타인의 자유와 양립 가능한 범위에서 최대한의 기본적 자유를 누려야 한다. 제2원칙은 '차등 원칙(difference principle)'과 '공정한 기회 균등의 원칙'으로 구성된다. 차등 원칙에 따르면 사회적·경제적 불평등은 오직 사회에서 가장 불리한 처지에 있는 구성원에게 최대의 이익이 되는 경우에만 정당화된다. 이는 능력주의적 분배가 초래하는 불평등을 사회 정의의 차원에서 교정하려는 시도이다.

로버트 노직(Robert Nozick)은 『무정부, 국가, 유토피아』(1974)에서 롤스를 정면으로 비판했다. 노직에 따르면 정의는 분배의 결과적 패턴이 아니라 소유의 과정에 있다. 그의 '소유 권원(entitlement) 이론'은 세 원칙으로 구성된다. 첫째, 최초 취득의 정의 — 누구의 것도 아닌 자원을 정당하게 취득한 것은 정의롭다. 둘째, 이전의 정의 — 정당하게 소유한 것을 자발적으로 교환·증여한 결과는 정의롭다. 셋째, 교정의 정의 — 앞의 두 원칙이 위반된 경우 그것을 바로잡는 것이 정의이다.

노직의 핵심 논점은 결과적 패턴에 기반한 재분배가 개인의 자유와 재산권을 침해한다는 것이다. 그는 '윌트 체임벌린 논증'을 제시한다. 농구 스타 체임벌린이 경기 관람료의 일부를 받기로 계약하고, 관중이 자발적으로 관람료를 지불하여 체임벌린이 거액을 벌었다고 하자. 모든 과정이 자발적이므로 결과는 정의롭다. 그런데 국가가 세금을 부과하여 이 소득을 재분배한다면, 이는 체임벌린과 관중 모두의 자유로운 선택을 침해하는 것이다. 노직에게 '결과의 평등'을 추구하는 재분배는 곧 개인 자유의 침해이며, 최소 국가(minimal state)만이 정당하다.

이에 대해 롤스주의자들은 반론한다. 첫째, 개인의 재능 자체가 '도덕적으로 임의적(morally arbitrary)'이다. 뛰어난 재능을 갖고 태어난 것은 개인의 공로가 아니므로, 재능에서 비롯된 불평등을 사회적으로 교정하는 것은 자유 침해가 아니라 공정의 실현이다. 둘째, 무지의 베일 아래에서 누구도 자신이 최하위 계층에 속할 가능성을 배제할 수 없으므로, 합리적 개인이라면 차등 원칙에 동의할 것이다.
}
\kfSecHeader{kfAreaNonlit}{문제}{}

\begin{multicols}{2}\raggedcolumns\setlength{\columnsep}{12pt}
\begin{kfQBlock}{1}{}
\kfQStem{윗글의 내용과 일치하지 않는 것은?}
\begin{kfChoices}
\kfChoice 롤스의 차등 원칙에 따르면 불평등은 사회에서 가장 불리한 처지의 구성원에게 최대의 이익이 될 때 정당화된다.
\kfChoice 노직은 정의가 분배의 결과적 패턴이 아니라 소유의 과정에 있다고 주장했다.
\kfChoice 노직의 소유 권원 이론은 최초 취득, 이전, 교정의 세 원칙으로 구성된다.
\kfChoice 노직은 롤스의 차등 원칙이 개인의 자유와 재산권을 충분히 보호한다고 동의했다.
\kfChoice 롤스주의자들은 개인의 재능이 도덕적으로 임의적이라고 주장한다.
\end{kfChoices}
\end{kfQBlock}

\begin{kfQBlock}{2}{}
\kfQStem{윗글에 따를 때, 노직이 '윌트 체임벌린 논증'을 통해 주장하려는 바로 가장 적절한 것은?}
\begin{kfChoices}
\kfChoice 자발적 거래에서 비롯된 불평등을 재분배로 바로잡는 것은 개인의 자유를 침해한다.
\kfChoice 농구 스타는 도덕적 의무에 따라 소득의 상당 부분을 사회에 환원해야 한다.
\kfChoice 관중이 지불하는 관람료는 진정으로 자발적이지 않으므로 결과가 정의롭지 않다.
\kfChoice 국가가 부과하는 세금은 롤스의 차등 원칙을 실현하는 수단이므로 정당하다.
\kfChoice 체임벌린의 소득은 도덕적으로 임의적인 결과이므로 마땅히 재분배되어야 한다.
\end{kfChoices}
\end{kfQBlock}

\begin{kfQBlock}{3}{}
\kfQStem{윗글의 롤스와 노직을 비교한 내용으로 적절하지 않은 것은?}
\begin{kfChoices}
\kfChoice 롤스는 분배의 결과적 패턴을, 노직은 소유의 과정을 정의의 기준으로 삼는다.
\kfChoice 롤스는 불평등의 사회적 교정을, 노직은 개인의 재산권 보호를 우선시한다.
\kfChoice 롤스는 무지의 베일이라는 가상적 장치를 사용하고, 노직은 윌트 체임벌린이라는 구체적 사례를 사용한다.
\kfChoice 롤스와 노직 모두 결과의 평등을 추구하는 재분배를 지지한다.
\kfChoice 롤스는 재능을 도덕적으로 임의적이라 보는 반면, 노직은 재능에서 비롯된 소득을 자유로운 선택의 결과로 본다.
\end{kfChoices}
\end{kfQBlock}

\begin{kfQBlock}{4}{}
\kfQStem{윗글을 바탕으로 <보기>에 대한 각 이론의 판단으로 적절하지 않은 것은?

<보기> 어떤 사회에서 상위 10\%가 전체 부의 80\%를 소유하고, 최하위 계층은 기본적 생활이 어렵다.}
\begin{kfChoices}
\kfChoice 롤스의 차등 원칙에 따르면, 이 불평등이 최하위 계층에게 최대 이익이 되지 않는 한 정당화될 수 없다.
\kfChoice 노직에 따르면, 상위 10\%의 부가 정당한 취득·이전의 과정을 거쳤다면 그 자체로 정의롭다.
\kfChoice 롤스주의자는 상위 10\%의 부를 재분배하여 최하위 계층의 처지를 개선하는 것이 공정하다고 볼 것이다.
\kfChoice 노직은 최하위 계층의 기본 생활을 보장하기 위해 국가가 대규모 재분배를 실시해야 한다고 주장할 것이다.
\kfChoice 롤스의 관점에서 재능에서 비롯된 불평등은 도덕적으로 임의적이므로 사회적 교정의 대상이 된다.
\end{kfChoices}
\end{kfQBlock}

\end{multicols}

\kfSecHeader{kfAreaNonlit}{지문}{레비나스의 타자 윤리학 (인문)}

\kfPassageNote{레비나스의 타자 윤리학 (인문)}{%
에마뉘엘 레비나스(Emmanuel Lévinas)는 서양 철학의 전통을 '동일자(le Même)의 철학'이라 진단하고, 이를 근본적으로 전복하려 했다. 플라톤 이래 서양 형이상학은 다양한 존재자를 하나의 보편적 원리(이데아, 실체, 주체 등)로 포섭하려 해 왔다. 레비나스는 이러한 전통이 타자(l'Autre)를 동일자의 틀 안으로 환원하여 타자의 진정한 타자성(altérité)을 말살한다고 비판한다.

레비나스에게 타자는 인식의 대상이 아니라, 나의 인식 능력을 넘어서는 무한한 존재이다. 타자의 '얼굴(visage)'은 이러한 무한의 현현(顯現)이다. 얼굴은 물리적 외모가 아니라, 나에게 윤리적 요청을 전달하는 표현이다. 타자의 얼굴은 「살인하지 말라」라는 명령을 스스로 발한다. 이 명령은 어떤 보편적 도덕 법칙에서 연역되는 것이 아니라, 타자의 얼굴과의 대면 자체에서 직접적으로 발생하는 것이다.

레비나스의 핵심 주장은 '윤리학이 제1철학이다(l'éthique est la philosophie première)'라는 것이다. 서양 철학의 전통에서 제1철학은 존재론이었다. 아리스토텔레스는 '존재하는 것으로서의 존재'를 탐구하는 학문을 제1철학이라 불렀다. 그러나 레비나스는 존재론이 타자를 동일자로 환원하는 폭력의 학문이라 비판하며, 타자에 대한 윤리적 책임이 존재론에 선행한다고 주장한다. 내가 '있다'는 것보다 타자에 대한 책임이 더 근원적이다.

이 책임은 비대칭적이고 무한하다. 비대칭적이란, 내가 타자에 대해 지는 책임은 타자가 나에 대해 지는 책임과 대칭적이지 않다는 뜻이다. 나는 타자의 응답을 기대하지 않고 일방적으로 책임을 진다. 무한하다는 것은, 이 책임에 끝이 없으며 완수될 수 없다는 의미이다. 레비나스는 이를 「나는 타자의 인질(otage)이다」라는 극단적 표현으로 요약한다. 나는 내가 선택하지 않은 책임에 의해 묶여 있으며, 이 책임으로부터 도피할 수 없다. 이것이 레비나스가 말하는 주체의 본래적 구조이다. 주체는 자유로운 자기 정초의 주인이 아니라, 타자의 부름에 응답하는 책임의 존재이다.
}
\kfSecHeader{kfAreaNonlit}{문제}{}

\begin{multicols}{2}\raggedcolumns\setlength{\columnsep}{12pt}
\begin{kfQBlock}{1}{}
\kfQStem{윗글의 내용과 일치하지 않는 것은?}
\begin{kfChoices}
\kfChoice 레비나스는 서양 철학의 전통이 타자를 동일자의 틀로 환원한다고 비판했다.
\kfChoice 타자의 얼굴은 물리적 외모가 아니라 윤리적 요청을 전달하는 표현이다.
\kfChoice 레비나스에게 존재론은 윤리학에 선행하는 제1철학이다.
\kfChoice 타자에 대한 책임은 비대칭적이고 무한하다.
\kfChoice 레비나스의 주체는 타자의 부름에 응답하는 책임의 존재이다.
\end{kfChoices}
\end{kfQBlock}

\begin{kfQBlock}{2}{}
\kfQStem{윗글에 따를 때, 레비나스가 존재론을 비판하는 이유로 가장 적절한 것은?}
\begin{kfChoices}
\kfChoice 존재론이 존재라는 주제 자체에 대해 전혀 탐구하지 않기 때문이다.
\kfChoice 존재론이 타자를 동일자의 틀로 환원해 그 타자성을 말살하기 때문이다.
\kfChoice 존재론이 타자에 대한 윤리적 책임만을 지나치게 과도하게 강조하기 때문이다.
\kfChoice 존재론이 사람의 물리적 외모와 표면적 특징에만 관심을 기울이기 때문이다.
\kfChoice 존재론이 타자의 자유와 권리를 침해하지 않고 그대로 보장하기 때문이다.
\end{kfChoices}
\end{kfQBlock}

\begin{kfQBlock}{3}{}
\kfQStem{윗글의 '나는 타자의 인질이다'라는 레비나스의 표현이 의미하는 바로 가장 적절한 것은?}
\begin{kfChoices}
\kfChoice 주체가 선택하지 않은 무한한 책임에 묶여 도피할 수 없다는 것
\kfChoice 주체가 타자에게 신체적·물리적으로 감금된 상태에 놓여 있다는 것
\kfChoice 타자가 주체를 향해 강력한 폭력적 지배권을 행사하고 있다는 것
\kfChoice 주체는 타자의 인질이 되어야만 비로소 진정한 자유를 얻을 수 있다는 것
\kfChoice 주체가 지는 책임과 타자가 지는 책임은 서로 완전히 대칭적이라는 것
\end{kfChoices}
\end{kfQBlock}

\begin{kfQBlock}{4}{}
\kfQStem{윗글을 바탕으로 <보기>를 이해한 내용으로 적절하지 않은 것은?

<보기> G 씨는 길에서 낯선 노숙인과 눈이 마주쳤다. 그 순간 G 씨는 말로 표현할 수 없는 윤리적 무게를 느꼈고, 그 자리를 떠나지 못했다.}
\begin{kfChoices}
\kfChoice G 씨가 노숙인과 눈이 마주친 것은 레비나스가 말하는 타자의 '얼굴'과의 대면에 해당하겠군.
\kfChoice G 씨가 느낀 윤리적 무게는 보편적 도덕 법칙에서 연역된 것이 아니라 대면 자체에서 발생한 것이겠군.
\kfChoice G 씨가 그 자리를 떠나지 못한 것은 타자에 대한 책임으로부터 도피할 수 없다는 점을 보여주겠군.
\kfChoice 레비나스의 관점에서 G 씨는 노숙인에 대해 대칭적 책임, 즉 동등한 책임을 기대할 수 있겠군.
\kfChoice G 씨의 경험은 윤리적 책임이 존재론적 물음보다 더 근원적임을 보여주는 사례이겠군.
\end{kfChoices}
\end{kfQBlock}

\end{multicols}

\kfStartWeekly{패턴 워크북 — 총 ?문항}


\kfSecHeader{kfAreaWeekly}{지문}{문항 1 지문 / 섞어치기}

\kfPassageNote{문항 1 지문 / 섞어치기}{%
개별적이고 특수한 작업을 수행하는 여러 대의 클라이언트와 연결되어 정보나 서비스를 제공하는 컴퓨터 시스템을 서버라고 한다. 서버는 클라이언트의 요청에 따라 작업을 수행하거나 정보를 검색하고, 그 결과를 클라이언트에 돌려준다. 서버를 구성하는 기계에서 작업을 신속하고 정확하게 처리하기 위해서는 데이터를 안전하게 관리해야 할 뿐 아니라 신속하게 데이터에 접근하여 데이터의 읽기나 쓰기를 실행할 필요가 있다. 이런 필요를 충족시키기 위해 개발된 것이 ㉠분산 파일 관리 체계이다. 분산 파일 관리 체계는 데이터의 불법 변조를 방지하고, 기계적 결함에 대응하여 데이터의 유실로부터 안전을 도모하며, 서버의 부하를 줄여 신속한 연산을 수행하도록 기계적 성능과 저장 용량을 확보하는 데 유리한 전략이다. 단일 기계로 구성된 서버에서 데이터를 관리할 경우에는 하드웨어적으로 램이나 CPU의 성능이 우수한 기계를 확보하는 데 큰 비용이 들지만, 분산 파일 관리 체계에서는 다수의 저렴한 저성능 기계를 연결하여

(중략)
}
\begin{kfQBlock}{1}{}
\kfQStem{㉮와 문맥상 의미가 가장 가까운 것은?}
\begin{kfChoices}
\kfChoice 합의가 깨져 협상이 중단되었다.
\kfChoice 유리컵이 바닥에 떨어져 산산이 깨졌다.
\kfChoice 넘어지면서 무릎이 깨져 피가 났다.
\end{kfChoices}
\end{kfQBlock}

\kfSecHeader{kfAreaWeekly}{지문}{문항 2 지문 / 부정상대어}

\kfPassageNote{문항 2 지문 / 부정상대어}{%
(가) 가상 공간에서 영생의 삶을 누린다는 흥미로운 이야기가 SF 영화의 소재로만 그치지 않는다고 주장하는 입장이 있다. 커즈와일이나 보스트룀은 기억, 가치관, 태도, 정서적 성향처럼 인간의 자아를 구성하는 것들을 특정 정보 패턴으로 만들어 보존할 수 있다면 신체적인 죽음 이후에도 인간이 여전히 생존할 수 있다고 주장한다. 이 입장은 인격 동일성에 관한 심리적 연속성 이론과 관련이 있다. 이 이론에서는 ‘나’가 누구인지를 규정하는 것은 ‘나’의 기억, 신념, 생각, 감정, 희망, 두려움 등의 묶음이라고 보며, 이러한 심리적인 특성들의 집합이 유지·보존되느냐에 ‘나’의 지속 여부가 달려 있다고 말한다.

(중략)

커즈와일이나 보스트룀은 마음이나 정신 상태를 의미하는 심적 상태의 본성을 계산 기능주의의 관점에서 이해한다. 계산 기능주의에서는 심적 상태의 독특한 인과적 역할이 두뇌의 물리적 상태에 의해서 이루어진다고 보기 때문에 정신 작용을 두뇌에 구현된 계산 체계의 작동으로 이해하며, 인간의 사고나

(중략)
}
\begin{kfQBlock}{2}{}
\kfQStem{윗글의 내용과 일치하는 것으로 가장 적절한 것은?}
\begin{kfChoices}
\kfChoice 복수 실현 가능성을 인정해도 물리적 동일성의 자동 보장은 성립하지 않는다고 본다.
\kfChoice 계산 기능주의는 정신 상태를 물질 구성과 무관한 순수 도덕 범주로 본다.
\kfChoice 체화된 인지는 신체 종류가 달라도 개념 이해 방식은 항상 같다고 본다.
\end{kfChoices}
\end{kfQBlock}

\kfSecHeader{kfAreaWeekly}{지문}{문항 3 지문 / 주체객체}

\kfPassageNote{문항 3 지문 / 주체객체}{%
일반적으로 경제학에서는 경제 주체들이 거래에 관한 완전한 정보를 지니고 있다고 가정한다. 완전 정보 상태가 성립하려면 모든 경제 주체들이 원하는 정보를 아무런 비용을 들이지 않고 얻을 수 있어야 한다. 그러나 현실에서는 정보 취득에 비용이 소요되므로 경제 주체들은 불완전한 정보를 지닐 수밖에 없다. 더욱이 각 경제 주체마다 비용 제약이 상이하기 때문에 취득하는 정보의 양과 질에 차이가 발생하는데, 이러한 상황을 [정보 비대칭 상황]이라고 한다.

(중략)

대표적인 정보 비대칭 상황은 거래의 한쪽이 상대방이나 상품의 특성에 관한 정보를 상대적으로 적게 가지고 있는 경우이다. 예를 들어 보험 회사는 보험에 가입하려는 사람에게 사고가 발생할 가능성, 즉 고객의 위험 수준을 알기 어렵다. 이러한 상황에서 정보가 부족한 쪽은 역선택 문제에 직면한다. 역선택이란 정보 비대칭으로 인해 정보가 부족한 경제 주체가 정보를 많이 가진 경제 주체 중 가장 바람직하지 않은 상대와 거래하게 되는 현상이다. 보험 회사가

(중략)
}
\begin{kfQBlock}{3}{}
\kfQStem{윗글을 이해한 내용으로 적절하지 않은 것은?}
\begin{kfChoices}
\kfChoice 품질 보증 제도는 소비자가 상품의 품질을 파악할 수 있도록 하기 위한 신호 발송 장치로 기능한다.
\kfChoice 리뷰 시스템과 평점 제도는 정보 우위에 있는 측이 정보가 적은 측보다 상대방의 숨겨진 특성을 더욱 용이하게 파악할 수 있게 해 준다.
\end{kfChoices}
\end{kfQBlock}

\kfSecHeader{kfAreaWeekly}{지문}{문항 4 지문 / 순서방향}

\kfPassageNote{문항 4 지문 / 순서방향}{%
육상 동물과 수생 동물은 각각의 생활 환경에 따라 호흡 과정에서 어려움을 겪는다. 육상 동물은 폐의 표면에 건조한 공기가 닿으면 폐가 건조해질 수 있는데, 이는 폐를 손상시키고 ⓐ탈수를 일으키며 기체를 용해하는 폐의 능력을 떨어뜨린다. 한편 수생 동물은 가용 산소가 적은 수중 환경과 수온 차에 의해 변화되는 산소 가용성에 적응해야 한다. 또한 수생 동물은 물속에서 몸을 움직이는 데 많은 근육 운동이 필요하고, 찬 바닷물의 높은 비열 때문에 아가미의 열을 빼앗기며, 삼투압 현상으로 인해 수분 균형에 문제가 발생하기도 한다.

(중략)

이러한 환경 속에서 수생 동물은 효과적인 호흡을 위해 아가미라는 특수한 기관을 사용한다. 아가미는 물속에 녹아 있는 산소를 흡수하고 체내에 있는 이산화 탄소를 배출하는 기능을 하는데, 외아가미와 내아가미로 나뉜다. 외아가미는 몸 밖으로 ⓑ돌출된 형태로, 표면적이 넓고 형태가 다양하며 유생기 도롱뇽이나 일부 무척추동물에서 나타난다. 외아가미는 주변 물과의 접촉을 극대화

(중략)
}
\begin{kfQBlock}{4}{}
\kfQStem{Ⓐ의 이유로 적절한 것은?}
\begin{kfChoices}
\kfChoice 물의 흐름으로 인해 구심성 혈관에서 원심성 혈관으로 흐르던 혈액의 방향이 반대 방향으로 바뀌기 때문에
\kfChoice 산소를 함유한 물이 혈액과 반대 방향으로 이동하며 계속해서 산소 농도가 낮은 혈액과 접촉하기 때문에
\end{kfChoices}
\end{kfQBlock}

\kfSecHeader{kfAreaWeekly}{지문}{문항 5 지문 / 인과도치}

\kfPassageNote{문항 5 지문 / 인과도치}{%
행정 대집행은 공법상 대체적 작위 의무를 의무자가 불이행한 경우에 행정청이 의무자 대신 직접 행위를 하거나 제삼자에게 이를 하게 하고, 그 비용을 의무자로부터 ⓐ징수하는 제도이다. ㉠대체적 작위 의무란 타인이 대신하여 이행할 수 있는 적극적 행위 의무를 의미하며, 건물의 철거, 불법 시설물의 제거, 오염 물질의 제거 의무 등이 해당한다. 이는 어떠한 행위를 하여서는 안 되는 소극적 의무인 ㉡부작위 의무나 오직 의무자만이 이행할 수 있는 의무인 ㉢비대체적 작위 의무와는 구별된다. 행정 대집행은 국민의 재산권 등 기본권에 대한 중대한 제한을 ⓑ수반하므로 반드시 법적 근거가 있어야 한다. 헌법상 법치 행정의 원리에 따라, 모든 행정 작용은 법률에 근거하여 이루어져야 하기 때문이다. 따라서 행정상 의무를 명할 수 있는 명령권의 근거가 되는 법이 동시에 강제로 행정을 집행하는 근거가 될 수는 없으며, 대집행을 위한 별도의 법적 근거가 필요하다.

(중략)

행정 대집행법상 대집행이 가능하기 위해서는 다음과

(중략)
}
\begin{kfQBlock}{5}{}
\kfQStem{㉮의 이유를 추론한 내용으로 가장 적절한 것은?}
\begin{kfChoices}
\kfChoice 퇴거 의무는 점유 이전이라는 인도 의무와는 성격이 다르므로, 대집행의 대상이 될 수 없는 비대체적 작위 의무에 해당하지 않기 때문이다.
\kfChoice 건물 철거를 위해서는 철거 의무자의 퇴거가 선행되어야 하므로, 퇴거 행위는 주된 행위인 철거 의무의 대집행에 포함되는 부수적 행위로 볼 수 있기 때문이다.
\end{kfChoices}
\end{kfQBlock}

\kfSecHeader{kfAreaWeekly}{지문}{문항 6 지문 / 의도/목적}

\kfPassageNote{문항 6 지문 / 의도/목적}{%
국민들의 다양한 행정 수요에 효과적으로 대응하기 위해 행정에는 국가, 지방 자치 단체, 공공 단체 등 다양한 주체들이 기능한다. 공공 단체에는 공공 조합이나 공법상 재단 등이 포함된다. 국가의 일반 행정 조직과는 달리, 법적 근거에 따라 특정 행정 분야를 전문적으로 수행하는 조직도 있는데, 이를 특별 공공 행정 조직이라고 부른다. 특별 공공 행정 조직은 국가가 직접 담당하기 \uline{ⓐ어려운} 행정 영역을 보완하거나 정치적 중립성과 기능적 효율성이 요구되는 업무를 수행하는 것을 목적으로 하며, 국가 기관의 형태로 설립되기도 하고 공공 단체의 형태로 설립되기도 한다.

특별 공공 행정 조직 중에는 헌법에 의해 설립된 국가 기관이 있는데, 선거 관리 위원회, 감사원이 이에 해당한다. 이들 기관은 국가 권력의 분립과 견제라는 헌법적 원리에 따라 설립된 것이다. 따라서 그 설립과 폐지는 단순한 법률 개정만으로는 불가능하고 반드시 헌법 개정을 통해서만 가능하며, 독립성과 공정성을 핵심 운영…(중략)
}
\begin{kfQBlock}{6}{}
\kfQStem{㉠의 이유로 가장 적절하지 않은 것은?}
\begin{kfChoices}
\kfChoice 해당 기관의 정치적 중립성을 확보하기 위해서
\kfChoice 해당 기관의 정치적 중립성을 확보하기 위해서
\kfChoice 해당 기관의 법적 지위를 보장하기 위해서
\kfChoice 해당 기관의 조직 규모를 확대하기 위해서
\end{kfChoices}
\end{kfQBlock}

\kfSecHeader{kfAreaWeekly}{지문}{문항 7 지문 / 상관관계 반전}

\kfPassageNote{문항 7 지문 / 상관관계 반전}{%
식물은 동물과 달리 생애 동안 지속적으로 성장하는 개체이며, 이러한 생장은 분열 조직의 활동에 의해 유지된다. 식물의 분열 조직은 크게 측생 분열 조직과 정단 분열 조직으로 나뉘는데, 측생 분열 조직은 줄기와 뿌리의 측면에서 부피 생장을, 정단 분열 조직은 줄기와 뿌리 끝에서 길이 생장을 담당한다.

측생 분열 조직의 대표로는 관다발 형성층과 코르크 형성층이 있다. 관다발 형성층은 물관과 체관 사이에서 새로운 관다발 조직을 만들고, 코르크 형성층은 표피 아래에서 보호 조직을 형성해 수분 손실을 \uline{ⓐ막는} 역할을 한다. 이들의 지속적 활동은 줄기와 뿌리를 굵고 단단하게 유지하며 나이테 형성과도 관련된다.

정단 분열 조직에는 뿌리 정단 분열 조직과 슈트 정단 분열 조직이 있다. 뿌리 끝은 뿌리골무, 분열대, 신장대, 성숙대로 구분되며 분열대에서 세포 분열이 가장 활발하다. 시원 세포는 스스로를 복제해 새로운 세포를 공급하고, 신장대에서는 세포가 길게 늘어나며, 성숙대에서는 세…(중략)
}
\begin{kfQBlock}{7}{}
\kfQStem{㉠의 함의로 가장 적절하지 않은 것은?}
\begin{kfChoices}
\kfChoice 초기 두 학설은 모든 식물에서 예외 없이 완전 검증되었다.
\kfChoice 초기 두 학설은 모든 식물에서 예외 없이 완전 검증되었다.
\kfChoice 초기 학설 일부는 특정 식물군에는 설명력이 있었으나 보편성 한계가 있었다.
\kfChoice 후속 학설은 이전 가설의 한계를 보완하는 방향으로 제시되었다.
\end{kfChoices}
\end{kfQBlock}

\kfSecHeader{kfAreaWeekly}{지문}{문항 8 지문 / 필요조건}

\kfPassageNote{문항 8 지문 / 필요조건}{%
개별적이고 특수한 작업을 수행하는 여러 대의 클라이언트와 연결되어 정보나 서비스를 제공하는 컴퓨터 시스템을 서버라고 한다. 서버는 클라이언트의 요청에 따라 작업을 수행하거나 정보를 검색하고, 그 결과를 클라이언트에 돌려준다. 서버를 구성하는 기계에서 작업을 신속하고 정확하게 처리하기 위해서는 데이터를 안전하게 관리해야 할 뿐 아니라 신속하게 데이터에 접근하여 데이터의 읽기나 쓰기를 실행할 필요가 있다. 이런 필요를 충족시키기 위해 개발된 것이 \uline{㉠분산 파일 관리 체계}이다. 분산 파일 관리 체계는 데이터의 불법 변조를 방지하고, 기계적 결함에 대응하여 데이터의 유실로부터 안전을 도모하며, 서버의 부하를 줄여 신속한 연산을 수행하도록 기계적 성능과 저장 용량을 확보하는 데 유리한 전략이다. 단일 기계로 구성된 서버에서 데이터를 관리할 경우에는 하드웨어적으로 램이나 CPU의 성능이 우수한 기계를 확보하는 데 큰 비용이 들지만, 분산 파일 관리 체계에서는 다수의 저렴한 저성능 기…(중략)
}
\begin{kfQBlock}{8}{}
\kfQStem{ⓐ를 구현할 때 <보기>의 ⓑ나 ⓒ의 활용에 대해 설명한 내용으로 적절한 것은?}
\begin{kfEvidence}<보기> 복제를 구현하려면 원본 데이터 세트와 그것 전체를 복사한 복제본 데이터 세트를 하나 이상 마련하여 저장해야 한다. 복제는 원본과 복제본을 구분하여 취급할 것이냐, 아니냐에 따라 ⓑ마스터 슬레이브 방법과 ⓒ피어 투 피어 방법으로 나뉜다. 마스터 슬레이브 방법은 원본과 복제본을 구분하여 취급하는 복제이다. 원본 데이터 세트를 마스터라고 부르고 복제본 데이터 세트를 슬레이브라고 부른다. 읽기 요청은 마스터와 슬레이브에서 모두 분산하여 처리할 수 있지만 쓰기 요청은 마스터에서만 처리한다. 이렇게 마스터에서 수정된 데이터 세트는 슬…(중략)\end{kfEvidence}
\begin{kfChoices}
\kfChoice ⓑ를 채택할 때에는 복수의 기계에 동일한 내용의 마스터가 샤드로 존재하게 된다.
\kfChoice ⓑ를 채택할 때에는 복수의 기계에 동일한 내용의 마스터가 샤드로 존재하게 된다.
\kfChoice ⓑ를 채택할 때에는 하나의 기계에 다른 내용의 마스터와 슬레이브가 샤드로 존재하게 된다.
\kfChoice ⓒ를 채택할 때에는 하나의 기계에 다른 내용의 샤드들이 피어로 존재하게 된다.
\end{kfChoices}
\end{kfQBlock}

\kfSecHeader{kfAreaWeekly}{지문}{문항 9 지문 / 결합/분리}

\kfPassageNote{문항 9 지문 / 결합/분리}{%
현대 경제는 과거와는 다른 양상으로 움직인다. 기술은 시시각각 변하고 세계 금융 시장은 얽히고설켜 예측하기 어려운 변동성을 보인다. 이러한 복잡한 경제 상황을 이해하기 위해 미국의 산타페 연구소를 중심으로 등장한 복잡계 경제학은 경제 주체가 언제나 가장 합리적 선택을 한다는 명제는 진리가 아니라고 주장한다. 그러면서 경제에 대한 근본적인 생각의 전환을 요구한다.

복잡계 경제학에서는 현대 경제를 수학적 계산이나 도식적 모델의 결과물이 아니라, 동적으로 변화하는 복잡한 시스템이라고 본다. 또한 현대 경제의 특징은 다양한 주체들이 제한된 정보 속에서 상호 작용하며 예측 불가능한 질서를 만들어 내는 비정형성이라고 설명한다. 어떤 기술이 표준이 될지, 어떤 기업이 시장을 주도할지, 어떤 정책이 장기적으로 성공할지는 누구도 확신할 수 없다는 것이다. 예를 들어 기술의 선택에서는 효율성보다는 경로 의존성이 더 크게 작용하기 때문에 새로운 기술 개발의 효과를 예측하는 것은 매우 어렵다. 여기서 …(중략)
}
\begin{kfQBlock}{9}{}
\kfQStem{[01] 현대 경제에 대한 '복잡계 경제학'의 관점에 부합하지 않는 것은?}
\begin{kfChoices}
\kfChoice 자발적으로 형성된 시장의 질서는 개별 주체들에서 보편적으로 발견되는 특성이다.
\kfChoice 자발적으로 형성된 시장의 질서는 개별 주체들에서 보편적으로 발견되는 특성이다.
\kfChoice 경로 고착은 정책이나 제도에서도 나타난다.
\kfChoice 경제 주체들의 상호 작용에 의해 창발이 나타난다.
\end{kfChoices}
\end{kfQBlock}

\kfSecHeader{kfAreaWeekly}{지문}{문항 10 지문 / 조건 왜곡}

\kfPassageNote{문항 10 지문 / 조건 왜곡}{%
인간은 고기와 가죽을 얻거나 실험에 이용하는 등의 목적으로 동물을 이용한다. 그리고 그런 이용은 일반적으로 동물에게 해악을 끼치는 방식으로 이루어지기에 잘못된 행동으로 인식된다. 동물에게 이렇게 해악을 끼치는 것을 정당화할 방법이 있을까? 한 가지 대답은 동물은 단지 동물이기에 이런 식으로 사용하는 것이 정당하다는 것이다. \uline{㉠‘종(種) 차별주의’}는 이렇게 우리가 속한 종의 이익을 다른 종의 이익보다 우선시하는 것을 말한다. 동물은 우리 종에 속하지 않기 때문에 인간에게 절대 하지 않는 방식으로 동물을 대할 수 있다는 것이다.

종 차별주의를 옹호하는 가장 직관적이고 상식적인 방법은 인간은 인간이라는 종에 속하므로 특별한 대우를 받아도 된다고 주장하는 것이다. 그런데 나와 같은 인간이므로 더 중요하게 대우받아야 한다는 주장이 허용된다면, 나와 같은 성별이기에 더 중요하게 대우받아야 한다는 ‘성차별주의’나, 나와 같은 인종이므로 더 중요하게 대우받아야 한다는 ‘인종 차별주의…(중략)
}
\begin{kfQBlock}{10}{}
\kfQStem{윗글의 핵심 주장으로 가장 적절한 것은?}
\begin{kfChoices}
\kfChoice 종 차별주의를 도덕적으로 옹호하기 위해서는 궁극적으로 종 중첩 논증이 해결되어야 한다.
\kfChoice 종 차별주의를 도덕적으로 옹호하기 위해서는 궁극적으로 종 중첩 논증이 해결되어야 한다.
\kfChoice 고차원적 정신 능력이 제시되지 않으면 종 차별주의는 허용되지 않는다.
\kfChoice 인간의 배타적 특성은 도덕적인 의미가 없으므로 종 차별주의는 옹호되지 않는다.
\end{kfChoices}
\end{kfQBlock}

\kfSecHeader{kfAreaWeekly}{지문}{문항 11 지문 / 긍부정 반전}

\kfPassageNote{문항 11 지문 / 긍부정 반전}{%
\#\#\# [32～34] 다음 글을 읽고 물음에 답하시오.

(가)   청강 녹초변에 소 먹이는 아이들이   석양에 흥이 겨워 피리를 빗기 부니   물 아래 잠긴 용이 잠 깨어 일어날 듯   내 기운에 나온 학이 제 깃을 던져 두고 반공에 솟아 뜰 듯   소선(蘇仙)* 적벽은 추칠월이 좋다 하되   팔월 십오야를 모두 어찌 칭찬하는가   구름이 걷히고 물결이 다 잔 적에   하늘에 돋은 달이 솔 위에 걸렸거든   잡다가 빠진 줄이 적선(謫仙)\textit{이 헌사할샤}   \textit{공산에 쌓인 잎을 삭풍이 거둬 불어}   \textit{떼구름 거느리고 눈조차 몰아오니}   \textit{천공이 호사로워 옥으로 꽃을 지어}   \textit{만수천림을 꾸며곰 낼세이고}   \textit{앞 여울 가리 얼어 독목교(獨木橋) 비꼈는데}   \textit{막대 멘 늙은 중이 어느 절로 간단 말고}   \textit{산옹의 이 부귀를 남더러 자랑 마오}   \textit{경요굴(瓊瑤窟)} 숨은 세계 찾을 이 있을세라   산중에 벗이 없어 서책을 쌓아 두고   만고 인물을 거슬러 혜여하니   성현도 많…(중략)
}
\begin{kfQBlock}{11}{}
\kfQStem{[A]에 대한 이해로 적절한 것은?}
\begin{kfChoices}
\kfChoice 하늘의 이치가 제대로 구현되지 못했음을 ‘시운’의 ‘흥망’에서 발견하고도 모를 일이 많다고 한 것에는, 인물의 담담한 태도가, 이상에 미치지 못하는 현실을 수용하는 것을 통해 드러난다.
\kfChoice 하늘의 이치가 제대로 구현되지 못했음을 ‘시운’의 ‘흥망’에서 발견하고도 모를 일이 많다고 한 것에는, 인물의 담담한 태도가, 이상에 미치지 못하는 현실을 수용하는 것을 통해 드러난다.
\kfChoice ‘삭풍’이 가을 잎을 쓸고 간 자리에 구름을 불러와 ‘공산’을 눈 세상으로 만들었다고 한 것에는, 인물이 거처한 공간의 아름다움에 대한 인식이 계절에 따른 자연의 변화를 통해 드러난다.
\kfChoice ‘앞 여울’을 건너가는 노승을 발견하고 ‘경요굴’이 들키지 않기를 바라는 것에는, 빼어난 경치를 소중하게 여기는 태도가, 숨어 있는 세계가 알려질 것에 대한 염려를 통해 드러난다.
\end{kfChoices}
\end{kfQBlock}

\kfSecHeader{kfAreaWeekly}{지문}{문항 12 지문 / 비교대상 뒤바꿈}

\kfPassageNote{문항 12 지문 / 비교대상 뒤바꿈}{%
[32\textasciitilde{}34] 다음 글을 읽고 물음에 답하시오.

(가)   동짓달에도 날씨가 며칠 푸근하면   철없는 개나리는 노란 얼굴 내민다   봄이 오면 꽃샘추위 아랑곳없이   진달래는 곳곳에 소담스럽게 피어난다   피어나는 꽃의 마음을   가냘프다고 / 억누를 수 있느냐   어두운 땅속으로 뻗어나가는 뿌리의 힘을   보이지 않는다고 / 업신여길 수 있느냐   땅에 깊숙이 뿌리내리고   하늘로 피어오르는 꿈을   드높은 가지 끝에 품은   나무처럼 젊은이들도   힘차게 위로 솟아오르고   ⓐ 조용히 아래로 깊어지며   밝고 넓게 퍼져 나가기를   그러나 행여 잊지 말기를   ⓑ 아무리 높다란 나뭇가지 끝에서   저 들판 너머를 볼 수 있어도   뿌리는 언제나 땅속에 있고 / 지하수가 수액이 되어   남모르게 줄기 속을 흐르지 않으면   바람결에 멀리 향냄새 풍기는   아카시아도 라일락도 / 절대로 피어날 수 없음을

* 김광규, ｢나무처럼 젊은이들도｣ -

(나)   ⓒ 사당역 4호선에서 2…(중략)
}
\begin{kfQBlock}{12}{}
\kfQStem{(가)와 (나)의 공통점으로 가장 적절한 것은?}
\begin{kfChoices}
\kfChoice 유사한 문장 구조를 반복하여 대상의 속성을 부각하고 있다.
\kfChoice 유사한 문장 구조를 반복하여 대상의 속성을 부각하고 있다.
\kfChoice 음성 상징어를 활용하여 대상의 역동성을 표현하고 있다.
\kfChoice 계절적 배경을 묘사하여 대상이 처한 상황을 드러내고 있다.
\end{kfChoices}
\end{kfQBlock}

\kfSecHeader{kfAreaWeekly}{지문}{문항 13 지문 / 주체/객체 혼동}

\kfPassageNote{문항 13 지문 / 주체/객체 혼동}{%
[31 \textasciitilde{} 34] 다음 글을 읽고 물음에 답하시오.   [앞부분 줄거리] 왕언의 딸 왕시는 홍관 땅의 김유령을 만나 혼인을 했지만 나라의 늙은 신하에 의해 이별하게 되었다.

김유령이 무릎을 꿇고 대답하였다.   “제 나이 스무 살 되었을 때 아내를 얻었는데, 나라의 노신하가 궁녀로 들이니 늘 서러워하며 지내고 있습니다. 세상일도 잊은 채, 다만 아내의 소식이나 한번 듣고 싶어 그것만을 희망하고 살고 있었습니다. 그런데 어느날 꿈에 선할아버님께서 이르시기를, ‘어찌 화산도사를 찾아가 보지 않는가? 그 도사가 못할 일이 없으니 네가 가보면 소원을 이룰 수 있으리라. 갈 때 돈 일만 관을 가져가라.’라고 하셨습니다. 그래서 꿈에서 깨어나자마자 돈을 장만하여 가지고 이렇게 온 것입니다.”   그러자 도사가 말했다.   “네 아내를 도로 밖으로 내어다 살고자 하느냐? 네 뜻을 자세히 말해라.”   김유령이 말했다.   “도로 내어다 살기야 바랄 수 있겠습니까? 그저 나와 하루만이라도 만나보…(중략)
}
\begin{kfQBlock}{13}{}
\kfQStem{윗글에 대한 이해로 적절하지 않은 것은?}
\begin{kfChoices}
\kfChoice 김유령은 도사에게 처음부터 숨김없이 소원을 말하였다.
\kfChoice 김유령은 도사에게 처음부터 숨김없이 소원을 말하였다.
\kfChoice 도사는 김유령에게 소원을 이루기 위한 과업을 제시하였다.
\kfChoice 김유령은 담당 관원에게 소청하여 왕시의 시신을 스무 날 안에 묻었다.
\end{kfChoices}
\end{kfQBlock}

\kfSecHeader{kfAreaWeekly}{지문}{문항 14 지문 / 내용과 방법 혼동}

\kfPassageNote{문항 14 지문 / 내용과 방법 혼동}{%
[31 \textasciitilde{} 34] 다음 글을 읽고 물음에 답하시오.

술이 알맞게 되었을 때, 청년 신사는 노래를 중지시키고, 예의 청산유수식 구변을 토하기 시작했다. ― 농촌 경제가 어떠니, 구태의연한 영농방법을 버리고 근대화를 해야 되느니, 그러기 위해서는 먼저 국민들의 비상한 각오가 필요하느니, 또 도시에 주택단지 공업단지가 서듯이 농촌에는 식량단지, 채소단지, 심지어 돼지단지까지 있어야 하느니 등, 그야말로 먼 앞날을 내다보는 ⓐ 유익한 얘기들이 꼬리를 물 듯 계속되었다.   옛날에는 권업계 서기요 지금은 산업계 서기들이 하는 말을 수타 들어왔기 때문에, 허 생원도 대강 짐작은 갔었지만, 결국 귀에 남는 것은 무슨 단지 단지 하는 새로운 말뿐이고, 청년이 말하는 〈먼 앞날〉보다 우선 코앞에 다가 있는 〈사는 문제〉가 더 절박했다.   “허 선생님은 이 고장 출신이시고, 또 누구보다 이곳 사정을 잘 아실 뿐 아니라 이해도 깊으실 터인 만큼 ―.”   드디어 청년 신사는 화제를 슬쩍 딴 데로 돌…(중략)
}
\begin{kfQBlock}{14}{}
\kfQStem{[A], [B]의 서술상 특징에 대한 설명으로 가장 적절하지 않은 것은?}
\begin{kfChoices}
\kfChoice [B]는 공간의 이동에 따른 인물의 행위를 제시하고 있다.
\kfChoice [B]는 공간의 이동에 따른 인물의 행위를 제시하고 있다.
\kfChoice [A]는 장면을 빈번하게 전환하여 긴박한 분위기를 조성하고 있다.
\kfChoice [B]는 내적 독백을 통해 사건의 흐름을 지연시키고 있다.
\end{kfChoices}
\end{kfQBlock}

\kfSecHeader{kfAreaWeekly}{지문}{문항 15 지문 / 내용 왜곡}

\kfPassageNote{문항 15 지문 / 내용 왜곡}{%
[22 \textasciitilde{} 27] 다음 글을 읽고 물음에 답하시오.

(가)   몰아라 어서 보자 총석정 어서 보자   총석정 좋단 말을 일찍이 들었거니   바람 불면 못 보려니 몰아라 어서 보자   벽해 위의 높은 집이 저것이 총석정인가   올라 보니 후면이라 전면으로 보오리라   배 대어라 사공들아 풍랑이 일지 않아   층파로 돌아 저어 총석 전면 보게 하라   배 띄워라 굽이마다 따라 저어 볼 양이면   영소전 태을궁*을 지으려고 경영턴가*   \textit{돌기둥 천백 개를 육모로 깎아 내어}   \textit{개개이 묶어 세워 몇 만 년이 되었던지}   \textit{황량한 데 벌였으니 배 없어 못 실린가}   \textit{(중략)}   \textit{하우씨 도끼뿔이 용문을 뚫었으나}   *이 돌*을 만났으면 이같이 깎을세며   영장*이 신묘하여 코끝의 것 찍었으나*   \textit{이 돌을 다듬는다고 이같이 곧을쏘냐}   \textit{어떠한 도끼로 용이히 깎았으며}   *어떠한 승묵*으로 천연히 골랐는고   끈 없이 묶었으되 틈 없이 묶었으며   풀 없이 붙였으되 흔적…(중략)
}
\begin{kfQBlock}{15}{}
\kfQStem{(나)에 대한 이해로 가장 적절한 것은?}
\begin{kfChoices}
\kfChoice <제6수>의 ‘술’은 자연과 어울리며 풍류를 즐기는 화자의 생활을 드러내는 것으로 볼 수 있다.
\kfChoice <제6수>의 ‘술’은 자연과 어울리며 풍류를 즐기는 화자의 생활을 드러내는 것으로 볼 수 있다.
\kfChoice <제1수>의 ‘신월’은 오래된 것보다는 새로운 것을 더 중시하는 삶의 자세를 강조하는 것으로 볼 수 있다.
\kfChoice <제4수>의 ‘남’은 화자의 삶을 지켜보며 그에 대해 정당한 판단을 내리는 인물로 볼 수 있다.
\end{kfChoices}
\end{kfQBlock}

\kfSecHeader{kfAreaWeekly}{지문}{문항 16 지문 / 과잉 해석}

\kfPassageNote{문항 16 지문 / 과잉 해석}{%
[43\textasciitilde{}45] 다음 글을 읽고 물음에 답하시오.

“도련님은 어디서 온 누구십니까? 지금 어디를 가시는 길인지 물어봐도 될까요?”   “네, 저는 하늘 옥황 문왕성 문 도령입니다. 지금 아랫마을 거무 선생님께 글공부 가는 길이오.”   자청비가 문 도령을 찬찬히 살펴보는데 인물이 단정하고 눈빛이 깊은 것이 마음에 들었다. 게다가 거무 선생께 글공부를 간다 하니 같이 글공부하러 가고 싶은 생각이 불쑥 솟아났다.   “도련님, 우리 집에도 나와 닮은 남동생이 있는데 마침 거무 선생께 글공부하러 가고 싶어 합니다. 이름은 자청 도령이라 하니 같이 벗하여 가는 것이 어떻겠습니까?”   조금이라도 자청비와 더 있고 싶은 문 도령은 선선히 그러겠다고 대답하고는 자청비를 따라갔다. 자청비는 문 도령을 집 앞 골목에 세워 놓고, 부모님 방으로 달려갔다.   “아버님, 어머님, 저도 다른 선비들처럼 글공부하러 가고 싶습니다.”   대감이 펄쩍 뛰었다.   “계집아이가 글을 배워 무엇에 쓴단 말인고?…(중략)
}
\begin{kfQBlock}{16}{}
\kfQStem{<보기>를 참고하여 윗글을 감상한 내용으로 적절하지 않은 것은? [3점]}
\begin{kfChoices}
\kfChoice 자청비가 무쇠 방석을 ‘우리 낭군이 깔고 앉았던 방석’이라고 말한 것은 상대방을 함정에 빠뜨려 자신의 편으로 만들기 위한 유인책으로 볼 수 있겠군.
\kfChoice 자청비가 무쇠 방석을 ‘우리 낭군이 깔고 앉았던 방석’이라고 말한 것은 상대방을 함정에 빠뜨려 자신의 편으로 만들기 위한 유인책으로 볼 수 있겠군.
\kfChoice 자청비가 ‘자청 도령’ 행세를 한 것은 문 도령과의 동질성을 획득하기 위한 속임수로 볼 수 있겠군.
\kfChoice 자청비가 ‘계집아이가 글을 배워 무엇에’ 쓰냐며 부모로부터 글공부를 제지당하는 것은 자청비가 받는 사회적 제약으로 볼 수 있겠군.
\end{kfChoices}
\end{kfQBlock}

\kfSecHeader{kfAreaWeekly}{지문}{문항 17 지문 / 범위 오류}

\kfPassageNote{문항 17 지문 / 범위 오류}{%
\#\#\# [42 \textasciitilde{} 45] 다음 글을 읽고 물음에 답하시오.

[앞부분 줄거리] 정 소저는 계모 박 씨의 모함을 의심 없이 받아들인 아버지 정공 때문에 위기에 처하고, 집에서 나와 숨어 다니던 중 도적을 만나 강물에 몸을 던진다. 이때, 정혼자 조무(용홍)와 동생 조성이 정 소저를 우연히 발견하여 구출한다.

소저가 매우 놀라며 말하였다.   “내가 외가로 가지 않고 구차하게 길가에서 분주하게 다닌 것은 조숙모에게 부끄럽고, 아버지의 허물을 드러내고 싶지 않아서였다. 뜻밖에 저 공자들을 만나니 내가 차마 사실을 말하여 부끄러움을 더하겠는가? 은인의 덕이 산과 바다 같으나 차마 근본을 아뢰게 되어 저 집에서 우리 집의 허물을 알게 되면 매우 부끄럽게 될 것이다. 모름지기 너는 다만 대답하기를 내가 타향에서 떠돌아다니다가 서울의 친척을 찾으러 왔다가 도적을 만나 물에 빠져 죽을 뻔했다고 말하여라. 조 공자가 이미 우리가 여자인 줄을 알았으니 남녀는 구별이 있는 것이다. 생명을 구해준 은혜…(중략)
}
\begin{kfQBlock}{17}{}
\kfQStem{[A]와 [B]에 대한 이해로 가장 적절하지 않은 것은?}
\begin{kfChoices}
\kfChoice [A]와 [B]는 모두 과거에 일어난 일을 상대에게 언급하고 있다.
\kfChoice [A]와 [B]는 모두 과거에 일어난 일을 상대에게 언급하고 있다.
\kfChoice [A]는 [B]와 달리 상대의 행동에 변화를 촉구하고 있다.
\kfChoice [B]는 [A]와 달리 상대에게 다른 인물의 말을 전하고 있다.
\end{kfChoices}
\end{kfQBlock}

\kfSecHeader{kfAreaWeekly}{지문}{문항 18 지문 / 시점/화자 혼동}

\kfPassageNote{문항 18 지문 / 시점/화자 혼동}{%
[26\textasciitilde{}28] 다음 글을 읽고 물음에 답하시오.

(가)   거미란 놈이 흉한 심보로 병원 뒤뜰 난간과 꽃밭 사이 사람 발이 잘 닿지 않는 곳에 그물을 쳐 놓았다. 옥외 요양을 받는 젊은 사나이가 누워서 치어다보기 바르게-

나비가 한 마리 꽃밭에 날아들다 그물에 걸리었다. 노오란 날개를 파득거려도 파득거려도 나비는 자꾸 감기우기만 한다. 거미가 쏜살같이 가더니 끝없는 끝없는 실을 뽑아 나비의 온몸을 감아 버린다. 사나이는 긴 한숨을 쉬었다.

나이보담 무수한 고생 끝에 때를 잃고 병을 얻은 이 사나이를 위로할 말이-거미줄을 헝클어 버리는 것밖에 위로의 말이 없었다.

* 윤동주, 「위로」 -

(나)   누가 와서 나를 부른다면   내 보여 주리라   저 얼은 들판 위에 내리는 달빛을.   얼은 들판을 걸어가는 한 그림자를   지금까지 내 생각해 온 것은 모두 무엇인가.   친구 몇몇 친구 몇몇 그들에게는   이제 내 것 가운데 그중 외로움이 아닌 길을   보여 주게 되리.   오…(중략)
}
\begin{kfQBlock}{18}{}
\kfQStem{<보기>를 바탕으로 (가), (나)를 감상한 내용으로 적절하지 않은 것은? [3점]}
\begin{kfChoices}
\kfChoice (가)에서 ‘나비’가 ‘꽃밭’으로 ‘날아’드는 것은 일제 강점기의 암울한 현실에 대응하려는 화자의 의지를 드러낸 것이겠군.
\kfChoice (가)에서 ‘나비’가 ‘꽃밭’으로 ‘날아’드는 것은 일제 강점기의 암울한 현실에 대응하려는 화자의 의지를 드러낸 것이겠군.
\kfChoice (가)에서 ‘한숨을 쉬’는 ‘사나이’를 위해 ‘거미줄을 헝클어 버리는’ 것은 무기력한 우리 민족의 상황을 위로하는 화자의 행동을 드러낸 것이겠군.
\kfChoice (나)에서 ‘달빛’을 ‘받’으며 ‘구름 개인 들판을 걸어’가는 것은 달의 밝은 이미지를 내면화하여 순수한 삶을 살겠다는 화자의 자세를 드러낸 것이겠군.
\end{kfChoices}
\end{kfQBlock}



\clearpage

\kfChapterCover{3}{비트겐슈타인3 ch03}{Chapter 3}{이번 챕터: 문법 → 문학 5작품 → 비문학 5지문 → 패턴 워크북.}

\kfStartGrammar{이번 챕터 문법 영역 — 음운 / 비음화·유음화 + 복합 음운 변동 분석}


\kfSecHeader{kfAreaGrammar}{문제}{}

\begin{multicols}{2}\raggedcolumns\setlength{\columnsep}{12pt}
\kfPassageFull{문항 1 지문 (concept)}{%
「비음화란 비음이 아닌 자음이 비음의 영향을 받아 비음으로 바뀌는 음운 변동이다. 파열음 'ㄱ, ㄷ, ㅂ'이 비음 'ㄴ, ㅁ' 앞에서 각각 'ㅇ, ㄴ, ㅁ'으로 교체된다. 예를 들어 '국물[궁물]', '받는[반는]', '합니다[함니다]'가 이에 해당한다. 비음화는 역행 동화로, 뒤에 오는 비음의 성질이 앞 자음에 영향을 미쳐 일어난다. 한편, 비음화의 결과로 만들어진 종성이 다시 뒤 자음에 영향을 줄 수 있다. '독립[동닙]'에서 'ㄱ'이 'ㄹ' 앞에서 바로 비음화되는 것이 아니라, 먼저 'ㄹ'이 'ㄴ'으로 바뀌고(유음의 비음화), 그 뒤에 'ㄱ'이 'ㄴ' 앞에서 비음화되어 'ㅇ'이 된다.」
}
\begin{kfQBlock}{1}{}
\kfQStem{윗글을 바탕으로 <보기>의 단어들의 발음 과정을 분석한 내용으로 적절한 것은?

<보기> ㉠ 막내[망내]   ㉡ 십리[심니]   ㉢ 학력[항녁] ㉣ 백리[뱅니]   ㉤ 먹는[멍는]}
\begin{kfChoices}
\kfChoice ㉠에서는 'ㄱ'이 'ㄴ' 앞에서 비음화만 일어나 [망내]가 된다.
\kfChoice ㉡에서는 비음화가 먼저 일어난 후 유음화가 일어난다.
\kfChoice ㉢에서는 비음화 한 번만 일어나 [항녁]이 된다.
\kfChoice ㉣에서는 'ㄱ'이 'ㄹ' 앞에서 직접 비음화되어 'ㅇ'이 된다.
\kfChoice ㉤에서는 'ㄱ'이 'ㄴ' 앞에서 유음화되어 [멀는]이 된다.
\end{kfChoices}
\end{kfQBlock}

\kfPassageFull{문항 2 지문 (concept)}{%
「유음화란 'ㄴ'이 'ㄹ'의 앞이나 뒤에서 'ㄹ'로 바뀌는 음운 변동이다. '신라[실라]', '칼날[칼랄]'이 대표적 예이다. '신라'에서는 'ㄴ'이 뒤의 'ㄹ'에 동화되어 'ㄹ'로 바뀌고(역행 동화), '칼날'에서는 'ㄴ'이 앞의 'ㄹ'에 동화되어 'ㄹ'로 바뀐다(순행 동화). 그런데 'ㄴ'과 'ㄹ'이 만나더라도 유음화가 아닌 비음화가 일어나는 경우도 있다. 합성어에서 뒷말의 첫소리 'ㄹ'이 'ㄴ'으로 바뀌는 경우가 있는데, '담력[담녁]', '침략[침냑]' 등이 이에 해당한다. 이 경우 'ㄹ'이 비음 뒤에서 'ㄴ'으로 바뀌는 것으로, 유음화와는 반대 방향의 변동이다.」
}
\begin{kfQBlock}{2}{}
\kfQStem{윗글을 바탕으로 <보기>의 ⓐ\textasciitilde{}ⓔ를 분석한 내용으로 적절하지 않은 것은?

<보기> ⓐ 관리[괄리]   ⓑ 난로[날로]   ⓒ 생산량[생산냥] ⓓ 설날[설랄]   ⓔ 음료[음뇨]}
\begin{kfChoices}
\kfChoice ⓐ에서 'ㄴ'이 뒤의 'ㄹ'에 동화되어 'ㄹ'로 바뀌는 역행 동화가 일어났다.
\kfChoice ⓑ에서 'ㄴ'이 뒤의 'ㄹ'에 동화되어 'ㄹ'로 바뀌는 역행 동화가 일어났다.
\kfChoice ⓒ에서 'ㄹ'이 'ㄴ' 뒤에서 'ㄴ'으로 바뀌는 비음화가 일어났다.
\kfChoice ⓓ에서 'ㄴ'이 앞의 'ㄹ'에 동화되어 'ㄹ'로 바뀌는 순행 동화가 일어났다.
\kfChoice ⓔ에서 'ㄹ'이 'ㅁ' 뒤에서 유음화되어 'ㄹ'이 유지된다.
\end{kfChoices}
\end{kfQBlock}

\kfPassageFull{문항 3 지문 (concept)}{%
「하나의 단어가 발음될 때 여러 음운 변동이 동시에 또는 순차적으로 일어날 수 있다. 음운 변동의 횟수를 셀 때에는 각각의 변동을 독립적으로 센다. 예를 들어 '꽃말[꼰말]'에서는 먼저 'ㅊ'이 음절 끝소리 규칙에 의해 [ㄷ]으로 교체되고, 이 [ㄷ]이 비음 'ㅁ' 앞에서 비음화되어 [ㄴ]이 된다. 따라서 교체 2회의 음운 변동이 일어난다. '닭볶음탕'처럼 복잡한 단어에서는 자음군 단순화, 경음화, 비음화 등이 연쇄적으로 적용되기도 한다.」
}
\begin{kfQBlock}{3}{}
\kfQStem{윗글을 바탕으로 <보기>의 밑줄 친 단어에서 일어나는 음운 변동의 횟수와 유형을 분석한 내용으로 적절한 것은?

<보기> ㉠ 닫는[단는]       ㉡ 흙냄새[흥냄새] ㉢ 없는[엄는]       ㉣ 꽃길[꼳낄→꼰낄]}
\begin{kfChoices}
\kfChoice ㉠에서는 교체 1회(비음화)가 일어난다.
\kfChoice ㉡에서는 교체 1회(비음화)만 일어난다.
\kfChoice ㉢에서는 탈락 1회와 교체 2회, 총 3회의 변동이 일어난다.
\kfChoice ㉣에서는 교체 1회만 일어나 [꼳길]이 된다.
\kfChoice ㉠\textasciitilde{}㉣에서 비음화는 모두 동일한 유형의 교체이다.
\end{kfChoices}
\end{kfQBlock}

\kfPassageFull{문항 4 지문 (concept)}{%
「음운 변동은 교체, 탈락, 첨가, 축약의 네 유형으로 분류된다. 교체는 하나의 음운이 다른 음운으로 바뀌는 것(ㄷ→ㄴ 등), 탈락은 음운이 사라지는 것(겹받침에서 하나 탈락 등), 첨가는 없던 음운이 새로 생기는 것(ㄴ 첨가 등), 축약은 두 음운이 합쳐져 하나로 줄어드는 것(ㅎ 축약 등)이다. 하나의 발음 과정에서 여러 유형이 복합적으로 나타날 수 있으며, 같은 유형이 반복되기도 한다. 음운 변동의 유형을 판별할 때는 변동 전후의 음운 개수를 비교하면 된다. 교체는 음운 개수가 변하지 않고, 탈락은 줄어들며, 첨가는 늘어나고, 축약은 줄어든다.」
}
\begin{kfQBlock}{4}{}
\kfQStem{윗글을 바탕으로 <보기>의 발음에서 일어나는 음운 변동의 유형을 바르게 분석한 것은?

<보기> ㉠ 놓는[논는]     ㉡ 닭장[닥짱]     ㉢ 한여름[한녀름] ㉣ 좋고[조코]     ㉤ 꽃잎[꼰닙]}
\begin{kfChoices}
\kfChoice ㉠: 축약 1회
\kfChoice ㉡: 탈락 1회, 교체 1회
\kfChoice ㉢: 교체 1회
\kfChoice ㉣: 교체 1회
\kfChoice ㉤: 교체 2회, 첨가 1회
\end{kfChoices}
\end{kfQBlock}

\kfPassageFull{문항 5 지문 (concept)}{%
「음운 변동이 여러 개 일어날 때 적용 순서가 중요하다. 순서가 바뀌면 표면형이 달라질 수 있기 때문이다. 일반적으로 음절 구조 관련 규칙(음절의 끝소리 규칙, 자음군 단순화)이 먼저 적용되고, 그 결과를 바탕으로 동화 규칙(비음화, 유음화)이나 경음화 등이 적용된다. 예를 들어 '읽는다'에서 먼저 겹받침 'ㄺ'의 자음군 단순화로 'ㄹ'이 탈락하여 [익]이 되고, 이 'ㄱ'이 'ㄴ' 앞에서 비음화되어 [ㅇ]이 되어 [잉는다]가 된다. 만약 비음화가 먼저 적용되면 잘못된 결과가 나올 수 있다.」
}
\begin{kfQBlock}{5}{}
\kfQStem{윗글을 바탕으로 <보기>의 단어에 적용되는 음운 변동의 순서를 바르게 배열한 것은?

<보기> '닭만' → [당만]

① 자음군 단순화 → ② 비음화 위 순서를 참고하여 아래 단어의 변동 순서를 분석하시오.

㉠ 흙냄새[흥냄새]  ㉡ 없는[엄는]  ㉢ 값나가다[감나가다]}
\begin{kfChoices}
\kfChoice ㉠\textasciitilde{}㉢ 모두 비음화 → 자음군 단순화 순서로 적용된다.
\kfChoice ㉠은 자음군 단순화 → 비음화, ㉡은 비음화 → 자음군 단순화 순서이다.
\kfChoice ㉠\textasciitilde{}㉢ 모두 자음군 단순화 → 비음화 순서로 적용된다.
\kfChoice ㉢은 자음군 단순화만 일어나고 비음화는 일어나지 않는다.
\kfChoice ㉠은 비음화만 일어나고 자음군 단순화는 일어나지 않는다.
\end{kfChoices}
\end{kfQBlock}

\end{multicols}


\kfStartLit{이번 챕터 문학 작품 5편}


\kfSecHeader{kfAreaLiterature}{지문}{성산별곡 — 정철 (가사(고전시가)) / 정철, 「성산별곡」, 핵심 발췌}

\kfPassageNote{성산별곡 — 정철 (가사(고전시가)) / 정철, 「성산별곡」, 핵심 발췌}{%
어주에 봄이 드니 경물이 새롭고야 \\[4pt]
도화행화(桃花杏花)는 석양리(夕陽裏)에 빛이 나고 \\[4pt]
벽초방초(碧草芳草)는 세우중(細雨中)에 프르로다 \\[4pt]
칼로 물을 베히건이 그칠 줄을 모르거니 \\[4pt]
사시풍경(四時風景)이 이 같으면 좋으련만 \\[14pt]
여름 하늘 높으니 바람이 절로 부니 \\[4pt]
오동반엽(梧桐半葉)이 가만히 떨어지거늘 \\[4pt]
어즈버 가을이냐 여름이 아직이라 \\[14pt]
낙엽 중에 기러기는 어디서 오는 것이 \\[4pt]
허공에 떠 있으니 세상을 다 잊었네 \\[14pt]
설월(雪月)에 혼자 앉아 소상강(瀟湘江) 밤비 소리 \\[4pt]
이 몸이 이렇거니 무엇을 부러하리
}
\kfSecHeader{kfAreaLiterature}{활동}{작품 정리표 — 성산별곡(정철)}

\noindent\textit{\small 빈칸에 알맞은 말을 채워 넣으시오. (초성 힌트 참고)}\par\medskip
\begin{kfActTable}{|>{\centering\arraybackslash}m{2.6cm}|m{6cm}|m{4cm}|}
\rowcolor{kfPrimaryLight!50}\textbf{\color{kfPrimary}항목} & \textbf{\color{kfPrimary}내용 (빈칸 채우기)} & \textbf{\color{kfPrimary}답안 작성}\\\hline
갈래 & ㄱㅅ(歌辭), 강호 ㅈㅇ ① & \kfTBlank \\\hline
작자 & ㅈㅊ(鄭澈) ② & \kfTBlank \\\hline
배경 & 전라도 ㅅㅅ(星山) ③ & \kfTBlank \\\hline
구조 & ㅅㄱ → ㅂ·ㅇㄹ·ㄱㅇ·ㄱㅇ → ㄱㅇ ④ & \kfTBlank \\\hline
봄 & ㄷㅎㅎㅎ, ㅂㅊㅂㅊ — 꽃과 풀의 아름다움 ⑤ & \kfTBlank \\\hline
여름→가을 & 'ㅇㅈㅂ ㄱㅇ이냐 ㅇㄹ이 아직이라' ⑥ & \kfTBlank \\\hline
가을 & ㄴㅇ 중에 ㄱㄹㄱ — 탈속적 심경 ⑦ & \kfTBlank \\\hline
겨울 & ㅅㅇ(雪月)에 혼자 앉아 ⑧ & \kfTBlank \\\hline
마무리 & '무엇을 ㅂㄹㅎㄹ' — ㅈㅈ감 ⑨ & \kfTBlank \\\hline
핵심 정서 & ㄱㅎ ㅈㅇ에 대한 ㄷㄱ과 ㅌㅅ ⑩ & \kfTBlank \\\hline
표현 기법 & ㄱㄱㅈ ㅁㅅ + 한자어 활용 ⑪ & \kfTBlank \\\hline
주제 & ㅈㅇ 속 ㅎㄱ 거주 + ㅎㅁ ㅅㅇ의 즐거움 ⑫ & \kfTBlank \\\hline
\end{kfActTable}
\kfSecHeader{kfAreaLiterature}{문제}{}

\begin{multicols}{2}\raggedcolumns\setlength{\columnsep}{12pt}
\begin{kfQBlock}{1}{}
\kfQStem{윗글에 대한 설명으로 적절하지 않은 것은?}
\begin{kfChoices}
\kfChoice 사계절의 변화에 따른 자연 풍경이 순차적으로 묘사되고 있다.
\kfChoice '도화행화'와 '벽초방초' 등 한자어와 자연물을 활용한 감각적 묘사가 나타나고 있다.
\kfChoice '허공에 떠 있으니 세상을 다 잊었네'에서 탈속적 심경이 드러나고 있다.
\kfChoice '이 몸이 이렇거니 무엇을 부러하리'에서 자족감이 나타나고 있다.
\kfChoice 화자는 자연을 떠나 관직에 복귀하겠다는 강한 의지를 드러내고 있다.
\end{kfChoices}
\end{kfQBlock}

\begin{kfQBlock}{2}{}
\kfQStem{윗글에서 '칼로 물을 베히건이 그칠 줄을 모르거니'가 함축하는 바로 가장 적절한 것은?}
\begin{kfChoices}
\kfChoice 관직에서 겪는 고단함과 잡무가 좀처럼 끝나지 않는다는 것
\kfChoice 화자가 흐르는 물을 베어 가며 실제로 농사를 짓고 있다는 것
\kfChoice 봄 경치가 별로 아름답지 않아 화자가 불만을 토로하고 있다는 것
\kfChoice 자연의 흐름처럼 봄 경치에 대한 감흥이 그치지 않는다는 것
\kfChoice 화자가 칼을 갈며 전쟁에 나갈 채비를 분주히 갖추고 있다는 것
\end{kfChoices}
\end{kfQBlock}

\begin{kfQBlock}{3}{}
\kfQStem{윗글의 마지막 행 '이 몸이 이렇거니 무엇을 부러하리'가 의미하는 바로 가장 적절한 것은?}
\begin{kfChoices}
\kfChoice 자연 속 생활에 만족하여 세속의 부귀영화를 부러워할 까닭이 없다는 자족의 태도
\kfChoice 자연 속 생활이 고단하여 관직에 대한 부러움이 극에 달하였음을 토로하는 것
\kfChoice 화자가 경제적으로 풍요로워 아쉬울 것이 전혀 없음을 자랑하는 진술
\kfChoice 자연의 아름다움이 사라져 허무감을 깊이 느끼는 정서를 드러낸 것
\kfChoice 타인의 성공과 부귀에 대한 화자의 질투심을 우회적으로 드러낸 것
\end{kfChoices}
\end{kfQBlock}

\begin{kfQBlock}{4}{}
\kfQStem{윗글을 바탕으로 <보기>를 감상한 내용으로 적절하지 않은 것은?

<보기> 정철의 가사에서 자연은 학문과 수양의 공간이자, 세속적 가치를 넘어선 이상적 삶의 터전이다. 사계절의 변화를 통해 화자는 자연의 질서 속에서 인간적 도리를 체득한다.}
\begin{kfChoices}
\kfChoice 사계절 풍경의 순차적 묘사는 자연의 질서를 체득하는 과정을 보여주겠군.
\kfChoice '세상을 다 잊었네'는 세속적 가치를 넘어선 태도를 드러내겠군.
\kfChoice '무엇을 부러하리'는 이상적 삶의 터전에서의 자족감이겠군.
\kfChoice 사계절 묘사는 단순한 배경이 아니라 학문·수양의 의미를 담고 있겠군.
\kfChoice 화자는 자연을 떠나 세속적 성공을 추구해야 한다는 교훈을 전달하겠군.
\end{kfChoices}
\end{kfQBlock}

\begin{kfQBlock}{5}{}
\kfQStem{윗글에서 '어즈버 가을이냐 여름이 아직이라'라는 구절의 표현 효과로 가장 적절한 것은?}
\begin{kfChoices}
\kfChoice 여름과 가을이라는 두 계절이 빚어내는 갈등을 통해 서사적 긴장을 만든다.
\kfChoice 화자가 계절의 변화를 도무지 구분하지 못하는 무지함을 우회적으로 드러낸다.
\kfChoice 과학적 계절 관측의 정확성과 정밀함을 강조하여 시의 객관성을 확보한다.
\kfChoice 화자가 다가오는 계절 변화를 깊이 두려워하고 있음을 직접적으로 보여 준다.
\kfChoice 계절 경계의 미묘한 순간을 포착해 자연 변화를 섬세히 표현한다.
\end{kfChoices}
\end{kfQBlock}

\end{multicols}

\kfSecHeader{kfAreaLiterature}{지문}{사미인곡 / 속사미인곡 — 정철 (가사(복합지문)) / (가) 정철, 「사미인곡」 핵심 발췌 / (나) 정철, 「속미인곡」 핵심 발췌}

\kfPassageNote{사미인곡 / 속사미인곡 — 정철 (가사(복합지문)) / (가) 정철, 「사미인곡」 핵심 발췌 / (나) 정철, 「속미인곡」 핵심 발췌}{%
(가) \\[4pt]
이 몸 삼기실 제 님을 좇아 삼기시니 \\[4pt]
한생연분이며 하늘 모를 일이런가 \\[4pt]
나 하나 젊어 있고 님 하나 날 괴시니 \\[4pt]
이 마음 이 사랑 견줄 데 없어라 \\[14pt]
평생에 원하오되 한데 녀자 하였더니 \\[4pt]
늙거야 무슨 일로 외오 두고 그리시는고 \\[4pt]
엊그제 님을 보니 님은 나를 잊으셨네 \\[4pt]
중문 나내 숨고 겉으로 괴어하시니 \\[4pt]
반기시는 낯이언마는 눈물이 앞을 가리온다 \\[14pt]
봄바람 건듯 불어 적은 비를 몰아오니 \\[4pt]
풍입송(風入松) 해 뫼신 소리 여 님의 소식 삼고 \\[4pt]
꽃 떨어진 뒤에 녹음이 가득한데 \\[4pt]
님 그린 상사몽이 춘풍을 못 이기네 \\[14pt]
추풍에 지는 잎을 보며 앉아 눈물짓고 \\[4pt]
어룬 달 비친 누각에 님 향한 마음 전할 데 없네 \\[4pt]
아마도 이 님 향한 일편단심이야 \\[4pt]
가실 줄이 있으랴 \\[14pt]
(나) \\[4pt]
님다히 늘그시니 나도 쫓아 늙새로다 \\[4pt]
세상에 나 같은 이 님 같은 이 또 있을까 \\[4pt]
하늘이 만드신 인연을 끊으려야 끊을 수 있으랴 \\[14pt]
님이야 나를 모르시나 나는 님을 몰라 하랴 \\[4pt]
님 향한 눈물이 바다 되어 님의 고기 되고지고 \\[4pt]
차라리 싀여 지여 나비나 되야이셔 \\[4pt]
꽃이야 님이신 양 꽃 위에 앉고지고 \\[14pt]
각시님 뉘라 하니 유덕한 이 이렇도다 \\[4pt]
봄바람 가벼운 양 임의 귀에 불고지고 \\[4pt]
가을 달 비친 양 임의 머리 위에 비치고지고 \\[14pt]
꿈에서나 님을 뵐까 잠이 올까 기다리니 \\[4pt]
잠은 아니 오고 바깥에 달만 밝았구나 \\[4pt]
기러기 저 기러기 님 계신 곳 아는가 \\[4pt]
소식이나 전하여다 전할 길이 없었구나
}
\kfSecHeader{kfAreaLiterature}{활동}{작품 정리표 — 사미인곡 / 속미인곡(정철)}

\noindent\textit{\small 빈칸에 알맞은 말을 채워 넣으시오. (초성 힌트 참고)}\par\medskip
\begin{kfActTable}{|>{\centering\arraybackslash}m{2.6cm}|m{6cm}|m{4cm}|}
\rowcolor{kfPrimaryLight!50}\textbf{\color{kfPrimary}항목} & \textbf{\color{kfPrimary}내용 (빈칸 채우기)} & \textbf{\color{kfPrimary}답안 작성}\\\hline
갈래 & ㄱㅅ(歌辭), ㅊㅊㅇㅈㅈㅅ 전통 ① & \kfTBlank \\\hline
작자 & ㅈㅊ(鄭澈, 1536\textasciitilde{}1593) ② & \kfTBlank \\\hline
(가) 화자 구성 & ㄷㄷ ㅎㅈ의 ㄷㅂㅊ ③ & \kfTBlank \\\hline
(가) 핵심 정서 & 님에 대한 ㅇㄱ + ㅇㅍㄷㅅ ④ & \kfTBlank \\\hline
(가) 계절 소재 & 봄 ㄲ, 가을 ㄷ, 겨울 ㅁㅎ ⑤ & \kfTBlank \\\hline
(나) 화자 구성 & ㄱㄴ(甲女)와 ㅇㄴ(乙女)의 ㄷㅎ ㅎㅅ ⑥ & \kfTBlank \\\hline
(나) 변신 소망 1 & 'ㄴㅂ나 되야이셔 ㄲ 위에 앉고지고' ⑦ & \kfTBlank \\\hline
(나) 변신 소망 2 & 'ㅂㅂㄹ 가벼운 양 임의 귀에 불고지고' ⑧ & \kfTBlank \\\hline
(나) 소식 전달 매개 & 'ㄱㄹㄱ 저 기러기 님 계신 곳 아는가' ⑨ & \kfTBlank \\\hline
공통 비유 구조 & 임금 → ㄴ, 충신 → ㅇㅅ ㅎㅈ ⑩ & \kfTBlank \\\hline
차이 — (가) & ㄷㅂ ㅌㄹ → ㅇㅈㅈ 다짐 ⑪ & \kfTBlank \\\hline
차이 — (나) & ㄷㅎ ㅎㅅ → ㅂㅅ ㅅㅁ의 극단적 간절함 ⑫ & \kfTBlank \\\hline
\end{kfActTable}
\kfSecHeader{kfAreaLiterature}{문제}{}

\begin{multicols}{2}\raggedcolumns\setlength{\columnsep}{12pt}
\begin{kfQBlock}{1}{}
\kfQStem{(가)와 (나)에 대한 설명으로 적절하지 않은 것은?}
\begin{kfChoices}
\kfChoice (가)에서 화자는 님과의 관계가 하늘이 맺어 준 것임을 강조하고 있다.
\kfChoice (나)에서 '나비나 되야이셔'는 님에게 다가가려는 간절한 소망을 표현한 것이다.
\kfChoice (가)와 (나) 모두 사계절의 자연물을 활용하여 그리움을 형상화하고 있다.
\kfChoice (가)의 화자는 님의 냉대에 분노하여 관계 단절을 선언하고 있다.
\kfChoice (나)에서 '기러기'는 님에게 소식을 전해 줄 매개체로 활용되고 있다.
\end{kfChoices}
\end{kfQBlock}

\begin{kfQBlock}{2}{}
\kfQStem{(나)에서 '차라리 싀여 지여 나비나 되야이셔 / 꽃이야 님이신 양 꽃 위에 앉고지고'에 담긴 함축적 의미로 가장 적절한 것은?}
\begin{kfChoices}
\kfChoice 님과의 관계를 포기하고 자연인으로 돌아가겠다는 결심
\kfChoice 님에 대한 원망이 극에 달하여 현실 도피를 추구하는 것
\kfChoice 나비와 꽃의 생태적 관계를 과학적으로 관찰한 것
\kfChoice 화자가 실제로 곤충학에 관심을 보이는 것
\kfChoice 자연물로 변신해서라도 님 곁에 머물고 싶은 극한의 그리움
\end{kfChoices}
\end{kfQBlock}

\begin{kfQBlock}{3}{}
\kfQStem{(가)의 '아마도 이 님 향한 일편단심이야 / 가실 줄이 있으랴'와 (나)의 '하늘이 만드신 인연을 끊으려야 끊을 수 있으랴'를 비교한 내용으로 가장 적절한 것은?}
\begin{kfChoices}
\kfChoice (가)는 화자의 의지적 결단으로, (나)는 천명의 불가분한 인연으로 변함없는 마음을 표현한다.
\kfChoice (가)와 (나) 모두 떠난 님에 대한 깊은 분노와 원한을 설의법으로 격하게 표출한다.
\kfChoice (가)는 님과의 재회에 대한 확신을, (나)는 영원한 이별의 담담한 수용을 나타낸다.
\kfChoice (가)는 화자가 타인의 조언을 수용한 결과이고, (나)는 화자 스스로의 자발적 깨달음이다.
\kfChoice (가)는 님을 잊겠다는 단호한 선언이고, (나)는 새로운 님을 찾겠다는 결심의 표명이다.
\end{kfChoices}
\end{kfQBlock}

\begin{kfQBlock}{4}{}
\kfQStem{(가)와 (나)를 바탕으로 <보기>를 이해한 내용으로 적절하지 않은 것은?

<보기> 정철의 가사에서 '연군(戀君)'은 충신이 임금을 그리워하는 정서를 말한다. 이를 여성 화자의 목소리에 실어 님에 대한 사랑으로 표현하는 것은 '충신연주지사(忠臣戀主之詞)'의 전통이다. (가)는 단독 화자의 독백체로, (나)는 두 여성의 대화체로 연군의 정서를 형상화한다.}
\begin{kfChoices}
\kfChoice (가)의 '님'은 <보기>에 따르면 임금을 가리키며, 화자의 그리움은 충신의 연군 정서에 해당하겠군.
\kfChoice (나)의 대화 형식은 갑녀의 그리움에 을녀의 위로가 더해져 연군의 정서에 객관적 거리를 부여하겠군.
\kfChoice (가)의 '일편단심이야 가실 줄이 있으랴'는 충신의 변함없는 충성을 여성 화자의 사랑으로 치환한 것이겠군.
\kfChoice (가)와 (나) 모두 여성 화자의 목소리를 빌려 정치적 정서를 표현하는 충신연주지사의 전통에 해당하겠군.
\kfChoice (나)의 대화 형식은 연군의 정서를 약화시켜 정치적 의미를 지워 낸 것이겠군.
\end{kfChoices}
\end{kfQBlock}

\begin{kfQBlock}{5}{}
\kfQStem{(나)에서 '봄바람 가벼운 양 임의 귀에 불고지고 / 가을 달 비친 양 임의 머리 위에 비치고지고'에 사용된 표현 기법과 그 효과로 가장 적절한 것은?}
\begin{kfChoices}
\kfChoice 의인법을 활용하여 봄바람과 가을 달이 지닌 내면의 심리를 세밀하게 묘사한다.
\kfChoice 자연물 비유로 화자가 님에게 감각적으로 다가가려는 소망을 드러낸다.
\kfChoice 역설법을 통해 사랑과 미움이 뒤섞인 화자의 모순된 감정을 직접적으로 드러낸다.
\kfChoice 직유법을 활용하여 님의 외모를 봄바람의 가벼움과 가을 달의 환함에 비교한다.
\kfChoice 반어법을 사용하여 님에 대한 깊은 원망을 부드러운 어조로 간접적으로 표현한다.
\end{kfChoices}
\end{kfQBlock}

\end{multicols}

\kfSecHeader{kfAreaLiterature}{지문}{강호사시가 — 맹사성 (시조(연시조 4수)) / 맹사성, 「강호사시가」 전문(4수)}

\kfPassageNote{강호사시가 — 맹사성 (시조(연시조 4수)) / 맹사성, 「강호사시가」 전문(4수)}{%
강호에 봄이 드니 미친 흥이 절로 난다 \\[4pt]
탁료계변(濁醪溪邊)에 금린어(錦鱗魚) 안주로다 \\[4pt]
이 몸이 한가하옴도 역군은(亦君恩)이샷다 \\[14pt]
강호에 녀름이 드니 초당(草堂)에 일이 업다 \\[4pt]
유신(有信)한 강파(江波)는 보내나니 바람이로다 \\[4pt]
이 몸이 서늘하옴도 역군은(亦君恩)이샷다 \\[14pt]
강호에 가을이 드니 고기마다 살져 있다 \\[4pt]
소정(小艇)에 그물 시러 흘려 띄워 보내노라 \\[4pt]
이 몸이 소일(消日)하옴도 역군은(亦君恩)이샷다 \\[14pt]
강호에 겨울이 드니 눈 기픠 자히 남다 \\[4pt]
삿갓 빗기 쓰고 누역(짚옷)을 두루니 \\[4pt]
이 몸이 칩지 아니하옴도 역군은(亦君恩)이샷다
}
\kfSecHeader{kfAreaLiterature}{활동}{작품 정리표 — 강호사시가(맹사성)}

\noindent\textit{\small 빈칸에 알맞은 말을 채워 넣으시오. (초성 힌트 참고)}\par\medskip
\begin{kfActTable}{|>{\centering\arraybackslash}m{2.6cm}|m{6cm}|m{4cm}|}
\rowcolor{kfPrimaryLight!50}\textbf{\color{kfPrimary}항목} & \textbf{\color{kfPrimary}내용 (빈칸 채우기)} & \textbf{\color{kfPrimary}답안 작성}\\\hline
갈래 & ㅇㅅㅈ(聯時調) 4수 ① & \kfTBlank \\\hline
작자 & ㅁㅅㅅ(孟思誠, 1360\textasciitilde{}1438) ② & \kfTBlank \\\hline
구조 & 4수 — ㅂ·ㅇㄹ·ㄱㅇ·ㄱㅇ 순서 ③ & \kfTBlank \\\hline
후렴구 & 'ㅇㄱㅇ(亦君恩)이샷다' ④ & \kfTBlank \\\hline
봄 & ㅌㄹㄱㅂ에 ㄱㄹㅇ 안주 — ㅁㅊ ㅎ ⑤ & \kfTBlank \\\hline
여름 & ㅊㄷ에 ㅇ이 없다 — ㄱㅍ가 ㅂㄹ 보냄 ⑥ & \kfTBlank \\\hline
가을 & ㄱㄱ마다 ㅅㅈ — ㅅㅈ에 ㄱㅁ 실어 ⑦ & \kfTBlank \\\hline
겨울 & ㄴ이 깊이 — ㅅㄱ·ㄴㅇ 차림 ⑧ & \kfTBlank \\\hline
핵심 정서 & ㄱㅎ ㅈㅇ의 ㅎㄱ함 + ㅇㄱ ㅇㅎ에 대한 감사 ⑨ & \kfTBlank \\\hline
주제 & ㅈㅇ ㅎㅇ와 ㅊㄱ ㅇㅅ의 조화 ⑩ & \kfTBlank \\\hline
표현 특징 & ㅎㄹㄱ 반복 + ㄱㄱㅈ ㅁㅅ ⑪ & \kfTBlank \\\hline
문학사적 의의 & 조선 전기 ㄱㅎ ㅅㅈ의 대표작 ⑫ & \kfTBlank \\\hline
\end{kfActTable}
\kfSecHeader{kfAreaLiterature}{문제}{}

\begin{multicols}{2}\raggedcolumns\setlength{\columnsep}{12pt}
\begin{kfQBlock}{1}{}
\kfQStem{윗글에 대한 설명으로 적절하지 않은 것은?}
\begin{kfChoices}
\kfChoice 사계절에 대응하는 4수의 연시조 형식으로, 각 수가 봄·여름·가을·겨울을 순차적으로 다루고 있다.
\kfChoice 매 수의 종장에 '역군은이샷다'라는 후렴구가 반복되어 통일성을 부여하고 있다.
\kfChoice 화자는 강호의 자연 속에서 한가한 삶을 누리며 자족감을 드러내고 있다.
\kfChoice 화자는 강호 생활이 임금의 은혜 덕분이라는 인식을 반복적으로 표현하고 있다.
\kfChoice 화자는 강호를 떠나 관직에 복귀해야 한다는 의무감을 강하게 토로하고 있다.
\end{kfChoices}
\end{kfQBlock}

\begin{kfQBlock}{2}{}
\kfQStem{윗글에서 '역군은(亦君恩)이샷다'에 담긴 함축적 의미로 가장 적절한 것은?}
\begin{kfChoices}
\kfChoice 임금에 대한 화자의 불만을 우회적으로 드러내는 반어적 진술
\kfChoice 강호의 평안한 삶이 임금의 은혜 덕분임을 감사하는 표현
\kfChoice 강호 생활이 불편하여 임금의 부름을 기다린다는 화자의 간청
\kfChoice 임금의 은혜가 강호 생활에까지 미치지 못한다는 화자의 불평
\kfChoice 화자가 임금에게 직접 은혜를 베풀었다는 자부심을 드러낸 표현
\end{kfChoices}
\end{kfQBlock}

\begin{kfQBlock}{3}{}
\kfQStem{윗글의 표현 기법에 대한 설명으로 적절하지 않은 것은?}
\begin{kfChoices}
\kfChoice 매 수 종장에 비유법을 활용하여 임금을 자연물에 비유하고 있다.
\kfChoice '유신한 강파는 보내나니 바람이로다'에서 의인법을 통해 자연에 인격적 속성을 부여하고 있다.
\kfChoice '삿갓 빗기 쓰고 누역을 두루니'에서 구체적 행위 묘사를 통해 겨울 강호의 생활상을 제시하고 있다.
\kfChoice 매 수 초장의 '강호에 \textasciitilde{}이 드니'라는 구절이 반복되어 계절 전환의 표지가 되고 있다.
\kfChoice '미친 흥이 절로 난다'에서 과장법을 통해 봄의 흥취를 강조하고 있다.
\end{kfChoices}
\end{kfQBlock}

\begin{kfQBlock}{4}{}
\kfQStem{윗글을 바탕으로 <보기>를 감상한 내용으로 적절하지 않은 것은?

<보기> 조선 전기 사대부 시조에서 자연과 정치는 분리되지 않는다. 강호(江湖)는 단순한 은거처가 아니라, 임금의 덕화(德化)가 미치는 이상적 공간이다. 사대부는 강호에서 도를 닦으면서도 군신 관계의 끈을 놓지 않는다.}
\begin{kfChoices}
\kfChoice '역군은이샷다'의 반복은 강호 생활이 임금의 덕화 안에 있음을 확인하는 행위이겠군.
\kfChoice 화자가 강호에서 한가하게 지내면서도 임금의 은혜를 말하는 것은 자연과 정치의 분리 불가능성을 보여주겠군.
\kfChoice '금린어 안주로다'나 '고기마다 살져 있다' 등은 임금의 덕화가 미치는 풍요로운 강호의 모습이겠군.
\kfChoice 화자의 강호 생활은 군신 관계를 단절한 채 자연에만 몰입하는 것이겠군.
\kfChoice 강호는 단순한 은거처가 아니라 임금의 통치가 미치는 이상적 공간으로 제시되어 있겠군.
\end{kfChoices}
\end{kfQBlock}

\begin{kfQBlock}{5}{}
\kfQStem{윗글에서 계절별 화자의 생활상과 정서를 바르게 짝지은 것은?}
\begin{kfChoices}
\kfChoice 가을 — 초당에서 무위의 생활, 겨울 — 강파에 배를 띄움
\kfChoice 봄 — 눈 속에서 누역을 두름, 여름 — 고기가 살찌는 것을 관찰
\kfChoice 봄 — 금린어를 안주 삼아 즐김, 여름 — 바람에서 서늘함을 느낌
\kfChoice 가을 — 탁료계변에서 술을 마심, 겨울 — 소정에 그물을 실음
\kfChoice 봄 — 초당에서 일이 없음, 여름 — 삿갓을 빗겨 씀
\end{kfChoices}
\end{kfQBlock}

\end{multicols}

\kfSecHeader{kfAreaLiterature}{지문}{만흥 — 윤선도 (시조(연시조 6수)) / 윤선도, 「만흥」 전문(6수)}

\kfPassageNote{만흥 — 윤선도 (시조(연시조 6수)) / 윤선도, 「만흥」 전문(6수)}{%
산수간(山水間) 바회 아래 뫼집을 짓노라 하니 \\[4pt]
그 모른 남이 웃는다 한들 어떠하리 \\[4pt]
어리고 우활한 이 내 생애 아닐런가 \\[14pt]
보리밥 풋나물을 알마츠 먹은 후에 \\[4pt]
바위 끝 물가에 슬근슬근 거러 보니 \\[4pt]
그 남은 여남은 일이야 부럴 것이 더욱 없다 \\[14pt]
잔 들고 혼자 앉아 먼 뫼를 바라보니 \\[4pt]
그리던 님이 오다 반가옴이 이러하랴 \\[4pt]
말씀도 우습도 아녀도 못내 좋아하노라 \\[14pt]
내 벗이 몇이냐 하니 수석(水石)과 송죽(松竹)이라 \\[4pt]
동산에 달 오르니 그 더욱 반갑고야 \\[4pt]
두어라 이 다섯 밖에 또 더하여 무엇하리 \\[14pt]
나무도 아닌 것이 풀도 아닌 것이 \\[4pt]
곧기는 뉘 시기며 속은 어이 볏겼는고 \\[4pt]
뉘라서 사시에 프르른 줄을 그를 몰라 하노라 \\[14pt]
더우면 꽃 피고 추우면 잎 지거늘 \\[4pt]
솔아 너는 어찌 눈서리를 모르느니 \\[4pt]
구태여 모진 형용을 못 말릴 줄 알면서 하노라
}
\kfSecHeader{kfAreaLiterature}{활동}{작품 정리표 — 만흥(윤선도)}

\noindent\textit{\small 빈칸에 알맞은 말을 채워 넣으시오. (초성 힌트 참고)}\par\medskip
\begin{kfActTable}{|>{\centering\arraybackslash}m{2.6cm}|m{6cm}|m{4cm}|}
\rowcolor{kfPrimaryLight!50}\textbf{\color{kfPrimary}항목} & \textbf{\color{kfPrimary}내용 (빈칸 채우기)} & \textbf{\color{kfPrimary}답안 작성}\\\hline
갈래 & ㅇㅅㅈ 6수 ① & \kfTBlank \\\hline
작자 & ㅇㅅㄷ(尹善道, 1587\textasciitilde{}1671) ② & \kfTBlank \\\hline
배경 & ㅂㄱㄷ(甫吉島) ③ & \kfTBlank \\\hline
주제 1 & ㅇㅂㄴㄷ — 가난해도 도를 즐김 ④ & \kfTBlank \\\hline
주제 2 & ㅈㅇ ㅎㅇ — 자연과 하나 됨 ⑤ & \kfTBlank \\\hline
소박한 삶 & 'ㅂㄹㅂ ㅍㄴㅁ을 알마츠 먹은 후에' ⑥ & \kfTBlank \\\hline
자연 벗 & ㅅㅅ과 ㅅㅈ + ㄷ — 다섯 벗 ⑦ & \kfTBlank \\\hline
무언의 교감 & 'ㅁㅆ도 ㅇㅅ도 아녀도 못내 ㅈㅇ하노라' ⑧ & \kfTBlank \\\hline
세속 초월 & 'ㅂㄹ 것이 더욱 없다' ⑨ & \kfTBlank \\\hline
절개 상징 1 & 대나무 — ㅅㅅ에 ㅍㄹㄹ ⑩ & \kfTBlank \\\hline
절개 상징 2 & 소나무 — ㄴㅅㄹ를 모르느니 ⑪ & \kfTBlank \\\hline
화자 태도 & 'ㅇㄹㄱ ㅇㅎㅎ 이 내 생애' — 자기 삶 ㄱㅈ ⑫ & \kfTBlank \\\hline
\end{kfActTable}
\kfSecHeader{kfAreaLiterature}{문제}{}

\begin{multicols}{2}\raggedcolumns\setlength{\columnsep}{12pt}
\begin{kfQBlock}{1}{}
\kfQStem{윗글에 대한 설명으로 적절하지 않은 것은?}
\begin{kfChoices}
\kfChoice 화자는 산수간에 집을 짓고 사는 삶을 스스로 선택한 것으로 제시하고 있다.
\kfChoice '보리밥 풋나물'은 소박한 음식으로, 안빈낙도의 가치관을 드러내고 있다.
\kfChoice 화자는 수석, 송죽, 달을 벗으로 삼아 자연과의 교감을 표현하고 있다.
\kfChoice 제5수와 제6수에서 화자는 대나무와 소나무의 변절을 비판하고 있다.
\kfChoice 화자는 타인의 비웃음에 개의치 않는 초연한 태도를 보이고 있다.
\end{kfChoices}
\end{kfQBlock}

\begin{kfQBlock}{2}{}
\kfQStem{제4수의 '두어라 이 다섯 밖에 또 더하여 무엇하리'에 담긴 함축적 의미로 가장 적절한 것은?}
\begin{kfChoices}
\kfChoice 인간 친구를 새로 다섯 명 더 사귀겠다는 의지의 표현
\kfChoice 벗이 너무 많아 그 수를 줄이고 싶다는 화자의 불만
\kfChoice 다섯 자연 벗만으로 충분하다는 자족감의 표현
\kfChoice 자연 벗을 다섯 가지로 제한해야 한다는 규범적인 주장
\kfChoice 자연 벗이 모자라 더 필요하다는 화자의 아쉬움의 토로
\end{kfChoices}
\end{kfQBlock}

\begin{kfQBlock}{3}{}
\kfQStem{제5수와 제6수에서 화자가 대나무와 소나무에 주목하는 공통적 이유로 가장 적절한 것은?}
\begin{kfChoices}
\kfChoice 사시사철 변치 않는 절개와 지조의 자연물이기 때문이다.
\kfChoice 경제적 가치가 높은 목재로서의 실용성 때문이다.
\kfChoice 다른 식물보다 크기가 크고 위용이 뛰어나기 때문이다.
\kfChoice 약초로서의 효능이 있어 건강에 도움이 되기 때문이다.
\kfChoice 계절에 따라 색이 변하는 아름다움이 있기 때문이다.
\end{kfChoices}
\end{kfQBlock}

\begin{kfQBlock}{4}{}
\kfQStem{윗글을 바탕으로 <보기>를 감상한 내용으로 적절하지 않은 것은?

<보기> 안빈낙도(安貧樂道)는 가난한 삶 속에서도 도(道)를 즐기는 유교적 이상이다. 이는 물질적 풍요를 추구하는 세속적 가치를 거부하고, 자연과의 합일 속에서 정신적 충만을 추구하는 삶의 태도이다. 조선 사대부 시조에서 이 가치는 자연 벗과의 교감, 소박한 생활, 세속에 대한 초월로 구체화된다.}
\begin{kfChoices}
\kfChoice '보리밥 풋나물을 알마츠 먹은 후에'는 소박한 생활을 통한 안빈낙도의 실천이겠군.
\kfChoice '그 남은 여남은 일이야 부럴 것이 더욱 없다'는 세속적 가치에 대한 초월을 보여주겠군.
\kfChoice '수석과 송죽'을 벗으로 삼는 것은 자연 벗과의 교감을 통한 정신적 충만이겠군.
\kfChoice '어리고 우활한 이 내 생애'는 자신의 삶을 부정하고 세속으로의 복귀를 열망하는 것이겠군.
\kfChoice 제5·6수에서 대나무와 소나무의 절개를 예찬하는 것은 도(道)를 즐기는 태도와 연결되겠군.
\end{kfChoices}
\end{kfQBlock}

\begin{kfQBlock}{5}{}
\kfQStem{제3수에서 '말씀도 우습도 아녀도 못내 좋아하노라'에 사용된 표현 기법과 그 효과로 가장 적절한 것은?}
\begin{kfChoices}
\kfChoice 대구법을 사용하여 화자와 님 사이의 대립적인 관계를 한층 부각하고 있다.
\kfChoice 과장법을 활용하여 님과의 만남에서 느끼는 화자의 슬픔을 극대화한다.
\kfChoice 반어법을 통해 님이 끝내 오지 않은 것에 대한 실망감을 우회적으로 드러낸다.
\kfChoice 부정의 반복적 열거로 자연과의 무언의 교감을 강조한다.
\kfChoice 직유법을 활용하여 무심한 자연을 인간의 모습에 빗대어 친근하게 표현한다.
\end{kfChoices}
\end{kfQBlock}

\end{multicols}

\kfSecHeader{kfAreaLiterature}{지문}{고공가 — 허전 (가사(교훈 가사)) / 허전, 「고공가」 핵심 발췌}

\kfPassageNote{고공가 — 허전 (가사(교훈 가사)) / 허전, 「고공가」 핵심 발췌}{%
고공(雇工)이야 고공이야 네 집주인 내 아니냐 \\[4pt]
주인의 분부(分付) 듣고 일을 하라 고공이야 \\[4pt]
너를 머여 쓰는 것이 농사를 위함이라 \\[14pt]
자네 밭을 갈려 하되 옥석(玉石)을 가리소냐 \\[4pt]
묵은밭 새밭이라도 기름지게 갈아 두소 \\[4pt]
짐승이 밟거든 울을 세워 막아 두고 \\[4pt]
벌레가 먹거든 잡아 죽여 없이 하소 \\[14pt]
씨를 뿌린 후에 김을 매되 고르게 하소 \\[4pt]
가물면 물을 대고 장마지면 물을 빼소 \\[4pt]
이웃에 험담하는 자 있거든 잘 달래소 \\[14pt]
고공의 녹(祿)은 주인의 은(恩)이라 \\[4pt]
곡식이 잘되면 고공의 공(功)이요 \\[4pt]
곡식이 못되면 고공의 죄(罪)로다 \\[14pt]
주인의 집이 흥하면 고공도 따라 흥하고 \\[4pt]
주인의 집이 망하면 고공이 뉘 밑에 가리 \\[4pt]
네 집이로다 네 집이로다 주인의 집이 네 집이로다
}
\kfSecHeader{kfAreaLiterature}{활동}{작품 정리표 — 고공가(허전)}

\noindent\textit{\small 빈칸에 알맞은 말을 채워 넣으시오. (초성 힌트 참고)}\par\medskip
\begin{kfActTable}{|>{\centering\arraybackslash}m{2.6cm}|m{6cm}|m{4cm}|}
\rowcolor{kfPrimaryLight!50}\textbf{\color{kfPrimary}항목} & \textbf{\color{kfPrimary}내용 (빈칸 채우기)} & \textbf{\color{kfPrimary}답안 작성}\\\hline
갈래 & ㄱㅅ(歌辭), ㄱㅎ 가사 ① & \kfTBlank \\\hline
작자 & ㅎㅈ(許筌) ② & \kfTBlank \\\hline
비유: 집주인 & = ㅇㄱ(君) ③ & \kfTBlank \\\hline
비유: 머슴 & = ㅅㅎ(臣下) ④ & \kfTBlank \\\hline
비유: 농사 & = ㄱㄱ ㄱㅇ(國家經營) ⑤ & \kfTBlank \\\hline
비유: 곡식 & = ㅂㅅ(百姓) ⑥ & \kfTBlank \\\hline
비유: 벌레·해충 & = ㄱㅅ·ㅇㅈ(奸臣·外敵) ⑦ & \kfTBlank \\\hline
핵심 구절 1 & '고공의 ㄴ은 주인의 ㅇ이라' ⑧ & \kfTBlank \\\hline
핵심 구절 2 & '주인의 집이 ㄴ ㅈ이로다' ⑨ & \kfTBlank \\\hline
농사 지시 & 밭 갈기, ㄱ 매기, ㅂㄹ 잡기, ㅁ 대기 ⑩ & \kfTBlank \\\hline
주제 & ㅇㅁ(爲民) 정치, 신하의 ㅊㅁ 강조 ⑪ & \kfTBlank \\\hline
표현 기법 & ㅇㅇㅈ ㅂㅇ(寓意的 比喩) — 농사로 정치 논리 전환 ⑫ & \kfTBlank \\\hline
\end{kfActTable}
\kfSecHeader{kfAreaLiterature}{문제}{}

\begin{multicols}{2}\raggedcolumns\setlength{\columnsep}{12pt}
\begin{kfQBlock}{1}{}
\kfQStem{윗글에 대한 설명으로 적절하지 않은 것은?}
\begin{kfChoices}
\kfChoice 집주인이 머슴에게 농사일을 지시하는 형식으로 전개되고 있다.
\kfChoice 농사의 비유를 통해 정치적 메시지를 전달하고 있다.
\kfChoice '벌레가 먹거든 잡아 죽여 없이 하소'에서 해충 방제는 간신 제거를 비유한 것이다.
\kfChoice 화자는 머슴의 자율성을 존중하여 일체의 간섭을 배제하고 있다.
\kfChoice '주인의 집이 네 집이로다'에서 군신 운명공동체 의식이 드러나고 있다.
\end{kfChoices}
\end{kfQBlock}

\begin{kfQBlock}{2}{}
\kfQStem{윗글에서 '고공의 녹(祿)은 주인의 은(恩)이라'에 담긴 함축적 의미로 가장 적절한 것은?}
\begin{kfChoices}
\kfChoice 신하의 봉록이 임금의 은혜에서 나오니 보답해야 한다는 군신 윤리
\kfChoice 머슴이 받는 임금이 너무 적어 큰 폭의 인상을 요구하는 발화
\kfChoice 집주인이 오히려 머슴에게 경제적으로 착취당하고 있다는 불만의 토로
\kfChoice 머슴이 집주인의 은혜를 굳이 갚을 필요가 없다는 거부의 주장
\kfChoice 녹과 은의 실제 화폐적 가치를 객관적으로 비교한 경제 논리
\end{kfChoices}
\end{kfQBlock}

\begin{kfQBlock}{3}{}
\kfQStem{윗글의 마지막 행 '네 집이로다 네 집이로다 주인의 집이 네 집이로다'가 의미하는 바로 가장 적절한 것은?}
\begin{kfChoices}
\kfChoice 임금과 신하는 운명공동체이므로 나라의 흥망이 곧 신하의 흥망
\kfChoice 머슴이 집주인의 재산을 빼앗아 자기 소유로 만들라는 도발적 권유
\kfChoice 집주인이 머슴에게 자기 집을 증여하겠다는 일종의 공식 약속
\kfChoice 머슴은 자기 소유의 집이 없으므로 결국 노숙해야 한다는 사실
\kfChoice 주인과 머슴이 법적으로 완전히 동등한 소유권을 가진다는 선언
\end{kfChoices}
\end{kfQBlock}

\begin{kfQBlock}{4}{}
\kfQStem{윗글을 바탕으로 <보기>를 감상한 내용으로 적절하지 않은 것은?

<보기> 「고공가」는 농사 비유를 통해 유교적 위민 정치론을 전개한 교훈 가사이다. 작품 속 비유 체계에서 집주인은 임금, 머슴은 신하, 농사는 국가 경영, 곡식은 백성, 벌레·해충은 간신·외적에 대응한다. 이 우의적 구조를 통해 신하의 책무를 강조한다.}
\begin{kfChoices}
\kfChoice '밭을 갈려 하되 옥석을 가리소냐'는 <보기>의 비유 체계에 따르면 신하가 백성을 차별 없이 다스려야 함을 의미하겠군.
\kfChoice '가물면 물을 대고 장마지면 물을 빼소'는 상황에 따른 유연한 국정 운영을 의미하겠군.
\kfChoice '곡식이 잘되면 고공의 공이요'는 백성이 잘 살게 되면 신하의 공로라는 뜻이겠군.
\kfChoice '이웃에 험담하는 자 있거든 잘 달래소'는 외교적 갈등 해소를 신하에게 당부하는 것이겠군.
\kfChoice 작품 전체는 임금이 신하를 칭찬하여 사기를 높이려는 축하의 글이겠군.
\end{kfChoices}
\end{kfQBlock}

\begin{kfQBlock}{5}{}
\kfQStem{윗글에 사용된 주된 표현 기법과 그 효과에 대한 설명으로 가장 적절한 것은?}
\begin{kfChoices}
\kfChoice 농사 비유로 추상적 정치 논리를 일상적 차원으로 전환하여 이해를 돕는다.
\kfChoice 역설법을 활용하여 머슴이 지닌 권리와 집주인이 져야 할 의무를 동시에 강조하고 있다.
\kfChoice 과장법을 활용하여 농사일이 안고 있는 본질적인 어려움을 극대화해 부각하고 있다.
\kfChoice 풍자적 어조를 사용하여 집주인이 보여 주는 임금으로서의 부당함을 신랄하게 비판한다.
\kfChoice 반어적 표현을 활용하여 머슴이 보이는 일상의 게으름과 무성의를 우회적으로 비웃고 있다.
\end{kfChoices}
\end{kfQBlock}

\end{multicols}

\kfStartNonlit{이번 챕터 비문학 지문 5편}


\kfSecHeader{kfAreaNonlit}{지문}{노자의 도와 한비자의 법 (인문)}

\kfPassageNote{노자의 도와 한비자의 법 (인문)}{%
(가) 노자(老子)의 도(道)는 천지만물의 근원이자 운행 원리이다. 「도가도비상도(道可道非常道)」, 즉 말로 표현할 수 있는 도는 영원한 도가 아니다. 도는 형체가 없고 이름 붙일 수 없으며, 감각으로 포착할 수 없는 것이다. 그러나 만물은 이 도에서 나와 도로 돌아간다. 노자는 이러한 도의 속성을 '무위자연(無爲自然)'이라 표현한다. 무위란 인위적 조작을 하지 않는 것이고, 자연이란 스스로 그러한 것이다.

노자의 정치 사상에서 이상적인 통치자는 '무위이치(無爲而治)', 즉 억지로 다스리지 않되 저절로 다스려지는 것을 실현하는 존재이다. 백성은 통치자의 존재를 의식하지 못하되, 그 덕에 자연스럽게 조화를 이룬다. 노자는 법률과 제도의 증가를 비판한다. 「법령이 번다해질수록 도적이 많아진다(法令滋彰 盜賊多有)」라는 표현은 인위적 규제가 오히려 혼란을 가중시킨다는 통찰이다.

(나) 한비자(韓非子)는 법가(法家)의 집대성자로서, 노자의 도를 전혀 다른 방향으로 전유(轉有)했다. 한비자는 노자의 '도'를 만물의 변화 법칙(理)으로 해석하고, 통치자가 이 법칙을 파악하여 '법(法)', '술(術)', '세(勢)' 삼위일체의 통치 기술로 활용해야 한다고 주장했다. 여기서 '법'은 명문화된 법률, '술'은 군주가 신하를 다루는 비밀스러운 통치술, '세'는 군주의 권위와 지위이다.

한비자는 인간의 본성을 이기적으로 파악한다. 사람은 이익을 좇고 해로움을 피하는 존재이므로, 통치자는 덕치(德治)나 인의(仁義)가 아니라 상벌(賞罰)의 이해(利害) 관계로 다스려야 한다. 노자가 법률의 증가를 비판한 반면, 한비자는 엄격한 법의 시행만이 질서를 유지할 수 있다고 보았다.

두 사상가의 관계는 '동원이류(同源異流)', 즉 같은 근원에서 나와 다른 흐름을 이룬 것이다. 노자의 도라는 형이상학적 원리에서 출발하되, 노자는 무위를, 한비자는 유위(有爲)의 법치를 도출했다. 이 차이의 핵심은 인간 본성에 대한 관점이다. 노자가 인간의 자연적 선함을 신뢰한다면, 한비자는 인간의 이기적 본성을 전제하여 제도적 통제를 강조한다.
}
\kfSecHeader{kfAreaNonlit}{문제}{}

\begin{multicols}{2}\raggedcolumns\setlength{\columnsep}{12pt}
\begin{kfQBlock}{1}{}
\kfQStem{(가)와 (나)의 내용과 일치하지 않는 것은?}
\begin{kfChoices}
\kfChoice 노자에게 도는 형체가 없고 감각으로 포착할 수 없는 것이다.
\kfChoice 노자는 법령이 번다해질수록 도적이 많아진다고 보았다.
\kfChoice 한비자는 노자의 도를 만물의 변화 법칙으로 재해석했다.
\kfChoice 한비자는 인간의 본성이 선하므로 덕치로 다스려야 한다고 주장했다.
\kfChoice 두 사상가의 관계는 같은 근원에서 나와 다른 흐름을 이룬 동원이류이다.
\end{kfChoices}
\end{kfQBlock}

\begin{kfQBlock}{2}{}
\kfQStem{(가)의 노자와 (나)의 한비자를 비교한 내용으로 가장 적절한 것은?}
\begin{kfChoices}
\kfChoice 한비자는 인간의 자연적 선함을 신뢰하고, 노자는 인간의 이기심을 전제한다.
\kfChoice 노자와 한비자 모두 엄격한 법치를 최선의 통치 방법으로 본다.
\kfChoice 노자는 무위자연을, 한비자는 법·술·세의 유위적 통치를 이상으로 삼는다.
\kfChoice 노자와 한비자 모두 도를 부정하고 새로운 원리를 제시한다.
\kfChoice 노자는 상벌을 강조하고, 한비자는 무위를 강조한다.
\end{kfChoices}
\end{kfQBlock}

\begin{kfQBlock}{3}{}
\kfQStem{윗글을 바탕으로 <보기>의 상황에 대해 두 사상가가 각각 어떻게 판단할지 추론한 것으로 적절하지 않은 것은?

<보기> 어떤 나라에서 범죄율이 높아지자, 통치자가 법률을 대폭 강화하고 엄벌을 시행했다.}
\begin{kfChoices}
\kfChoice 노자는 법령의 강화가 오히려 혼란을 가중시킬 수 있다고 비판할 것이다.
\kfChoice 한비자는 엄격한 상벌 체계가 질서 유지에 효과적이라고 지지할 것이다.
\kfChoice 노자는 무위이치를 통해 백성이 자연스럽게 조화를 이루게 해야 한다고 주장할 것이다.
\kfChoice 한비자는 통치자가 법·술·세를 활용하여 신하와 백성을 통제해야 한다고 볼 것이다.
\kfChoice 노자와 한비자 모두 법률 강화를 최선의 해결책이라 동의할 것이다.
\end{kfChoices}
\end{kfQBlock}

\begin{kfQBlock}{4}{}
\kfQStem{윗글의 '동원이류(同源異流)'가 의미하는 바로 가장 적절한 것은?}
\begin{kfChoices}
\kfChoice 같은 근원에서 출발해 다른 방향으로 흘러간 사상적 관계
\kfChoice 서로 다른 근원에서 출발하여 결국 같은 결론에 도달하게 된 관계
\kfChoice 한비자가 노자의 사상을 정면에서 완전히 부정하고 폐기한 관계
\kfChoice 노자와 한비자가 정치 운영에서 완전히 동일한 결론을 도출한 관계
\kfChoice 두 사상이 역사적·이론적으로 서로 전혀 무관함을 가리키는 표현
\end{kfChoices}
\end{kfQBlock}

\end{multicols}

\kfSecHeader{kfAreaNonlit}{지문}{장재·정이·주희의 이기론 심화 (인문)}

\kfPassageNote{장재·정이·주희의 이기론 심화 (인문)}{%
중국 송대 성리학에서 세계의 근본 원리를 설명하는 이기론(理氣論)은 장재(張載), 정이(程頤), 주희(朱熹)를 거치며 점차 정교해졌다. 장재의 기일원론(氣一元論)은 세계를 기(氣)라는 하나의 원리로 설명한다. 그에게 태허(太虛)는 기가 응결되기 이전의 투명하고 고요한 상태이며, 기가 모이면 만물이 되고 흩어지면 태허로 돌아간다. 장재에게 리(理)는 기의 운동 법칙에 불과하며 별도의 독립적 실체가 아니다.

정이는 장재와 달리 리를 기와 구별되는 독립적 원리로 격상시켰다. 이것이 이기이원론(理氣二元論)의 시작이다. 정이에 따르면 리는 형이상(形而上)의 원리이고 기는 형이하(形而下)의 재료이다. 리는 기에 의존하지 않고 독자적으로 존재하며, 만물의 존재 근거가 된다. 정이는 「리가 먼저이고 기가 뒤이다(理先氣後)」라는 원칙을 제시하여 리의 존재론적 우위를 확립했다.

주희는 정이의 이원론을 계승하면서도 '이기불상리불상잡(理氣不相離不相雜)'이라는 정식을 통해 리와 기의 관계를 더 정밀하게 규정했다. '불상리(不相離)'란 리와 기가 결코 떨어질 수 없다는 것이고, '불상잡(不相雜)'이란 리와 기가 서로 섞이지 않는다는 것이다. 리는 형체도 조작도 없는 순수한 원리이지만, 기가 없으면 리가 붙어 있을 곳이 없다. 반대로 기는 리가 없으면 운동의 법칙이 없어진다. 그러나 리와 기는 본질이 다르므로 뒤섞여서는 안 된다.

주희의 체계에서 리와 기의 관계는 인간의 도덕성에도 적용된다. 인간의 성(性)은 리의 차원에서 보면 순선(純善)하지만, 기질(氣質)의 차원에서 보면 맑고 탁함, 순수하고 잡박함의 차이가 있다. 성인과 범인의 차이는 리에 있는 것이 아니라 기질에 있다. 따라서 도덕 수양의 핵심은 기질의 탁함을 맑혀 본래의 리(天理)를 회복하는 것, 즉 '존천리거인욕(存天理去人欲)'이다.
}
\kfSecHeader{kfAreaNonlit}{문제}{}

\begin{multicols}{2}\raggedcolumns\setlength{\columnsep}{12pt}
\begin{kfQBlock}{1}{}
\kfQStem{윗글의 내용과 일치하지 않는 것은?}
\begin{kfChoices}
\kfChoice 장재에게 리는 기의 운동 법칙에 불과하며 별도의 독립적 실체가 아니다.
\kfChoice 정이는 리를 형이상의 원리, 기를 형이하의 재료로 구분했다.
\kfChoice 주희에 따르면 리와 기는 떨어질 수 없으면서도 섞이지 않는다.
\kfChoice 주희에게 성인과 범인의 차이는 리의 차이에서 비롯된다.
\kfChoice 주희의 도덕 수양은 기질의 탁함을 맑혀 천리를 회복하는 것이다.
\end{kfChoices}
\end{kfQBlock}

\begin{kfQBlock}{2}{}
\kfQStem{윗글의 장재, 정이, 주희의 입장을 비교한 내용으로 적절하지 않은 것은?}
\begin{kfChoices}
\kfChoice 장재는 리를 기에 종속시키지만, 정이는 리를 기와 독립된 원리로 격상시킨다.
\kfChoice 정이는 리의 존재론적 우위를 주장하며, 주희는 이를 계승했다.
\kfChoice 주희는 리와 기가 떨어질 수 없다고 보아, 장재의 기일원론으로 회귀했다.
\kfChoice 장재에서 정이, 주희로 갈수록 리의 독립성이 강화된다.
\kfChoice 주희는 리와 기의 불가분리와 불가혼합을 동시에 주장한다.
\end{kfChoices}
\end{kfQBlock}

\begin{kfQBlock}{3}{}
\kfQStem{윗글의 '이기불상리불상잡'이 의미하는 바로 가장 적절한 것은?}
\begin{kfChoices}
\kfChoice 리와 기가 떨어질 수 없으면서도 본질이 달라 섞이지 않는다는 것
\kfChoice 리와 기가 본질에서나 작용에서나 완전히 동일한 하나의 실체라는 것
\kfChoice 리가 기에 전적으로 종속되어 별도로는 존재할 수 없다는 것
\kfChoice 기가 리에 비해 존재론적으로 한층 우위에 있다는 것
\kfChoice 리와 기가 서로 무관하게 완전히 독립적으로 존재한다는 것
\end{kfChoices}
\end{kfQBlock}

\begin{kfQBlock}{4}{}
\kfQStem{윗글에 따를 때, 주희가 '존천리거인욕'을 수양의 핵심으로 제시한 이유로 가장 적절한 것은?}
\begin{kfChoices}
\kfChoice 성이 리에서는 순선하나 기질의 탁함으로 가려지기 때문이다.
\kfChoice 리와 기가 본래 완전히 분리 가능하므로 기를 제거하면 되기 때문이다.
\kfChoice 인간의 본성 자체가 악하므로 그 본성을 통째로 교체해야 하기 때문이다.
\kfChoice 장재의 기일원론에 따라 기질의 맑고 탁함만 다스리면 되기 때문이다.
\kfChoice 정이가 리선기후의 명제를 정면으로 부정한 입장을 따르기 때문이다.
\end{kfChoices}
\end{kfQBlock}

\end{multicols}

\kfSecHeader{kfAreaNonlit}{지문}{왕필의 귀무론과 곽상의 독화론 (인문)}

\kfPassageNote{왕필의 귀무론과 곽상의 독화론 (인문)}{%
(가) 위진 시대의 현학(玄學)은 유교 경전에 대한 새로운 해석을 통해 형이상학적 세계관을 정립하려 한 지적 운동이다. 그 가운데 왕필(王弼, 226\textasciitilde{}249)의 '귀무론(貴無論)'은 현학의 한 갈래를 대표한다. 왕필에 따르면 우주의 근본은 '무(無)'이다. 「노자」의 「천하만물은 유(有)에서 생기고, 유는 무에서 생긴다(天下萬物生於有 有生於無)」를 해석하면서, 왕필은 무를 유의 근거이자 본체(本體)로 파악했다.

왕필의 무는 단순한 '없음'이 아니라, 모든 유(존재하는 것)를 가능하게 하는 근원이다. 마치 수레바퀴의 빈 구멍이 바퀴의 기능을 가능하게 하듯, 무는 만유의 작동 근거이다. 왕필의 정치 사상에서 이 논리는 '성인무위론(聖人無爲論)'으로 이어진다. 이상적 통치자는 자기를 비워(無) 만물이 스스로 그 자리를 찾게 한다.

(나) 곽상(郭象, 252\textasciitilde{}312)은 왕필의 귀무론을 정면으로 비판하고, '독화론(獨化論)'을 제시했다. 곽상에 따르면 왕필의 '무'는 사실상 또 하나의 형이상학적 실체를 상정하는 것이며, 이는 노자의 정신에 오히려 어긋난다. 만약 무가 유의 근거라면, 무 자체도 '무엇(something)'이 되어 버리는 모순에 빠진다.

곽상의 대안은 만물이 스스로 생겨난다(自生)는 것이다. 어떤 외부의 근거(무이든 유이든)도 만물의 존재를 설명하는 원리가 될 수 없으며, 각각의 사물은 다른 것에 의존하지 않고 '홀로 스스로 변화한다(獨化)'. 이것이 독화론의 핵심이다. 곽상에게 만물의 존재는 원인 없이 자발적으로 발생하는 것이며, 외부의 제1원인이나 궁극적 근거를 찾으려는 시도 자체가 무의미하다.

곽상의 정치 사상에서 독화론은 '각자가 자기 분수에 맞게 자족한다(各適其性)'는 원칙으로 이어진다. 왕필이 통치자의 '비움'을 통한 무위를 강조한 반면, 곽상은 통치자와 백성 각각이 자기 자리에서 자기 몫을 다하는 것이 자연스러운 질서라 보았다. 두 사상가 모두 노자의 무위를 계승하되, 왕필은 '무'라는 본체에 무위를 기초시키고, 곽상은 본체 자체를 부정함으로써 무위를 급진화한 것이다.
}
\kfSecHeader{kfAreaNonlit}{문제}{}

\begin{multicols}{2}\raggedcolumns\setlength{\columnsep}{12pt}
\begin{kfQBlock}{1}{}
\kfQStem{(가)와 (나)의 내용과 일치하지 않는 것은?}
\begin{kfChoices}
\kfChoice 왕필에게 무는 단순한 '없음'이 아니라 만유의 작동 근거이다.
\kfChoice 곽상은 왕필의 무가 또 하나의 형이상학적 실체를 상정하는 것이라 비판했다.
\kfChoice 곽상의 독화론에 따르면 만물은 외부 근거 없이 스스로 생겨난다.
\kfChoice 왕필과 곽상 모두 노자의 무위를 전면 부정하고 유위 정치를 주장했다.
\kfChoice 곽상은 외부의 제1원인이나 궁극적 근거를 찾으려는 시도 자체가 무의미하다고 보았다.
\end{kfChoices}
\end{kfQBlock}

\begin{kfQBlock}{2}{}
\kfQStem{(가)의 왕필과 (나)의 곽상을 비교한 내용으로 가장 적절한 것은?}
\begin{kfChoices}
\kfChoice 왕필은 무를 만유의 본체로, 곽상은 본체를 부정하고 자생을 주장한다.
\kfChoice 왕필과 곽상 모두 무를 만유의 궁극적 근거로 받아들이는 데 의견을 같이한다.
\kfChoice 곽상은 왕필이 정립한 귀무론의 핵심 논리를 그대로 계승하여 한층 발전시켰다.
\kfChoice 왕필은 적극적인 유위 정치를 옹호한 반면 곽상은 철저한 무위 정치를 주장한다.
\kfChoice 두 사상가 모두 만유의 궁극적 제1원인을 반드시 찾아야 한다고 함께 보았다.
\end{kfChoices}
\end{kfQBlock}

\begin{kfQBlock}{3}{}
\kfQStem{윗글에 따를 때, 곽상이 왕필의 '무' 개념을 비판하는 논리적 근거로 가장 적절한 것은?}
\begin{kfChoices}
\kfChoice 무가 유의 근거라면 무 자체도 '무엇'이 되어 또 하나의 실체를 상정하는 자기모순에 빠진다는 점
\kfChoice 왕필의 무 개념은 유교 경전의 본래 정신과 양립할 수 없어 해석상의 충돌이 발생한다는 점
\kfChoice 왕필은 정치적으로 유위(有爲)의 적극적 통치를 옹호함으로써 노자의 본래 입장을 배반한다는 점
\kfChoice 왕필의 무 개념은 어떠한 경험으로도 검증할 수 없어 학문적 객관성을 갖추지 못한다는 점
\kfChoice 노자 자신이 저술 안에서 무 개념을 명시적으로 부정한 사례가 다수 존재한다는 점
\end{kfChoices}
\end{kfQBlock}

\begin{kfQBlock}{4}{}
\kfQStem{윗글의 '독화(獨化)'가 의미하는 바로 가장 적절한 것은?}
\begin{kfChoices}
\kfChoice 만물이 외부 근거 없이 홀로 스스로 변화하여 존재한다는 것
\kfChoice 무(無)라는 본체가 만물을 변화시킨다는 것
\kfChoice 통치자 한 사람이 만물을 독점적으로 다스린다는 것
\kfChoice 만물이 하나의 궁극적 원인에서 비롯된다는 것
\kfChoice 자연의 변화가 인간의 인위적 조작에 의해 이루어진다는 것
\end{kfChoices}
\end{kfQBlock}

\end{multicols}

\kfSecHeader{kfAreaNonlit}{지문}{퇴계 이황의 이발설과 율곡 이이의 기발설 (인문)}

\kfPassageNote{퇴계 이황의 이발설과 율곡 이이의 기발설 (인문)}{%
조선 성리학의 핵심 논쟁인 사단칠정(四端七情) 논쟁은 퇴계 이황(1501\textasciitilde{}1570)과 고봉 기대승(1527\textasciitilde{}1572) 사이에서 시작되어, 이후 율곡 이이(1536\textasciitilde{}1584)에 의해 새로운 방향으로 전개되었다. 이 논쟁의 핵심은 인간의 감정(情)이 리(理)에서 발하는가, 기(氣)에서 발하는가 하는 문제이다.

퇴계 이황은 '이기호발설(理氣互發說)'을 제시했다. 그에 따르면 사단(측은·수오·사양·시비)은 리가 발하고 기가 따르는 것(理發而氣隨之)이며, 칠정(희·노·애·락·애·오·욕)은 기가 발하고 리가 타는 것(氣發而理乘之)이다. 즉 사단은 리가 주도하여 발현한 순선한 감정이고, 칠정은 기가 주도하여 발현한 감정으로서 선악이 뒤섞일 수 있다. 퇴계는 리에 능동성(자발성)을 부여함으로써 리의 도덕적 역할을 강조했다.

율곡 이이는 퇴계의 이기호발설을 비판하고, '기발이승일도설(氣發理乘一途說)'을 제시했다. 율곡에 따르면 실제로 발하는 것은 오직 기뿐이며, 리는 스스로 움직이지 않고 기를 타고 있을 뿐이다. 리는 형체도 조작도 없는 순수한 원리이므로, 스스로 발할 수 없다. 따라서 사단과 칠정은 모두 기가 발하고 리가 타는 하나의 길(一途)이며, 사단은 칠정의 선한 부분이지 별도의 발출 경로를 갖는 것이 아니다.

두 입장의 핵심 차이는 리의 능동성 여부에 있다. 퇴계는 리가 스스로 발할 수 있다고 봄으로써 도덕의 초월적 근거를 확보하려 했다. 만약 리가 스스로 발하지 못한다면, 도덕의 근거가 기의 우발성에 의존하게 되어 도덕의 절대성이 위태로워진다는 우려가 있기 때문이다. 반면 율곡은 리에 능동성을 부여하면 주희의 '리는 형이상, 기는 형이하'라는 원칙이 무너진다고 비판했다. 형체도 없는 리가 스스로 운동한다면, 리는 이미 기적 성격을 갖게 되므로 리와 기의 구분이 불분명해진다는 것이다. 율곡은 리의 능동성 대신 기의 자발성을 인정하면서, 리는 기에 탑승하여 기의 방향을 규정하는 역할을 한다고 보았다.
}
\kfSecHeader{kfAreaNonlit}{문제}{}

\begin{multicols}{2}\raggedcolumns\setlength{\columnsep}{12pt}
\begin{kfQBlock}{1}{}
\kfQStem{윗글의 내용과 일치하지 않는 것은?}
\begin{kfChoices}
\kfChoice 퇴계는 사단이 리가 발하고 기가 따르는 것이라 보았다.
\kfChoice 율곡은 실제로 발하는 것은 오직 기뿐이라 주장했다.
\kfChoice 퇴계와 율곡 모두 사단과 칠정이 별도의 발출 경로를 갖는다고 동의했다.
\kfChoice 율곡에 따르면 사단은 칠정의 선한 부분이지 별도의 경로가 아니다.
\kfChoice 퇴계는 리의 능동성을 부여함으로써 도덕의 초월적 근거를 확보하려 했다.
\end{kfChoices}
\end{kfQBlock}

\begin{kfQBlock}{2}{}
\kfQStem{윗글에 따를 때, 율곡이 퇴계의 이기호발설을 비판하는 핵심 논거로 가장 적절한 것은?}
\begin{kfChoices}
\kfChoice 퇴계가 주희의 이기론 체계를 정면에서 완전히 부정해 버렸기 때문이다.
\kfChoice 퇴계가 칠정을 본래 선한 감정의 영역으로만 한정해 좁혔기 때문이다.
\kfChoice 퇴계의 이기호발설이 불교의 형이상학 사상에 깊이 의존하기 때문이다.
\kfChoice 사단이 리가 아닌 기에서 발한다고 퇴계가 일관되게 주장했기 때문이다.
\kfChoice 형체 없는 리가 스스로 운동하면 리·기 구분이 불분명해지기 때문이다.
\end{kfChoices}
\end{kfQBlock}

\begin{kfQBlock}{3}{}
\kfQStem{윗글의 퇴계와 율곡을 비교한 내용으로 적절하지 않은 것은?}
\begin{kfChoices}
\kfChoice 퇴계는 리에 능동성을 부여하지만, 율곡은 리의 능동성을 부정한다.
\kfChoice 퇴계는 사단과 칠정의 발출 경로가 다르다고 보지만, 율곡은 하나의 길이라 본다.
\kfChoice 퇴계는 도덕의 초월적 근거 확보를, 율곡은 리기 구분의 일관성을 우선시한다.
\kfChoice 퇴계와 율곡 모두 기가 발할 수 있음을 인정한다.
\kfChoice 퇴계와 율곡 모두 리의 능동성을 동일하게 인정한다.
\end{kfChoices}
\end{kfQBlock}

\begin{kfQBlock}{4}{}
\kfQStem{윗글의 '기발이승일도설'이 의미하는 바로 가장 적절한 것은?}
\begin{kfChoices}
\kfChoice 발하는 것은 기뿐, 리는 기에 타며 사단·칠정은 하나의 경로라는 것
\kfChoice 리와 기가 각각 독립적으로 발하여 두 가지 경로가 동시에 존재한다는 것
\kfChoice 기가 리에 비해 존재론적으로 열등하므로 기를 마땅히 제거해야 한다는 것
\kfChoice 사단은 기에서 발하고 칠정은 리에서 각각 따로 발한다는 것
\kfChoice 리와 기가 완전히 같은 것이라는 일원론적 주장
\end{kfChoices}
\end{kfQBlock}

\end{multicols}

\kfSecHeader{kfAreaNonlit}{지문}{다산 정약용의 실학적 경학 (인문)}

\kfPassageNote{다산 정약용의 실학적 경학 (인문)}{%
정약용(1762\textasciitilde{}1836)은 조선 후기 실학의 집대성자로서, 성리학의 경전 해석을 근본적으로 재검토했다. 그의 실학적 경학(經學)은 주희(朱熹)의 해석을 절대시하는 당대 학계에 대한 전복적 도전이었다.

정약용의 경학 방법론의 핵심은 '고증(考證)'과 '실사구시(實事求是)'이다. 주희의 주석을 무조건 따르는 대신, 경전의 원래 맥락과 고대 문자의 본뜻으로 돌아가 재해석한다. 정약용은 주희가 불교·도교의 영향을 받아 유학의 본래 의미를 왜곡했다고 비판하며, 공자·맹자의 원뜻을 복원하려 했다.

대표적 재해석이 '인(仁)' 개념이다. 주희는 인을 '심(心)의 덕(德), 사랑의 리(理)'로 정의하여, 마음이 갖추고 있는 본래적 덕성으로 파악했다. 이에 대해 정약용은 인을 '두 사람 사이의 관계에서 실현되는 것(二人之際)'으로 재정의했다. 인은 마음속에 이미 갖추어진 고정된 덕이 아니라, 사람과 사람 사이의 구체적 관계에서 행위를 통해 실현되는 것이다. 이 해석은 도덕을 내면의 본성 문제에서 사회적 관계와 실천의 문제로 전환시킨다.

또한 정약용은 '예(禮)'에 대해서도 기존 해석을 전환했다. 성리학에서 예는 천리(天理)의 표현이자 마음의 절제로 이해되었지만, 정약용에게 예는 사회적 질서를 유지하기 위한 구체적 제도이자 규범이다. 예의 근거를 형이상학적 천리가 아니라 인간 사회의 현실적 필요에서 찾은 것이다.

정약용의 실학적 경학은 조선 후기 지식 세계에 두 가지 의의를 갖는다. 첫째, 주희 해석의 절대성을 해체하고 경전 해석의 다원성을 열었다. 둘째, 도덕과 예를 형이상학에서 실천과 제도의 영역으로 전환함으로써, 학문이 현실 개혁에 기여할 수 있는 토대를 마련했다.
}
\kfSecHeader{kfAreaNonlit}{문제}{}

\begin{multicols}{2}\raggedcolumns\setlength{\columnsep}{12pt}
\begin{kfQBlock}{1}{}
\kfQStem{윗글의 내용과 일치하지 않는 것은?}
\begin{kfChoices}
\kfChoice 정약용은 주희가 불교·도교의 영향으로 유학의 본래 의미를 왜곡했다고 비판했다.
\kfChoice 주희는 인을 마음이 갖추고 있는 본래적 덕성으로 파악했다.
\kfChoice 정약용은 인을 마음속에 이미 갖추어진 고정된 덕이라고 보았다.
\kfChoice 정약용에게 예는 사회적 질서를 위한 구체적 제도이자 규범이다.
\kfChoice 정약용의 경학은 주희 해석의 절대성을 해체하고 경전 해석의 다원성을 열었다.
\end{kfChoices}
\end{kfQBlock}

\begin{kfQBlock}{2}{}
\kfQStem{윗글의 주희와 정약용을 비교한 내용으로 가장 적절한 것은?}
\begin{kfChoices}
\kfChoice 주희와 정약용 모두 예를 천리의 표현으로 전적으로 동일하게 이해하였다.
\kfChoice 주희는 인을 마음의 본래적 덕으로, 정약용은 관계의 실현으로 보았다.
\kfChoice 정약용은 주희의 경전 해석을 절대시하여 그대로 계승해 학파를 세웠다.
\kfChoice 주희가 실천을, 정약용이 오히려 내면 성찰을 우선적으로 강조하였다.
\kfChoice 정약용이 경전 고증보다는 권위 있는 주석의 해석을 최우선으로 내세웠다.
\end{kfChoices}
\end{kfQBlock}

\begin{kfQBlock}{3}{}
\kfQStem{윗글에 따를 때, 정약용이 인(仁)을 '두 사람 사이의 관계'로 재정의한 것이 갖는 의의로 가장 적절한 것은?}
\begin{kfChoices}
\kfChoice 도덕을 내면의 본성 문제에서 사회적 관계와 실천의 문제로 전환시킨 것이다.
\kfChoice 인을 개인의 내면에 국한시켜 사회적 맥락을 배제한 것이다.
\kfChoice 주희의 인 해석을 그대로 계승하여 정통성을 확보한 것이다.
\kfChoice 인의 실현이 행위와 무관하고 마음의 상태로 충분하다는 것을 강조한 것이다.
\kfChoice 불교의 자비 개념을 유학에 도입한 것이다.
\end{kfChoices}
\end{kfQBlock}

\begin{kfQBlock}{4}{}
\kfQStem{윗글의 '실사구시(實事求是)'가 의미하는 바로 가장 적절한 것은?}
\begin{kfChoices}
\kfChoice 실제 사실에 근거하여 참된 뜻을 구한다는 것
\kfChoice 주희의 주석을 무조건 따르는 태도
\kfChoice 형이상학적 원리에서 도덕의 근거를 찾는 것
\kfChoice 불교와 도교의 교리를 유학에 통합하는 것
\kfChoice 경전 해석의 다원성을 거부하고 하나의 해석만 인정하는 것
\end{kfChoices}
\end{kfQBlock}

\end{multicols}

\kfStartWeekly{패턴 워크북 — 총 ?문항}


\kfSecHeader{kfAreaWeekly}{지문}{문항 1 지문 / 섞어치기}

\kfPassageNote{문항 1 지문 / 섞어치기}{%
무선 통신에 대한 수요가 폭발적으로 증가함에 따라 미국 정부는 ⓐ경매를 통해 통신 사업자들에게 통신 주파수 대역들을 판매하게 되었다. 경매가 시행되기 전에는 ⓑ심사나 ⓒ제비뽑기와 같은 방식을 활용하였다. 심사 방식은 여러 통신 사업자들이 사업 계획서를 제출하면 평가자들이 사업 계획서에 근거하여 무선 통신 사업을 가장 잘 운영할 만한 사업자를 선정하고, 그렇게 선정된 사업자가 주파수를 할당받는 것이었다. 그러나 사업자가 고용한 법률가들과 로비스트들이 서로 자기들의 사업 계획이 최선이라고 주장하는 상황에서 사업 계획서에 대한 평가 절차와 이의 제기 과정을 전부 거쳐 주파수를 할당하기까지 수년이 걸리기도 하였다.

(중략)

1982년 무렵에는 더 이상 심사 방식으로 수많은 주파수를 할당하는 것이 어려운 수준에 이르렀고, 결국 의회는 제비뽑기로 주파수를 판매하는 것에 동의하게 되었다. 제비뽑기 방식은 주파수를 할당하는 데 드는 시간과 노력을 대폭 줄여 주었지만 새로운 문제를 낳았다. 아무 심사도 없이 

(중략)
}
\begin{kfQBlock}{1}{}
\kfQStem{ⓐ\textasciitilde{}ⓒ에 대한 설명으로 가장 적절한 것은?}
\begin{kfChoices}
\kfChoice 미국에서 ⓑ를 통해 주파수를 판매할 때에는 시간의 문제가 발생하였고, ⓒ를 통해 판매할 때에는 투기의 문제가 발생하였다.
\kfChoice 미국에서 ⓐ를 통해 주파수를 판매하기 전에는 ⓑ와 ⓒ를 통한 판매 방식이 동시에 활용되었다.
\kfChoice 미국 정부는 ⓒ를 통해 주파수를 통신 사업자들에게 효율적으로 분배할 수 있었지만, 정부 수입을 극대화하기 어려웠다.
\end{kfChoices}
\end{kfQBlock}

\kfSecHeader{kfAreaWeekly}{지문}{문항 2 지문 / 부정상대어}

\kfPassageNote{문항 2 지문 / 부정상대어}{%
(가) 가상 공간에서 영생의 삶을 누린다는 흥미로운 이야기가 SF 영화의 소재로만 그치지 않는다고 주장하는 입장이 있다. 커즈와일이나 보스트룀은 기억, 가치관, 태도, 정서적 성향처럼 인간의 자아를 구성하는 것들을 특정 정보 패턴으로 만들어 보존할 수 있다면 신체적인 죽음 이후에도 인간이 여전히 생존할 수 있다고 주장한다. 이 입장은 인격 동일성에 관한 심리적 연속성 이론과 관련이 있다. 이 이론에서는 ‘나’가 누구인지를 규정하는 것은 ‘나’의 기억, 신념, 생각, 감정, 희망, 두려움 등의 묶음이라고 보며, 이러한 심리적인 특성들의 집합이 유지·보존되느냐에 ‘나’의 지속 여부가 달려 있다고 말한다.

(중략)

커즈와일이나 보스트룀은 마음이나 정신 상태를 의미하는 심적 상태의 본성을 계산 기능주의의 관점에서 이해한다. 계산 기능주의에서는 심적 상태의 독특한 인과적 역할이 두뇌의 물리적 상태에 의해서 이루어진다고 보기 때문에 정신 작용을 두뇌에 구현된 계산 체계의 작동으로 이해하며, 인간의 사고나

(중략)
}
\begin{kfQBlock}{2}{}
\kfQStem{(가)와 (나)를 바탕으로 [업로딩]을 이해한 내용으로 적절한 것은?}
\begin{kfChoices}
\kfChoice 업로딩은 신체 폐기 후 생존 가능성을 열어 주지만, 비판적 관점에서는 경험의 동일성 문제를 남긴다.
\kfChoice 업로딩은 심리적 특성 일부만 남아도 인격 동일성이 자동 보장된다는 뜻이다.
\kfChoice 업로딩은 물리적 동일성 보존을 유일 조건으로 삼는 접근이다.
\end{kfChoices}
\end{kfQBlock}

\kfSecHeader{kfAreaWeekly}{지문}{문항 3 지문 / 주체객체}

\kfPassageNote{문항 3 지문 / 주체객체}{%
일반적으로 혼인은 남자와 여자가 부부가 되는 일을 말한다. 과거에는 특정한 법적 절차 없이도 결혼식과 같은 사회적 관습에 의한 의식을 치르면 혼인이 성립한 것으로 보았다. 하지만 현재 대부분의 국가는 법에서 규정한 절차를 따라야만 혼인이 성립되는 법률혼 제도를 두고 있으며, 사회적 관습보다는 혼인 신고와 같은 법적 절차가 더 중요하다. 혼인은 사회를 구성하는 기본 단위인 가족을 형성하는 일이기 때문에 국가의 근간을 이루는 중요한 과정이고 이에 따라 법을 통해 관계를 보호하는 것이다. 따라서 현대 사회에서 혼인은 법적 절차에 의해 이루어지게 되고 여기에 참여한 남녀 간에는 법적 효력이 발생하게 된다.

(중략)

혼인이 성립하기 위해서는 법에서 규정하고 있는 실질적 요건과 형식적 요건을 모두 만족하여야 한다. 실질적 요건에서 가장 중요한 것은 혼인 의사의 합치이다. 혼인 의사는 부부 관계를 성립하겠다는 의사로 반드시 혼인 당사자 간의 합의가 필요하다. 또 다른 실질적 요건은 연령과 혼인 자격에 대한

(중략)
}
\begin{kfQBlock}{3}{}
\kfQStem{㉮와 관련한 이해로 가장 적절한 것은?}
\begin{kfChoices}
\kfChoice 협의상 이혼에서도 법원은 재판상 이혼처럼 이혼 사유 존재를 본안 심리한다.
\kfChoice 협의상 이혼에서 법원이 이혼 사유를 심사하지 않는 이유는 당사자 자율 합의를 중심으로 절차가 설계되어 있기 때문이다.
\end{kfChoices}
\end{kfQBlock}

\kfSecHeader{kfAreaWeekly}{지문}{문항 4 지문 / 순서방향}

\kfPassageNote{문항 4 지문 / 순서방향}{%
육상 동물과 수생 동물은 각각의 생활 환경에 따라 호흡 과정에서 어려움을 겪는다. 육상 동물은 폐의 표면에 건조한 공기가 닿으면 폐가 건조해질 수 있는데, 이는 폐를 손상시키고 ⓐ탈수를 일으키며 기체를 용해하는 폐의 능력을 떨어뜨린다. 한편 수생 동물은 가용 산소가 적은 수중 환경과 수온 차에 의해 변화되는 산소 가용성에 적응해야 한다. 또한 수생 동물은 물속에서 몸을 움직이는 데 많은 근육 운동이 필요하고, 찬 바닷물의 높은 비열 때문에 아가미의 열을 빼앗기며, 삼투압 현상으로 인해 수분 균형에 문제가 발생하기도 한다.

(중략)

이러한 환경 속에서 수생 동물은 효과적인 호흡을 위해 아가미라는 특수한 기관을 사용한다. 아가미는 물속에 녹아 있는 산소를 흡수하고 체내에 있는 이산화 탄소를 배출하는 기능을 하는데, 외아가미와 내아가미로 나뉜다. 외아가미는 몸 밖으로 ⓑ돌출된 형태로, 표면적이 넓고 형태가 다양하며 유생기 도롱뇽이나 일부 무척추동물에서 나타난다. 외아가미는 주변 물과의 접촉을 극대화

(중략)
}
\begin{kfQBlock}{4}{}
\kfQStem{윗글의 역류 관계에 대한 이해로 가장 적절한 것은?}
\begin{kfChoices}
\kfChoice 역류 관계에서는 물의 흐름 때문에 혈액이 구심성에서 원심성으로 흐르지 못하게 된다.
\kfChoice 역류 관계에서는 물이 이동하며 계속 산소가 낮은 혈액과 접촉하므로 확산이 계속될 수 있다.
\end{kfChoices}
\end{kfQBlock}

\kfSecHeader{kfAreaWeekly}{지문}{문항 5 지문 / 인과도치}

\kfPassageNote{문항 5 지문 / 인과도치}{%
㉠시민 참여에 대한 긍정적 관점은 고전적인 민주주의 이론가들을 중심으로 주창된 이래 1960년대에 이르러서는 참여 민주주의 이론으로 발전하였다. 이들은 민주적인 사회 체제가 제대로 기능하기 위해서는 시민에게 일정한 시민적 자질이 요구된다고 보았다. 이러한 자질에는 정치 문제에 관한 지식과 민주주의에 대한 신념이 있는데, 높은 수준의 시민적 자질을 바탕으로 시민 참여가 이루어질 때 민주주의 사회는 성공적으로 기능할 수 있다.

(중략)

민주주의 사회에서 시민 참여에 관해 기본적인 가정을 제공한 대표적인 학자로는 고전적인 민주주의 이론가로 분류되는 루소와 밀이 있다. 루소에게 참여는 의사 결정에의 참여를 의미하였다. 루소에 따르면, 인간은 공동의 이익을 추구하는 ‘일반 의지’를 가지고 있으며, 이 일반 의지는 타인에 의해 대표될 수 없기 때문에 선거를 통하여 선출된 의원이 시민들의 의사를 대표하는 대의제하에서의 시민은 단지 의회의 구성원에 대한 선거가 이루어지는 동안만 자유로울 뿐이다. 따라서 루소

(중략)
}
\begin{kfQBlock}{5}{}
\kfQStem{윗글을 바탕으로 <보기>를 이해한 반응으로 적절하지 않은 것은?}
\begin{kfEvidence}<보기> A국은 선거권의 평등을 보장하는 민주적인 선거가 이루어지고 있다. B 씨는 자신이 사는 ○○구 △△동의 하천 오염 문제를 해결하기 위해 주민 자치 위원회 대표 자격으로 최근 선거에서 비례 대표로 당선된 국회 의원을 찾아가 면담했다.\end{kfEvidence}
\begin{kfChoices}
\kfChoice 루소는 A국의 선거를 통해 뽑힌 국회 의원이 B 씨가 가지고 있는 일반 의지를 대표할 수 없다고 보겠군.
\kfChoice 페이트먼은 A국의 국회 의원 선거나 B 씨의 국회 의원 면담은 체제의 존속과 유지에 필요한 자질을 발전시키는 기능이 없다고 보겠군.
\end{kfChoices}
\end{kfQBlock}

\kfSecHeader{kfAreaWeekly}{지문}{문항 6 지문 / 의도/목적}

\kfPassageNote{문항 6 지문 / 의도/목적}{%
미니멀리즘 음악은 1960년대 이후 미국에서 태동한 현대 음악의 한 흐름으로, 라 몬테 영, 테리 라일리, 스티브 라이히, 필립 글래스 등이 대표적 작곡가로 꼽힌다. 이들은 소수의 \uline{음형}(몇 개의 음이 연속하여 어떤 가락이나 악곡의 요소가 되는 음의 모양)이나 단순한 \uline{패턴}을 반복하는 방식으로 기존 음악의 복잡한 서사나 화성 전개를 과감히 축소하여, 감상자가 소리 그 자체를 경험하도록 유도한다.

전통적인 서양 음악에서 반복은 주제를 강조하여 이를 서사적으로 강화하기 위해 사용되었다. 예컨대 제시부-발전부-재현부로 구성된 소나타 형식은 제시부에서 등장한 주제가 발전부에서 여러 차례 변주되다가 재현부에서 제시부에 등장한 형태로 회귀함으로써, 작곡가가 설정한 주제를 확인 및 강화하는 구조를 이룬다. 철학자 \uline{들뢰즈}의 표현을 빌리면, 이와 같은 ㉠ \uline{‘같은 것의 반복’}에서는 특정 주제가 중심에 놓이고 음악의 각 부분은 그 주제를 재확인하는 역할만…(중략)
}
\begin{kfQBlock}{6}{}
\kfQStem{미니멀리즘 음악에 대해 이해한 내용으로 적절한 것은?}
\begin{kfChoices}
\kfChoice 미니멀리즘 음악은 ‘무엇을 의미하는가’보다 ‘그 자체로 존재함’을 중시하며 특정 상징이나 주제를 전달하기 위한 개방적 구조를 지향한다.
\kfChoice 미니멀리즘 음악은 ‘무엇을 의미하는가’보다 ‘그 자체로 존재함’을 중시하며 특정 상징이나 주제를 전달하기 위한 개방적 구조를 지향한다.
\kfChoice 미니멀리즘 음악은 작품의 주제가 감상자에게 전달되는 것보다 감상자의 감상 경험 자체를 중시한다.
\kfChoice 미니멀리즘 음악에서 단순한 음형이나 패턴의 반복은 소리 자체의 순간적 변화를 드러내기 위한 작곡 전략이다.
\end{kfChoices}
\end{kfQBlock}

\kfSecHeader{kfAreaWeekly}{지문}{문항 7 지문 / 상관관계 반전}

\kfPassageNote{문항 7 지문 / 상관관계 반전}{%
근대 경제학은 국민 소득, 물가, 고용, 경기 변동 등 경제의 \uline{거시적 현상}을 설명하기 위해 다양한 \uline{모형}을 \uline{ⓐ발전시켜} 왔다. \uline{고전학파}는 시장의 \uline{자율 조정} 능력을 강조하며, 가격과 임금의 \uline{유연성}을 전제로 \uline{완전 고용}이 자연스럽게 달성된다고 보았다. 그러나 20세기 초 \uline{대공황}을 계기로 시장의 자율 조정에 한계가 있다는 비판이 제기되었고, 이는 \uline{케인스 경제학}의 등장으로 이어졌다. 케인스는 \uline{총수요}(한 경제 내에서 주어진 기간과 물가 수준하에 최종 재화와 서비스에 대한 수요의 총합)가 자동으로 \uline{총공급}(주어진 기간과 물가 수준하에 국가 경제에서 재화와 서비스에 대한 공급의 총합)에 맞춰진다고 보는 고전학파의 이론을 비판하며, 낮은 총수요가 소득의 감소 및 실업률 상승을 초래할 수 있다고 보았다. 따라서 그는 정부 지출 확대 등 정책 개입을 통해 총수요를 관리해야…(중략)
}
\begin{kfQBlock}{7}{}
\kfQStem{다음 <보기>의 상황을 IS-MP 모형으로 이해한 내용으로 가장 적절하지 않은 것은?}
\begin{kfEvidence}다음 그래프는 현재 경제의 균형 상태(E)를 보여 주는 IS-MP 모형이다. 이후 중앙은행이 경기 부양을 위해 완화적 정책 기조로 전환하였다.

[그래프] 세로축: 이자율, 가로축: 국민 소득. 우하향하는 IS 곡선과 우상향하는 MP 곡선이 교차하며, 교차점이 E이다.\end{kfEvidence}
\begin{kfChoices}
\kfChoice 중앙은행이 완화적 기조로 전환하면 MP 곡선은 아래쪽으로 이동하여 새로운 균형이 형성될 수 있다.
\kfChoice 중앙은행이 완화적 기조로 전환하면 MP 곡선은 아래쪽으로 이동하여 새로운 균형이 형성될 수 있다.
\kfChoice 중앙은행이 완화적 기조로 전환하면 MP 곡선은 위쪽으로 이동할 것이다.
\kfChoice 중앙은행의 정책 기조 변화는 내생적 변화이므로 MP 곡선 위에서 점만 이동할 것이다.
\end{kfChoices}
\end{kfQBlock}

\kfSecHeader{kfAreaWeekly}{지문}{문항 8 지문 / 필요조건}

\kfPassageNote{문항 8 지문 / 필요조건}{%
개별적이고 특수한 작업을 수행하는 여러 대의 클라이언트와 연결되어 정보나 서비스를 제공하는 컴퓨터 시스템을 서버라고 한다. 서버는 클라이언트의 요청에 따라 작업을 수행하거나 정보를 검색하고, 그 결과를 클라이언트에 돌려준다. 서버를 구성하는 기계에서 작업을 신속하고 정확하게 처리하기 위해서는 데이터를 안전하게 관리해야 할 뿐 아니라 신속하게 데이터에 접근하여 데이터의 읽기나 쓰기를 실행할 필요가 있다. 이런 필요를 충족시키기 위해 개발된 것이 \uline{㉠분산 파일 관리 체계}이다. 분산 파일 관리 체계는 데이터의 불법 변조를 방지하고, 기계적 결함에 대응하여 데이터의 유실로부터 안전을 도모하며, 서버의 부하를 줄여 신속한 연산을 수행하도록 기계적 성능과 저장 용량을 확보하는 데 유리한 전략이다. 단일 기계로 구성된 서버에서 데이터를 관리할 경우에는 하드웨어적으로 램이나 CPU의 성능이 우수한 기계를 확보하는 데 큰 비용이 들지만, 분산 파일 관리 체계에서는 다수의 저렴한 저성능 기…(중략)
}
\begin{kfQBlock}{8}{}
\kfQStem{㉢의 취지를 고려할 때 적절한 것은?}
\begin{kfChoices}
\kfChoice 샤드마다 서로 다른 데이터만 두고 복제 없이 운영해도 ㉢이 자연히 높아진다.
\kfChoice 샤드마다 서로 다른 데이터만 두고 복제 없이 운영해도 ㉢이 자연히 높아진다.
\kfChoice 복제본을 함께 두어 장애 발생 시 정상 데이터로 복구할 수 있게 하는 것은 ㉢을 높이는 방식이다.
\kfChoice 읽기 요청이 몰릴 때 복제본으로 트래픽을 분산하는 것은 ㉢과 연관된 가용성 향상 전략이 될 수 있다.
\end{kfChoices}
\end{kfQBlock}

\kfSecHeader{kfAreaWeekly}{지문}{문항 9 지문 / 결합/분리}

\kfPassageNote{문항 9 지문 / 결합/분리}{%
근대 경제학은 국민 소득, 물가, 고용, 경기 변동 등 경제의 \uline{거시적 현상}을 설명하기 위해 다양한 \uline{모형}을 \uline{ⓐ발전시켜} 왔다. \uline{고전학파}는 시장의 \uline{자율 조정} 능력을 강조하며, 가격과 임금의 \uline{유연성}을 전제로 \uline{완전 고용}이 자연스럽게 달성된다고 보았다. 그러나 20세기 초 \uline{대공황}을 계기로 시장의 자율 조정에 한계가 있다는 비판이 제기되었고, 이는 \uline{케인스 경제학}의 등장으로 이어졌다. 케인스는 \uline{총수요}(한 경제 내에서 주어진 기간과 물가 수준하에 최종 재화와 서비스에 대한 수요의 총합)가 자동으로 \uline{총공급}(주어진 기간과 물가 수준하에 국가 경제에서 재화와 서비스에 대한 공급의 총합)에 맞춰진다고 보는 고전학파의 이론을 비판하며, 낮은 총수요가 소득의 감소 및 실업률 상승을 초래할 수 있다고 보았다. 따라서 그는 정부 지출 확대 등 정책 개입을 통해 총수요를 관리해야…(중략)
}
\begin{kfQBlock}{9}{}
\kfQStem{윗글을 바탕으로 <보기>를 이해한 내용으로 적절한 것은?}
\begin{kfEvidence}다음 그래프는 현재 경제의 균형 상태(E)를 보여 주는 IS-MP 모형이다.

[그래프] 세로축: 이자율, 가로축: 국민 소득. 우하향하는 IS 곡선과 우상향하는 MP 곡선이 교차하며, 교차점이 E이다.\end{kfEvidence}
\begin{kfChoices}
\kfChoice 정부가 재정 지출을 확대하면 IS 곡선은 오른쪽으로 이동하여 총수요가 늘어날 것이며, 이에 따라 높아진 국민 소득에 반응하여 중앙은행이 기준 금리를 올리게 되므로 MP 곡선 자체도 위쪽으로 이동하여 새로운 균형을 이룰 것이다.
\kfChoice 정부가 재정 지출을 확대하면 IS 곡선은 오른쪽으로 이동하여 총수요가 늘어날 것이며, 이에 따라 높아진 국민 소득에 반응하여 중앙은행이 기준 금리를 올리게 되므로 MP 곡선 자체도 위쪽으로 이동하여 새로운 균형을 이룰 것이다.
\kfChoice MP 곡선이 우상향하는 것은 국민 소득이 증가할수록 중앙은행이 기준 금리를 인상하는 방향으로 반응하는 경향을 반영하기 때문일 것이다.
\kfChoice 국민 소득이 감소하면 중앙은행은 기준 금리를 인하하는 방향으로 반응하므로, 이는 MP 곡선상에서 금리가 내생적으로 조정되는 과정으로 나타날 것이다.
\end{kfChoices}
\end{kfQBlock}

\kfSecHeader{kfAreaWeekly}{지문}{문항 10 지문 / 조건 왜곡}

\kfPassageNote{문항 10 지문 / 조건 왜곡}{%
어느 시대에서나 기적에 대한 보고는 꾸준히 있었고, 오늘날에도 기적을 목격했다는 사람은 부지기수로 많다. 죽은 사람이 살아난다든가 바위가 스스로 움직여 떠다니는 것을 보았다는 증언이 그런 예이다. 어떤 사건을 \uline{㉠기적적 사건}이라고 해야 할까?

기적이 특이한 사건이라야 한다는 점에는 동의할 수 있다. 그러나 일어나기 힘든 일이긴 하지만 자연법칙으로 설명될 수 있는 한 그것은 기적적 사건으로 간주될 수 없다. 이순신 장군의 명량 대첩처럼 13척의 배로 133척의 왜적을 물리친 일은 일어나기에 아주 힘든 일일 뿐이지 자연법칙에 어긋나지는 않는다. 그렇다면 우리에게 알려진 어떤 자연법칙으로도 설명할 수 없는 사건은 기적일까? 사진 감광판을 완전히 어두운 곳에 줄곧 보관했는데도 감광판이 빛에 노출된다는 사실을 맨 처음 알게 되었을 때, 이 사건은 당시에 알려진 어떤 자연법칙으로도 설명할 수 없었다. 그러나 머지않아 방사능에 관한 과학이 정립되어 이 진기한 현상에 대한 과학적 설명…(중략)
}
\begin{kfQBlock}{10}{}
\kfQStem{<보기>에서 ㉠에 해당하는 진술들만 묶은 것은?}
\begin{kfEvidence}ㄱ. 벼락에 맞았는데도 이렇게 살아 있다니, 기적임이 틀림없다. ㄴ. 마법사는 초자연적 힘으로 돌멩이를 황금으로 바꾸는 기적을 일으켰다. ㄷ. 6·25 전쟁을 겪은 대한민국이 경제를 빠르게 재건한 것은 한강의 기적이다. ㄹ. 당첨 확률이 백만분의 1인 로또에 두 번 연속 당첨되다니 기적이 아닐 수 없다. ㅁ. 쇠막대가 하늘로 날아오르는 기적 같은 일이 일어났는데, 사실은 헛것을 본 것이다.\end{kfEvidence}
\begin{kfChoices}
\kfChoice ㄴ
\kfChoice ㄴ
\kfChoice ㄱ, ㄴ
\kfChoice ㄱ, ㄹ, ㅁ
\end{kfChoices}
\end{kfQBlock}

\kfSecHeader{kfAreaWeekly}{지문}{문항 11 지문 / 긍부정 반전}

\kfPassageNote{문항 11 지문 / 긍부정 반전}{%
\# 
}
\begin{kfQBlock}{11}{}
\kfQStem{㉠을 중심으로 (나), (다)를 이해한 내용으로 가장 적절하지 않은 것은?}
\begin{kfChoices}
\kfChoice (나)와 (다)는 모두 음성상징어를 활용하여 시적 분위기를 부각하고 있다.
\kfChoice (나)와 (다)는 모두 음성상징어를 활용하여 시적 분위기를 부각하고 있다.
\kfChoice (나)는 청유형 어미를 활용하여 주제 의식을 부각하고 있다.
\kfChoice (다)는 수미상관의 형식을 활용하여 시적 상황을 부각하고 있다.
\end{kfChoices}
\end{kfQBlock}

\kfSecHeader{kfAreaWeekly}{지문}{문항 12 지문 / 비교대상 뒤바꿈}

\kfPassageNote{문항 12 지문 / 비교대상 뒤바꿈}{%
\#\#\# [43 \textasciitilde{} 45] 다음 글을 읽고 물음에 답하시오.

(가)   여기저기서 단풍잎 같은 슬픈 가을이 뚝뚝 떨어진다. 단풍잎 떨어져 나온 자리마다 봄을 마련해 놓고 나뭇가지 위에 하늘이 펼쳐 있다. 가만히 하늘을 들여다 보려면 눈썹에 파란 물감이 든다. 두 손으로 따뜻한 볼을 씃어* 보면 손바닥에도 파란 물감이 묻어 난다. 다시 손바닥을 들여다 본다. 손금에는 맑은 강물이 흐르고, 맑은 강물이 흐르고, 강물 속에는 사랑처럼 슬픈 얼굴 ― 아름다운 순이의 얼굴이 어린다. 소년은 황홀히 눈을 감아 본다. 그래도 맑은 강물은 흘러 사랑처럼 슬픈 얼굴 ― 아름다운 순이의 얼굴은 어린다.

* 윤동주, ｢ 소년 ｣ -   * 씃어 : 씻어.

(나)   할머니들이 아파트 앞에 모여 햇볕을 쪼이고 있다. 굵은 주름 잔주름 하나도 놓치지 않고   [A] 꼼꼼하게 햇볕을 채워넣고 있다.   겨우내 얼었던 뼈와 관절들 다 녹도록   온몸을 노곤노곤하게 지지고 있다.   [B] 마른버짐 사이로 …(중략)
}
\begin{kfQBlock}{12}{}
\kfQStem{(가)와 (나)의 공통점으로 가장 적절한 것은?}
\begin{kfChoices}
\kfChoice 현재 시제를 활용하여 시적 상황을 제시하고 있다.
\kfChoice 현재 시제를 활용하여 시적 상황을 제시하고 있다.
\kfChoice 연쇄법을 활용하여 역동적인 분위기를 형성하고 있다.
\kfChoice 다양한 음성 상징어를 사용하여 대상을 묘사하고 있다.
\end{kfChoices}
\end{kfQBlock}

\kfSecHeader{kfAreaWeekly}{지문}{문항 13 지문 / 주체/객체 혼동}

\kfPassageNote{문항 13 지문 / 주체/객체 혼동}{%
\#\#\# [22～27] 다음 글을 읽고 물음에 답하시오.

(가)   첩첩산중에도 없는 마을이 여긴 있습니다. 잎 진 사잇길 저 모랫둑, 그 너머 강기슭에서도 보이진 않습니다. 허방다리* 들어내면 보이는 마을.   갱 속 같은 마을. ㉠꼴깍, 해가, 노루꼬리 해가 지면 집집마다 봉당에 불을 켜지요. 콩깍지, 콩깍지처럼 후미진 외딴집, 외딴집에도 불빛은 앉아 이슥토록 창문은 모과빛입니다.   기인 밤입니다. 외딴집 노인은 홀로 잠이 깨어 출출한 나머지 무우를 깎기도 하고 고구마를 깎다, 문득 바람도 없는데 시나브로 풀려 풀려 내리는 짚단, 짚오라기의 설레임을 듣습니다. 귀를 모으고 듣지요. ㉡후루룩 후루룩 처마 깃에 나래 묻는 이름 모를 새, 새들의 온기를 생각합니다. 숨을 죽이고 생각하지요.   참 오래오래, 노인의 자리맡에 밭은기침 소리도 없을 양이면 벽 속에서 겨울 귀뚜라미는 울지요. 떼를 지어 웁니다, 벽이 무너지라고 웁니다.   어느덧 밖에는 눈발이라도 치는지, 펄펄 함박눈이라…(중략)
}
\begin{kfQBlock}{13}{}
\kfQStem{㉠～㉤에 대한 설명으로 적절하지 않은 것은?}
\begin{kfChoices}
\kfChoice ㉢ : 높이 날아오른 연을 동경하는 심리를 드러내고 있다.
\kfChoice ㉢ : 높이 날아오른 연을 동경하는 심리를 드러내고 있다.
\kfChoice ㉠ : 아주 짧은 순간에 해가 지는 모습을 나타낸 말로, 시간의 변화를 함축하고 있다.
\kfChoice ㉡ : 소리를 통해 연상되는 새의 모습을 감각적으로 형상화하고 있다.
\end{kfChoices}
\end{kfQBlock}

\kfSecHeader{kfAreaWeekly}{지문}{문항 14 지문 / 내용과 방법 혼동}

\kfPassageNote{문항 14 지문 / 내용과 방법 혼동}{%
[27 \textasciitilde{} 30] 다음 글을 읽고 물음에 답하시오.

2424 혹은 5454번의 전화번호를 보디에 커다랗게 써 붙인 삼륜차 또는 픽업이 대충 비슷비슷한 내용물들을 실은 채 속속들이 닿고 있었고, 감색 유니폼의 관리인들이 요소요소마다 늘어선 채 똑같은 말들을 외쳐 대고 있었다. 일테면,   “차는 현관 옆으로 바짝 붙여 주십시오!”   “호실 키는 임시 관리 사무소에서 입주증과 교환해 드리고 있습니다. 관리 사무소는 217동과 219동 사이에 위치하고 있습니다…….”   “계단이 혼잡하오니 도착순대로 짐을 올리시고, 화장실 및 주방의 부착물은 248동과 249동 간에 위치하고 있는…….”   삼륜차 위에서 나는 한동안 멍청하게 흔들리고만 있었다. 수백 수천의 똑같은 5층짜리 콘크리트 건물군과 그리고 그 협곡 사이사이마다 출렁이고 있는 입주자들의 행렬……. 그것은 실로 기이한 대조였다. 나는 무거운 압박감과 마음 붙일 곳 없는 황량함을 동시에 의식하지 않을 수 없었다. 차가움, 견고함,…(중략)
}
\begin{kfQBlock}{14}{}
\kfQStem{[A]와 [B]의 서술상 특징으로 가장 적절하지 않은 것은?}
\begin{kfChoices}
\kfChoice [A]는 장면에 대한 관찰을 중심으로, [B]는 인물의 복잡한 내면을 중심으로 서술하고 있다.
\kfChoice [A]는 장면에 대한 관찰을 중심으로, [B]는 인물의 복잡한 내면을 중심으로 서술하고 있다.
\kfChoice [A]는 사건 해결의 실마리를 중심으로, [B]는 인물의 행위에 담긴 의미를 중심으로 서술하고 있다.
\kfChoice [A]는 인물들 간에 심화되는 갈등을 중심으로, [B]는 인물이 겪는 내적 갈등을 중심으로 서술하고 있다.
\end{kfChoices}
\end{kfQBlock}

\kfSecHeader{kfAreaWeekly}{지문}{문항 15 지문 / 내용 왜곡}

\kfPassageNote{문항 15 지문 / 내용 왜곡}{%
[22 \textasciitilde{} 27] 다음 글을 읽고 물음에 답하시오.

(가)   몰아라 어서 보자 총석정 어서 보자   총석정 좋단 말을 일찍이 들었거니   바람 불면 못 보려니 몰아라 어서 보자   벽해 위의 높은 집이 저것이 총석정인가   올라 보니 후면이라 전면으로 보오리라   배 대어라 사공들아 풍랑이 일지 않아   층파로 돌아 저어 총석 전면 보게 하라   배 띄워라 굽이마다 따라 저어 볼 양이면   영소전 태을궁*을 지으려고 경영턴가*   \textit{돌기둥 천백 개를 육모로 깎아 내어}   \textit{개개이 묶어 세워 몇 만 년이 되었던지}   \textit{황량한 데 벌였으니 배 없어 못 실린가}   \textit{(중략)}   \textit{하우씨 도끼뿔이 용문을 뚫었으나}   *이 돌*을 만났으면 이같이 깎을세며   영장*이 신묘하여 코끝의 것 찍었으나*   \textit{이 돌을 다듬는다고 이같이 곧을쏘냐}   \textit{어떠한 도끼로 용이히 깎았으며}   *어떠한 승묵*으로 천연히 골랐는고   끈 없이 묶었으되 틈 없이 묶었으며   풀 없이 붙였으되 흔적…(중략)
}
\begin{kfQBlock}{15}{}
\kfQStem{(다)에 대한 설명으로 가장 적절한 것은?}
\begin{kfChoices}
\kfChoice 연상을 통해 다양한 대상을 열거하며 공간에 대한 애정을 드러내고 있다.
\kfChoice 연상을 통해 다양한 대상을 열거하며 공간에 대한 애정을 드러내고 있다.
\kfChoice 상황에 따라 의성어를 다채롭게 구사하여 현장감을 부각하고 있다.
\kfChoice 말줄임표를 통해 과거의 연인과의 재회에 대한 회의감을 표현하고 있다.
\end{kfChoices}
\end{kfQBlock}

\kfSecHeader{kfAreaWeekly}{지문}{문항 16 지문 / 과잉 해석}

\kfPassageNote{문항 16 지문 / 과잉 해석}{%
[18 \textasciitilde{} 23] 다음 글을 읽고 물음에 답하시오.

(가)   삼 년을 임을 떠나 해도(海島)에 유배되니   ㉠ 내 언제 무심하여 임에게 득죄했나   임이 언제 박정(薄情)하여 날 대접 소홀히 했나   내 얼굴 고왔던지 질투하는 건 뭇 여자로다   유한한* 이내 몸을 음란하다 이르로세   (중략)   긴 소매 들고 앉아 옛 잘못을 헤아리니   우직하기 본성이오 망령됨도 내 죄로되   근본을 생각하면 임 위한 정성일세   일월 같은 우리 임이 거의 아니 굽어볼까   날 살리신 이 은혜를 결초(結草)하기 생각하나   광주리의 가을 부채 어느 날 다시 날꼬   황금을 못 얻으니 장문부*를 어이 사리*   \textit{마름과 연(蓮)으로 옷을 짓고 부용(芙蓉)으로 치마 지어}   \textit{상자 안에 두어신들 눌 위하여 단장할꼬}   \textit{고향에 돌아갈 꿈 벽해(碧海)를 밟아 건너}   \textit{옥루(玉樓) 높은 곳에 밤마다 임을 모셔}   \textit{일당우불에 수답이 여향하니}   가까이 다가앉아 귀신을 묻던 가태부 이러한가…(중략)
}
\begin{kfQBlock}{16}{}
\kfQStem{<보기>를 바탕으로 (가)를 감상한 내용으로 적절하지 않은 것은? [3점]}
\begin{kfChoices}
\kfChoice ‘이내 몸을’ ‘일월 같은 우리 임이 거의 아니 굽어볼까’라고 한 것은 작가가 유배지에서 생활하고 있는 자신의 일상에 관심을 보이는 임금에 대한 감사함을 드러낸 것이겠군.
\kfChoice ‘이내 몸을’ ‘일월 같은 우리 임이 거의 아니 굽어볼까’라고 한 것은 작가가 유배지에서 생활하고 있는 자신의 일상에 관심을 보이는 임금에 대한 감사함을 드러낸 것이겠군.
\kfChoice ‘뭇 여자’가 ‘질투하’여 ‘음란하다 이르’었다고 한 것은 작가가 반대파의 모함을 받아 유배되었다고 생각하고 있음을 나타낸 것이겠군.
\kfChoice ‘옛 잘못’에 대해 ‘근본을 생각하면 임 위한 정성일세’라고 한 것은 작가가 자신의 시련이 임금을 위한 충정에서 비롯되었다고 생각하고 있음을 나타낸 것이겠군.
\end{kfChoices}
\end{kfQBlock}

\kfSecHeader{kfAreaWeekly}{지문}{문항 17 지문 / 범위 오류}

\kfPassageNote{문항 17 지문 / 범위 오류}{%
\#\#\# 2024년 11월 고3 수능

[22-27] 다음 글을 읽고 물음에 답하시오.

(가)   배를 민다   배를 밀어보는 것은 아주 드문 경험   희번덕이는 잔잔한 가을 바닷물 위에   배를 밀어넣고는   온몸이 아주 추락하지 않을 순간의 한 허공에서   밀던 힘을 한껏 더해 밀어주고는   아슬아슬히 배에서 떨어진 손, 순간 환해진 손을   허공으로부터 거둔다

사랑은 참 부드럽게도 떠나지   뵈지도 않는 길을 부드럽게도

배를 한껏 세게 밀어내듯이 슬픔도   그렇게 밀어내는 것이지

배가 나가고 남은 빈 물 위의 흉터   잠시 머물다 가라앉고

그런데 오, 내 안으로 들어오는 배여   아무 소리 없이 밀려들어오는 배여   - 장석남, <배를 밀며>

(나)   당신……, 당신이라는 말 참 좋지요, 그래서 불러봅니다 킥킥거리며 한때 적요로움의 울음이 있었던 때, 한 슬픔이 문을 닫으면 또 한 슬픔이 문을 여는 것을 이만큼 살아옴의 상처에 기대, 나 킥킥……, 당신을 부릅니다 단…(중략)
}
\begin{kfQBlock}{17}{}
\kfQStem{문항별}
\begin{kfChoices}
\kfChoice 번은 '통제할 수 없는 익명의 욕구'를 상대를 향한 사랑이 운명적이어서 멈출 수 없다는 뜻으로 해석했습니다.
\kfChoice 번은 '통제할 수 없는 익명의 욕구'를 상대를 향한 사랑이 운명적이어서 멈출 수 없다는 뜻으로 해석했습니다.
\kfChoice 번은 배가 들어오는 것을 '대상과의 재회가 예상대로 이루어짐'이라고 해석했습니다.
\kfChoice 번은 '당신'을 화자의 내면에 사는 '병자'로서 연민의 대상이라고 설명했습니다.
\end{kfChoices}
\end{kfQBlock}

\kfSecHeader{kfAreaWeekly}{지문}{문항 18 지문 / 시점/화자 혼동}

\kfPassageNote{문항 18 지문 / 시점/화자 혼동}{%
[18 \textasciitilde{} 23] 다음 글을 읽고 물음에 답하시오.

(가)   삼 년을 임을 떠나 해도(海島)에 유배되니   ㉠ 내 언제 무심하여 임에게 득죄했나   임이 언제 박정(薄情)하여 날 대접 소홀히 했나   내 얼굴 고왔던지 질투하는 건 뭇 여자로다   유한한* 이내 몸을 음란하다 이르로세   (중략)   긴 소매 들고 앉아 옛 잘못을 헤아리니   우직하기 본성이오 망령됨도 내 죄로되   근본을 생각하면 임 위한 정성일세   일월 같은 우리 임이 거의 아니 굽어볼까   날 살리신 이 은혜를 결초(結草)하기 생각하나   광주리의 가을 부채 어느 날 다시 날꼬   황금을 못 얻으니 장문부*를 어이 사리*   \textit{마름과 연(蓮)으로 옷을 짓고 부용(芙蓉)으로 치마 지어}   \textit{상자 안에 두어신들 눌 위하여 단장할꼬}   \textit{고향에 돌아갈 꿈 벽해(碧海)를 밟아 건너}   \textit{옥루(玉樓) 높은 곳에 밤마다 임을 모셔}   \textit{일당우불에 수답이 여향하니}   가까이 다가앉아 귀신을 묻던 가태부 이러한가…(중략)
}
\begin{kfQBlock}{18}{}
\kfQStem{(나)의 시상 전개에 대한 설명으로 적절하지 않은 것은?}
\begin{kfChoices}
\kfChoice <제3수>의 초장에서는 <제2수>의 종장에 제시된 소망이 실현될 것이라는 화자의 믿음이 드러난다.
\kfChoice <제3수>의 초장에서는 <제2수>의 종장에 제시된 소망이 실현될 것이라는 화자의 믿음이 드러난다.
\kfChoice <제1수>에서는 화자에게 일어난 일이 시간의 순서에 따라 제시된다.
\kfChoice <제2수>의 중장에서는 초장에 제시된 상황과 관련된 화자의 정서가 드러난다.
\end{kfChoices}
\end{kfQBlock}



\clearpage

\kfChapterCover{4}{비트겐슈타인3 ch04}{Chapter 4}{이번 챕터: 문법 → 문학 5작품 → 비문학 5지문 → 패턴 워크북.}

\kfStartGrammar{이번 챕터 문법 영역 — 음운 / 교체·탈락·첨가·축약 종합 + 음운 변동 규칙의 적용 순서}


\kfSecHeader{kfAreaGrammar}{문제}{}

\begin{multicols}{2}\raggedcolumns\setlength{\columnsep}{12pt}
\kfPassageFull{문항 1 지문 (concept)}{%
「구개음화는 끝소리가 'ㄷ, ㅌ'인 형태소 뒤에 모음 'ㅣ'나 반모음 'ㅣ'로 시작하는 형식 형태소가 결합할 때, 'ㄷ'은 'ㅈ'으로, 'ㅌ'은 'ㅊ'으로 바뀌는 현상이다. '같이[가치]', '굳이[구지]', '맏이[마지]' 등이 이에 해당한다. 구개음화는 잇몸소리가 센입천장소리로 바뀌는 교체 현상이다. 한편, 'ㅎ' 축약은 'ㅎ'이 뒤 음절 초성의 예사소리(ㄱ, ㄷ, ㅈ, ㅂ)와 결합하여 거센소리(ㅋ, ㅌ, ㅊ, ㅍ)가 되는 현상이다. '놓다[노타]', '좋고[조코]' 등이 그 예이다. 축약은 두 음운이 하나로 합쳐지므로 음운의 개수가 줄어든다.」
}
\begin{kfQBlock}{1}{}
\kfQStem{윗글을 바탕으로 <보기>의 발음에서 일어나는 음운 변동을 분석한 내용으로 적절하지 않은 것은?

<보기> ㉠ 해돋이[해도지]   ㉡ 닫히다[다치다] ㉢ 놓지[노치]       ㉣ 쌓이다[싸이다] ㉤ 밭이[바치]}
\begin{kfChoices}
\kfChoice ㉠에서 'ㄷ'이 'ㅣ' 앞에서 'ㅈ'으로 교체되는 구개음화가 일어났다.
\kfChoice ㉡에서 'ㅌ'과 'ㅎ'이 축약되어 'ㅊ'이 된 것이 아니라, 'ㄷ'이 구개음화된 것이다.
\kfChoice ㉢에서 'ㅎ'과 'ㅈ'이 축약되어 'ㅊ'이 된다.
\kfChoice ㉣에서 'ㅎ'이 모음 사이에서 탈락하여 [싸이다]가 된다.
\kfChoice ㉤에서 'ㅌ'이 접미사 '-이' 앞에서 구개음화되어 [바치]가 된다.
\end{kfChoices}
\end{kfQBlock}

\kfPassageFull{문항 2 지문 (concept)}{%
「탈락은 음운이 일정한 환경에서 사라지는 현상이다. 국어에서 대표적인 탈락 현상으로는 자음군 단순화, 'ㅎ' 탈락, 'ㅡ' 탈락, 동일 모음 탈락 등이 있다. 자음군 단순화는 겹받침에서 자음 하나가 탈락하는 것이고, 'ㅎ' 탈락은 'ㅎ'이 모음 앞에서 사라지는 것이다. 'ㅡ' 탈락은 '쓰-+-어→써', '끄-+-어→꺼'처럼 'ㅡ'로 끝나는 어간에 모음 어미가 결합할 때 'ㅡ'가 탈락하는 현상이다. 동일 모음 탈락은 '가-+-아→가'처럼 같은 모음이 연속될 때 하나가 탈락하는 것이다.」
}
\begin{kfQBlock}{2}{}
\kfQStem{윗글을 바탕으로 <보기>의 활용형에서 일어나는 탈락 현상을 분석한 내용으로 적절한 것은?

<보기> ㉠ 놓아[노아]   ㉡ 크-+-어→커   ㉢ 서-+-어→서 ㉣ 담그-+-어→담가   ㉤ 넣어[너어]}
\begin{kfChoices}
\kfChoice ㉠에서는 'ㅎ' 축약이 일어났다.
\kfChoice ㉡에서는 동일 모음 탈락이 일어났다.
\kfChoice ㉢에서는 'ㅡ' 탈락이 일어났다.
\kfChoice ㉣에서는 'ㅡ' 탈락이 일어났다.
\kfChoice ㉤에서는 'ㅎ' 탈락과 동일 모음 탈락이 모두 일어났다.
\end{kfChoices}
\end{kfQBlock}

\kfPassageFull{문항 3 지문 (concept)}{%
「첨가는 형태소의 결합 과정에서 원래 없던 음운이 새로 끼어드는 현상이다. 대표적으로 'ㄴ' 첨가와 사이시옷 현상이 있다. 'ㄴ' 첨가는 합성어나 파생어에서 앞 형태소가 자음으로 끝나고 뒷 형태소가 'ㅣ, ㅑ, ㅕ, ㅛ, ㅠ'로 시작할 때 'ㄴ'이 첨가되는 현상이다. '꽃잎[꼰닙]', '솜이불[솜니불]' 등이 그 예이다. 사이시옷 현상은 합성어에서 앞말이 모음으로 끝나고 뒷말의 첫소리가 된소리로 발음되거나, 'ㄴ' 소리가 첨가되는 경우에 나타난다. '냇가[내:까]', '나뭇잎[나무닙]' 등이 이에 해당한다.」
}
\begin{kfQBlock}{3}{}
\kfQStem{윗글을 바탕으로 <보기>의 발음에서 일어나는 음운 변동을 분석한 내용으로 적절하지 않은 것은?

<보기> ㉠ 색연필[생년필]   ㉡ 나뭇잎[나문닙] ㉢ 한여름[한녀름]   ㉣ 식용유[시굥뉴]}
\begin{kfChoices}
\kfChoice ㉠에서 'ㄴ' 첨가 후 비음화가 일어나 [생년필]이 된다.
\kfChoice ㉡에서 사이시옷에 의해 'ㄴ'이 첨가되고 비음화가 일어난다.
\kfChoice ㉢에서 'ㄴ' 첨가만 일어나 [한녀름]이 된다.
\kfChoice ㉣에서 'ㄴ' 첨가가 두 번 일어나 [시굥뉴]가 된다.
\kfChoice ㉠과 ㉢에서 첨가되는 'ㄴ'은 모두 뒷 형태소가 'ㅣ, ㅑ, ㅕ, ㅛ, ㅠ'로 시작하는 환경에서 일어난다.
\end{kfChoices}
\end{kfQBlock}

\kfPassageFull{문항 4 지문 (concept)}{%
「경음화(된소리되기)는 예사소리가 된소리로 바뀌는 교체 현상이다. 주요 환경은 다음과 같다. ① 받침 'ㄱ, ㄷ, ㅂ'(파열음) 뒤에서 예사소리가 된소리로 바뀐다: '학교[학꾜]', '국밥[국빱]'. ② 어간 받침 'ㄴ, ㅁ' 뒤에서 어미의 첫소리 'ㄱ, ㄷ, ㅅ, ㅈ'이 된소리로 바뀐다: '신고[신꼬]', '감다[감따]'. ③ 관형사형 어미 '-(으)ㄹ' 뒤에서 예사소리가 된소리로 바뀐다: '갈 곳[갈꼳]'. ④ 한자어에서 'ㄹ' 받침 뒤에 'ㄷ, ㅅ, ㅈ'이 오면 된소리로 바뀐다: '발달[발딸]', '갈등[갈뜽]'. ⑤ 사이시옷 뒤에서 된소리가 된다: '깃발[기빨]'.」
}
\begin{kfQBlock}{4}{}
\kfQStem{윗글을 바탕으로 <보기>의 밑줄 친 부분의 발음에 적용된 경음화 규칙의 환경을 바르게 짝지은 것은?

<보기> ㄱ. 먹고[먹꼬]     ㄴ. 할 수[할쑤] ㄷ. 안다[안따]     ㄹ. 물질[물찔] ㅁ. 손가락[손까락]}
\begin{kfChoices}
\kfChoice ㄱ-①, ㄴ-③, ㄷ-②, ㄹ-④, ㅁ-②
\kfChoice ㄱ-①, ㄴ-③, ㄷ-②, ㄹ-④, ㅁ-①
\kfChoice ㄱ-①, ㄴ-②, ㄷ-②, ㄹ-④, ㅁ-②
\kfChoice ㄱ-①, ㄴ-③, ㄷ-②, ㄹ-④, ㅁ-⑤
\kfChoice ㄱ-①, ㄴ-③, ㄷ-②, ㄹ-④
\end{kfChoices}
\end{kfQBlock}

\kfPassageFull{문항 5 지문 (concept)}{%
「다음은 수업 시간에 선생님이 제시한 탐구 과제이다. 선생님은 학생들에게 네 가지 음운 변동 유형(교체, 탈락, 첨가, 축약)을 구분하고, 각 단어에서 일어나는 음운 변동의 유형과 횟수를 분석해 보라고 하였다. 이때 음운 변동의 결과를 정확히 파악하려면 변동 전후의 음운 개수를 비교해야 한다. 교체는 음운 개수가 유지되고, 탈락은 줄어들며, 첨가는 늘어나고, 축약도 줄어든다.」
}
\begin{kfQBlock}{5}{}
\kfQStem{<보기>의 각 단어에 대해 음운 변동의 유형별 횟수를 바르게 분석한 것은?

<보기> ㉠ 꽃힘[꼬침]   ㉡ 넓적하다[넙쩌카다] ㉢ 솜이불[솜니불]}
\begin{kfChoices}
\kfChoice ㉠에서는 축약 1회가 일어난다.
\kfChoice ㉡에서는 탈락 1회, 교체 2회가 일어난다.
\kfChoice ㉢에서는 첨가 1회만 일어난다.
\kfChoice ㉠에서는 교체 1회, 축약 1회가 일어난다.
\kfChoice ㉠\textasciitilde{}㉢에서 교체는 총 3회 일어난다.
\end{kfChoices}
\end{kfQBlock}

\end{multicols}


\kfStartLit{이번 챕터 문학 작품 5편}


\kfSecHeader{kfAreaLiterature}{지문}{춘향전 — 작자 미상 (판소리계 소설(고전소설)) / 「춘향전」 핵심 발췌(수청 강요\textasciitilde{}암행어사 출도)}

\kfPassageNote{춘향전 — 작자 미상 (판소리계 소설(고전소설)) / 「춘향전」 핵심 발췌(수청 강요\textasciitilde{}암행어사 출도)}{%
이때 변 사또 하는 행사 못 할 일이 없으니, 온갖 형벌을 다 쓰며 수청 들라 호령하되 춘향이 굽히지 아니하니라.

「네 이 년, 관장의 분부를 어기다니 당장 죽여 마땅하나, 고이 타이르노니 수청을 드니 아니 드니?」

춘향이 여쭈되,

「소녀 비록 천한 기생의 딸이오나 한 번 정한 남편을 저버리지 못하겠나이다. 일부종사가 여자의 도리어늘, 죽기로 맹세하고 절개를 지키오리다. 사또께서 백성의 어버이가 되시어 형벌로 수청을 시키시오니, 이것이 어찌 목민관의 도리이리까?」

변 사또 대로하여 형방에 분부하되,

「저년을 끌어다 곤장을 쳐라!」

형리들이 춘향을 끌어내어 형틀에 올리니, 춘향이 하늘을 우러러 통곡하되,

「충신은 불사이군이요 열녀는 불경이부라. 이 몸이 죽어도 임 향한 일편단심은 가실 줄이 없사옵나이다.」

이윽고 어사또 출도하는 날, 본관 수령 변학도가 생일잔치를 크게 벌이거늘, 어사또 거지 꼴로 잔치에 나아가 앉아 음식을 얻어먹으며 좌중을 살피더니, 이윽고 좌우를 호령하여 마패를 내어 놓으니,

「암행어사 출도야!」

역졸들이 우르르 달려들어 변 사또를 끌어내리니, 그 위세가 하늘을 찌를 듯하더라. 어사또 좌정하여 봉고파직을 선언하되,

「남원부사 변학도는 탐학무도하여 백성의 재물을 빼앗고, 무고한 여인에게 수청을 강요하여 형벌을 남용하였으니, 봉고파직에 처하노라.」

옥에 갇힌 춘향을 불러내어 석방하니, 춘향이 어사또를 우러러 보며 서로 부둥켜 안고 울음을 토하니, 보는 이 눈물 아니 흘리는 자 없더라.
}
\kfSecHeader{kfAreaLiterature}{활동}{작품 정리표 — 춘향전}

\noindent\textit{\small 빈칸에 알맞은 말을 채워 넣으시오. (초성 힌트 참고)}\par\medskip
\begin{kfActTable}{|>{\centering\arraybackslash}m{2.6cm}|m{6cm}|m{4cm}|}
\rowcolor{kfPrimaryLight!50}\textbf{\color{kfPrimary}항목} & \textbf{\color{kfPrimary}내용 (빈칸 채우기)} & \textbf{\color{kfPrimary}답안 작성}\\\hline
갈래 & ㅍㅅㄹ계 ㅅㅅ, 고전소설 ① & \kfTBlank \\\hline
핵심 갈등 & ㅂ 사또의 ㅅㅊ 강요 vs ㅊㅎ의 ㅈㄱ ② & \kfTBlank \\\hline
춘향의 논거 & 'ㅇㅂㅈㅅ가 여자의 도리' — 유교 윤리 ③ & \kfTBlank \\\hline
춘향의 비판 & 'ㅁㅁㄱ의 도리'에 어긋남 — 권력 ㄴㅇ 비판 ④ & \kfTBlank \\\hline
춘향의 선언 & 'ㅊㅅ은 ㅂㅅㅇㄱ, ㅇㄴ는 ㅂㄱㅇㅂ' ⑤ & \kfTBlank \\\hline
극적 반전 & ㄱㅈ 꼴 → ㅁㅍ 제시 → 'ㅇㅁㅇㅅ ㅊㄷ야!' ⑥ & \kfTBlank \\\hline
변 사또 처벌 & ㅂㄱㅍㅈ — 탐학무도 ⑦ & \kfTBlank \\\hline
춘향 결말 & ㅇ에서 ㅅㅂ ⑧ & \kfTBlank \\\hline
법적 장치 & ㅇㅅ ㅈㄷ — 부패 관리 감시·단죄 ⑨ & \kfTBlank \\\hline
주제 1 & ㄱㄹ ㄴㅇ에 대한 ㅂㅍ ⑩ & \kfTBlank \\\hline
주제 2 & ㅈㅇ ㅅㅎ과 ㅁㅈ ㅈㅇ ⑪ & \kfTBlank \\\hline
문학사적 의의 & ㅅㅂ ㅈㅅ ㅁㅅ과 ㅁㅈ ㅇㅅ이 반영된 대표작 ⑫ & \kfTBlank \\\hline
\end{kfActTable}
\kfSecHeader{kfAreaLiterature}{문제}{}

\begin{multicols}{2}\raggedcolumns\setlength{\columnsep}{12pt}
\begin{kfQBlock}{1}{}
\kfQStem{윗글에 대한 설명으로 적절하지 않은 것은?}
\begin{kfChoices}
\kfChoice 춘향은 기생의 딸이라는 신분에도 불구하고 절개를 지키겠다는 의지를 드러내고 있다.
\kfChoice 변 사또는 관장의 권력을 이용하여 춘향에게 수청을 강요하고 있다.
\kfChoice 춘향은 '일부종사'라는 유교적 여성 윤리를 근거로 수청을 거부하고 있다.
\kfChoice 암행어사는 변 사또의 탐학을 단죄하고 춘향을 석방하고 있다.
\kfChoice 춘향은 변 사또의 요구에 처음엔 거부하다가 결국 수청에 응하고 있다.
\end{kfChoices}
\end{kfQBlock}

\begin{kfQBlock}{2}{}
\kfQStem{윗글에서 춘향이 '이것이 어찌 목민관의 도리이리까?'라고 말한 것에 담긴 함축적 의미로 가장 적절한 것은?}
\begin{kfChoices}
\kfChoice 변 사또에게 더 높은 직위를 마땅히 주어야 한다는 건의의 표현
\kfChoice 목민관의 도리를 미처 모르는 자기 자신의 무지를 한탄하는 발화
\kfChoice 관리가 형벌로 수청을 강요함은 본분에 어긋난다는 비판
\kfChoice 관청에서 처리하는 행정 절차가 지나치게 복잡하다는 불만의 토로
\kfChoice 목민관에게는 형벌을 행사할 권한이 본래 없다는 법률적 주장
\end{kfChoices}
\end{kfQBlock}

\begin{kfQBlock}{3}{}
\kfQStem{윗글에서 춘향이 '충신은 불사이군이요 열녀는 불경이부라'라고 말한 것에 대한 해석으로 가장 적절한 것은?}
\begin{kfChoices}
\kfChoice 충신·열녀 윤리에 근거하여 수청 거부의 정당성을 주장한 것
\kfChoice 충신과 열녀의 차이점을 설명하여 자신의 학식을 과시하려 한 것
\kfChoice 변 사또가 충신이자 열녀라는 점을 인정하며 복종의 뜻을 표한 것
\kfChoice 나라에 충성하겠으니 수청 대신 관직을 달라고 요청한 발화이다
\kfChoice 충신과 열녀가 더 이상 존재하지 않는 세상을 안타깝게 한탄한 것
\end{kfChoices}
\end{kfQBlock}

\begin{kfQBlock}{4}{}
\kfQStem{윗글을 바탕으로 <보기>를 감상한 내용으로 적절하지 않은 것은?

<보기> 「춘향전」은 판소리계 소설로, 조선 후기 사회의 신분 질서 모순과 권력 남용에 대한 민중의 비판 의식이 반영되어 있다. 춘향의 수청 거부는 단순한 사랑의 절개를 넘어 부당한 권력에 대한 저항으로, 암행어사의 출도는 정의 실현의 기대를 형상화한 것이다.}
\begin{kfChoices}
\kfChoice 변 사또의 수청 강요는 <보기>의 '권력 남용'에 해당하며, 관리의 폭력적 권력 행사를 보여주겠군.
\kfChoice 춘향의 수청 거부는 단순한 사랑의 절개를 넘어 부당한 권력에 대한 저항이겠군.
\kfChoice 암행어사의 봉고파직 선언은 정의 실현의 기대를 형상화한 것이겠군.
\kfChoice 춘향이 '기생의 딸'임에도 절개를 주장하는 것은 신분 질서의 모순에 대한 도전이겠군.
\kfChoice 이 작품은 당대 사회의 신분 질서를 전면 긍정하고 관리의 권위를 찬양하는 것이겠군.
\end{kfChoices}
\end{kfQBlock}

\begin{kfQBlock}{5}{}
\kfQStem{윗글에서 어사또의 출도 장면에 사용된 극적 기법과 그 효과로 가장 적절한 것은?}
\begin{kfChoices}
\kfChoice 어사또가 미리 자신의 정체를 밝힌 후 점진적으로 권위를 보여주고 있다.
\kfChoice 거지 꼴로 잔치에 나타나는 반전으로 정의 실현의 극적 쾌감을 부각하고 있다.
\kfChoice 내면 독백을 통해 어사또의 심리적 갈등을 세밀하게 묘사하고 있다.
\kfChoice 시간의 역전을 통해 과거와 현재를 교차하며 서사적 긴장을 만들고 있다.
\kfChoice 대화 없이 지문만으로 상황을 전달하여 객관적 거리를 유지하고 있다.
\end{kfChoices}
\end{kfQBlock}

\end{multicols}

\kfSecHeader{kfAreaLiterature}{지문}{유충렬전 — 작자 미상 (고전소설(영웅소설·군담소설)) / 「유충렬전」 핵심 발췌(간신 처단\textasciitilde{}천자 구출)}

\kfPassageNote{유충렬전 — 작자 미상 (고전소설(영웅소설·군담소설)) / 「유충렬전」 핵심 발췌(간신 처단\textasciitilde{}천자 구출)}{%
이때 간신 정한담이 오랑캐와 내통하여 황성을 함락하고 천자를 가두니, 온 조정이 어찌할 줄 몰라 통곡하더라.

유충렬이 기병을 거느리고 황성으로 향하며 하늘을 우러러 맹세하되,

「소장이 비록 미약하오나, 선친의 원수를 갚고 천자를 구하여 충효를 다하리이다. 정한담의 역적 무리를 소탕하지 못하면 이 칼을 이 몸에 꽂으리라.」

대군이 황성에 이르러 적진을 향해 진격하니, 유충렬이 선봉에 서서 적장을 베고 성문을 열어 함락된 성을 되찾으니, 그 위세가 천지를 진동케 하더라.

정한담을 사로잡아 끌어내니, 유충렬이 꾸짖어 이르되,

「네 이놈 정한담아, 네가 조정의 녹을 먹으며 군부(君父)의 은혜를 입었거늘, 도리어 오랑캐와 결탁하여 종사(宗社)를 위태롭게 하고 충신의 집안을 멸하였으니, 그 죄가 만 번 죽어도 마땅하다. 천하에 간신의 끝이 이러한 줄 알라!」

천자를 옥좌에 모시고 유충렬이 엎드려 아뢰되,

「신의 선친 유심이 정한담의 모함으로 죽임을 당하였사오니, 원통한 충신의 넋을 위로하여 주시옵소서.」

천자 크게 감동하여 이르시되,

「유충렬의 충효가 천지를 감동시켰으니, 선친 유심의 충렬을 추증하고, 유충렬에게는 일등 공신의 직첩을 내리노라.」

이에 정한담 일당은 처형되고, 유심의 충렬이 복권되니, 충과 역의 시비가 밝혀진지라.
}
\kfSecHeader{kfAreaLiterature}{활동}{작품 정리표 — 유충렬전}

\noindent\textit{\small 빈칸에 알맞은 말을 채워 넣으시오. (초성 힌트 참고)}\par\medskip
\begin{kfActTable}{|>{\centering\arraybackslash}m{2.6cm}|m{6cm}|m{4cm}|}
\rowcolor{kfPrimaryLight!50}\textbf{\color{kfPrimary}항목} & \textbf{\color{kfPrimary}내용 (빈칸 채우기)} & \textbf{\color{kfPrimary}답안 작성}\\\hline
갈래 & ㅇㅇ ㅅㅅ(영웅소설·군담소설) ① & \kfTBlank \\\hline
주인공 & ㅇㅊㄹ — 충신 유심의 아들 ② & \kfTBlank \\\hline
적대자 & ㅈㅎㄷ — 간신, ㅇㅈ과 결탁 ③ & \kfTBlank \\\hline
핵심 갈등 & ㅊ(忠) vs ㅇ(逆)의 대립 ④ & \kfTBlank \\\hline
유충렬의 맹세 & 선친 ㅇㅅ + ㅊㅈ 구출 + ㅊㅎ 실현 ⑤ & \kfTBlank \\\hline
정한담의 죄상 & ㅇㄹㅋ 결탁, ㅈㅅ 위태, ㅊㅅ 멸문 ⑥ & \kfTBlank \\\hline
처벌 & 정한담 일당 ㅊㅎ ⑦ & \kfTBlank \\\hline
포상 & 유심 ㅊㅈ + 유충렬 ㅇㄷ ㄱㅅ ⑧ & \kfTBlank \\\hline
결말 유형 & 충·역 시비 밝혀진 ㄷㅎ ㄱㅁ ⑨ & \kfTBlank \\\hline
영웅소설 구조 & 고귀한 ㅎㅌ → ㅅㄹ → ㅈㄹ → ㅅㄹ → ㅍㅅ ⑩ & \kfTBlank \\\hline
주제 & ㅊㅇ ㄷㄹ의 ㅂㅈ·ㅇㄹㅈ 구조 ⑪ & \kfTBlank \\\hline
문학사적 의의 & 유교적 ㅈㅊ ㅇㄹ를 형상화한 ㄱㄷ ㅅㅅ ⑫ & \kfTBlank \\\hline
\end{kfActTable}
\kfSecHeader{kfAreaLiterature}{문제}{}

\begin{multicols}{2}\raggedcolumns\setlength{\columnsep}{12pt}
\begin{kfQBlock}{1}{}
\kfQStem{윗글에 대한 설명으로 적절하지 않은 것은?}
\begin{kfChoices}
\kfChoice 정한담은 오랑캐와 결탁하여 황성을 함락하고 천자를 가두었다.
\kfChoice 유충렬은 선친의 원수를 갚고 천자를 구하겠다고 맹세하고 있다.
\kfChoice 유충렬은 정한담을 사로잡아 그 죄를 꾸짖고 있다.
\kfChoice 천자는 유충렬의 공을 인정하여 관직을 수여하고 있다.
\kfChoice 정한담은 처벌을 면하고 유충렬과 화해하였다.
\end{kfChoices}
\end{kfQBlock}

\begin{kfQBlock}{2}{}
\kfQStem{윗글에서 유충렬이 정한담에게 '천하에 간신의 끝이 이러한 줄 알라!'라고 말한 것에 담긴 함축적 의미로 가장 적절한 것은?}
\begin{kfChoices}
\kfChoice 전쟁이 가져오는 무익함을 강조하여 평화의 회복을 호소하는 발화
\kfChoice 정한담을 향한 화자 개인의 복수심만을 강하게 드러낸 발화
\kfChoice 간신이라 해도 때로는 용서받을 수 있다는 관용의 마음을 표현
\kfChoice 간신의 배신과 역모는 반드시 응보를 받는다는 정의의 선언
\kfChoice 충신과 간신의 구분 자체가 무의미하다는 상대주의적 인식의 표명
\end{kfChoices}
\end{kfQBlock}

\begin{kfQBlock}{3}{}
\kfQStem{윗글을 바탕으로 <보기>를 감상한 내용으로 적절하지 않은 것은?

<보기> 「유충렬전」은 영웅소설의 전형적 구조(고귀한 혈통→시련→조력→승리→포상)를 따르면서, 충·역 대립이라는 유교적 정치 윤리를 핵심 주제로 삼는다. 간신의 처단과 충신의 복권은 선악의 법적·윤리적 구분이 최종적으로 확립되는 과정이다.}
\begin{kfChoices}
\kfChoice 유충렬의 고난과 승리는 영웅소설의 '시련→승리' 구조에 해당하겠군.
\kfChoice 이 작품은 충·역 대립의 결과가 모호하게 남겨져 열린 결말을 형성하겠군.
\kfChoice 천자가 일등 공신 직첩을 내리는 것은 영웅소설의 '포상' 단계이겠군.
\kfChoice 충·역 대립이 핵심 주제라는 <보기>의 관점에서, 유충렬의 맹세는 충의 실천 의지를 보여주겠군.
\kfChoice 정한담의 처형과 유심의 복권은 선악의 법적·윤리적 구분이 확립되는 과정이겠군.
\end{kfChoices}
\end{kfQBlock}

\begin{kfQBlock}{4}{}
\kfQStem{윗글에서 유충렬이 정한담의 죄를 꾸짖을 때 거론한 세 가지 핵심 죄상을 바르게 짝지은 것은?}
\begin{kfChoices}
\kfChoice 오랑캐와의 결탁, 종사를 위태롭게 함, 충신 집안의 멸문
\kfChoice 국가 세금의 횡령, 민간인의 학대, 대외 외교의 거듭된 실패
\kfChoice 오랑캐와의 결탁, 천자의 강제 퇴위, 충신에 대한 매수
\kfChoice 왕위 찬탈 음모, 백성 학살, 종교 탄압의 자행
\kfChoice 군사 반란 시도, 세금 인상 강행, 충신과의 화해 거부
\end{kfChoices}
\end{kfQBlock}

\begin{kfQBlock}{5}{}
\kfQStem{윗글에서 '소장이 비록 미약하오나'에 사용된 표현 기법과 그 효과로 가장 적절한 것은?}
\begin{kfChoices}
\kfChoice 겸양의 표현으로 영웅의 결의를 역설적으로 부각하고 있다.
\kfChoice 과장법을 활용하여 유충렬의 무력과 위세를 극도로 극대화하여 보여 준다.
\kfChoice 반어법을 통해 유충렬이 내면에 품은 자만심을 우회적으로 드러내고 있다.
\kfChoice 의인법을 활용하여 유충렬의 칼에 살아 있는 인격적 속성을 부여하고 있다.
\kfChoice 직유법을 통해 유충렬을 위엄 있는 자연물의 모습에 비유하여 표현하고 있다.
\end{kfChoices}
\end{kfQBlock}

\end{multicols}

\kfSecHeader{kfAreaLiterature}{지문}{이춘풍전 / 이대봉전 — 작자 미상 (고전소설(복합지문)) / (가) 「이춘풍전」 핵심 발췌 / (나) 「이대봉전」 핵심 발췌}

\kfPassageNote{이춘풍전 / 이대봉전 — 작자 미상 (고전소설(복합지문)) / (가) 「이춘풍전」 핵심 발췌 / (나) 「이대봉전」 핵심 발췌}{%
(가) 이춘풍의 처 춘풍 아내가 남편이 평양 기생 추월에게 빠져 재산을 탕진하였다는 소식을 듣고, 스스로 남장을 하고 평양감사의 비장이 되어 평양으로 향하니라.

비장으로 위장한 춘풍 처가 추월의 집에 이르러 추월을 불러 앉히고 이르되,

「네 이년, 양반의 재물을 빼앗아 놓고 그 남편을 종으로 부리고 있으니, 관가에서 다스리리라.」

추월이 놀라 엎드려 빌거늘, 춘풍 처가 빼앗긴 재물을 모두 도로 찾고, 이춘풍을 불러내어 이르되,

「여보시오 서방님, 기생에게 재산 다 빼앗기고 종노릇 하는 꼴이 양반의 체면이오? 이제 집에 돌아가 다시는 이런 짓 마시오.」

이춘풍이 머리를 조아리며 부끄러워하니, 보는 이들이 모두 웃더라.

(나) 이대봉이 남장을 하고 과거에 응시하여 장원 급제하니, 천자가 크게 기뻐하여 높은 관직을 내렸더라. 때마침 오랑캐가 쳐들어오거늘, 이대봉이 원수(元帥)가 되어 출전하니라.

이대봉이 대군을 이끌고 적진에 나아가 진을 치고 적장을 향해 이르되,

「네 이놈, 대국의 위엄을 모르고 감히 침범하였으니, 이 칼에 목을 놓을지어다.」

단번에 적장의 머리를 베니, 적군이 크게 패하여 물러가더라. 천자가 이대봉의 공을 높이 사 일등 공신에 봉하고 재상의 직을 내리니, 조정 대신이 모두 칭찬하되, 이대봉이 여자인 줄은 아무도 모르더라.

후에 이대봉의 여자임이 밝혀지매, 천자 크게 놀라 이르시되,

「여자의 몸으로 이토록 큰 공을 세우니, 참으로 기이한 일이로다.」
}
\kfSecHeader{kfAreaLiterature}{활동}{작품 정리표 — 이춘풍전 / 이대봉전}

\noindent\textit{\small 빈칸에 알맞은 말을 채워 넣으시오. (초성 힌트 참고)}\par\medskip
\begin{kfActTable}{|>{\centering\arraybackslash}m{2.6cm}|m{6cm}|m{4cm}|}
\rowcolor{kfPrimaryLight!50}\textbf{\color{kfPrimary}항목} & \textbf{\color{kfPrimary}내용 (빈칸 채우기)} & \textbf{\color{kfPrimary}답안 작성}\\\hline
공통 기법 & 여성의 ㄴㅈ을 통한 ㅇㅎ ㅈㅂ ① & \kfTBlank \\\hline
(가) 인물 & ㅇㅊㅍ(무능 남편) vs ㅊㅍ ㅊ(유능 아내) ② & \kfTBlank \\\hline
(가) 남장 역할 & 평양감사의 ㅂㅈ(裨將) ③ & \kfTBlank \\\hline
(가) 해결 범위 & ㄱㅈ 문제 — ㅈㅅ 회수·남편 ㅈㅊ ④ & \kfTBlank \\\hline
(가) 서사 방식 & ㅍㅈ — 남성 무능을 비웃음 ⑤ & \kfTBlank \\\hline
(나) 인물 & ㅇㄷㅂ — ㄴㅈ ㅇㅅ 영웅 ⑥ & \kfTBlank \\\hline
(나) 남장 역할 & ㅈㅇ ㄱㅈ + ㅇㅅ(元帥) ⑦ & \kfTBlank \\\hline
(나) 해결 범위 & ㄱㄱ적 문제 — ㅇㅈ 퇴치 ⑧ & \kfTBlank \\\hline
(나) 서사 방식 & ㅇㅇ ㅅㅅ — 여성 능력 긍정 ⑨ & \kfTBlank \\\hline
공통 주제 & ㅅㅂ ㅈㅅ에 대한 문학적 ㅈㅂ ⑩ & \kfTBlank \\\hline
차이 — (가) & ㅍㅈ적 접근, ㄱㅈ 차원 ⑪ & \kfTBlank \\\hline
차이 — (나) & ㅇㅇ적 접근, ㄱㄱ 차원 ⑫ & \kfTBlank \\\hline
\end{kfActTable}
\kfSecHeader{kfAreaLiterature}{문제}{}

\begin{multicols}{2}\raggedcolumns\setlength{\columnsep}{12pt}
\begin{kfQBlock}{1}{}
\kfQStem{(가)와 (나)에 대한 설명으로 적절하지 않은 것은?}
\begin{kfChoices}
\kfChoice (가)에서 춘풍 처는 남장을 하고 비장이 되어 재산을 되찾고 있다.
\kfChoice (나)에서 이대봉은 남장으로 과거에 급제하고 전쟁에서 공을 세우고 있다.
\kfChoice (가)와 (나) 모두 여성 인물이 남장을 통해 남성의 사회적 역할을 수행하고 있다.
\kfChoice (가)의 이춘풍은 유능한 남편으로서 가정을 이끌어가고 있다.
\kfChoice (나)에서 이대봉이 여자임이 밝혀지자 천자가 놀라고 있다.
\end{kfChoices}
\end{kfQBlock}

\begin{kfQBlock}{2}{}
\kfQStem{(가)에서 춘풍 처가 '양반의 체면이오?'라고 말한 것에 담긴 함축적 의미로 가장 적절한 것은?}
\begin{kfChoices}
\kfChoice 양반 제도 자체를 즉시 폐지해야 한다는 혁명적이고 급진적 주장
\kfChoice 양반의 체면을 지켜 준 기생 추월에 대한 깊은 감사의 표현
\kfChoice 양반 신분을 유지하려면 더 많은 재산이 반드시 필요하다는 주장
\kfChoice 기생과의 관계를 공식적으로 인정받게 해 달라는 사적인 요청
\kfChoice 재산을 빼앗기고 종노릇 하는 양반의 추태에 대한 풍자적 비판
\end{kfChoices}
\end{kfQBlock}

\begin{kfQBlock}{3}{}
\kfQStem{(가)와 (나)를 바탕으로 <보기>를 감상한 내용으로 적절하지 않은 것은?

<보기> 조선 후기 소설에서 여성이 남장을 하고 남성의 사회적 역할을 수행하는 것은 기존의 성별 질서에 대한 문학적 전복이다. (가)는 풍자를 통해 남성의 무능을 비웃고, (나)는 영웅 서사를 통해 여성의 능력을 긍정한다. 두 작품 모두 역할 전복을 통해 기존 질서의 모순을 드러낸다.}
\begin{kfChoices}
\kfChoice (가)에서 춘풍 처가 비장이 되어 문제를 해결하는 것은 성별 질서의 전복이겠군.
\kfChoice (가)의 풍자는 이춘풍의 무능을 비웃음으로써 남성 중심 질서의 모순을 드러내겠군.
\kfChoice (나)에서 이대봉이 장원 급제하고 전쟁에서 공을 세우는 것은 여성의 능력을 긍정하는 것이겠군.
\kfChoice (가)와 (나) 모두 기존 성별 질서를 강화하고 남성의 우월성을 확인하는 것이겠군.
\kfChoice 두 작품에서 여성의 남장은 기존 역할 구분에 대한 문학적 도전이겠군.
\end{kfChoices}
\end{kfQBlock}

\begin{kfQBlock}{4}{}
\kfQStem{(가)와 (나)에서 남장 여성이 수행하는 역할의 차이로 가장 적절한 것은?}
\begin{kfChoices}
\kfChoice (가) 춘풍 처는 가정 차원, (나) 이대봉은 국가 차원의 문제를 해결한다.
\kfChoice (가) 춘풍 처와 (나) 이대봉이 모두 전쟁에 출정하여 적장을 베고 공을 세우는 활약을 한다.
\kfChoice (가) 춘풍 처는 남장이 도중에 밝혀져 본모습이 드러나고, (나) 이대봉은 끝까지 남장 신분을 유지한다.
\kfChoice (가) 춘풍 처는 결국 왕비의 자리에 오르고, (나) 이대봉은 본래의 자리로 내려와 기생이 되어 살아간다.
\kfChoice (가)와 (나) 모두 남장한 여성 주인공이 임무 수행에 실패한 뒤 남성 인물의 구원을 받는다.
\end{kfChoices}
\end{kfQBlock}

\begin{kfQBlock}{5}{}
\kfQStem{(가)에서 '보는 이들이 모두 웃더라'라는 서술이 지닌 서사적 기능으로 가장 적절한 것은?}
\begin{kfChoices}
\kfChoice 춘풍 처의 남장 작전이 결국 들통나 무위로 돌아갔음을 구경꾼의 웃음을 통해 확인시키는 기능을 한다.
\kfChoice 이춘풍의 부끄러워하는 모습을 긍정적으로 평가하여 독자에게 회개의 교훈을 전달하는 효과를 거두고 있다.
\kfChoice 작중 구경꾼의 웃음을 통해 독자에게도 풍자적 웃음을 유도하며 이춘풍의 행위가 조롱의 대상임을 확인시킨다.
\kfChoice 기생 추월이 양반 부부를 굴복시킨 승리의 순간을 구경꾼의 환호와 웃음으로 축하하는 장면이 된다.
\kfChoice 양반의 권위가 무너지는 비극적 정황을 구경꾼의 웃음과 대비해 더욱 무겁게 부각하는 장치로 기능한다.
\end{kfChoices}
\end{kfQBlock}

\end{multicols}

\kfSecHeader{kfAreaLiterature}{지문}{저당 잡힌 사내 — 이동하 (현대소설) / 이동하, 「저당 잡힌 사내」 핵심 발췌}

\kfPassageNote{저당 잡힌 사내 — 이동하 (현대소설) / 이동하, 「저당 잡힌 사내」 핵심 발췌}{%
그는 돈을 빌렸다. 아니, 빌릴 수밖에 없었다. 아내의 수술비, 아이의 학비, 밀린 월세. 삶이 그를 채무자로 만들었다.

처음에 그 사람은 친절했다. 「언제든 갚으시오. 급할 것 없소.」 그러나 이자가 이자를 낳고, 원금은 줄지 않고, 어느새 빚은 눈덩이처럼 불어났다. 그러자 그 사람의 얼굴이 달라졌다.

「이봐, 자네 그렇게 살 거야? 내가 자네한테 돈을 빌려준 것은 사업하라고 빌려준 거야. 나한테 갚을 생각이 있기나 한 거야?」

그 사람은 이제 그의 집에 마음대로 들어왔다. 냉장고를 열어 보고, 아내가 입는 옷을 훑어보고, 아이들의 장난감까지 손으로 만져 보며 고개를 저었다.

「이런 것들은 다 팔아야지. 빚쟁이 주제에 이런 사치를 하고 있으니.」

그는 아무 말도 할 수 없었다. 항의하면 「그래, 당장 갚아」 할 것이고, 법적으로 그 사람이 옳았다. 빚을 진 쪽이 약자였다. 아니, 약자가 아니라 아예 사람이 아니었다. 그는 어느새 사람이 아니라 저당물이 되어 있었다.

밤이면 그는 생각했다. 내가 빌린 것은 돈이었는데, 내가 잃어버린 것은 무엇인가. 돈이 아니었다. 자존심이었다. 자유였다. 인간으로서의 존엄이었다.

아내가 말했다. 「여보, 우리 사람답게 살 수는 없는 건가요?」

그는 대답하지 못했다. 사람답게 산다는 것, 그것이 무엇인지 이미 잊어버렸기 때문이었다.
}
\kfSecHeader{kfAreaLiterature}{활동}{작품 정리표 — 저당 잡힌 사내(이동하)}

\noindent\textit{\small 빈칸에 알맞은 말을 채워 넣으시오. (초성 힌트 참고)}\par\medskip
\begin{kfActTable}{|>{\centering\arraybackslash}m{2.6cm}|m{6cm}|m{4cm}|}
\rowcolor{kfPrimaryLight!50}\textbf{\color{kfPrimary}항목} & \textbf{\color{kfPrimary}내용 (빈칸 채우기)} & \textbf{\color{kfPrimary}답안 작성}\\\hline
갈래 & ㅎㄷ ㅅㅅ(단편) ① & \kfTBlank \\\hline
작자 & ㅇㄷㅎ ② & \kfTBlank \\\hline
제목의 비유 & 'ㅈㄷ'의 법적 개념 → 인간 존엄 상실의 ㅇㄹㄱㄹ ③ & \kfTBlank \\\hline
빚의 원인 & 아내 ㅅㅅ비, 아이 ㅎㅂ, 밀린 ㅇㅅ ④ & \kfTBlank \\\hline
채권자 변화 & ㅊㅈ → ㅇㅂ → ㅊㅇ적 지배 ⑤ & \kfTBlank \\\hline
주인공의 상실 & ㅈㅊㅅ + ㅈㅇ + ㅇㄱ으로서의 ㅈㅇ ⑥ & \kfTBlank \\\hline
법과 정의 괴리 & '법적으로 그 사람이 ㅇㅇ' — 법적 ㅈㄷㅅ vs 인간 ㅈㅇ ⑦ & \kfTBlank \\\hline
핵심 비유 & '사람이 아니라 ㅈㄷㅁ이 되어 있었다' ⑧ & \kfTBlank \\\hline
아내의 질문 & 'ㅅㄹ답게 살 수는 없는 건가요?' ⑨ & \kfTBlank \\\hline
주제 & ㄱㅈㅈ ㄱㄱ가 인간 존엄을 ㅇㄷ ⑩ & \kfTBlank \\\hline
사회적 배경 & ㅈㅂㅈㅇ ㅅㅎ의 ㄱㅈㅈ 폭력 ⑪ & \kfTBlank \\\hline
문학적 기법 & ㅊㄱ-ㅊㅁ 관계의 ㅇㄹㄱㄹ ⑫ & \kfTBlank \\\hline
\end{kfActTable}
\kfSecHeader{kfAreaLiterature}{문제}{}

\begin{multicols}{2}\raggedcolumns\setlength{\columnsep}{12pt}
\begin{kfQBlock}{1}{}
\kfQStem{윗글에 대한 설명으로 적절하지 않은 것은?}
\begin{kfChoices}
\kfChoice 주인공은 아내의 수술비, 아이의 학비 등 생활의 필요로 빚을 지게 되었다.
\kfChoice 채권자는 처음에는 친절했으나 빚이 불어나자 태도가 변하였다.
\kfChoice 채권자는 주인공의 집에 마음대로 들어와 물건을 살피며 압박하고 있다.
\kfChoice 주인공은 법적으로 채권자에게 항의할 수 있는 강한 입장에 있다.
\kfChoice 주인공은 빚으로 인해 자존심, 자유, 존엄을 잃어버렸다고 느끼고 있다.
\end{kfChoices}
\end{kfQBlock}

\begin{kfQBlock}{2}{}
\kfQStem{윗글에서 '그는 어느새 사람이 아니라 저당물이 되어 있었다'에 담긴 함축적 의미로 가장 적절한 것은?}
\begin{kfChoices}
\kfChoice 채무 관계가 인간 존엄을 침식해 사람이 물건처럼 전락했다는 비유
\kfChoice 주인공이 실제로 부동산을 담보로 설정하여 비로소 안전을 확보했다는 것
\kfChoice 저당권이라는 근대 법제도의 정교함과 우수성을 우회적으로 칭찬하는 표현
\kfChoice 주인공이 누구의 강요도 받지 않고 채권자에게 스스로를 바친 자발적 행위
\kfChoice 채권자와 채무자가 결국 평등하고 대등한 관계임을 새삼 확인해 주는 표현
\end{kfChoices}
\end{kfQBlock}

\begin{kfQBlock}{3}{}
\kfQStem{윗글을 바탕으로 <보기>를 감상한 내용으로 적절하지 않은 것은?

<보기> 이동하의 소설에서 채권-채무 관계는 단순한 경제적 거래가 아니라, 인간의 존엄과 자유를 침식하는 구조적 폭력의 알레고리이다. 법적으로 채권자의 권리 행사는 정당하지만, 그것이 인간의 존엄을 파괴하는 데까지 이르렀을 때, 법과 정의 사이의 괴리가 드러난다.}
\begin{kfChoices}
\kfChoice 채권자가 주인공의 집에 마음대로 들어와 물건을 살피는 것은 구조적 폭력의 양상이겠군.
\kfChoice '법적으로 그 사람이 옳았다'는 서술은 법적 정당성과 인간 존엄 사이의 괴리를 보여주겠군.
\kfChoice 이 작품은 채권자의 권리 행사가 언제나 정의롭다는 것을 증명하는 작품이겠군.
\kfChoice 아내의 '사람답게 살 수는 없는 건가요?'라는 질문은 경제적 관계가 인간의 존엄을 압도하는 현실에 대한 문제 제기이겠군.
\kfChoice '저당물이 되어 있었다'는 표현은 채무 관계가 인간 존엄을 침식한 결과의 알레고리이겠군.
\end{kfChoices}
\end{kfQBlock}

\begin{kfQBlock}{4}{}
\kfQStem{윗글에서 주인공이 '내가 빌린 것은 돈이었는데, 내가 잃어버린 것은 무엇인가'라고 성찰한 내용으로 적절하지 않은 것은?}
\begin{kfChoices}
\kfChoice 빌린 것은 돈이지만, 잃어버린 것은 자존심이다.
\kfChoice 빌린 것은 돈이지만, 잃어버린 것은 자유이다.
\kfChoice 빌린 것은 돈이지만, 잃어버린 것은 인간으로서의 존엄이다.
\kfChoice 빌린 것은 돈이지만, 잃어버린 것은 채권자와의 우정이다.
\kfChoice 돈이라는 물질적 대상이 비물질적 가치(자존심·자유·존엄)를 잠식한 역설이다.
\end{kfChoices}
\end{kfQBlock}

\begin{kfQBlock}{5}{}
\kfQStem{윗글에서 채권자의 태도 변화 과정을 서술한 방식의 특징으로 가장 적절한 것은?}
\begin{kfChoices}
\kfChoice 채권자의 변화하는 내면 심리를 시점 이동과 독백을 통해 세밀하게 분석하여, 독자가 그에게 공감하도록 유도하는 서술이 이어지고 있다.
\kfChoice 채권자의 폭력적 태도가 처음부터 끝까지 일관되게 유지되어, 인물의 평면적 성격과 단조로운 행동 양상을 명확히 부각해 보여주고 있다.
\kfChoice 처음의 친절한 말투에서 점차 명령적이고 침입적인 태도로 변화하는 과정을 대비하여, 채무 관계의 본질적 불평등을 점층적으로 드러내고 있다.
\kfChoice 채권자와 채무자의 관계가 처음부터 끝까지 대등한 거래 관계로 유지되어, 두 인물 사이에 어떤 권력의 위계도 발생하지 않는 모습을 보여주고 있다.
\kfChoice 현재의 위협적 태도에서 과거의 친절했던 시기로 거슬러 올라가는 시간의 역순 구성을 통해 서사적 긴장을 인위적으로 고조시키고 있다.
\end{kfChoices}
\end{kfQBlock}

\end{multicols}

\kfSecHeader{kfAreaLiterature}{지문}{골목 안 — 박태원 (현대소설(세태소설·모더니즘)) / 박태원, 「골목 안」 핵심 발췌}

\kfPassageNote{골목 안 — 박태원 (현대소설(세태소설·모더니즘)) / 박태원, 「골목 안」 핵심 발췌}{%
이 골목에는 세 채의 집이 있다.

왼편 첫째 집에 이발사 박 씨가 살고 있다. 아내는 폐병이 있다. 병원에 가야 하지만 돈이 없다. 이발 수입으로는 약값도 대기 어렵다. 아내는 아침마다 기침을 한다. 마른기침. 그 소리가 골목 끝까지 들린다. 아이들은 밥도 제대로 먹지 못한 얼굴로 골목에서 놀고 있다.

가운데 집에 만년 필경사 김 씨가 살고 있다. 석 달째 월세를 밀렸다. 집주인이 찾아와 문을 두드린다. 김 씨는 안에서 숨을 죽이고 있다. 집주인이 돌아가면 김 씨는 창밖을 내다보며 한숨을 쉰다. 아무 일도 일어나지 않는다.

오른편 끝 집에 최 씨 노인이 혼자 살고 있다. 아들이 만주로 떠난 뒤로 소식이 없다. 노인은 매일 골목 어귀에 앉아 기다린다. 누가 들어오나 고개를 내민다. 아무도 오지 않는다. 해가 지면 노인은 빈 방으로 돌아간다.

골목 위로 전깃줄이 지나간다. 전기는 이 골목 안의 어떤 집에도 들어오지 않는다. 밤이면 세 채의 집에 호롱불이 켜진다. 불빛이 약하여 얼굴도 제대로 보이지 않는다.

봄비가 내린다. 골목에 물이 고인다. 아무도 치우지 않는다.
}
\kfSecHeader{kfAreaLiterature}{활동}{작품 정리표 — 골목 안(박태원)}

\noindent\textit{\small 빈칸에 알맞은 말을 채워 넣으시오. (초성 힌트 참고)}\par\medskip
\begin{kfActTable}{|>{\centering\arraybackslash}m{2.6cm}|m{6cm}|m{4cm}|}
\rowcolor{kfPrimaryLight!50}\textbf{\color{kfPrimary}항목} & \textbf{\color{kfPrimary}내용 (빈칸 채우기)} & \textbf{\color{kfPrimary}답안 작성}\\\hline
갈래 & ㅎㄷ ㅅㅅ(ㅅㅌ소설·ㅁㄷㄴㅈ) ① & \kfTBlank \\\hline
작자 & ㅂㅌㅇ(1909\textasciitilde{}1986) ② & \kfTBlank \\\hline
배경 & 1930년대 ㄱㅅ(서울)의 ㄱㅁ ③ & \kfTBlank \\\hline
서술 기법 & ㅋㅁㄹ ㅇㅇ — ㄱㅊㅈ 시점, ㅇㅁㅈ 행위 기록 ④ & \kfTBlank \\\hline
인물 1 & 이발사 ㅂ 씨 — 아내 ㅍㅂ, ㅂㅇ비 부담 ⑤ & \kfTBlank \\\hline
인물 2 & 필경사 ㄱ 씨 — ㅇㅅ 밀려 집주인 피함 ⑥ & \kfTBlank \\\hline
인물 3 & ㅊ 씨 노인 — 아들 ㅁㅈ행, 혼자 기다림 ⑦ & \kfTBlank \\\hline
전깃줄 상징 & ㅈㄱㅈ 지나가지만 집에 안 들어옴 → ㅅㅇ ⑧ & \kfTBlank \\\hline
마지막 문장 & 'ㅇㅁㄷ ㅊㅇㅈ ㅇㄴㄷ' → ㅁㄱㅅ과 ㅂㅊ ⑨ & \kfTBlank \\\hline
주제 & ㄷㅅ ㅎㅊㅁ의 ㄱㅍ + ㅈㄷ ㅅㄱㅈㄷ ⑩ & \kfTBlank \\\hline
문체 특징 & ㅉ은 ㄷㅁ + ㅅㄹ ㅁㅅ ㅊㅅ + ㄱㅈ적 거리 ⑪ & \kfTBlank \\\hline
문학사적 의의 & ㅁㄷㄴㅈ ㄱㅂ의 한국 ㅅㅌ 소설 ⑫ & \kfTBlank \\\hline
\end{kfActTable}
\kfSecHeader{kfAreaLiterature}{문제}{}

\begin{multicols}{2}\raggedcolumns\setlength{\columnsep}{12pt}
\begin{kfQBlock}{1}{}
\kfQStem{윗글에 대한 설명으로 적절하지 않은 것은?}
\begin{kfChoices}
\kfChoice 골목 안 세 채의 집에 사는 인물들의 일상이 담담하게 서술되고 있다.
\kfChoice 이발사 박 씨는 아내의 병원비를 감당하지 못하는 경제적 궁핍에 처해 있다.
\kfChoice 필경사 김 씨는 밀린 월세 때문에 집주인을 피해 숨어 있다.
\kfChoice 최 씨 노인은 만주로 떠난 아들의 소식을 기다리고 있다.
\kfChoice 골목 안 주민들은 서로 적극적으로 돕고 연대하여 어려움을 극복하고 있다.
\end{kfChoices}
\end{kfQBlock}

\begin{kfQBlock}{2}{}
\kfQStem{윗글에서 '전기는 이 골목 안의 어떤 집에도 들어오지 않는다'라는 서술이 함축하는 의미로 가장 적절한 것은?}
\begin{kfChoices}
\kfChoice 근대 문명의 혜택이 이 골목의 하층민에게는 미치지 못한다는 소외의 상징
\kfChoice 골목 주민들이 자발적으로 전기를 거부하여 전통적 삶의 양식을 고수한다는 것
\kfChoice 당시 전기 기술이 아직 발명되지 않았음을 보여 주는 시대적 배경 설명
\kfChoice 골목이 산비탈에 있어 전선 설치 자체가 물리적으로 불가능하다는 사실
\kfChoice 비싼 전기 요금을 아끼려는 주민들의 절약 지향적 경제 판단
\end{kfChoices}
\end{kfQBlock}

\begin{kfQBlock}{3}{}
\kfQStem{윗글을 바탕으로 <보기>를 감상한 내용으로 적절하지 않은 것은?

<보기> 박태원의 소설은 모더니즘 기법, 특히 '카메라 아이(camera eye)' 기법으로 알려져 있다. 작가는 인물의 심리를 분석하거나 가치 판단을 내리지 않고, 관찰자적 시점에서 인물들의 외면적 행위와 환경을 기록한다. 이를 통해 독자 스스로 현실의 의미를 인식하게 한다.}
\begin{kfChoices}
\kfChoice '아내는 아침마다 기침을 한다. 마른기침.'이라는 서술은 관찰자적 시점에서 외면적 행위를 기록한 것이겠군.
\kfChoice '아무 일도 일어나지 않는다'라는 서술은 작가의 가치 판단 없이 상황을 기록한 것이겠군.
\kfChoice 세 가구의 일상을 병렬적으로 나열하는 것은 카메라가 공간을 비추듯 인물들을 기록하는 기법이겠군.
\kfChoice '봄비가 내린다. 골목에 물이 고인다.'라는 서술은 환경을 객관적으로 기록한 것이겠군.
\kfChoice 작가는 인물들의 비참한 삶에 대해 직접적으로 분노를 표출하며 사회 변혁을 촉구하고 있겠군.
\end{kfChoices}
\end{kfQBlock}

\begin{kfQBlock}{4}{}
\kfQStem{윗글의 마지막 문장 '아무도 치우지 않는다'가 상징하는 바로 가장 적절한 것은?}
\begin{kfChoices}
\kfChoice 하층민의 고통을 외면하는 제도와 사회의 무관심과 방치
\kfChoice 골목 주민들이 청소를 싫어하여 의도적으로 방치하고 있는 게으른 태도
\kfChoice 봄비의 풍경이 아름다워 흘러내리는 물조차 치우고 싶지 않다는 정서
\kfChoice 골목에 배수 시설이 잘 갖추어져 있어 별달리 걱정할 필요가 없다는 사실
\kfChoice 주민들이 한마음으로 단결하여 공동으로 물을 치울 계획을 세우고 있다는 것
\end{kfChoices}
\end{kfQBlock}

\begin{kfQBlock}{5}{}
\kfQStem{윗글의 서술 방식에 대한 설명으로 가장 적절한 것은?}
\begin{kfChoices}
\kfChoice 짧은 단문 연속으로 심리를 절제하며 객관적 거리를 유지한다.
\kfChoice 긴 복합문을 활용하여 인물의 내면 심리를 한 문단에 걸쳐 세밀하게 분석하고 있다.
\kfChoice 1인칭 서술자가 자기 체험을 독자에게 고백하는 회고적 형식으로 작품을 전개하고 있다.
\kfChoice 과거와 현재의 장면을 교차로 배치하는 시간 역전 기법을 일관되게 사용하고 있다.
\kfChoice 풍부하고 장식적인 비유와 상징을 동원하여 화려하고 미려한 문체를 구사하고 있다.
\end{kfChoices}
\end{kfQBlock}

\end{multicols}

\kfStartNonlit{이번 챕터 비문학 지문 5편}


\kfSecHeader{kfAreaNonlit}{지문}{법률행위의 무효와 취소 (법학(복합지문))}

\kfPassageNote{법률행위의 무효와 취소 (법학(복합지문))}{%
(가) 법률행위의 무효(無效)란 법률행위가 성립 요건은 갖추었으나 효력 발생 요건에 흠이 있어 처음부터 법적 효력이 발생하지 않는 것을 말한다. 무효인 법률행위는 당사자가 별도의 의사 표시를 하지 않아도 처음부터 자동적으로 효력이 없으며, 누구나 무효를 주장할 수 있다. 무효의 원인은 다양하다. 의사 표시의 핵심 내용이 확정되지 않은 경우, 원시적으로 이행이 불가능한 급부를 내용으로 하는 경우, 선량한 풍속 기타 사회질서에 위반하는 내용인 경우(민법 제103조), 불공정한 법률행위인 경우(민법 제104조), 비진의 의사 표시에서 상대방이 표의자의 진의 아님을 알았거나 알 수 있었을 경우(민법 제107조 제1항 단서), 통정 허위 표시인 경우(민법 제108조) 등이 있다.

무효인 법률행위의 효과는 원칙적으로 처음부터 효력이 없으므로, 이미 이행된 것이 있다면 부당이득으로 반환해야 한다. 그러나 거래 안전을 위해 일정한 제한이 존재한다. 통정 허위 표시의 무효는 '선의의 제3자'에게 대항할 수 없다(민법 제108조 제2항). 여기서 '선의'란 그 법률행위가 허위라는 사실을 모른다는 뜻이며, '제3자'란 당사자 및 그 포괄 승계인 이외의 자로서 허위 표시에 기초하여 새로운 법률상 이해관계를 맺은 자를 말한다. 이 규정은 외관을 신뢰한 제3자를 보호하여 거래의 안전을 도모하기 위한 것이다.

무효인 법률행위는 원칙적으로 추인(追認)에 의해 유효로 전환될 수 없다. 무효 행위를 추인하더라도 그것은 새로운 법률행위를 한 것으로 볼 뿐이며, 처음부터 유효였던 것으로 되지는 않는다. 다만 무효 행위의 전환(轉換)이라 하여, 무효인 법률행위가 다른 법률행위의 요건을 갖추고 당사자가 무효임을 알았다면 다른 법률행위를 하였을 것이라고 인정될 때에는 그 다른 법률행위로서 효력을 가질 수 있다(민법 제138조).

(나) 법률행위의 취소(取消)란 일단 유효하게 성립한 법률행위를 일정한 사유가 있을 때 취소권자의 의사 표시에 의해 소급(遡及)하여 효력을 소멸시키는 것을 말한다. 취소 사유에는 제한능력자(미성년자, 피성년후견인 등)의 행위, 착오에 의한 의사 표시(민법 제109조), 사기·강박에 의한 의사 표시(민법 제110조) 등이 있다.

착오에 의한 취소에서 '착오'란 의사와 표시가 일치하지 않는 것을 표의자 스스로 모르는 경우를 말한다. 착오가 '법률행위의 내용의 중요 부분'에 해당해야 취소할 수 있으며, 표의자에게 중대한 과실이 있을 때에는 취소하지 못한다. 사기에 의한 취소에서 제3자의 사기인 경우에는 상대방이 그 사실을 알았거나 알 수 있었을 때에만 취소할 수 있다. 사기·강박에 의한 취소는 선의의 제3자에게 대항할 수 없다(민법 제110조 제3항).

취소의 효과는 소급 무효이다. 즉, 취소된 법률행위는 처음부터 무효였던 것으로 된다(민법 제141조). 이미 이행된 것이 있다면 원상회복 의무가 발생하며, 제한능력자의 경우 현존 이익의 한도에서만 반환 의무를 부담한다. 취소와 무효의 결정적 차이는, 취소는 취소 전에는 유효한 법률행위로 존재하며, 취소권자만이 취소를 주장할 수 있다는 점이다.

취소할 수 있는 법률행위는 추인에 의해 확정적으로 유효한 것으로 될 수 있다(민법 제143조). 추인은 취소의 원인이 소멸한 후에 해야 효력이 있으며, 추인하면 다시 취소하지 못한다. 또한 취소권은 추인할 수 있는 날로부터 3년 내, 법률행위를 한 날로부터 10년 내에 행사해야 한다(민법 제146조). 이 기간이 지나면 취소권이 소멸하여 법률행위가 확정적으로 유효하게 된다.
}
\kfSecHeader{kfAreaNonlit}{문제}{}

\begin{multicols}{2}\raggedcolumns\setlength{\columnsep}{12pt}
\begin{kfQBlock}{1}{}
\kfQStem{(가)와 (나)의 내용과 일치하지 않는 것은?}
\begin{kfChoices}
\kfChoice 무효인 법률행위는 처음부터 자동적으로 효력이 없으며, 누구나 무효를 주장할 수 있다.
\kfChoice 취소할 수 있는 법률행위는 취소 전에는 유효한 것으로 존재하며, 취소권자만이 취소를 주장할 수 있다.
\kfChoice 통정 허위 표시의 무효는 선의의 제3자에게도 언제나 대항할 수 있다.
\kfChoice 착오에 의한 취소는 법률행위 내용의 중요 부분에 착오가 있어야 한다.
\kfChoice 취소할 수 있는 법률행위는 추인에 의해 확정적으로 유효하게 될 수 있다.
\end{kfChoices}
\end{kfQBlock}

\begin{kfQBlock}{2}{}
\kfQStem{(가)와 (나)를 바탕으로, 무효와 취소에서 '추인'의 효력이 다른 이유로 가장 적절한 것은?}
\begin{kfChoices}
\kfChoice 무효는 처음부터 효력이 없으나, 취소는 유효한 행위를 확정할 수 있기 때문이다.
\kfChoice 무효의 경우에는 제3자가 관여할 수 없지만, 취소에는 제3자가 적극 개입할 수 있기 때문이다.
\kfChoice 무효는 형사법이 다루는 영역이고 취소는 민사법이 관할하는 영역으로 구분되기 때문이다.
\kfChoice 무효인 행위는 반드시 위법한 것이지만, 취소할 수 있는 행위는 어떠한 경우에도 합법이기 때문이다.
\kfChoice 무효 추인에는 반드시 법원의 사전 허가가 필요하지만, 취소 추인에는 그것이 전혀 필요하지 않기 때문이다.
\end{kfChoices}
\end{kfQBlock}

\begin{kfQBlock}{3}{}
\kfQStem{(가)와 (나)를 바탕으로 <보기>를 이해한 내용으로 적절하지 않은 것은?

<보기> 甲은 자신의 토지를 乙에게 가장 매매(통정 허위 표시)하였고, 乙은 이 사실을 모르는 丙에게 그 토지를 매도하였다. 한편 丁은 17세 미성년자로서 법정대리인의 동의 없이 戊에게 자전거를 매도하였으나, 성년이 된 후에도 아무런 행동을 하지 않은 채 5년이 지났다.}
\begin{kfChoices}
\kfChoice 甲乙 간의 가장 매매는 통정 허위 표시이므로 무효이다.
\kfChoice 丙이 甲乙 간 매매가 허위라는 사실을 몰랐다면, 甲은 丙에게 무효를 주장할 수 없다.
\kfChoice 丁의 자전거 매매는 제한능력자의 행위이므로 취소할 수 있는 법률행위이다.
\kfChoice 丁이 성년이 된 후 5년이 지났다면, 추인할 수 있는 날로부터 3년이 경과하여 취소권이 소멸하였다.
\kfChoice 丁의 매매가 취소되면 처음부터 유효였던 것으로 확정된다.
\end{kfChoices}
\end{kfQBlock}

\begin{kfQBlock}{4}{}
\kfQStem{(가)에서 설명한 '무효 행위의 전환(轉換)'에 대한 이해로 적절하지 않은 것은?}
\begin{kfChoices}
\kfChoice 무효인 법률행위가 다른 법률행위의 요건을 갖추고 있어야 한다.
\kfChoice 당사자가 무효임을 알았더라면 다른 법률행위를 하였을 것이라고 인정되어야 한다.
\kfChoice 전환이 인정되면 원래의 법률행위가 소급하여 유효하게 된다.
\kfChoice 전환은 무효인 행위를 다른 법률행위로서 효력을 갖게 하는 것이다.
\kfChoice 전환은 민법 제138조에 근거한다.
\end{kfChoices}
\end{kfQBlock}

\end{multicols}

\kfSecHeader{kfAreaNonlit}{지문}{형법의 고의와 과실 심화 (법학)}

\kfPassageNote{형법의 고의와 과실 심화 (법학)}{%
형법에서 범죄가 성립하려면 구성요건해당성, 위법성, 책임(유책성)이라는 세 요소가 충족되어야 한다. 이 가운데 '고의'와 '과실'은 구성요건의 주관적 요소이자 책임의 요소로 기능하며, 행위자의 심리 상태에 따라 형사책임의 범위가 달라진다.

고의(故意)란 범죄 사실을 인식하고 그 결과 발생을 의욕(意欲)하거나 용인(容認)하는 심리 상태를 말한다. 고의는 다시 확정적 고의와 미필적 고의로 나뉜다. 확정적 고의는 행위의 결과가 반드시 발생할 것임을 확실히 인식하고 이를 의욕하는 경우이다. 예를 들어 상대방을 살해하겠다고 결심하고 칼로 찌르는 행위가 이에 해당한다. 미필적 고의는 결과 발생이 불확실하지만 '발생하더라도 어쩔 수 없다'고 용인하는 경우이다. 예를 들어 폭발물을 설치하면서 행인이 다칠 수 있음을 알면서도 '다쳐도 상관없다'고 생각하는 경우가 미필적 고의이다.

과실(過失)이란 주의의무를 위반하여 범죄 사실을 인식하지 못하거나, 인식하였더라도 결과 발생을 의욕하지 않은 심리 상태를 말한다. 과실은 인식 있는 과실과 인식 없는 과실로 나뉜다. 인식 있는 과실은 결과 발생 가능성을 인식하였으나 '그런 일은 발생하지 않을 것'이라고 신뢰하여 주의를 기울이지 않은 경우이다. 인식 없는 과실은 결과 발생 가능성을 전혀 인식하지 못한 경우이다. 미필적 고의와 인식 있는 과실의 경계는 '용인 여부'에 달려 있다. 결과 발생을 인식하면서도 이를 용인한 경우는 미필적 고의, 결과가 발생하지 않을 것이라고 신뢰한 경우는 인식 있는 과실이다.

구성요건에 해당하는 행위라도 위법성이 조각(阻却)되면 범죄가 성립하지 않는다. 위법성 조각 사유에는 정당방위, 긴급피난, 자구행위, 피해자의 승낙, 정당행위 등이 있다. 정당방위는 자기 또는 타인의 법익에 대한 현재의 부당한 침해를 방위하기 위한 행위로서 상당한 이유가 있을 때 성립한다. 긴급피난은 자기 또는 타인의 법익에 대한 현재의 위난(危難)을 피하기 위한 행위로서, 이로 인한 피해가 피하려는 피해보다 크지 않아야 한다. 정당방위는 '부당한 침해'에 대한 방어인 반면, 긴급피난은 '위난'에 대한 회피라는 점에서 구별된다.

위법한 행위라도 행위자에게 책임을 물을 수 없으면 범죄가 되지 않는다. 책임의 요소에는 책임능력, 위법성의 인식(가능성), 기대가능성이 있다. 책임능력은 사물을 변별(辨別)하고 그에 따라 의사를 결정할 수 있는 능력으로, 14세 미만의 자(형사미성년자)와 심신상실자에게는 책임능력이 인정되지 않는다. 위법성의 인식은 자기 행위가 법에 어긋난다는 점을 인식하는 것이며, 정당한 이유로 위법성을 인식하지 못한 경우에는 벌하지 않는다(형법 제16조). 기대가능성은 행위 당시의 구체적 사정 하에서 적법한 행위를 기대할 수 있었느냐의 문제로, 기대가능성이 없으면 책임이 조각된다.
}
\kfSecHeader{kfAreaNonlit}{문제}{}

\begin{multicols}{2}\raggedcolumns\setlength{\columnsep}{12pt}
\begin{kfQBlock}{1}{}
\kfQStem{윗글의 내용과 일치하지 않는 것은?}
\begin{kfChoices}
\kfChoice 확정적 고의는 결과 발생을 확실히 인식하고 의욕하는 경우이다.
\kfChoice 미필적 고의는 결과 발생이 불확실하지만 발생을 용인하는 경우이다.
\kfChoice 인식 있는 과실은 결과 발생 가능성을 인식하였으나 발생하지 않을 것이라 신뢰한 경우이다.
\kfChoice 미필적 고의와 인식 있는 과실의 구별 기준은 결과 발생의 인식 여부이다.
\kfChoice 인식 없는 과실은 결과 발생 가능성을 전혀 인식하지 못한 경우이다.
\end{kfChoices}
\end{kfQBlock}

\begin{kfQBlock}{2}{}
\kfQStem{윗글에 따를 때, 정당방위와 긴급피난의 차이로 가장 적절한 것은?}
\begin{kfChoices}
\kfChoice 정당방위는 부당한 침해에 대한 방어이고, 긴급피난은 위난에 대한 회피이다.
\kfChoice 정당방위는 타인의 법익만 보호하고, 긴급피난은 자기의 법익만 보호한다.
\kfChoice 정당방위에는 상당성 요건이 없고, 긴급피난에만 있다.
\kfChoice 정당방위는 과실범에만 적용되고, 긴급피난은 고의범에만 적용된다.
\kfChoice 정당방위는 책임 조각 사유이고, 긴급피난은 위법성 조각 사유이다.
\end{kfChoices}
\end{kfQBlock}

\begin{kfQBlock}{3}{}
\kfQStem{윗글을 바탕으로 <보기>의 사례를 판단한 내용으로 적절하지 않은 것은?

<보기> (ㄱ) A는 B를 죽이려고 결심하고 B에게 독약을 먹였다. (ㄴ) C는 건물을 폭파하면서 안에 사람이 있을 수 있음을 알았지만, '죽어도 상관없다'고 생각하고 실행했다. (ㄷ) D는 운전 중 과속하면 사고 날 수 있다고 생각했지만, '내 운전 실력이면 괜찮을 것'이라 믿고 과속하다가 사고를 냈다. (ㄹ) E(12세)가 친구의 물건을 훔쳤다. (ㅁ) F는 길을 가다가 칼로 공격해 오는 상대를 막기 위해 상대를 밀쳐 넘어뜨렸다.}
\begin{kfChoices}
\kfChoice (ㄱ)은 결과를 확실히 인식하고 의욕하였으므로 확정적 고의에 해당한다.
\kfChoice (ㄴ)은 결과 발생을 용인하였으므로 미필적 고의에 해당한다.
\kfChoice (ㄷ)은 결과가 발생하지 않을 것이라 신뢰하였으므로 인식 있는 과실에 해당한다.
\kfChoice (ㄹ)은 14세 미만이므로 책임능력이 없어 범죄가 성립하지 않는다.
\kfChoice (ㅁ)은 긴급피난에 해당하여 위법성이 조각된다.
\end{kfChoices}
\end{kfQBlock}

\begin{kfQBlock}{4}{}
\kfQStem{윗글에서 설명한 '책임'의 세 요소에 대한 이해로 적절하지 않은 것은?}
\begin{kfChoices}
\kfChoice 책임능력은 사물 변별 능력과 의사 결정 능력을 포함한다.
\kfChoice 14세 미만의 자에게는 책임능력이 인정되지 않는다.
\kfChoice 위법성의 인식이 없더라도 그 이유가 정당하지 않으면 처벌할 수 있다.
\kfChoice 기대가능성은 적법 행위를 기대할 수 있었느냐의 문제이다.
\kfChoice 심신상실자에게는 책임능력이 인정되지만 형이 감경된다.
\end{kfChoices}
\end{kfQBlock}

\end{multicols}

\kfSecHeader{kfAreaNonlit}{지문}{행정 대집행과 이행강제금 (법학(복합지문))}

\kfPassageNote{행정 대집행과 이행강제금 (법학(복합지문))}{%
(가) 행정 대집행(代執行)이란 행정법상 의무자가 대체적 작위의무를 이행하지 않을 때, 행정청이 스스로 또는 제3자로 하여금 그 의무를 대신 이행하게 하고 그 비용을 의무자로부터 징수하는 행정상 강제집행 수단이다. 여기서 '대체적 작위의무'란 의무자 본인이 아닌 제3자가 대신 이행할 수 있는 성질의 의무를 말한다. 예를 들어 위법 건축물의 철거 의무는 대체적 작위의무이다. 건축물 철거는 의무자가 아닌 다른 사람도 할 수 있기 때문이다.

대집행의 요건은 세 가지이다. 첫째, 법률에 의해 직접 명령되거나 법률에 의거하여 행정청이 명한 행위로서 타인이 대신 할 수 있는 행위여야 한다(대체적 작위의무). 둘째, 다른 수단으로는 그 이행을 확보하기 곤란해야 한다(보충성). 셋째, 불이행을 방치하면 공익을 크게 해칠 것이 인정되어야 한다. 대집행의 절차는 계고(戒告) → 대집행 영장의 통지 → 대집행 실행 → 비용 징수의 4단계로 진행된다. 계고란 상당한 이행 기간을 정하여 그 기한까지 이행하지 않으면 대집행한다는 뜻을 미리 문서로 알리는 것이다.

대집행에 대한 권리구제로는 행정심판과 행정소송이 가능하다. 대집행의 계고, 영장 통지, 실행 행위 각각에 대해 취소소송을 제기할 수 있으며, 긴급한 경우에는 집행정지를 신청할 수도 있다.

(나) 이행강제금(履行强制金)이란 행정법상 의무를 이행하지 않는 자에 대해 일정 기한까지 의무를 이행하지 않으면 금전적 부담을 부과하겠다고 미리 알림으로써 심리적 압박을 가하여 자발적 이행을 유도하는 간접적 강제 수단이다. 대집행이 행정청이 직접 또는 제3자를 통해 의무를 대신 이행하는 '직접적' 수단인 반면, 이행강제금은 금전적 불이익이라는 '간접적' 수단으로 의무 이행을 확보한다.

이행강제금은 대체적 작위의무뿐 아니라 비대체적 작위의무와 부작위의무에도 부과할 수 있다는 점에서 대집행보다 적용 범위가 넓다. 비대체적 작위의무란 의무자 본인만이 이행할 수 있는 의무를 말한다. 예를 들어 건축법상 시정 명령 불이행 시 이행강제금을 부과할 수 있다. 이행강제금의 가장 큰 특징은 반복 부과가 가능하다는 점이다. 의무자가 의무를 이행할 때까지 반복하여 이행강제금을 부과할 수 있으므로, 지속적인 심리적 압박 효과가 있다. 다만 이행강제금은 과거의 의무 위반에 대한 제재가 아니라 장래의 의무 이행을 확보하기 위한 수단이므로, 행정형벌이나 과태료와는 성격이 다르다.

이행강제금에 대한 권리구제는 이행강제금 부과 처분에 대한 취소소송을 제기할 수 있다. 다만 이행강제금 부과 후 의무를 이행하면 새로운 이행강제금은 부과되지 않으나, 이미 부과된 이행강제금의 납부 의무는 소멸하지 않는다.
}
\kfSecHeader{kfAreaNonlit}{문제}{}

\begin{multicols}{2}\raggedcolumns\setlength{\columnsep}{12pt}
\begin{kfQBlock}{1}{}
\kfQStem{(가)와 (나)의 내용과 일치하지 않는 것은?}
\begin{kfChoices}
\kfChoice 대집행은 대체적 작위의무의 불이행에 대한 직접적 강제 수단이다.
\kfChoice 이행강제금은 금전적 부담을 통한 간접적 강제 수단이다.
\kfChoice 이행강제금은 대체적 작위의무에만 부과할 수 있다.
\kfChoice 대집행의 절차는 계고, 영장 통지, 실행, 비용 징수 순으로 진행된다.
\kfChoice 이행강제금은 의무 이행 시까지 반복 부과가 가능하다.
\end{kfChoices}
\end{kfQBlock}

\begin{kfQBlock}{2}{}
\kfQStem{(가)와 (나)를 바탕으로, 대집행과 이행강제금의 차이로 가장 적절한 것은?}
\begin{kfChoices}
\kfChoice 대집행은 의무를 대신 이행, 이행강제금은 금전 압박으로 이행을 유도한다.
\kfChoice 대집행은 비대체적 의무에만 적용되고, 이행강제금은 대체적 의무에만 적용된다.
\kfChoice 대집행은 권리구제 절차가 봉쇄되어 있지만, 이행강제금은 권리구제가 가능하다.
\kfChoice 대집행은 동일 사안에 반복 실행이 가능하지만, 이행강제금은 단 한 번만 부과한다.
\kfChoice 대집행은 간접적인 행정 수단이고, 이행강제금은 의무자에게 직접적인 수단이다.
\end{kfChoices}
\end{kfQBlock}

\begin{kfQBlock}{3}{}
\kfQStem{(가)와 (나)를 바탕으로 <보기>를 이해한 내용으로 적절하지 않은 것은?

<보기> 甲은 무허가 건축물을 소유하고 있어 관할 구청으로부터 철거 명령을 받았으나 이행하지 않았다. 구청은 甲에게 '2주 내에 철거하지 않으면 대집행한다'는 문서를 보냈다. 한편 乙은 건축법 위반으로 시정 명령을 받았으나 이행하지 않아 이행강제금 500만 원이 부과되었고, 이후 의무를 이행하였다.}
\begin{kfChoices}
\kfChoice 甲의 건축물 철거 의무는 타인이 대신 이행할 수 있으므로 대체적 작위의무에 해당한다.
\kfChoice 구청이 甲에게 보낸 문서는 대집행 절차의 '계고'에 해당한다.
\kfChoice 乙에게 부과된 이행강제금은 과거 위반에 대한 제재 성격의 행정형벌이다.
\kfChoice 乙이 의무를 이행하면 새로운 이행강제금은 부과되지 않는다.
\kfChoice 乙이 의무를 이행하더라도 이미 부과된 500만 원의 납부 의무는 소멸하지 않는다.
\end{kfChoices}
\end{kfQBlock}

\begin{kfQBlock}{4}{}
\kfQStem{(가)에서 설명한 대집행의 '보충성' 요건에 대한 이해로 가장 적절한 것은?}
\begin{kfChoices}
\kfChoice 다른 수단으로는 의무 이행을 확보하기 곤란한 경우에만 대집행이 가능하다는 것
\kfChoice 대집행은 다른 강제 수단을 모두 시행한 후 보충적으로만 사용할 수 있다는 것
\kfChoice 대집행 비용이 부족할 때 보충 재원을 확보해야 한다는 것
\kfChoice 대집행 후 추가 행정 처분으로 보충할 수 있다는 것
\kfChoice 대집행은 의무자의 동의를 보충적으로 받아야 한다는 것
\end{kfChoices}
\end{kfQBlock}

\end{multicols}

\kfSecHeader{kfAreaNonlit}{지문}{혼인의 법적 효력 (법학)}

\kfPassageNote{혼인의 법적 효력 (법학)}{%
우리 민법은 법률혼주의(法律婚主義)를 채택하고 있다. 이는 법률이 정한 요건과 절차를 갖추어 혼인 신고를 해야만 법적으로 유효한 혼인이 성립한다는 원칙이다. 혼인의 성립 요건은 실질적 요건과 형식적 요건으로 나뉜다. 실질적 요건으로는 당사자 간의 혼인 의사의 합치, 혼인 적령(만 18세), 근친혼 금지, 중혼(重婚) 금지 등이 있다. 형식적 요건은 「가족관계의 등록 등에 관한 법률」에 따른 혼인 신고이다. 혼인 신고가 수리(受理)되면 법률상 혼인이 성립한다.

혼인이 성립하면 부부 사이에 여러 법적 효력이 발생한다. 부부는 동거·협조·부양의 의무를 지며, 이는 혼인의 본질적 효력이다. 동거 의무란 부부가 함께 생활할 의무이고, 협조 의무란 혼인 생활을 함께 영위하기 위해 서로 돕는 의무이며, 부양 의무란 부부 일방이 스스로의 자력이나 근로에 의해 생활을 유지할 수 없는 경우 상대방이 생활을 보장해 주는 의무이다. 또한 혼인 중 취득한 재산에 대해서는 부부 각자의 특유재산(고유재산)을 제외하고 부부 공동 생활에 기여한 정도에 따라 재산분할청구권이 인정된다.

일상가사대리권도 혼인의 중요한 효력이다. 부부 일방이 일상의 가사에 관하여 제3자와 법률행위를 한 경우, 다른 일방도 이에 대해 연대책임을 진다(민법 제832조). 여기서 '일상가사'란 부부 공동생활에 통상 필요한 사항, 예를 들어 식료품 구입, 의료비 지출, 자녀 교육비 등을 말한다. 이 규정은 거래 상대방을 보호하기 위한 것이다.

사실혼(事實婚)이란 혼인 신고를 하지 않았으나 실질적으로 부부 공동생활을 영위하는 관계를 말한다. 사실혼은 법률혼의 형식적 요건(혼인 신고)을 갖추지 못한 것이므로 법률상 혼인으로 인정되지 않는다. 그러나 판례와 법률은 사실혼 배우자에게도 일정한 법적 보호를 부여한다. 사실혼 배우자는 동거·협조·부양 의무, 일상가사대리권, 재산분할청구권 등을 인정받을 수 있다. 다만 사실혼 배우자에게는 상속권이 인정되지 않는다는 점이 법률혼과의 가장 큰 차이이다. 즉, 사실혼 배우자는 상대방이 사망하더라도 법률상 상속인이 되지 못한다.

혼인의 해소에는 이혼과 혼인 무효·취소가 있다. 이혼은 부부 쌍방의 합의에 의한 협의이혼과 재판에 의한 재판상 이혼으로 나뉜다. 이혼 시에는 재산분할, 위자료, 자녀의 친권·양육권 등의 문제가 발생한다. 재산분할청구권은 이혼한 날로부터 2년 이내에 행사해야 한다.
}
\kfSecHeader{kfAreaNonlit}{문제}{}

\begin{multicols}{2}\raggedcolumns\setlength{\columnsep}{12pt}
\begin{kfQBlock}{1}{}
\kfQStem{윗글의 내용과 일치하지 않는 것은?}
\begin{kfChoices}
\kfChoice 우리 민법은 혼인 신고를 해야 법적으로 유효한 혼인이 성립하는 법률혼주의를 채택하고 있다.
\kfChoice 혼인의 실질적 요건에는 혼인 의사의 합치, 혼인 적령, 근친혼 금지, 중혼 금지 등이 있다.
\kfChoice 사실혼 배우자에게는 동거·협조·부양 의무와 재산분할청구권이 인정될 수 있다.
\kfChoice 사실혼 배우자에게도 법률혼 배우자와 동일하게 상속권이 인정된다.
\kfChoice 재산분할청구권은 이혼한 날로부터 2년 이내에 행사해야 한다.
\end{kfChoices}
\end{kfQBlock}

\begin{kfQBlock}{2}{}
\kfQStem{윗글에 따를 때, '일상가사대리권'에 대한 설명으로 가장 적절한 것은?}
\begin{kfChoices}
\kfChoice 일상가사 법률행위에 다른 일방이 연대책임을 지도록 한 거래 보호 규정
\kfChoice 부부 일방이 상대방이 보유한 모든 재산을 자유롭게 처분할 수 있는 포괄적 권한
\kfChoice 부부 일방이 상대방을 대리하여 소송을 단독으로 수행할 수 있도록 부여된 권한
\kfChoice 이혼한 후에도 계속해서 유지되는 부부 사이의 법적·도덕적 의무
\kfChoice 부부의 동거 의무를 효과적으로 실현하기 위해 마련된 법적 보충 수단
\end{kfChoices}
\end{kfQBlock}

\begin{kfQBlock}{3}{}
\kfQStem{윗글을 바탕으로 <보기>를 이해한 내용으로 적절하지 않은 것은?

<보기> 甲과 乙은 5년간 혼인 신고 없이 사실혼 관계로 동거하며 공동생활을 해왔다. 이 기간 동안 甲은 근로 소득으로, 乙은 가사노동으로 각각 기여하였다. 어느 날 甲이 사망하여 재산 상속 문제가 발생하였다.}
\begin{kfChoices}
\kfChoice 甲과 乙의 관계는 혼인 신고가 없으므로 사실혼에 해당한다.
\kfChoice 사실혼 관계에서 乙에게는 동거·협조·부양 의무가 인정될 수 있다.
\kfChoice 사실혼 관계가 부당하게 파기된 경우 乙은 재산분할을 청구할 수 있다.
\kfChoice 甲이 사망한 경우 乙은 법률상 상속인으로서 甲의 재산을 상속받을 수 있다.
\kfChoice 甲과 乙이 혼인 신고를 했다면 乙은 法律상 상속인이 될 수 있었을 것이다.
\end{kfChoices}
\end{kfQBlock}

\begin{kfQBlock}{4}{}
\kfQStem{윗글에서 설명한 혼인의 '실질적 요건'과 '형식적 요건'에 대한 이해로 적절하지 않은 것은?}
\begin{kfChoices}
\kfChoice 실질적 요건은 혼인 의사의 합치 등 혼인의 내용에 관한 요건이다.
\kfChoice 형식적 요건은 혼인 신고라는 절차에 관한 요건이다.
\kfChoice 혼인 적령인 만 18세에 이르지 못한 자의 혼인은 실질적 요건을 갖추지 못한 것이다.
\kfChoice 중혼은 실질적 요건의 위반에 해당한다.
\kfChoice 혼인 신고가 수리되면 실질적 요건의 충족 여부와 관계없이 혼인이 영구히 유효하다.
\end{kfChoices}
\end{kfQBlock}

\end{multicols}

\kfSecHeader{kfAreaNonlit}{지문}{법적 의제와 법적 추정 (법학)}

\kfPassageNote{법적 의제와 법적 추정 (법학)}{%
법학에서 '의제(擬制)'와 '추정(推定)'은 일정한 법적 효과를 발생시키기 위해 사실 관계를 법적으로 확정하는 기법이다. 이 두 개념은 유사해 보이지만, 반증(反證)의 허용 여부에서 결정적으로 다르다.

법적 의제란 실제로는 존재하지 않거나 확인할 수 없는 사실을 법률이 존재하는 것으로 확정하여 취급하는 것이다. 의제는 반증에 의해 번복할 수 없다. 법률이 '∼으로 본다(看做)'라는 표현을 사용하면 의제에 해당하는 경우가 많다. 대표적인 예로 실종선고(失踪宣告)에 의한 사망 의제가 있다. 부재자의 생사가 5년간 분명하지 않으면 법원은 실종선고를 할 수 있으며, 이 경우 실종 기간이 만료한 때에 사망한 것으로 '본다'(민법 제28조). 이 사망 의제는 반증을 들어 번복할 수 없고, 오직 실종선고의 취소라는 별도의 절차에 의해서만 번복된다. 또 다른 예로 민법 제187조는 상속·공용징수 등에 의한 부동산 물권 취득은 등기 없이도 그 효력이 생긴다고 규정하는데, 이 역시 등기 없는 물권 변동을 유효한 것으로 '보는' 의제이다.

법적 추정이란 일정한 사실이 존재하면 다른 사실의 존재를 잠정적으로 인정하되, 반증에 의해 번복할 수 있는 것이다. 법률이 '∼으로 추정한다(推定)'라는 표현을 사용하면 추정에 해당한다. 대표적인 예로 친생추정(親生推定)이 있다. 민법 제844조는 '아내가 혼인 중에 포태(胞胎)한 자녀는 남편의 자녀로 추정한다'고 규정한다. 이 추정은 남편 또는 아내가 친생부인(親生否認)의 소를 제기하여 반증을 제출하면 번복될 수 있다. 또한 민법 제197조는 '물건의 점유자는 소유의 의사로 점유한 것으로 추정한다'고 규정한다. 이 점유 추정 역시 상대방이 반증을 제출하여 타주점유(他主占有)임을 증명하면 번복될 수 있다.

의제와 추정은 증명책임(證明責任)의 전환이라는 공통된 기능을 수행한다. 일반적으로 어떤 사실을 주장하는 자가 그 사실을 증명할 책임을 부담한다. 그런데 추정이 적용되면, 추정되는 사실의 존재를 다투는 자가 반증의 책임을 부담하게 된다. 예를 들어 점유 추정에서는, 점유자가 소유 의사를 증명할 필요 없이 상대방이 타주점유임을 증명해야 한다. 의제의 경우에는 증명책임의 전환을 넘어서, 해당 사실 관계를 확정적으로 고정하여 더 이상 다툴 수 없게 만든다.

이 두 기법의 구별이 중요한 이유는 실무적 결과가 크게 다르기 때문이다. 의제의 경우 당사자가 아무리 반대 증거를 제출하더라도 그 법적 효과를 번복할 수 없다. 반면 추정의 경우 반대 증거를 통해 추정된 사실을 뒤집을 수 있으므로, 실제 사실 관계에 부합하는 결론에 도달할 가능성이 열려 있다. 법률이 의제를 선택하느냐 추정을 선택하느냐는 법적 안정성과 구체적 타당성 사이의 균형에 대한 입법적 결단이다.
}
\kfSecHeader{kfAreaNonlit}{문제}{}

\begin{multicols}{2}\raggedcolumns\setlength{\columnsep}{12pt}
\begin{kfQBlock}{1}{}
\kfQStem{윗글의 내용과 일치하지 않는 것은?}
\begin{kfChoices}
\kfChoice 의제는 반증에 의해 번복할 수 없고, 추정은 반증에 의해 번복할 수 있다.
\kfChoice 법률이 '∼으로 본다'라고 표현하면 의제에, '∼으로 추정한다'라고 표현하면 추정에 해당하는 경우가 많다.
\kfChoice 실종선고에 의한 사망 의제는 반증을 통해 번복할 수 있다.
\kfChoice 친생추정은 친생부인의 소를 통해 번복될 수 있다.
\kfChoice 의제와 추정은 모두 증명책임의 전환 기능을 수행한다.
\end{kfChoices}
\end{kfQBlock}

\begin{kfQBlock}{2}{}
\kfQStem{윗글에 따를 때, 법률이 '의제'를 선택하는 것과 '추정'을 선택하는 것의 실무적 차이로 가장 적절한 것은?}
\begin{kfChoices}
\kfChoice 의제는 반대 증거를 제출해도 법적 효과가 번복되지 않지만, 추정은 반대 증거로 뒤집을 수 있어 사실 관계에 부합하는 결론 가능성이 열린다.
\kfChoice 의제는 사실심인 하급심에서만 적용되어 효력을 가지고, 추정은 법률심인 상급심에서만 적용되어 별도의 효력을 발생시키는 차이가 있다.
\kfChoice 의제는 민사 법률 관계에만 한정되어 적용되고, 추정은 형사 절차에서만 활용되어 두 기법은 적용 영역 자체가 완전히 분리되어 있다.
\kfChoice 의제는 당사자 사이의 합의에 의해 그 효력을 번복할 수 있으나, 추정은 어떠한 합의에 의해서도 번복이 절대 불가능한 차이가 있다.
\kfChoice 의제는 당사자의 증명책임 자체를 완전히 면제해 주는 효과를 가지고, 추정은 오히려 당사자의 증명책임을 한층 가중하는 효과를 가진다.
\end{kfChoices}
\end{kfQBlock}

\begin{kfQBlock}{3}{}
\kfQStem{윗글을 바탕으로 <보기>를 이해한 내용으로 적절하지 않은 것은?

<보기> (ㄱ) 甲은 10년 전 실종되어 법원의 실종선고를 받았다. 그 후 甲의 친척이 '甲이 외국에서 살아 있다'는 증거를 제출했다. (ㄴ) 乙은 혼인 중 출생한 자녀에 대해, 남편 丙이 '자신의 자녀가 아니다'라는 유전자 검사 결과를 법원에 제출하였다. (ㄷ) 丁은 토지를 점유하고 있으며, 戊가 丁에게 '당신은 소유 의사 없이 점유하고 있다'고 주장하였다.}
\begin{kfChoices}
\kfChoice (ㄱ)에서 甲의 사망 의제는 친척의 증거 제출만으로는 번복되지 않는다.
\kfChoice (ㄱ)에서 甲의 사망 의제를 번복하려면 실종선고 취소 절차를 밟아야 한다.
\kfChoice (ㄴ)에서 乙의 출생 자녀에 대한 친생추정은 유전자 검사 결과라는 반증으로 번복될 수 있다.
\kfChoice (ㄷ)에서 丁의 점유는 소유 의사에 의한 점유로 추정되므로, 丁이 소유 의사를 증명해야 한다.
\kfChoice (ㄷ)에서 戊가 丁의 점유가 타주점유임을 증명하면 추정이 번복될 수 있다.
\end{kfChoices}
\end{kfQBlock}

\begin{kfQBlock}{4}{}
\kfQStem{윗글에서 의제와 추정의 관계를 '법적 안정성'과 '구체적 타당성'의 관점에서 설명한 내용으로 가장 적절한 것은?}
\begin{kfChoices}
\kfChoice 의제는 법적 안정성을, 추정은 구체적 타당성을 더 중시하는 결단이다.
\kfChoice 의제와 추정 모두 오로지 법적 안정성만을 일관되게 추구하는 기법이다.
\kfChoice 의제는 구체적 타당성을, 추정은 오히려 법적 안정성을 더 중시한다.
\kfChoice 의제와 추정의 선택은 법적 안정성이나 구체적 타당성과 본질적으로 무관하다.
\kfChoice 의제와 추정 모두 구체적 타당성만을 일관되게 추구하는 입법적 기법이다.
\end{kfChoices}
\end{kfQBlock}

\end{multicols}

\kfStartWeekly{패턴 워크북 — 총 ?문항}


\kfSecHeader{kfAreaWeekly}{지문}{문항 1 지문 / 섞어치기}

\kfPassageNote{문항 1 지문 / 섞어치기}{%
무선 통신에 대한 수요가 폭발적으로 증가함에 따라 미국 정부는 ⓐ경매를 통해 통신 사업자들에게 통신 주파수 대역들을 판매하게 되었다. 경매가 시행되기 전에는 ⓑ심사나 ⓒ제비뽑기와 같은 방식을 활용하였다. 심사 방식은 여러 통신 사업자들이 사업 계획서를 제출하면 평가자들이 사업 계획서에 근거하여 무선 통신 사업을 가장 잘 운영할 만한 사업자를 선정하고, 그렇게 선정된 사업자가 주파수를 할당받는 것이었다. 그러나 사업자가 고용한 법률가들과 로비스트들이 서로 자기들의 사업 계획이 최선이라고 주장하는 상황에서 사업 계획서에 대한 평가 절차와 이의 제기 과정을 전부 거쳐 주파수를 할당하기까지 수년이 걸리기도 하였다.

(중략)

1982년 무렵에는 더 이상 심사 방식으로 수많은 주파수를 할당하는 것이 어려운 수준에 이르렀고, 결국 의회는 제비뽑기로 주파수를 판매하는 것에 동의하게 되었다. 제비뽑기 방식은 주파수를 할당하는 데 드는 시간과 노력을 대폭 줄여 주었지만 새로운 문제를 낳았다. 아무 심사도 없이 

(중략)
}
\begin{kfQBlock}{1}{}
\kfQStem{㉡에 대한 설명으로 적절하지 않은 것은?}
\begin{kfChoices}
\kfChoice 경매 참여자들 간의 무리한 경쟁을 예방하기 위해 경매 참여자가 이전 라운드에 입찰한 주파수 개수보다 더 많은 주파수에 입찰할 수 없게 하였다.
\kfChoice 기존의 경매 방식으로 주파수를 경매했을 때 발생했던 문제를 해결하기 위해 폴 밀그럼 등이 제안하였다.
\kfChoice 여러 주파수를 사고 싶은 사업자가 입찰을 주저하게 되는 상황을 막기 위해 모든 주파수를 동시에 경매하였다.
\end{kfChoices}
\end{kfQBlock}

\kfSecHeader{kfAreaWeekly}{지문}{문항 2 지문 / 부정상대어}

\kfPassageNote{문항 2 지문 / 부정상대어}{%
(가) 가상 공간에서 영생의 삶을 누린다는 흥미로운 이야기가 SF 영화의 소재로만 그치지 않는다고 주장하는 입장이 있다. 커즈와일이나 보스트룀은 기억, 가치관, 태도, 정서적 성향처럼 인간의 자아를 구성하는 것들을 특정 정보 패턴으로 만들어 보존할 수 있다면 신체적인 죽음 이후에도 인간이 여전히 생존할 수 있다고 주장한다. 이 입장은 인격 동일성에 관한 심리적 연속성 이론과 관련이 있다. 이 이론에서는 ‘나’가 누구인지를 규정하는 것은 ‘나’의 기억, 신념, 생각, 감정, 희망, 두려움 등의 묶음이라고 보며, 이러한 심리적인 특성들의 집합이 유지·보존되느냐에 ‘나’의 지속 여부가 달려 있다고 말한다.

(중략)

커즈와일이나 보스트룀은 마음이나 정신 상태를 의미하는 심적 상태의 본성을 계산 기능주의의 관점에서 이해한다. 계산 기능주의에서는 심적 상태의 독특한 인과적 역할이 두뇌의 물리적 상태에 의해서 이루어진다고 보기 때문에 정신 작용을 두뇌에 구현된 계산 체계의 작동으로 이해하며, 인간의 사고나

(중략)
}
\begin{kfQBlock}{2}{}
\kfQStem{(가)와 (나)를 바탕으로 <보기>의 주장에 대해 보인 반응으로 적절하지 않은 것은?}
\begin{kfEvidence}<보기> C는 인간 정신을 완전한 계산 절차로 보고, 동일한 알고리즘이 다른 매체에서 실행되면 인격적 연속성도 유지된다고 주장한다. D는 감각 경험이 몸의 감각운동 조건에 깊게 의존하므로, 기능이 같아 보여도 경험 내용은 달라질 수 있다고 주장한다.\end{kfEvidence}
\begin{kfChoices}
\kfChoice 체화된 인지 입장에서는 C의 주장을 지지하며, 신체 차이가 커질수록 경험 동일성이 강화된다고 본다.
\kfChoice 계산 기능주의 입장에서는 C의 주장 중 ‘동일 알고리즘이면 매체가 달라도 기능 재현 가능’이라는 부분에 동의할 여지가 있다.
\kfChoice 체화된 인지 입장에서는 D의 주장처럼 기능 유사성이 경험 내용 동일성을 자동 보장하지 않는다고 볼 수 있다.
\end{kfChoices}
\end{kfQBlock}

\kfSecHeader{kfAreaWeekly}{지문}{문항 3 지문 / 주체객체}

\kfPassageNote{문항 3 지문 / 주체객체}{%
㉠아시시의 프란체스코는 성서 속 예수와 그 제자들을 본받아 지극히 검소한 삶의 방식을 고수하는 수도자들의 모임을 만들었다. 그는 자신을 따르는 수도자들로 하여금 헌금을 받기보다 하층민의 일을 맡아 하거나 차라리 탁발(종교의 수행자가 집집마다 다니며 동냥하는 일)을 하며 살도록 가르쳤다. 그러나 프란체스코가 세상을 떠난 후 상황은 예기치 못한 방향으로 전개되었다. 수도회의 명성이 높아지면서 수도자들을 위한 기부가 이어졌는데, 설립자를 기리는 사업과 늘어난 구성원들을 돌보는 일들에 재정적 어려움을 겪고 있던 수도회는 결국 기부받은 재산을 사용하기는 하되 그것을 소유하지는 않는다는 입장을 취하기에 이르렀다. 이에 호응하여 교황 그레고리우스 9세는 만약 수도자들이 불가피하게 그 재산을 사용해야 한다면, 그에 대한 소유권은 편의상 교황에게 귀속시킬 수 있게 해 주었다.

(중략)

기부에는 음식과 의복의 제공뿐 아니라 숙소나 성당의 건축도 포함되었으며, 이러한 물질적 풍요는 곧 내부적 반발과 외부적 비판을

(중략)
}
\begin{kfQBlock}{3}{}
\kfQStem{㉮에 대한 추론으로 적절하지 않은 것은?}
\begin{kfChoices}
\kfChoice 수도회 자체가 법적 인격을 가지고 재산을 소유하는 것으로 보면 수도자들도 청빈 서약의 위반에 대한 우려를 덜 수 있었을 것이다.
\kfChoice 베르가모의 보나그라티아는 수도회 자체가 법적 인격을 가지므로 수도자들은 단순한 사용을 할 수 있다고 보았을 것이다.
\end{kfChoices}
\end{kfQBlock}

\kfSecHeader{kfAreaWeekly}{지문}{문항 4 지문 / 순서방향}

\kfPassageNote{문항 4 지문 / 순서방향}{%
1908년 조지 하디와 빌헬름 바인베르크는 멘델 유전 법칙에 기초한 수학적 모델을 각각 독립적으로 제시했다. 하디는 집단 내 대립 유전자 A와 a의 빈도를 p와 q로 두고 무작위 교배에서 다음 세대의 유전자형 빈도를 분석했다. 바인베르크도 무작위 결합 조건에서 유전자 빈도가 세대를 거쳐 안정될 수 있음을 통계 자료와 함께 제시했다. 이 두 접근이 결합되어 하디-바인베르크 평형 법칙이 정립되었다.

하디-바인베르크 평형은 ㉠집단이 충분히 크고, ㉡무작위 교배가 일어나며, ㉢새 돌연변이가 없고, ㉣이입·이출 같은 이주가 없고, ㉤자연 선택이 작용하지 않을 때 유지된다. [A] 예를 들어 p=0.7, q=0.3이면 p+q=1이고 유전자형 빈도는 AA=p²=0.49, Aa=2pq=0.42, aa=q²=0.09로 계산되며 p²+2pq+q²=1이 성립한다. 유전자 빈도가 유지되면 다음 세대에서도 유전자형 빈도는 일정하게 유지된다. [A]

(중략)
}
\begin{kfQBlock}{4}{}
\kfQStem{[A]를 바탕으로 <보기>를 이해한 내용으로 적절하지 않은 것은?}
\begin{kfEvidence}<보기> X-연관 열성 질환에서 남성 환자 빈도가 9\%(q=0.09)일 때, p=0.91이며 여성 보인자 빈도는 2pq, 여성 환자 빈도는 q²로 계산된다.\end{kfEvidence}
\begin{kfChoices}
\kfChoice 남성 표현형 빈도를 q로 둘 수 있다.
\kfChoice 여성 환자 빈도가 남성 환자 빈도보다 높다.
\end{kfChoices}
\end{kfQBlock}

\kfSecHeader{kfAreaWeekly}{지문}{문항 5 지문 / 인과도치}

\kfPassageNote{문항 5 지문 / 인과도치}{%
대중적 정부의 형태를 채택하는 국가는 항상 시민 사회의 분열과 파벌 형성이라는 위험에 노출되어 있다. 여기서 파벌이란 공동체 전체의 이익과 일치하지 않는, 일부 구성원들만을 위한 이해관계에 기초하여 결합하고 행동하는 시민 집단을 뜻한다. 이들이 비교적 소수의 인원들로 구성되어 있을 경우에는 다수의 시민들이 선거에 의하여 그들의 준동을 좌절시킬 수 있을 것이나, 그와 반대의 상황이라면 파벌이 공공선 및 파벌에 속하지 않은 시민들의 권리를 자신들의 정념(情念)과 이익을 위해 희생시키는 결과가 초래될 수 있다.

(중략)

미국 건국의 아버지 중 한 사람으로 알려진 제임스 매디슨은 바로 이와 같은 문제가 사회적 다수의 지배를 그 핵심으로 하는 ㉮민주주의에 내재해 있는 치명적인 약점이라고 생각했다. 누구도 자기 자신의 송사에서 재판관이 될 수 없다는 것은 정의의 관점에서 지극히 당연한데, 이 말이 의미하는 바를 곱씹어 볼 때 사회적 다수가 자신들의 이익이 걸려 있는 상황에서 정치적 권력을 중립적으로 행사

(중략)
}
\begin{kfQBlock}{5}{}
\kfQStem{㉮에 대한 설명으로 가장 적절한 것은?}
\begin{kfChoices}
\kfChoice 시민들이 자신의 정치적 힘을 적절한 한도 내에서 절제하며 발휘할 수 없다.
\kfChoice 당파적 이익을 우선시하는 세력이 정치적 다수를 형성하면 그 폐단을 막을 수 없다.
\end{kfChoices}
\end{kfQBlock}

\kfSecHeader{kfAreaWeekly}{지문}{문항 6 지문 / 의도/목적}

\kfPassageNote{문항 6 지문 / 의도/목적}{%
미니멀리즘 음악은 1960년대 이후 미국에서 태동한 현대 음악의 한 흐름으로, 라 몬테 영, 테리 라일리, 스티브 라이히, 필립 글래스 등이 대표적 작곡가로 꼽힌다. 이들은 소수의 \uline{음형}(몇 개의 음이 연속하여 어떤 가락이나 악곡의 요소가 되는 음의 모양)이나 단순한 \uline{패턴}을 반복하는 방식으로 기존 음악의 복잡한 서사나 화성 전개를 과감히 축소하여, 감상자가 소리 그 자체를 경험하도록 유도한다.

전통적인 서양 음악에서 반복은 주제를 강조하여 이를 서사적으로 강화하기 위해 사용되었다. 예컨대 제시부-발전부-재현부로 구성된 소나타 형식은 제시부에서 등장한 주제가 발전부에서 여러 차례 변주되다가 재현부에서 제시부에 등장한 형태로 회귀함으로써, 작곡가가 설정한 주제를 확인 및 강화하는 구조를 이룬다. 철학자 \uline{들뢰즈}의 표현을 빌리면, 이와 같은 ㉠ \uline{‘같은 것의 반복’}에서는 특정 주제가 중심에 놓이고 음악의 각 부분은 그 주제를 재확인하는 역할만…(중략)
}
\begin{kfQBlock}{6}{}
\kfQStem{㉠과 ㉡을 이해한 내용으로 가장 적절하지 않은 것은?}
\begin{kfChoices}
\kfChoice ㉠은 주제를 강조함으로써 통일성을 강화하는 반면, ㉡은 표면적 동일성 속에서 미세한 음향적 차이를 지속적으로 생성한다.
\kfChoice ㉠은 주제를 강조함으로써 통일성을 강화하는 반면, ㉡은 표면적 동일성 속에서 미세한 음향적 차이를 지속적으로 생성한다.
\kfChoice ㉠은 작곡가가 의도한 주제를 강화하는 반복이지만, ㉡은 작곡가의 의도를 감상자가 수동적으로 재확인하게 만드는 반복이다.
\kfChoice ㉠은 반복 과정에서 주제의 변화가 일어나지 않는 반면, ㉡은 반복 속에서 주제가 점진적으로 변화되어 새로운 주제로 발전한다.
\end{kfChoices}
\end{kfQBlock}

\kfSecHeader{kfAreaWeekly}{지문}{문항 7 지문 / 상관관계 반전}

\kfPassageNote{문항 7 지문 / 상관관계 반전}{%
(가) 생수병이나 찻주전자의 내용물을 따를 때 액체가 용기 바깥 면을 타고 흘러 당황한 경험이 있다. 왈칵 따를 때에는 잘 쏟아지던 액체가 속력이 줄면 바깥 면을 타고 흘러내리는데, 이런 현상을 \uline{드리블링} 또는 찻주전자 효과라고 부른다. 찻주전자 효과를 체계적으로 연구한 대표 사례는 2010년 뒤에즈 연구 팀의 실험이다. 이들은 원인으로 주둥이의 젖음성에 주목했고, 이 현상은 점성과 무관하다고 \uline{ⓐ밝혔다}. 젖음성은 고체 표면에서 액체가 퍼지는 정도를 뜻하며, 젖음성이 클수록 액체는 더 얇게 퍼진다. 물에 대한 젖음성이 큰 경우를 친수성, 작은 경우를 소수성이라고 한다. 연구 팀은 유속, 주둥이 젖음성, 가장자리 곡률, 점성에 따라 이탈 각이 어떻게 달라지는지 실험으로 확인했다. 주둥이에서 이탈하는 액체가 연직선과 이루는 각을 이탈 각이라 하며, 이탈 각이 작을수록 찻주전자 효과가 커진다. 유속이 충분히 크면 이탈 각은 90°에 가깝게 유지되지만 유속이 줄어들면 …(중략)
}
\begin{kfQBlock}{7}{}
\kfQStem{(가)와 (나)에서 알 수 있는 내용이 아닌 것은?}
\begin{kfChoices}
\kfChoice 공기 중에서 고체 면에 부착된 물방울에서 고체-액체 간 표면 에너지가 고체-기체 간 표면 에너지보다 커지면 물의 접촉각이 0°에 가까워진다.
\kfChoice 공기 중에서 고체 면에 부착된 물방울에서 고체-액체 간 표면 에너지가 고체-기체 간 표면 에너지보다 커지면 물의 접촉각이 0°에 가까워진다.
\kfChoice 물에 대한 젖음성이 작아지면 소수성이 커진다.
\kfChoice 찻주전자 효과는 액체의 유속이 충분히 크지 않을 때 일어난다.
\end{kfChoices}
\end{kfQBlock}

\kfSecHeader{kfAreaWeekly}{지문}{문항 8 지문 / 필요조건}

\kfPassageNote{문항 8 지문 / 필요조건}{%
(가) 가상 공간에서 영생의 삶을 누린다는 흥미로운 이야기가 SF 영화의 소재로만 그치지 않는다고 주장하는 입장이 있다. 커즈와일이나 보스트룀은 기억, 가치관, 태도, 정서적 성향처럼 인간의 자아를 구성하는 것들을 특정 정보 패턴으로 만들어 보존할 수 있다면 신체적인 죽음 이후에도 인간이 여전히 생존할 수 있다고 주장한다. 이 입장은 인격 동일성에 관한 심리적 연속성 이론과 관련이 있다. 이 이론에서는 ‘나’가 누구인지를 규정하는 것은 ‘나’의 기억, 신념, 생각, 감정, 희망, 두려움 등의 묶음이라고 보며, 이러한 심리적인 특성들의 집합이 유지·보존되느냐에 ‘나’의 지속 여부가 달려 있다고 말한다.

커즈와일이나 보스트룀은 마음이나 정신 상태를 의미하는 심적 상태의 본성을 계산 기능주의의 관점에서 이해한다. 계산 기능주의에서는 심적 상태의 독특한 인과적 역할이 두뇌의 물리적 상태에 의해서 이루어진다고 보기 때문에 정신 작용을 두뇌에 구현된 계산 체계의 작동으로 이해하며, 인간의 사…(중략)
}
\begin{kfQBlock}{8}{}
\kfQStem{[05] 윗글의 내용과 일치하지 않는 것은?}
\begin{kfChoices}
\kfChoice 인간의 뇌를 생물학적 컴퓨터로 보면 물리적 기반이 없더라도 다양한 방식으로 정신의 구현이 가능하다.
\kfChoice 인간의 뇌를 생물학적 컴퓨터로 보면 물리적 기반이 없더라도 다양한 방식으로 정신의 구현이 가능하다.
\kfChoice 커즈와일이나 보스트룀은 심리적인 특성들의 집합을 인간의 자아를 구성하는 요소로 여긴다.
\kfChoice 계산 기능주의의 관점에서는 두뇌가 수행하는 모종의 기능적 역할이 심적 상태를 만든다고 말한다.
\end{kfChoices}
\end{kfQBlock}

\kfSecHeader{kfAreaWeekly}{지문}{문항 9 지문 / 결합/분리}

\kfPassageNote{문항 9 지문 / 결합/분리}{%
근대 경제학은 국민 소득, 물가, 고용, 경기 변동 등 경제의 \uline{거시적 현상}을 설명하기 위해 다양한 \uline{모형}을 \uline{ⓐ발전시켜} 왔다. \uline{고전학파}는 시장의 \uline{자율 조정} 능력을 강조하며, 가격과 임금의 \uline{유연성}을 전제로 \uline{완전 고용}이 자연스럽게 달성된다고 보았다. 그러나 20세기 초 \uline{대공황}을 계기로 시장의 자율 조정에 한계가 있다는 비판이 제기되었고, 이는 \uline{케인스 경제학}의 등장으로 이어졌다. 케인스는 \uline{총수요}(한 경제 내에서 주어진 기간과 물가 수준하에 최종 재화와 서비스에 대한 수요의 총합)가 자동으로 \uline{총공급}(주어진 기간과 물가 수준하에 국가 경제에서 재화와 서비스에 대한 공급의 총합)에 맞춰진다고 보는 고전학파의 이론을 비판하며, 낮은 총수요가 소득의 감소 및 실업률 상승을 초래할 수 있다고 보았다. 따라서 그는 정부 지출 확대 등 정책 개입을 통해 총수요를 관리해야…(중략)
}
\begin{kfQBlock}{9}{}
\kfQStem{다음 <보기>를 IS-MP 모형의 곡선 이동으로 연결한 것으로 가장 적절하지 않은 것은?}
\begin{kfEvidence}물가 상승 압력이 커지자 중앙은행은 물가 안정을 최우선 목표로 삼아 긴축적 정책 기조로 선회했고, 정부는 재정 지출을 축소하였다.\end{kfEvidence}
\begin{kfChoices}
\kfChoice IS 곡선은 왼쪽으로, MP 곡선은 위쪽으로 이동한다.
\kfChoice IS 곡선은 왼쪽으로, MP 곡선은 위쪽으로 이동한다.
\kfChoice IS 곡선은 오른쪽으로, MP 곡선은 아래쪽으로 이동한다.
\kfChoice IS 곡선은 오른쪽으로, MP 곡선은 위쪽으로 이동한다.
\end{kfChoices}
\end{kfQBlock}

\kfSecHeader{kfAreaWeekly}{지문}{문항 10 지문 / 조건 왜곡}

\kfPassageNote{문항 10 지문 / 조건 왜곡}{%
어느 시대에서나 기적에 대한 보고는 꾸준히 있었고, 오늘날에도 기적을 목격했다는 사람은 부지기수로 많다. 죽은 사람이 살아난다든가 바위가 스스로 움직여 떠다니는 것을 보았다는 증언이 그런 예이다. 어떤 사건을 \uline{㉠기적적 사건}이라고 해야 할까?

기적이 특이한 사건이라야 한다는 점에는 동의할 수 있다. 그러나 일어나기 힘든 일이긴 하지만 자연법칙으로 설명될 수 있는 한 그것은 기적적 사건으로 간주될 수 없다. 이순신 장군의 명량 대첩처럼 13척의 배로 133척의 왜적을 물리친 일은 일어나기에 아주 힘든 일일 뿐이지 자연법칙에 어긋나지는 않는다. 그렇다면 우리에게 알려진 어떤 자연법칙으로도 설명할 수 없는 사건은 기적일까? 사진 감광판을 완전히 어두운 곳에 줄곧 보관했는데도 감광판이 빛에 노출된다는 사실을 맨 처음 알게 되었을 때, 이 사건은 당시에 알려진 어떤 자연법칙으로도 설명할 수 없었다. 그러나 머지않아 방사능에 관한 과학이 정립되어 이 진기한 현상에 대한 과학적 설명…(중략)
}
\begin{kfQBlock}{10}{}
\kfQStem{<보기>에서 2문단의 ‘기적적 사건(㉠)’ 정의에 해당하는 진술만을 고른 것은?}
\begin{kfEvidence}ㄱ. 매우 드문 자연재해가 일어나 많은 사람이 ‘기적’이라고 불렀다. ㄴ. 어떤 사람이 초자연적 힘으로 돌을 순간이동시켰다고 주장했다. ㄷ. 과학이 발전하면서 설명할 수 있게 된 이전의 ‘미지의 현상’이 있었다. ㄹ. 어떤 사람이 아무 도구도 없이 공중에 떠올랐다고 한다. ㅁ. 사람들이 ‘기적 같은 일’이라며 떠들었지만, 사실은 착시였다.\end{kfEvidence}
\begin{kfChoices}
\kfChoice ㄴ, ㄹ
\kfChoice ㄴ, ㄹ
\kfChoice ㄱ, ㄴ
\kfChoice ㄴ, ㄷ, ㄹ
\end{kfChoices}
\end{kfQBlock}

\kfSecHeader{kfAreaWeekly}{지문}{문항 11 지문 / 긍부정 반전}

\kfPassageNote{문항 11 지문 / 긍부정 반전}{%
\#\#\# [34 \textasciitilde{} 38] 다음 글을 읽고 물음에 답하시오.

(가)   산아. 우뚝 솟은 푸른 산아. 철철철 흐르듯 짙푸른 산아. 숱한 나무들, 무성히 무성히 우거진 산마루에, 금빛 기름진 햇살은 내려오고, 둥둥 산을 넘어, 흰 구름 건넌 자리 씻기는 하늘. 사슴도 안 오고 바람도 안 불고, 넘엇 골 골짜기서 울어 오는 뻐꾸기…….

산아. 푸른 산아. 네 가슴 향기로운 풀밭에 엎드리면, 나는 가슴이 울어라. 흐르는 골짜기 스며드는 물소리에, 내사 줄줄줄 가슴이 울어라. 아득히 가 버린 것 잊어버린 하늘과, 아른아른 오지 않는 보고 싶은 하늘에, 어찌면 만나도질 볼이 고운 사람이, 난 혼자 그리워라. 가슴으로 그리워라.

티끌 부는 세상에도 벌레 같은 세상에도 눈 맑은, 가슴 맑은, 보고지운 나의 사람. 달밤이나 새벽녘, 홀로 서서 눈물 어릴 볼이 고운 나의 사람. 달 가고, 밤 가고, 눈물도 가고, 틔어 올 밝은 하늘 빛난 아침 이르면, 향기로운 이슬밭 푸른 언덕을, 총총총 달려도…(중략)
}
\begin{kfQBlock}{11}{}
\kfQStem{(다)의 구절과 관련지어 지지(止止)에 대해 설명한 내용으로 가장 적절하지 않은 것은?}
\begin{kfChoices}
\kfChoice ‘거처를 이렇게 이름한 것’과 관련되어 자신의 본성에 맞는 삶을 실천하고자 하는 글쓴이의 지향을 반영한다.
\kfChoice ‘거처를 이렇게 이름한 것’과 관련되어 자신의 본성에 맞는 삶을 실천하고자 하는 글쓴이의 지향을 반영한다.
\kfChoice ‘어찌 해칠 자가 있겠는가?’를 통해 위험을 감수하는 삶의 태도가 중요하다는 글쓴이의 믿음을 강조한다.
\kfChoice ‘그렇지 않으면 사람들이 재앙으로 여기고’를 통해 어쩔 수 없이 그쳐야 하는 상황에 대한 글쓴이의 두려움을 나타낸다.
\end{kfChoices}
\end{kfQBlock}

\kfSecHeader{kfAreaWeekly}{지문}{문항 12 지문 / 비교대상 뒤바꿈}

\kfPassageNote{문항 12 지문 / 비교대상 뒤바꿈}{%
\#\#\# [22～27] 다음 글을 읽고 물음에 답하시오.

(가)   아득한 옛날에 나는 떠났다   ㉠ 부여를 숙신을 발해를 여진을 요를 금을   흥안령을 음산을 아무우르를 숭가리를   범과 사슴과 너구리를 배반하고   송어와 메기와 개구리를 속이고 나는 떠났다

나는 그때   ㉡ 자작나무와 이깔나무의 슬퍼하던 것을 기억한다   갈대와 장풍의 붙드던 말도 잊지 않았다   ㉢ 오로촌이 멧돝을 잡아 나를 잔치해 보내던 것도   쏠론이 십릿길을 따라 나와 울던 것도 잊지 않았다

나는 그때   ㉣ 아무 이기지 못할 슬픔도 시름도 없이   다만 게을리 먼 앞대로 떠나 나왔다   그리하여 따사한 햇귀에서 하이얀 옷을 입고 매끄러운 밥을 먹고 단 샘을 마시고 낮잠을 잤다   밤에는 먼 개소리에 놀라나고   아침에는 지나가는 사람마다에게 절을 하면서도   나는 나의 부끄러움을 알지 못했다

그동안 돌비는 깨어지고 많은 은금보화는 땅에 묻히고 가마귀도 긴 족보를 이루었는데   이리하여 또 한 아득한…(중략)
}
\begin{kfQBlock}{12}{}
\kfQStem{(가)～(다)의 공통점으로 가장 적절한 것은?}
\begin{kfChoices}
\kfChoice 빗대어 표현하는 방식으로 ‘나’의 인식을 드러내고 있다.
\kfChoice 빗대어 표현하는 방식으로 ‘나’의 인식을 드러내고 있다.
\kfChoice 비판적 태도로 현실의 부정적 측면을 부각하고 있다.
\kfChoice 역사적 상황을 묘사하여 비극적 현실을 부각하고 있다.
\end{kfChoices}
\end{kfQBlock}

\kfSecHeader{kfAreaWeekly}{지문}{문항 13 지문 / 주체/객체 혼동}

\kfPassageNote{문항 13 지문 / 주체/객체 혼동}{%
\#\#\# [31\textasciitilde{}34] 다음 글을 읽고 물음에 답하시오.

동방이 차차 밝아 오매 마침 영릉골 관비 한 사람이 외촌에 가다가 돌아오는 길에 청수 가에 다다르니 어떤 여자가 물가에서 통곡하며 물에 빠져 죽고자 하거늘 급히 쫓아와 강 낭자를 붙들어 물가에 앉히고 이유를 물으니라. 그 후에 제 집으로 가자 하나 낭자 한사코 죽으려 하거늘 관비 여러 가지로 타일러 데리고 와서 수양딸로 정한 후에 자색과 태도를 살펴보니 천상 선녀 같은지라. 이 고을 동리마다 수청을 드리면 천금의 재산이 부럽지 않으며, 만 량 가진 태수를 원하겠느냐. 만 가지로 달래어 다른 데로 못 가게 하더라.   각설. 이때에 유충렬이 강 승상의 집을 떠나서 서쪽 하늘을 바라보고 정처 없이 가며 신세를 생각하니, 속절없고 하릴없다. 이제는 아무것도 할 수 없구나. 산중에 들어가 삭발하고 중이 되어 훗날의 도를 닦으리라 하고 청산을 바라보고 종일토록 가더니 한 곳에 다다르더라. 앞에 ㉠ 큰 산이 있으되 천 개의 봉우리와 만 …(중략)
}
\begin{kfQBlock}{13}{}
\kfQStem{[A]와 [B]에 대한 설명으로 가장 적절한 것은?}
\begin{kfChoices}
\kfChoice [A]는 이전 사건에 대한 정보를 전달하고, [B]는 변화된 현재 상황에 대한 심리를 드러내고 있다.
\kfChoice [A]는 이전 사건에 대한 정보를 전달하고, [B]는 변화된 현재 상황에 대한 심리를 드러내고 있다.
\kfChoice [A]는 [B]와 달리 과거 사건을 근거로 들며 문제 해결을 유보하고 있다.
\kfChoice [B]는 [A]와 달리 불가능한 상황을 설정하여 상대를 설득하고 있다.
\end{kfChoices}
\end{kfQBlock}

\kfSecHeader{kfAreaWeekly}{지문}{문항 14 지문 / 내용과 방법 혼동}

\kfPassageNote{문항 14 지문 / 내용과 방법 혼동}{%
\#\#\# [32～34] 다음 글을 읽고 물음에 답하시오.

(가)   장풍에 돛을 달고 육선이 함께 떠나   삼현과 군악 소리 해산을 진동하니   물속의 어룡들이 응당히 놀라리라   해구를 얼른 나서 오륙도를 뒤 지우고   고국을 돌아보니 야색이 아득하여   아무것도 아니 뵈고 연해 각진포에   불빛 두어 점이 구름 밖에 뵐 만하다   배 방에 누워 있어 내 신세를 생각하니   가뜩이 심란한데 대풍이 일어나서   태산 같은 성난 물결 천지에 자욱하니   크나큰 만곡주가 나뭇잎 불리이듯   하늘에 올랐다가 지함에 내려지니   열두 발 쌍돛대는 차아처럼 굽어 있고   쉰두 폭 초석 돛은 반달처럼 배불렀네   (중략)   날이 마침 극열하고 석양이 비치어서   끓는 땅에 엎디어서 말씀을 여쭈오니   속에서 불이 나고 관대에 땀이 배어   물 흐르듯 하는지라 나라께서 보시고서   너희 더위 어려우니 먼저 나가 쉬라시니   곡배하고 사퇴하니 천은이 망극하다   더위를 장히 먹어 막힐 듯하는지라…(중략)
}
\begin{kfQBlock}{14}{}
\kfQStem{(가), (나)의 표현상 특징에 대한 설명으로 가장 적절하지 않은 것은?}
\begin{kfChoices}
\kfChoice (가)는 사물의 형태가 변화한 모습을 묘사하여 외부 환경의 영향력을 부각하고 있다.
\kfChoice (가)는 사물의 형태가 변화한 모습을 묘사하여 외부 환경의 영향력을 부각하고 있다.
\kfChoice (가)는 과거를 회상하는 표현을 통해 현재 상황에 대한 아쉬움을 드러내고 있다.
\kfChoice (나)는 계절을 나타내는 어휘를 활용해 애달픈 정서를 부각하고 있다.
\end{kfChoices}
\end{kfQBlock}

\kfSecHeader{kfAreaWeekly}{지문}{문항 15 지문 / 내용 왜곡}

\kfPassageNote{문항 15 지문 / 내용 왜곡}{%
[16 \textasciitilde{} 20] 다음 글을 읽고 물음에 답하시오.

(가)   잠깐 초록을 본 마음이 돌아가지 않는다.   초록에 붙잡힌 마음이   초록에 붙어 바람에 세차게 흔들리는 마음이   종일 떨어지지 않는다   여리고 연하지만 불길처럼 이글이글 휘어지는 초록   땅에 박힌 심지에서 끝없이 솟구치는 초록   나무들이 온몸의 진액을 다 쏟아내는 초록   ㉠ 지금 저 초록 아래에서는   얼마나 많은 잔뿌리들이 발끝에 힘주고 있을까   초록은 수많은 수직선 사이에 있다   수직선들을 조금씩 지우며 번져가고 있다   직선과 사각에 밀려 꺼졌다가는 다시 살아나고 있다   흙이란 흙은 도로와 건물로 모조리 딱딱하게 덮인 줄 알았는데   이렇게 많은 초록이 갑자기 일어날 줄은 몰랐다   아무렇게나 버려지고 잘리고 갇힌 것들이   자투리땅에서 이렇게 크게 세상을 덮을 줄은 몰랐다   콘크리트 갈라진 틈에서도 솟아나고 있는   저 저돌적인 고요   단단하고 건조한 것들에게 옮겨 붙고 있는   저 촉촉한 불…(중략)
}
\begin{kfQBlock}{15}{}
\kfQStem{(나)에 대한 이해로 적절하지 않은 것은?}
\begin{kfChoices}
\kfChoice <제3수>의 ‘이제야 아무리 애달픈들’과 <제6수>의 ‘내 못하여 애달프니’에는 세월의 무상감에서 벗어나고자 하는 심리가 드러나는군.
\kfChoice <제3수>의 ‘이제야 아무리 애달픈들’과 <제6수>의 ‘내 못하여 애달프니’에는 세월의 무상감에서 벗어나고자 하는 심리가 드러나는군.
\kfChoice <제1수>의 ‘어져 내 일이야’에 담긴 한탄은, <제2수>의 ‘장부의 허다 사업’을 못 다 한 데서 비롯되는군.
\kfChoice <제1수>의 ‘노대에 도상비로다’에 담긴 애상감은, <제4수>의 ‘늙어지면 거짓 것이’로 이어지는군.
\end{kfChoices}
\end{kfQBlock}

\kfSecHeader{kfAreaWeekly}{지문}{문항 16 지문 / 과잉 해석}

\kfPassageNote{문항 16 지문 / 과잉 해석}{%
\#\#\# 2025년 11월 고3 수능 

\#\#\# [18\textasciitilde{}21] 다음 글을 읽고 물음에 답하시오.

[중모리] 그때에 사슴이 발론하되 근래 인간이 하 무서워 짐승을 잡아먹기 온갖 꾀가 다 생기고 산중에 수목이 없어 은신할 곳 없어지니 각기 의견 들어 보면 방책이 있을런가 이 모임을 했사오니 수령님의 좋은 꾀를 일러 주옵소서

[아니리] 호랑이가 수령 말을 듣더니마는 거드름을 피우며 오늘은 노소고하를 막론하고 자세히 말해 보라 토끼가 여짜오되

[자진모리] 사냥개라 허는 것은 같은 우리 모족(毛族)으로 사람 집에 기식허니 제 무슨 아첨으로 내 잘 맡는 자랑허여 심산궁곡 층암절벽 찾고 찾어 들어와 동제 간 살해만 허니 수령님 이후로는 사냥개를 있는 대로 다 잡어 잡수오면 그 덕이 모든 금수에게 미치오리다

[아니리] 호랑이 듣더니만 다 잡어 먹었으면 네 원통함도 풀고 나도 배부른 꼴을 보련마는 일등 포수가 따러다녀 어설피 물랴다가 조총에 불이 번듯 탄환이 쑥 나오면 거 내 신세는 어쩔 것이…(중략)
}
\begin{kfQBlock}{16}{}
\kfQStem{다음에 제시된 선생님의 설명을 참고하여 윗글을 감상한 내용으로 적절하지 않은 것은? [3점]}
\begin{kfChoices}
\kfChoice 자라가 '동의보감'을 떠올린 데서, 현실의 의서를 우화 속 상황에 중첩함으로써 명약을 탐하는 속내를 지식을 내세워 숨기는 위선적 인간에 대한 비판이 드러남을 알 수 있군.
\kfChoice 자라가 '동의보감'을 떠올린 데서, 현실의 의서를 우화 속 상황에 중첩함으로써 명약을 탐하는 속내를 지식을 내세워 숨기는 위선적 인간에 대한 비판이 드러남을 알 수 있군.
\kfChoice '사냥개'에 대한 토끼의 평가에서, 현실에서 사냥개가 사람에게 길들여진 것을 우화 속 상황에 중첩함으로써 강자의 환심을 사 이익을 얻는 인간에 대한 비판이 드러남을 알 수 있군.
\kfChoice '포수'에 대한 호랑이의 태도에서, 현실의 인간이 지닌 힘을 우화 속 인물들의 위계질서에 중첩함으로써 권력자가 상대에 대한 두려움을 보여 위신을 잃는 상황이 강조됨을 알 수 있군.
\end{kfChoices}
\end{kfQBlock}

\kfSecHeader{kfAreaWeekly}{지문}{문항 17 지문 / 범위 오류}

\kfPassageNote{문항 17 지문 / 범위 오류}{%
\#\#\# 2024년 11월 고3 수능

[22-27] 다음 글을 읽고 물음에 답하시오.

(가)   배를 민다   배를 밀어보는 것은 아주 드문 경험   희번덕이는 잔잔한 가을 바닷물 위에   배를 밀어넣고는   온몸이 아주 추락하지 않을 순간의 한 허공에서   밀던 힘을 한껏 더해 밀어주고는   아슬아슬히 배에서 떨어진 손, 순간 환해진 손을   허공으로부터 거둔다

사랑은 참 부드럽게도 떠나지   뵈지도 않는 길을 부드럽게도

배를 한껏 세게 밀어내듯이 슬픔도   그렇게 밀어내는 것이지

배가 나가고 남은 빈 물 위의 흉터   잠시 머물다 가라앉고

그런데 오, 내 안으로 들어오는 배여   아무 소리 없이 밀려들어오는 배여   - 장석남, <배를 밀며>

(나)   당신……, 당신이라는 말 참 좋지요, 그래서 불러봅니다 킥킥거리며 한때 적요로움의 울음이 있었던 때, 한 슬픔이 문을 닫으면 또 한 슬픔이 문을 여는 것을 이만큼 살아옴의 상처에 기대, 나 킥킥……, 당신을 부릅니다 단…(중략)
}
\begin{kfQBlock}{17}{}
\kfQStem{문항별}
\begin{kfChoices}
\kfChoice 번은 '통제할 수 없는 익명의 욕구'를 상대를 향한 사랑이 운명적이어서 멈출 수 없다는 뜻으로 해석했습니다.
\kfChoice 번은 '통제할 수 없는 익명의 욕구'를 상대를 향한 사랑이 운명적이어서 멈출 수 없다는 뜻으로 해석했습니다.
\kfChoice 번은 배가 들어오는 것을 '대상과의 재회가 예상대로 이루어짐'이라고 해석했습니다.
\kfChoice 번은 '당신'을 화자의 내면에 사는 '병자'로서 연민의 대상이라고 설명했습니다.
\end{kfChoices}
\end{kfQBlock}

\kfSecHeader{kfAreaWeekly}{지문}{문항 18 지문 / 시점/화자 혼동}

\kfPassageNote{문항 18 지문 / 시점/화자 혼동}{%
[18 \textasciitilde{} 23] 다음 글을 읽고 물음에 답하시오.

(가)   삼 년을 임을 떠나 해도(海島)에 유배되니   ㉠ 내 언제 무심하여 임에게 득죄했나   임이 언제 박정(薄情)하여 날 대접 소홀히 했나   내 얼굴 고왔던지 질투하는 건 뭇 여자로다   유한한* 이내 몸을 음란하다 이르로세   (중략)   긴 소매 들고 앉아 옛 잘못을 헤아리니   우직하기 본성이오 망령됨도 내 죄로되   근본을 생각하면 임 위한 정성일세   일월 같은 우리 임이 거의 아니 굽어볼까   날 살리신 이 은혜를 결초(結草)하기 생각하나   광주리의 가을 부채 어느 날 다시 날꼬   황금을 못 얻으니 장문부*를 어이 사리*   \textit{마름과 연(蓮)으로 옷을 짓고 부용(芙蓉)으로 치마 지어}   \textit{상자 안에 두어신들 눌 위하여 단장할꼬}   \textit{고향에 돌아갈 꿈 벽해(碧海)를 밟아 건너}   \textit{옥루(玉樓) 높은 곳에 밤마다 임을 모셔}   \textit{일당우불에 수답이 여향하니}   가까이 다가앉아 귀신을 묻던 가태부 이러한가…(중략)
}
\begin{kfQBlock}{18}{}
\kfQStem{(가)와 (나)의 시어에 대한 이해로 가장 적절한 것은?}
\begin{kfChoices}
\kfChoice (가)와 (나)의 ‘꿈’에는 모두 현재 상황에서 화자가 갖는 소망이 투영되어 있다.
\kfChoice (가)와 (나)의 ‘꿈’에는 모두 현재 상황에서 화자가 갖는 소망이 투영되어 있다.
\kfChoice (가)의 ‘닭’은 (나)의 ‘꾀꼬리’와 달리 꿈속에서의 임과의 만남을 방해하는 존재이다.
\kfChoice (나)의 ‘진짜 진주’에는 (가)의 ‘치마’와 달리 임에 대한 화자의 애정이 담겨 있다.
\end{kfChoices}
\end{kfQBlock}



\end{document}
